% Generated by build_chapter.py from the GY1979 files. Do not edit by hand; edit the files and rebuild.
\chapter{The academic year 1979--80}\label{ch:1979}
\noindent\textbf{Scope and editorial basis.}
This chapter covers all 27 MCCME records 462--488 in manuscript source order, with
dated assessments as of 10 November 2026. The Russian transcription is the textual
authority; the scans arbitrate doubtful formulas and attribution. The supplied
79-page envelope includes three boundary pages: P000840 is already represented by
1978-460.12--1978-460.20, and P000917--P000918 belong to the session dated
22 September 1980. Neither is counted again as a 1979--80 occurrence. The present
year's content occupies P000841--P000916. Record 481 is retained as a catalog
heading. The critical date of record 473 is 29 October 1979, against the catalog's
29 September; on 19 November, Gelfand's record 477 precedes Beilinson's record 476
in the manuscript.

The occurrence index contains 252 distinct mathematical items, 10
problem threads, and 104 bibliography entries. Item boundaries distinguish
definitions, constructions, mathematical claims, explicit questions, and identifiable
interventions. A class records historical fate; literal correctness and strength of
publication match are separate assessments.
\begin{center}
\begin{tabular}{lrrrrrr}
\toprule
Class & A & B & C & D & E & F\\
\midrule
Items & 12 & 7 & 3 & 0 & 203 & 27\\
\bottomrule
\end{tabular}
\end{center}

\noindent\textbf{The main lines.}
Much of the autumn is an extended mathematical encounter with quantum field theory:
Kallosh's reports on gravity and supergravity; Batalin's careful progression through
constraints, gauge slices, ghosts and BRST identities; and Kobzarev's account of
gauge fields, confinement and particle masses. Formal functional integrals are kept
separate from convergent finite-dimensional formulas. The early BRST material is
not retrospectively identified wholesale with the later BV formalism.

The surrounding mathematics is unusually broad. Zakharov treats resonance and
integrability in September and long-wave hydrodynamics in February. Gindikin's
twistor geometry and Berezin's supermathematics lead into projective bundles and
exterior algebras. Shabat's double-translation surfaces require the missing
nondegeneracy hypotheses in the classical Jacobian theorem. H\o egh-Krohn's energy
representations retain a delicate two-dimensional question. In November, Gelfand's
dilogarithm list is followed by Beilinson's specific regulator construction; the
latter has a close 1980 publication match. January brings Migdal's loop equations,
Zelevinsky's Kazhdan--Lusztig material and MacPherson's topology. Olshanskii's
representations and MacPherson's combinatorial characteristic classes distinguish
pre-talk results from genuinely subsequent developments. Their precise source
locators and chronology are given in the corresponding entries.

\noindent\textbf{What remains and what is developed here.}
The chapter does not turn a broad research programme into a definite new conjecture.
Continuum confinement is linked to the older static-potential problem; a closed
large-rank equation is not itself an area-law proof. The general two-dimensional
energy-representation question is distinguished from restricted cases and from a
published proof claim with a documented gap. Negative literature searches are reported as retrieval
limits, not proof that an idea was abandoned.

Berezin's two explicit matrix questions support a bounded mathematical treatment.
The normal-form question has an algorithmic answer through a finite-dimensional
matrix-algebra action; an explicit matrix-pair reduction explains the limits of a
simple Jordan analogy. For the even commutant of an arbitrary odd two-by-two matrix
over four Grassmann generators, the development gives a complete coefficient
parametrization, all dimension strata, and full proofs. Its exact computations were
implemented in two independent ways and agree in all 36 cases. These are proofs
and checks supplied here, with no priority claim. The literal finite commutant task
is elementary linear algebra, classified E; no open problem is manufactured to
justify its explanatory development. A separate compactly supported SU(2) example
checks the documented flaw in the energy-representation argument, without claiming
to resolve irreducibility.

\noindent\textbf{Continuity and corrections.}
The earlier occurrence index supplies links to exterior-module descriptions of
projective bundles, combinatorial Pontryagin classes, continuum confinement, and
the preceding year's integrable-system discussions. A common topic is not enough
to merge two problems with different quantifiers. The 106 corrections distinguish
notebook errors faithfully transcribed, unresolved readings, and provenance
discrepancies. In particular the adjacent-year
spillover, the October date, regulator signs, normalized Wilson traces, and the
Pontryagin construction's low-dimensional cases must remain visible to a reader.
The report gives the counting rules, evidence limits, search record and website
presentation recommendations.

\GStalk{462}{1979-09-24}{Renata Kallosh}{P000841-P000842}{Quantum gravity, higher-spin indices and superselfduality}
The actual new talk begins on P000841 and ends on P000842. Kallosh reports on the Cambridge relativity meeting, then describes higher-spin zero-mode formulae and self-dual supergravity. P000840 is a carry-in from the preceding notebook material: its vertex-model, Yang--Baxter and Ising observations were already indexed as 1978-460.12--20. They are preserved as provenance, not reassigned to Kallosh or duplicated as 1979 occurrences. The overlapping catalogue page interval is retained in the provenance crosswalk; the source-page field gives the actual talk span. The programme of a superspace index theorem is distinguished from the concrete index computations and from any claim of unrestricted ultraviolet finiteness.

\begin{GSitem}{1979-462.1}{E}{Spacetime foam as a quantum-gravity proposal}
\GSfield{Kind; source}{construction; \GSpage{P000841}}
\GSfield{The notes}{Small-scale spacetime topology should include many little handles; Hawking's paper on spacetime foam is named.}
\GSfield{Statement}{The identifiable proposal is to include nontrivial four-dimensional topology among Euclidean gravitational configurations. The notes specify neither a measure on metrics modulo diffeomorphisms nor a convergent sum over topologies, and therefore state no rigorous existence theorem for a quantum theory.}
\GSfield{Reconstruction}{This is a physical construction programme. A theorem claiming convergence of the gravitational functional integral would be an unjustified addition.}
\GSfield{Fate}{Hawking's named paper is already published in Nuclear Physics B 144 (1978), pp.~349--362 \citep{hawking1978foam}. The author's bibliography verifies the title, journal and date; its full argument was not retrieved.}
\GSfield{Present status}{As of November 2026, the notebook is identifiable as exposition of that preexisting proposal, not as a precise unresolved mathematical problem.}
\GSfield{Assessment}{literal statement: not determinable; publication match: exact; verification: publisher metadata checked; development: annotated; confidence: high}
\end{GSitem}
\begin{GSitem}{1979-462.2}{F}{The Page--Atiyah graviton report}
\GSfield{Kind; source}{claim; \GSpage{P000841}}
\GSfield{The notes}{Page and Atiyah are associated with something involving \ru{изв.} or \ru{изл.} gravitons; the abbreviation is unclear.}
\GSfield{Statement}{The report does not identify an equation, geometric class or assertion. In particular, the adjective cannot securely be expanded to a unique type of gravitational instanton.}
\GSfield{Reconstruction}{The original PNG and both transcriptions retain the ambiguity. No theorem has been manufactured from the two names.}
\GSfield{Fate}{Targeted searches for the Cambridge meeting and Page--Atiyah gravitational instantons did not verify an exact result matching this fragment.}
\GSfield{Present status}{As of November 2026, the specific oral result is unidentifiable from the retained text.}
\GSfield{Assessment}{literal statement: not determinable; publication match: none found; verification: unresolved lead; development: annotated; confidence: low}
\end{GSitem}
\begin{GSitem}{1979-462.3}{F}{Coleman and Jackiw on topology and instanton methods}
\GSfield{Kind; source}{claim; \GSpage{P000841}}
\GSfield{The notes}{Coleman and Jackiw are said to discuss gravity with topological numbers and Yang--Mills; a survey of instanton methods is mentioned.}
\GSfield{Statement}{No particular topological invariant, gauge group, action or conclusion is recorded. The passage supplies conference provenance for instanton methods but not a uniquely reconstructable mathematical proposition.}
\GSfield{Reconstruction}{Instanton-number quantization and tunnelling are different claims; neither is silently substituted for the fragment.}
\GSfield{Fate}{The topic is consistent with the period's instanton literature, but the cited survey and exact oral result were not identified. The appropriate classification is indeterminate, not an open problem.}
\GSfield{Present status}{As of November 2026, exact publication identification remains unresolved in this search.}
\GSfield{Assessment}{literal statement: not determinable; publication match: none found; verification: unresolved lead; development: annotated; confidence: low}
\end{GSitem}
\begin{GSitem}{1979-462.4}{F}{The four-manifold and bundle-classification report}
\GSfield{Kind; source}{claim; \GSpage{P000841}}
\GSfield{The notes}{Milnor: classification of four-dimensional manifolds; somebody else: bundles on four-manifolds.}
\GSfield{Statement}{The category (smooth, topological or homotopy), compactness, simple connectivity, orientation and type of bundle are all absent. No definite classification theorem follows from this topic line.}
\GSfield{Reconstruction}{Later theorems with restricted hypotheses cannot be counted as resolving an unspecified request to classify all four-manifolds. The unnamed bundle contribution is retained in the same conference-topic item.}
\GSfield{Fate}{No exact theorem or publication was verified for this report. It remains provenance, not a claimed precursor to a particular later classification theorem.}
\GSfield{Present status}{As of November 2026, the recorded mathematical scope is insufficient for a unique fate.}
\GSfield{Assessment}{literal statement: not determinable; publication match: none found; verification: unresolved lead; development: annotated; confidence: low}
\end{GSitem}
\begin{GSitem}{1979-462.5}{E}{Signature and characteristic numbers in dimension four}
\GSfield{Kind; source}{claim; \GSpage{P000841}}
\GSfield{The notes}{Singer is reported to have spoken on the signature of a four-manifold; Euler and Pontryagin numbers then occur in the index formula.}
\GSfield{Statement}{For a closed oriented smooth four-manifold \(M\), the signature of its middle-dimensional intersection pairing is \(\sigma(M)=\frac13\langle p_1(TM),[M]\rangle\). This is the standard signature theorem relevant to the report, not a reconstruction of any unstated new claim by Singer.}
\GSfield{Reconstruction}{Closedness and orientation are necessary. Manifolds with boundary require additional boundary data and correction terms; the notebook supplies none.}
\GSfield{Fate}{The signature theorem is preexisting background, treated in \citet[Ch.~3]{gilkey1984invariance}. The exact Cambridge lecture may have contained more; its novelty is not recoverable here.}
\GSfield{Present status}{As of November 2026, the stated, qualified result is established; the notebook wording is interpreted only with the hypotheses specified above.}
\GSfield{Assessment}{literal statement: not determinable; publication match: partial; verification: read the result; development: annotated; confidence: high}
\end{GSitem}
\begin{GSitem}{1979-462.6}{F}{The recorded higher-spin zero-mode index formula}
\GSfield{Kind; source}{claim; \GSpage{P000841-P000842}}
\GSfield{The notes}{With \(S=A+B,T=A-B\), formula~(2) reads \(n[S;T]-n[s;T-1]=4^{s-1}[(T-\frac12)\chi+(s-\frac13)p/2]\), where \(n\) counts zero modes.}
\GSfield{Statement}{The meaningful index setting is an elliptic operator or elliptic complex on a closed oriented Riemannian four-manifold, with the required spin structure and specified higher-spin bundles. Its analytic index is an integral of characteristic classes. The exact displayed coefficient cannot be certified without fixing the operators and whether brackets denote unsymmetrized tensor powers.}
\GSfield{Reconstruction}{P000853, formula~(39), instead has \(s-\frac23\). Both scans were inspected; the two expressions must not be silently identified. Bare zero-mode differences for unspecified Laplacians need not be topological indices.}
\GSfield{Fate}{The exact nearby publication is \citet{christensenduff1979indices}, July 1979, whose publisher abstract explicitly covers arbitrary-spin indices and heat coefficients. The full paper was inaccessible, so numerical agreement is unverified. Duff's later primary exposition confirms the programme \citep[Sec.~4]{duff1982ultraviolet}.}
\GSfield{Present status}{As of November 2026, the underlying higher-spin index theory is published, but this notebook formula's normalization and discrepancy remain unresolved.}
\GSfield{Assessment}{literal statement: not determinable; publication match: partial; verification: indirect; development: annotated; confidence: medium}
\end{GSitem}
\begin{GSitem}{1979-462.7}{F}{Boson--fermion index identities on self-dual backgrounds}
\GSfield{Kind; source}{claim; \GSpage{P000842}}
\GSfield{The notes}{Formula~(3) relates \(n[A,B]\) and \(2n[A-\frac12,B]\) to \(4^{A+B-1}(A-\frac7{12})(\chi+p/2)\); the minus sign is overwritten. Self-duality is explicitly required.}
\GSfield{Statement}{On appropriately chosen self-dual Riemannian spin backgrounds, differential intertwiners can pair nonzero bosonic and fermionic spectra; their unmatched kernels produce an index relation. The concrete bundles, gauge constraints and topological normalizations must be supplied before a numerical formula is stated.}
\GSfield{Reconstruction}{The overwritten sign and absent operator definitions prevent certifying the exact identity. The self-dual hypothesis is present in the notebook and is retained, rather than generalized to arbitrary metrics.}
\GSfield{Fate}{\citet{christensenduff1979indices} explicitly announces such supertheorems in its primary abstract. \citet[Secs.~2.3 and~4]{duff1982ultraviolet} supplies independently checked context; this is a partial match, not verification of the written coefficient.}
\GSfield{Present status}{As of November 2026, a published family of results is identified, but this exact formula remains underdetermined.}
\GSfield{Assessment}{literal statement: not determinable; publication match: partial; verification: indirect; development: annotated; confidence: medium}
\end{GSitem}
\begin{GSitem}{1979-462.8}{E}{Nonzero-mode cancellation and the heat coefficient}
\GSfield{Kind; source}{claim; \GSpage{P000842}}
\GSfield{The notes}{\(m[S;T]=m[S]=4^s m[0]\), \(B_4=n+m\); nonzero modes cancel in supergravity, leaving zero modes.}
\GSfield{Statement}{For a nonnegative elliptic Laplace-type operator \(P\) on a closed four-manifold, \(\zeta_P(s)=\sum_{\lambda>0}\lambda^{-s}\) has a continuation regular at zero, and the constant heat coefficient satisfies \(a_4(P)=\dim\ker P+\zeta_P(0)\). A verified spectral pairing can cancel the corresponding nonzero contributions in an appropriately weighted supertrace.}
\GSfield{Reconstruction}{The term \(m\) must be a regularized spectral contribution, not a finite number of positive eigenvalues or a common multiplicity. P000852 explicitly rejects the latter interpretation. The power \(4^s\) needs the tensor-space convention before it can be checked.}
\GSfield{Fate}{The general heat/zeta identity is established \citep[Secs.~1.7 and~1.10]{gilkey1984invariance}; the self-dual supergravity application is discussed in \citet{duff1982ultraviolet}. The numerical spin factor remains only a partial publication match.}
\GSfield{Present status}{As of November 2026, the stated, qualified result is established; the notebook wording is interpreted only with the hypotheses specified above.}
\GSfield{Assessment}{literal statement: correct with missing hypotheses; publication match: partial; verification: read the result; development: annotated; confidence: high}
\end{GSitem}
\begin{GSitem}{1979-462.9}{A}{Finiteness in a self-dual supergravity sector}
\GSfield{Kind; source}{claim; \GSpage{P000842}}
\GSfield{The notes}{They proved that the functional integral has no divergences; the speaker obtained related results.}
\GSfield{Statement}{The identified perturbative claim concerns local ultraviolet counterterms evaluated on backgrounds obeying specified self-duality conditions. It is not a construction of a nonperturbative gravitational measure or a claim that every supergravity amplitude on every background is finite.}
\GSfield{Reconstruction}{The preceding sentence restricts the calculation to a self-dual sector. Gauge fixing, zero modes, regularization, loop order and the precise supergravity theory are absent from the shorthand.}
\GSfield{Fate}{\citet{christensenduff1979indices} states a beyond-one-loop self-dual result; \citet{kallosh1980selfduality} concerns superspace self-duality. Duff's checked 1982 lectures, Sec.~2.3, explain both the self-dual cancellation and allowed counterterms outside that restriction \citep{duff1982ultraviolet}.}
\GSfield{Present status}{As of November 2026, the qualified sector-specific publication is identified; the unrestricted reading is unsupported.}
\GSfield{Assessment}{literal statement: correct with missing hypotheses; publication match: partial; verification: indirect; development: annotated; confidence: high}
\end{GSitem}
\begin{GSitem}{1979-462.10}{F}{An unspecified super-index theorem}
\GSfield{Kind; source}{problem; \GSpage{P000842}}
\GSfield{The notes}{Interesting problem: the same on a supermanifold; Singer is working on it; make a super index theorem.}
\GSfield{Statement}{No super-elliptic operator, symbol class, function space, compactness condition or index target is given. Since odd-coordinate differential calculi differ from ordinary elliptic theory, these choices determine the theorem.}
\GSfield{Reconstruction}{This is a research programme, not a single well-posed conjecture. It is kept distinct from the concrete self-dual boson--fermion identities.}
\GSfield{Fate}{The later superanalysis framework \citep{berezin1987superanalysis} supplies language, but no exact result can be matched to the unspecified request.}
\GSfield{Present status}{As of November 2026, the notebook does not determine a unique problem whose resolution can be recorded.}
\GSfield{Assessment}{literal statement: not determinable; publication match: thematic only; verification: publisher metadata checked; development: annotated; confidence: high}
\end{GSitem}
\begin{GSitem}{1979-462.11}{E}{Exceptional symmetry of maximal supergravity}
\GSfield{Kind; source}{claim; \GSpage{P000842}}
\GSfield{The notes}{The French group replaces an unsatisfactory \(SO(8)\) by \(E_7\longrightarrow SU(8)\); \(\dim E_7=133\), with an abbreviation possibly meaning compact.}
\GSfield{Statement}{In ungauged four-dimensional maximal supergravity the relevant continuous exceptional symmetry is the split real group \(E_{7(7)}\), with maximal compact subgroup locally \(SU(8)\); the scalar manifold has local form \(E_{7(7)}/SU(8)\), of dimension \(133-63=70\).}
\GSfield{Reconstruction}{The arrow is not a specified homomorphism or symmetry-breaking theorem. The compact real form of \(E_7\) would be the wrong real form for this scalar coset. Gauged \(SO(8)\) supergravity is a related but distinct construction.}
\GSfield{Fate}{The preexisting Cremmer--Julia work of 1978 is identified in \citet[supergravity field-content and symmetry discussion]{duff1982ultraviolet}. The independently opened historical source verifies the context, but no exact meaning is assigned to the notebook arrow.}
\GSfield{Present status}{As of November 2026, the stated, qualified result is established; the notebook wording is interpreted only with the hypotheses specified above.}
\GSfield{Assessment}{literal statement: correct with missing hypotheses; publication match: partial; verification: indirect; development: annotated; confidence: high}
\end{GSitem}
\begin{GSitem}{1979-462.12}{E}{The exceptional Jordan algebra and compact F4}
\GSfield{Kind; source}{theorem; \GSpage{P000842}}
\GSfield{The notes}{Take Hermitian \(3\times3\) octonionic matrices with product \((AB+BA)/2\); their automorphism group is \(F_4\).}
\GSfield{Statement}{For the real division octonions \(\mathbb O\), the Albert algebra \(J=\operatorname{Herm}_3(\mathbb O)\), with \(A\circ B=(AB+BA)/2\), is a 27-dimensional real Jordan algebra. Its real automorphism group is the compact real form of \(F_4\).}
\GSfield{Reconstruction}{Only products of two octonions occur in the displayed matrix product, so it is well defined despite nonassociativity. Split octonions or complexification change the real form.}
\GSfield{Fate}{This is a classical result, not new in 1979. \citet[Sec.~4.2, Theorem~5]{baez2002octonions} gives the precise statement and credits the 1950 Chevalley--Schafer construction.}
\GSfield{Present status}{As of November 2026, the stated, qualified result is established; the notebook wording is interpreted only with the hypotheses specified above.}
\GSfield{Assessment}{literal statement: correct with missing hypotheses; publication match: exact; verification: read the result; development: annotated; confidence: high}
\end{GSitem}

\GStalk{463}{1979-09-24}{V. E. Zakharov}{P000843-P000847}{Degenerate dispersion laws, resonances and integrability}
The talk begins on P000843, after Kallosh's report. Zakharov studies additional additive functions on wave-resonance manifolds, difference-curve parametrizations, dressing and a necessary resonance test for extra Hamiltonian integrals. Arnold connects this with translation surfaces, Poincaré obstructions and formal Birkhoff normal form. The subsequent programme appears in Zakharov--Schulman (1980, 1988) and the later Balk--Ferapontov web classification (1994); one additional collision invariant and infinitely many are distinct hypotheses. None of these statements identifies every notion of nonlinear integrability with degenerate scalar dispersion. P000847 ends with Gelfand's reading programme.

\begin{GSitem}{1979-463.1}{A}{Additional additive functions on a resonance manifold}
\GSfield{Kind; source}{definition; \GSpage{P000843}}
\GSfield{The notes}{\(\Gamma=\{(k_1,k_2):\omega(k_1+k_2)=\omega(k_1)+\omega(k_2)\}\); seek another \(f\) additive there and not differing from \(\omega\) by a linear function.}
\GSfield{Statement}{For a smooth real dispersion \(\omega\) on an open momentum domain, restrict to a nonempty regular component of the three-wave resonance equation. An additional collision invariant is a smooth \(f\) satisfying the same additivity relation and lying outside \(\{a\omega+b\cdot k:a\in\mathbb R,b\in\mathbb R^N\}\).}
\GSfield{Reconstruction}{Nonlinearity alone does not guarantee a smooth hypersurface or even a nonempty resonance set. The quotient by all scalar multiples of \(\omega\) is essential: otherwise \(2\omega\) is spuriously additional.}
\GSfield{Fate}{The contemporary definition and additive equation are in \citet[Sec.~1, pp.~192--193, eqs.~(1), (3)--(5)]{zakharovschulman1980degenerative}. The paper was received on 1 September 1979 and published in June 1980; receipt is not asserted to be a public circulation date. The later systematic formulation is \citet[Sec.~3, eqs.~(3.8)--(3.9)]{zakharovschulman1988invariants}.}
\GSfield{Present status}{As of November 2026, the stated, qualified result is established; the notebook wording is interpreted only with the hypotheses specified above.}
\GSfield{Assessment}{literal statement: false as written; publication match: exact; verification: read the result; development: annotated; confidence: high}
\end{GSitem}
\begin{GSitem}{1979-463.2}{A}{Difference curves yield families of collision invariants}
\GSfield{Kind; source}{construction; \GSpage{P000843-P000844}}
\GSfield{The notes}{\(p=\xi_1-\xi_2,\ q=a(\xi_1)-a(\xi_2),\ \omega=b(\xi_1)-b(\xi_2)\); replace \(b\) by arbitrary \(c\) to obtain \(f\). Three parameters telescope along \((\xi_1,\xi_2,\xi_3)\).}
\GSfield{Statement}{If \(a,b,c\) are smooth and \(a'(\xi_1)\ne a'(\xi_2)\), the first two expressions define local momentum coordinates. The associated functions \(\omega,f\) obey the three-wave identity on composable pairs \((\xi_1,\xi_2),(\xi_2,\xi_3)\), wherever the chosen local branches are compatible.}
\GSfield{Reconstruction}{The determinant of the coordinate map is \(a'(\xi_1)-a'(\xi_2)\). Telescoping proves the identity on the constructed component; it does not prove additivity on every disconnected resonance component or global single-valuedness.}
\GSfield{Fate}{The exact difference-curve parametrization appears in \citet[Sec.~1, p.~193, eqs.~(12)--(13)]{zakharovschulman1980degenerative}, published in June 1980. That paper credits the general example to S.~V.~Manakov, private communication. The arbitrary-function construction is written explicitly in \citet[Sec.~3, eqs.~(3.10)--(3.12)]{zakharovschulman1988invariants}, followed by the distinction between partial and complete degeneracy. The 1980 formula pages were read and visually checked in the original.}
\GSfield{Present status}{As of November 2026, the stated, qualified result is established; the notebook wording is interpreted only with the hypotheses specified above.}
\GSfield{Assessment}{literal statement: correct with missing hypotheses; publication match: exact; verification: read the result; development: annotated; confidence: high}
\end{GSitem}
\begin{GSitem}{1979-463.3}{B}{Is the difference-curve construction universal?}
\GSfield{Kind; source}{problem; \GSpage{P000843}}
\GSfield{The notes}{The page first asks for one additional additive function, then frames the construction as obtaining equations with infinitely many first integrals, and finally asks whether it is general.}
\GSfield{Statement}{In the infinitely-many-invariants reading, ask whether a regular local two-dimensional three-wave dispersion law admitting infinitely many independent smooth collision invariants can be represented by \((p,q,\omega)=(\xi_1-\xi_2,a(\xi_1)-a(\xi_2),b(\xi_1)-b(\xi_2))\), with an invertible momentum parametrization and compatible resonance branches.}
\GSfield{Reconstruction}{The initial one-extra-invariant question and the subsequent infinitely-many-invariants programme are different quantifiers. The former cannot silently replace the latter. Additivity on a selected local component must also be distinguished from additivity on the entire resonance set and from nonlinear conserved quantities of the full wave equation.}
\GSfield{Fate}{The speaker's construction is published in \citet[Sec.~3]{zakharovschulman1988invariants}. Later \citet[abstract]{balkferapontov1994invariants} explicitly classifies three-wave interactions with infinitely many independent invariants using web geometry; the original theorem and its hypotheses were not retrieved, so this remains a partial publication match. The earlier none-or-infinitely-many expectation is contradicted by the finite-invariant examples announced in \citet[Secs.~4--6]{balknazarenkozakharov1991drift}. Their uniqueness proof is deferred there; \citet[Sec.~2.1]{balk2014rossby} confirms the Rossby uniqueness result and cites \citet{balkferapontov1998web}.}
\GSfield{Present status}{As of November 2026, the infinitely-many-invariants classification has an identified later publication. An exact hypothesis-by-hypothesis comparison with the handwritten parametrization remains unverified because the 1994 theorem was inaccessible. The broader one-invariant reading is distinct; finite-extra-invariant systems occur in the published literature. No current open problem is inferred from the access limit.}
\GSfield{Assessment}{literal statement: not determinable; publication match: partial; verification: indirect; development: annotated; confidence: medium; problem \hyperref[prob:GP-1979-05]{GP-1979-05}}
\end{GSitem}
\begin{GSitem}{1979-463.4}{E}{Zero curvature and the dispersion of a dressed scalar pair}
\GSfield{Kind; source}{construction; \GSpage{P000844}}
\GSfield{The notes}{\(L_t-A_y+[L,A]=0\), with scalar differential operators in \(x\); linearization at \(L_0,A_0\) has constant coefficients and difference-symbol dispersion.}
\GSfield{Statement}{Take constant-coefficient scalar operators \(L_0=\ell(\partial_x)\), \(A_0=a(\partial_x)\), and linearize a compatible pair \(\partial_y-L,\partial_t-A\). Plane-wave or exponential perturbations give algebraic relations expressible using differences of the two symbols in the standard dressing construction.}
\GSfield{Reconstruction}{The zero-curvature sign depends on the chosen auxiliary operators. Linearization has constant coefficients only about a constant background; variable background coefficients do not disappear merely by discarding nonlinear terms.}
\GSfield{Fate}{The preexisting dressing framework is \citet{zakharovshabat1974scheme}. Its exact bibliographic metadata was checked; the notebook's particular symbols and linearization were verified algebraically rather than attributed to a separately read theorem.}
\GSfield{Present status}{As of November 2026, the stated, qualified result is established; the notebook wording is interpreted only with the hypotheses specified above.}
\GSfield{Assessment}{literal statement: correct with missing hypotheses; publication match: exact; verification: publisher metadata checked; development: annotated; confidence: high}
\end{GSitem}
\begin{GSitem}{1979-463.5}{E}{The linearized Gel'fand--Levitan--Marchenko equation}
\GSfield{Kind; source}{construction; \GSpage{P000844}}
\GSfield{The notes}{\(F+K+K\oplus F=0\), followed by \(K=-F\) and equations for the \(t,y\) dependence of \(F\).}
\GSfield{Statement}{Interpret \(K\oplus F\) as the integral-operator composition appropriate to a chosen triangular dressing factorization. For a small kernel \(F\), \(K=-F+O(F^2)\). The equality \(K=-F\) is exact only at first order or when the composition term vanishes.}
\GSfield{Reconstruction}{The displayed differential equations attach \(t\) to \(L_0\) and \(y\) to \(A_0\), while the exponential attaches \(y\) to \(\ell\) and \(t\) to \(a\); their signs also depend on action on the second kernel variable. The transcript is retained and the inconsistent assignments are not silently reconciled.}
\GSfield{Fate}{The GLM/dressing method predates this lecture \citep{zakharovshabat1974scheme}; the first-order identity follows directly by retaining degree one in the kernel equation.}
\GSfield{Present status}{As of November 2026, the stated, qualified result is established; the notebook wording is interpreted only with the hypotheses specified above.}
\GSfield{Assessment}{literal statement: correct with missing hypotheses; publication match: partial; verification: publisher metadata checked; development: annotated; confidence: high}
\end{GSitem}
\begin{GSitem}{1979-463.6}{A}{The resonance obstruction for an additional quadratic-leading integral}
\GSfield{Kind; source}{theorem; \GSpage{P000845-P000846}}
\GSfield{The notes}{If the equation is integrable and its dispersion nondegenerate, the interaction vanishes on the resonance manifold; the notes also question integrals beginning at higher orders.}
\GSfield{Statement}{For an analytic translation-invariant Hamiltonian wave system admitting an additional analytic integral whose leading term is \(\int f(k)|a_k|^2dk\), each on-shell scattering amplitude \(W\) satisfies \((\sum_{\rm in}f-\sum_{\rm out}f)W=0\) on its regular resonance set. Thus either the additional additive identity holds there or that amplitude vanishes.}
\GSfield{Reconstruction}{The hypotheses concern a specified analytic expansion with a quadratic leading integral and regular coefficients. They do not identify all notions of integrability, and do not automatically cover integrals whose first nonzero term has higher degree.}
\GSfield{Fate}{The kinetic form of the resonance obstruction was published soon after the seminar in \citet[Sec.~1, p.~195, final paragraph]{zakharovschulman1980degenerative}, excluding singular delta-type kinetic kernels. The general statement for analytic Hamiltonian invariants with a quadratic leading term is proved in \citet[Theorem~3.1, pp.~291--293]{zakharovschulman1988invariants}; its introduction, pp.~285--286, also identifies a brief 1985 announcement. Class A records the early kinetic result; the full modern reconstruction is supported by the later theorem.}
\GSfield{Present status}{As of November 2026, the stated, qualified result is established; the notebook wording is interpreted only with the hypotheses specified above.}
\GSfield{Assessment}{literal statement: correct with missing hypotheses; publication match: partial; verification: read the result; development: annotated; confidence: high}
\end{GSitem}
\begin{GSitem}{1979-463.7}{E}{The KP dispersion from quadratic and cubic symbols}
\GSfield{Kind; source}{construction; \GSpage{P000844-P000845}}
\GSfield{The notes}{\(\ell(\xi)=\xi^2,\ a(\xi)=\xi^3\) is first called KdV and then KP; \(\omega=\frac14(p^3+3q^2/p)\).}
\GSfield{Statement}{With \(p=\xi_1-\xi_2\), \(q=\xi_1^2-\xi_2^2\) and \(p\ne0\), one has \(\xi_1=(q/p+p)/2,\ \xi_2=(q/p-p)/2\), hence \(\xi_1^3-\xi_2^3=(p^3+3q^2/p)/4\). This is the two-spatial-variable KP symbol in the lecture's exponential convention; a \(y\)-independent reduction yields KdV after scaling.}
\GSfield{Reconstruction}{The rational expression is undefined at \(p=0\). The written KP evolution lacks its final equality to zero, and switching from real exponentials to oscillatory Fourier modes changes displayed signs.}
\GSfield{Fate}{Scalar dressing by these symbols is part of the established Zakharov--Shabat framework \citep{zakharovshabat1974scheme}; the exact algebraic identity is verified directly. The resonance geometry is discussed in \citet[Sec.~3]{zakharovschulman1988invariants}.}
\GSfield{Present status}{As of November 2026, the stated, qualified result is established; the notebook wording is interpreted only with the hypotheses specified above.}
\GSfield{Assessment}{literal statement: correct with missing hypotheses; publication match: exact; verification: read the result; development: annotated; confidence: high}
\end{GSitem}
\begin{GSitem}{1979-463.8}{E}{A matrix dressing pair and the Davey--Stewartson system}
\GSfield{Kind; source}{construction; \GSpage{P000845}}
\GSfield{The notes}{Dressing \(\operatorname{diag}(1,-1)\partial_x\) and \(I\partial_x^2\) yields a two-component operator with off-diagonal \(\psi,\bar\psi\), then the displayed hyperbolic Schrödinger/elliptic potential system.}
\GSfield{Statement}{The recorded system is \(i\psi_t-\psi_{xx}+\psi_{yy}+u\psi=0,\quad \Delta u=(\partial_x^2-\partial_y^2)|\psi|^2\), with real \(u\). In a decaying setting the second equation determines \(u\) up to an imposed harmonic component. It belongs, after coefficient and sign conventions are fixed, to the Davey--Stewartson family generated by matrix dressing.}
\GSfield{Reconstruction}{The exact form of the dressed \(A\), reduction signs, boundary conditions and nonlinear coefficient normalization are omitted. Consequently the displayed pair is not by itself a verified complete Lax representation.}
\GSfield{Fate}{The family is explicitly discussed in the primary later account \citep[Secs.~1 and~4]{zakharovschulman1988invariants}; the earlier dressing source is \citet{zakharovshabat1974scheme}. This is a partial match to the particular shorthand system.}
\GSfield{Present status}{As of November 2026, the stated, qualified result is established; the notebook wording is interpreted only with the hypotheses specified above.}
\GSfield{Assessment}{literal statement: correct with missing hypotheses; publication match: partial; verification: read the result; development: annotated; confidence: high}
\end{GSitem}
\begin{GSitem}{1979-463.9}{E}{Fourier Hamiltonians and the complex variational derivative}
\GSfield{Kind; source}{construction; \GSpage{P000845}}
\GSfield{The notes}{The Fourier equation is written \(ia_t=\delta H/\delta a\); the quartic interaction integral displays only two integration differentials.}
\GSfield{Statement}{For complex canonical wave variables with the usual Poisson convention, \(i\partial_t a(k)=\delta H/\delta\overline{a(k)}\). A quartic translation-invariant Hamiltonian integrates \(\bar a_1\bar a_2a_3a_4\) against a kernel and \(\delta(k_1+k_2-k_3-k_4)\), with all four momenta integrated, or three after eliminating the delta constraint.}
\GSfield{Reconstruction}{The missing conjugation is a mathematical error in the displayed canonical equation. Omitting integration variables is shorthand, but must not be interpreted as a two-variable integral with free momenta.}
\GSfield{Fate}{The precise convention is explicit in \citet[Sec.~1, equation~(1.1)]{zakharovschulman1988invariants}; it also follows from varying the quadratic term \(H_0=\int\omega|a|^2\).}
\GSfield{Present status}{As of November 2026, the stated, qualified result is established; the notebook wording is interpreted only with the hypotheses specified above.}
\GSfield{Assessment}{literal statement: false as written; publication match: exact; verification: read the result; development: annotated; confidence: high}
\end{GSitem}
\begin{GSitem}{1979-463.10}{F}{The displayed four-wave coefficient does not vanish on resonance}
\GSfield{Kind; source}{claim; \GSpage{P000845-P000846}}
\GSfield{The notes}{The two rational fractions printed for \(w\) are followed by \(w|_\Gamma=0\), with energy and momentum conservation defining \(\Gamma\).}
\GSfield{Statement}{\oursComputed\ For the printed expression and \(\omega(p,q)=-p^2+q^2\), take \(k_1=(1,0),k_2=(0,1),k_3=(0,0),k_4=(1,1)\). Both conservation laws hold and all denominators are nonzero, but the two fractions give \(1+1/5=6/5\), so the written vanishing claim is false.}
\GSfield{Reconstruction}{The original PNG confirms plus signs inside the second denominator. Replacing those signs by minus signs alone still fails: it gives \(2\) at this quartet. A correct symmetrized scattering amplitude requires additional index or term changes, which were not reconstructed from the notes.}
\GSfield{Fate}{The later scattering-amplitude theorem \citep{zakharovschulman1988invariants} provides the appropriate context but does not certify this erroneous displayed kernel. The explicit counterexample is checked in this chapter.}
\GSfield{Present status}{As of November 2026, the printed formula is disproved; an exact source-level restoration of the intended coefficient remains unverified.}
\GSfield{Assessment}{literal statement: false as written; publication match: thematic only; verification: read the result; development: annotated; confidence: high}
\end{GSitem}
\begin{GSitem}{1979-463.11}{E}{Poincaré's first-order homological obstruction}
\GSfield{Kind; source}{claim; \GSpage{P000846}; V. I. Arnold}
\GSfield{The notes}{For \(H=H_0(p)+\varepsilon\sum h_k(p)e^{ikq}\), a nondegenerate frequency map is assumed and the notes state \(h_k=0\) on \(k\cdot\lambda=0\); infinitely many violations supposedly exclude an independent integral.}
\GSfield{Statement}{If \(F=F_0(p)+\varepsilon F_1(p,q)+\cdots\) is an analytic or formal integral, the first-order homological equation gives \((k\cdot\nabla H_0)F_{1,k}=(k\cdot\nabla F_0)h_k\), up to a common sign convention. Therefore \((k\cdot\nabla F_0)h_k=0\) on each resonance.}
\GSfield{Reconstruction}{The factor involving \(F_0\) cannot be dropped. A nonintegrability theorem additionally needs a sufficiently rich Poincaré set of independent resonances and a nondegenerate Hessian; merely counting infinitely many nonzero coefficients is insufficient.}
\GSfield{Fate}{The classical perturbative context is the one explicitly compared with the wave theorem by \citet[Introduction]{zakharovschulman1988invariants}. Arnold's textbook metadata and appendix context were checked \citep{arnold1989mechanics}; the displayed necessary condition is derived here, not quoted from an unread proof.}
\GSfield{Present status}{As of November 2026, the stated, qualified result is established; the notebook wording is interpreted only with the hypotheses specified above.}
\GSfield{Assessment}{literal statement: false as written; publication match: partial; verification: read the result; development: annotated; confidence: high}
\end{GSitem}
\begin{GSitem}{1979-463.12}{E}{Translation-surface interpretation of the resonance parametrization}
\GSfield{Kind; source}{construction; \GSpage{P000846}; V. I. Arnold and V. E. Zakharov}
\GSfield{The notes}{Arnold interprets the difference graph as translating one curve along another; Zakharov writes a three-branch resonance law \(\omega_1(k_1+k_2)=\omega_2(k_1)+\omega_3(k_2)\).}
\GSfield{Statement}{The graph parametrized by \((p,q,\omega)=(\xi_1,a(\xi_1),b(\xi_1))-(\xi_2,a(\xi_2),b(\xi_2))\) is the local difference of two curve arcs, hence a translation surface where its differential has rank two. Distinct branches lead to analogous sums of possibly different curves.}
\GSfield{Reconstruction}{This is a geometric interpretation and local construction. It does not classify every three-branch functional equation without transversality and compatibility hypotheses.}
\GSfield{Fate}{The single-branch construction is verified in \citet[Sec.~3]{zakharovschulman1988invariants}; the general translation-geometry language is classical \citep[Sec.~2]{little1983translation}.}
\GSfield{Present status}{As of November 2026, the stated, qualified result is established; the notebook wording is interpreted only with the hypotheses specified above.}
\GSfield{Assessment}{literal statement: correct with missing hypotheses; publication match: partial; verification: read the result; development: annotated; confidence: high}
\end{GSitem}
\begin{GSitem}{1979-463.13}{F}{Integrals with higher-order leading terms}
\GSfield{Kind; source}{problem; \GSpage{P000846}}
\GSfield{The notes}{What if the integrals start directly at higher orders?}
\GSfield{Statement}{The question asks whether a resonance test can detect conserved quantities with no quadratic leading term, but gives no order, regularity class, Hamiltonian class or desired criterion. The quadratic-leading theorem does not answer it.}
\GSfield{Reconstruction}{Once those choices are fixed, expanding \(\{H,F\}=0\) gives degree-specific homological equations; different choices produce different problems.}
\GSfield{Fate}{\citet[Sec.~1, equation~(1.3)]{zakharovschulman1988invariants} explicitly starts with a quadratic leading invariant. Its theorem is therefore not a solution of this broader, unspecified request.}
\GSfield{Present status}{As of November 2026, this fragment is a legitimate question about the theorem's scope but is insufficiently specified for a unique global publication fate.}
\GSfield{Assessment}{literal statement: not determinable; publication match: partial; verification: read the result; development: annotated; confidence: medium}
\end{GSitem}
\begin{GSitem}{1979-463.14}{E}{Formal Birkhoff normal form at a nonresonant elliptic equilibrium}
\GSfield{Kind; source}{theorem; \GSpage{P000847}; V. I. Arnold}
\GSfield{The notes}{Arnold: \(H=\sum\omega_k z_k\bar z_k+\cdots\) reduces to \(H(\Pi_k)\), \(\Pi_k=z_k\bar z_k\), if the frequencies are independent; explicitly at the level of formal series.}
\GSfield{Statement}{For a formal Hamiltonian near an elliptic equilibrium with rationally independent frequencies, a formal symplectic change of variables puts it into a power series in the actions \(I_k=z_k\bar z_k\). Rational independence means \(m\cdot\omega\ne0\) for every nonzero \(m\in\mathbb Z^n\).}
\GSfield{Reconstruction}{At each degree the homological operator is diagonal on monomials; nonresonant terms can be divided by their nonzero frequency combinations. The remaining monomials have equal powers of \(z_k,\bar z_k\). No convergence or analytic integrability follows.}
\GSfield{Fate}{This is the classical Birkhoff theorem summarized in Arnold's mechanics text, Appendix~7 \citep{arnold1989mechanics}. Publisher metadata was checked; the finite-degree elimination explanation above independently verifies the formal assertion.}
\GSfield{Present status}{As of November 2026, the stated, qualified result is established; the notebook wording is interpreted only with the hypotheses specified above.}
\GSfield{Assessment}{literal statement: correct; publication match: exact; verification: publisher metadata checked; development: annotated; confidence: high}
\end{GSitem}

\GStalk{464}{1979-10-01}{I. M. Gelfand}{P000848}{Spinors, quaternions and the Dirac equation}
P000848 contains Gelfand's introductory discussion of spinors, the projective action of two-component vectors, Hamilton's square-root idea and physical consequences of the Dirac theory. The quaternion conclusion requires an extra quotient relation: three anticommuting square roots of minus one alone generate a Clifford algebra. Kallosh's subsequent programme begins at the foot of the same page.

\begin{GSitem}{1979-464.1}{E}{Electron--positron annihilation and its inverse}
\GSfield{Kind; source}{claim; \GSpage{P000848}}
\GSfield{The notes}{Dirac theory predicts positrons; an electron and positron can annihilate into two photons, and two photons can produce an electron--positron pair.}
\GSfield{Statement}{These are allowed QED scattering processes when energy, momentum and angular momentum permit them. In units \(c=1\), two-photon pair production requires \(s=(k_1+k_2)^2=2E_1E_2(1-\cos\theta)\ge4m_e^2\). A free electron--positron pair cannot annihilate into a single real photon because its total four-momentum is timelike.}
\GSfield{Reconstruction}{The notebook describes possible processes, not deterministic outcomes of every collision; a Dirac equation without quantized electromagnetic interactions is not alone the full scattering theory.}
\GSfield{Fate}{This is preexisting physical exposition, not an open mathematical assertion. The kinematic restriction is reconstructed directly from four-momentum conservation; the seminar's Dirac and QED background recurs from 1977-421 and 1978-452.}
\GSfield{Present status}{As of November 2026, the stated, qualified result is established; the notebook wording is interpreted only with the hypotheses specified above.}
\GSfield{Assessment}{literal statement: correct with missing hypotheses; publication match: exact; verification: indirect; development: annotated; confidence: high}
\end{GSitem}
\begin{GSitem}{1979-464.2}{E}{Hamilton's square-root relations and the missing quaternion quotient}
\GSfield{Kind; source}{claim; \GSpage{P000848}}
\GSfield{The notes}{Three symbols square to \(-1\) and anticommute; adding the identity is said to give the quaternions.}
\GSfield{Statement}{The universal unital real algebra with \(a^2=b^2=c^2=-1\) and pairwise anticommutation is \(\operatorname{Cl}_{0,3}\), of dimension eight, isomorphic to \(\mathbb H\oplus\mathbb H\). Adding \(ab=c\), equivalently choosing \(abc=-1\), gives a quaternion quotient. Either sign of the central volume element selects one summand.}
\GSfield{Reconstruction}{The square-root identity determines the Clifford relations but does not itself impose \(ab=c\). The four-dimensional span of \(1,a,b,c\) need not be closed without the missing relation.}
\GSfield{Fate}{Classical Clifford/quaternion algebra; the relation with normed division algebras is treated in \citet[Secs.~2.2--2.3]{baez2002octonions}. The dimension and quotient distinction follow directly from the monomial basis.}
\GSfield{Present status}{As of November 2026, the stated, qualified result is established; the notebook wording is interpreted only with the hypotheses specified above.}
\GSfield{Assessment}{literal statement: false as written; publication match: exact; verification: read the result; development: annotated; confidence: high}
\end{GSitem}
\begin{GSitem}{1979-464.3}{E}{Möbius transformations as projectivized spinor transformations}
\GSfield{Kind; source}{construction; \GSpage{P000848}}
\GSfield{The notes}{A fractional linear transformation of \(z=z_1/z_2\) comes from a linear transformation of the two complex coordinates; their complex conjugates are conjugate spinors.}
\GSfield{Statement}{An invertible \(2\times2\) complex matrix acts on \(\mathbb P^1\) by \(z\mapsto(az+b)/(cz+d)\). Restricting to \(SL_2(\mathbb C)\) gives a double cover of \(PSL_2(\mathbb C)\), also the connected proper Lorentz group through its action \(X\mapsto gXg^*\) on Hermitian \(2\times2\) matrices.}
\GSfield{Reconstruction}{The denominator may vanish, giving the point at infinity. A two-component spinor belongs to the linear representation before projectivization, so its sign is retained even though it disappears on \(\mathbb P^1\).}
\GSfield{Fate}{This is standard preexisting spinor geometry; \citet[Sec.~3.1]{woit2021euclidean} gives the real-form distinction used here.}
\GSfield{Present status}{As of November 2026, the stated, qualified result is established; the notebook wording is interpreted only with the hypotheses specified above.}
\GSfield{Assessment}{literal statement: correct with missing hypotheses; publication match: exact; verification: read the result; development: annotated; confidence: high}
\end{GSitem}

\GStalk{465}{1979-10-01}{Renata Kallosh}{P000848-P000853}{Spinor calculus and gravitational index theorems}
The eight-part programme on P000848 is followed on P000849--P000853 by two-spinor calculus, self-duality, Laplace-type operators, heat coefficients and higher-spin spectral pairings. The notes themselves correct a proposed interpretation of the nonzero-mode term as a common eigenvalue multiplicity. They also contain incompatible numerical versions of an index formula on P000841 and P000853; these are recorded without silently repairing an unverified normalization. The dated heading on P000854 starts the next meeting.

\begin{GSitem}{1979-465.1}{E}{The scope of the cited self-dual cancellation results}
\GSfield{Kind; source}{claim; \GSpage{P000849}}
\GSfield{The notes}{Hawking--Pope 1978 is cited for one-loop zero-mode contributions; Kallosh and Christensen--Duff for higher-order self-dual results. Four references are listed.}
\GSfield{Statement}{The identified claim is cancellation of paired nonzero modes in appropriate supergravity one-loop combinations on self-dual gravitational-instanton backgrounds, followed by stronger restrictions on counterterms in self-dual superspace. Boundary terms and zero modes remain relevant.}
\GSfield{Reconstruction}{The comparison with Yang--Mills is contextual, not a claim that all Yang--Mills spectra or all supergravities share one determinant formula. The Russian page numbers 192 and 493 refer to two Kallosh letters, not an erratum.}
\GSfield{Fate}{The exact first reference is \citet{hawkingpope1978symmetry}. Its publisher abstract and author bibliography were checked. The later, fully read \citet[Secs.~2.3 and~4.4]{duff1982ultraviolet} explains the limitations and identifies the two Kallosh works.}
\GSfield{Present status}{As of November 2026, the stated, qualified result is established; the notebook wording is interpreted only with the hypotheses specified above.}
\GSfield{Assessment}{literal statement: correct with missing hypotheses; publication match: partial; verification: indirect; development: annotated; confidence: high}
\end{GSitem}
\begin{GSitem}{1979-465.2}{E}{Two-spinor notation and the Lorentz/Euclidean distinction}
\GSfield{Kind; source}{construction; \GSpage{P000849-P000850}}
\GSfield{The notes}{\(SL_2(\mathbb C)\) acts on \(V\), conjugately on \(\bar V\); \(V\otimes\bar V\) has a real form and an invariant quadratic form. The text alternates Lorentz and \(SO(4)\).}
\GSfield{Statement}{For Lorentz signature, Hermitian matrices in \(V\otimes\bar V\) identify with Minkowski vectors and determinant is the Lorentz quadratic form. For Euclidean signature, \(\operatorname{Spin}(4)=SU(2)\times SU(2)\) has two independent complex two-dimensional spin modules \(S_+,S_-\); they are not related by the same Lorentz conjugation rule.}
\GSfield{Reconstruction}{The dotted and undotted notation survives complexification, but the real structures differ. Spinors transform under the spin double cover, not under \(SO(4)\) itself without qualification.}
\GSfield{Fate}{Classical background checked in \citet[Sec.~3.1]{woit2021euclidean} and \citet[Sec.~3.2]{gilkey1984invariance}. The change of real form is the missing hypothesis behind several following signs.}
\GSfield{Present status}{As of November 2026, the stated, qualified result is established; the notebook wording is interpreted only with the hypotheses specified above.}
\GSfield{Assessment}{literal statement: correct with missing hypotheses; publication match: exact; verification: read the result; development: annotated; confidence: high}
\end{GSitem}
\begin{GSitem}{1979-465.3}{E}{Pauli multiplication and spin-generator conventions}
\GSfield{Kind; source}{claim; \GSpage{P000849-P000850}}
\GSfield{The notes}{Equation~(3) writes \(\sigma_i\sigma_j=\delta_{ij}+\epsilon_{ijk}\sigma^k\), and (4) writes \(\Sigma_{ab}=\sigma_a\bar\sigma_b-\bar\sigma_b\sigma_a\). The next page scales the standard Hermitian Pauli matrices by \(1/\sqrt2\).}
\GSfield{Statement}{For the usual unscaled Hermitian Pauli matrices, \(\sigma_i\sigma_j=\delta_{ij}I+i\epsilon_{ijk}\sigma_k\). If all are scaled by \(1/\sqrt2\), the right side becomes \(\delta_{ij}I/2+i\epsilon_{ijk}\sigma_k/\sqrt2\). Spin generators use an antisymmetrized pair, proportional to \(\sigma_a\bar\sigma_b-\sigma_b\bar\sigma_a\), with normalization fixed consistently.}
\GSfield{Reconstruction}{The original scan lacks \(i\) in (3); neither the unscaled nor the explicitly scaled matrices satisfy that formula. Equation~(4) reverses matrix order instead of antisymmetrizing the spacetime labels and needs its convention restored.}
\GSfield{Fate}{These are elementary matrix identities, checked by direct multiplication and consistent with the classical spinor calculus \citep[Sec.~3.1]{woit2021euclidean}. They are faithful errors, not transcription repairs.}
\GSfield{Present status}{As of November 2026, the stated, qualified result is established; the notebook wording is interpreted only with the hypotheses specified above.}
\GSfield{Assessment}{literal statement: false as written; publication match: partial; verification: read the result; development: annotated; confidence: high}
\end{GSitem}
\begin{GSitem}{1979-465.4}{E}{Irreducible multispinors and epsilon contractions}
\GSfield{Kind; source}{theorem; \GSpage{P000849-P000850}}
\GSfield{The notes}{Symmetrizing \(2A\) undotted and \(2B\) dotted indices gives an irreducible representation; epsilon lowers indices.}
\GSfield{Statement}{For the complexified spin group \(SL_2(\mathbb C)\times SL_2(\mathbb C)\), the irreducible module labelled \((A,B)\) is \(\operatorname{Sym}^{2A}S_+\otimes\operatorname{Sym}^{2B}S_-\), of dimension \((2A+1)(2B+1)\), for nonnegative half-integers \(A,B\). The invariant alternating form identifies each standard module with its dual.}
\GSfield{Reconstruction}{Symmetrization is separate in the two index types. Raising and lowering requires a stated order of epsilon indices; writing the same numerical epsilon entries above and below is consistent only with the corresponding reversed contraction convention.}
\GSfield{Fate}{This is preexisting representation theory, checked in \citet[Sec.~3.2]{gilkey1984invariance}. The notebook's epsilon convention is ambiguous but is not automatically declared erroneous.}
\GSfield{Present status}{As of November 2026, the stated, qualified result is established; the notebook wording is interpreted only with the hypotheses specified above.}
\GSfield{Assessment}{literal statement: correct with missing hypotheses; publication match: exact; verification: read the result; development: annotated; confidence: high}
\end{GSitem}
\begin{GSitem}{1979-465.5}{E}{Two-form decomposition into symmetric spinors}
\GSfield{Kind; source}{theorem; \GSpage{P000850}}
\GSfield{The notes}{Equation~(10) splits an antisymmetric tensor into two symmetric two-spinors; a marginal tensor identity explains the split.}
\GSfield{Statement}{For two-dimensional spaces \(S,S'\), \(\Lambda^2(S\otimes S')=(\operatorname{Sym}^2S\otimes\Lambda^2S')\oplus(\Lambda^2S\otimes\operatorname{Sym}^2S')\). Thus \(F_{\alpha\dot\beta\,\gamma\dot\delta}=\epsilon_{\alpha\gamma}f_{\dot\beta\dot\delta}+\epsilon_{\dot\beta\dot\delta}f_{\alpha\gamma}\), with both \(f\)'s symmetric.}
\GSfield{Reconstruction}{The left-hand indices in the printed equation are ordered differently from the right-hand indices. The invariant tensor-product identity fixes the intended pairing without changing the underlying theorem.}
\GSfield{Fate}{Classical spinor linear algebra, directly verified by decomposition into symmetry types and the dimension count \(6=3+3\); compare \citet[Sec.~3.1]{woit2021euclidean}.}
\GSfield{Present status}{As of November 2026, the stated, qualified result is established; the notebook wording is interpreted only with the hypotheses specified above.}
\GSfield{Assessment}{literal statement: correct with missing hypotheses; publication match: exact; verification: read the result; development: annotated; confidence: high}
\end{GSitem}
\begin{GSitem}{1979-465.6}{E}{Hodge self-duality and signature}
\GSfield{Kind; source}{definition; \GSpage{P000850}}
\GSfield{The notes}{\(*F_{ab}=\frac12\epsilon_{abcd}F^{cd}\), eigenvalues \(\pm1\), with either symmetric two-spinor zero.}
\GSfield{Statement}{On an oriented Riemannian four-manifold, \(*^2=1\) on real two-forms, giving \(\Lambda^2=\Lambda^2_+\oplus\Lambda^2_-\). After complexification these correspond to the two symmetric-spinor summands. In Lorentz signature \(*^2=-1\) on two-forms, so the complex eigenvalues are \(\pm i\).}
\GSfield{Reconstruction}{The \(\pm1\) statement presupposes the Euclidean real form, unlike the preceding Lorentz conjugation notation.}
\GSfield{Fate}{Established Hodge/spinor theory \citep[Secs.~1.5 and~3.2]{gilkey1984invariance}; the signature qualification is essential.}
\GSfield{Present status}{As of November 2026, the stated, qualified result is established; the notebook wording is interpreted only with the hypotheses specified above.}
\GSfield{Assessment}{literal statement: correct with missing hypotheses; publication match: exact; verification: read the result; development: annotated; confidence: high}
\end{GSitem}
\begin{GSitem}{1979-465.7}{E}{Self-dual full curvature implies the vacuum Einstein equation}
\GSfield{Kind; source}{theorem; \GSpage{P000850}}
\GSfield{The notes}{Einstein equations follow from the self-duality condition on \(R_{abcd}\).}
\GSfield{Statement}{For an oriented Riemannian four-manifold, if the full Riemann curvature operator has image entirely in one Hodge summand, equivalently \(*R=\pm R\) on one antisymmetric index pair, then the metric is Ricci-flat.}
\GSfield{Reconstruction}{In the block decomposition of curvature on \(\Lambda^2_+\oplus\Lambda^2_-\), off-diagonal blocks encode traceless Ricci and each diagonal trace is scalar curvature. A vanishing Hodge block and symmetry force both to vanish. Self-duality of the Weyl tensor alone does not imply Ricci-flatness.}
\GSfield{Fate}{This is the gravitational-instanton setting of \citet{hawkingpope1978symmetry}; the block argument supplies the exact missing meaning of self-duality. Primary abstract checked, full original paper not read.}
\GSfield{Present status}{As of November 2026, the stated, qualified result is established; the notebook wording is interpreted only with the hypotheses specified above.}
\GSfield{Assessment}{literal statement: correct with missing hypotheses; publication match: exact; verification: indirect; development: annotated; confidence: high}
\end{GSitem}
\begin{GSitem}{1979-465.8}{F}{The proposed universal higher-spin Laplacian}
\GSfield{Kind; source}{construction; \GSpage{P000851-P000852}}
\GSfield{The notes}{Equations~(15), (23)--(25) define \(\Delta=-\Box+X\), with spin-generator curvature contractions and different integer/half-integer formulae.}
\GSfield{Statement}{A Laplace-type operator on a spin-associated bundle is determined by a connection and a curvature endomorphism \(X\), once the representation, metric signature and normalizations are fixed. A common principal symbol does not identify the endomorphism or the physical gauge-constrained operator.}
\GSfield{Reconstruction}{The barred field in (15), generator normalization, symmetrization factor and \(R_+\) normalization are unresolved. The notebook changes from irreducible symmetric spinors to unsymmetrized brackets later. Consequently its universal numerical formula is not certified here.}
\GSfield{Fate}{\citet{christensenduff1979indices} is an exact bibliographic match for arbitrary-spin operators; only its primary abstract was accessible. \citet[Sec.~4]{duff1982ultraviolet} corroborates the operator/heat-coefficient programme but does not repair every coefficient in the notebook.}
\GSfield{Present status}{As of November 2026, the intended published framework is identified; the exact notebook operator remains underdetermined.}
\GSfield{Assessment}{literal statement: not determinable; publication match: partial; verification: indirect; development: annotated; confidence: medium}
\end{GSitem}
\begin{GSitem}{1979-465.9}{E}{Formal symmetry versus self-adjoint spectral theory}
\GSfield{Kind; source}{claim; \GSpage{P000851}}
\GSfield{The notes}{\((\phi,\Delta\phi)=(\Delta\phi,\phi)\) is listed as the first property.}
\GSfield{Statement}{For a metric connection and fibrewise self-adjoint endomorphism \(X\) on a closed Riemannian manifold, \(\nabla^*\nabla+X\) is formally symmetric. Its standard elliptic self-adjoint realization has discrete real spectrum and finite-dimensional eigenspaces. Boundary or noncompact settings require domains and additional hypotheses.}
\GSfield{Reconstruction}{The single diagonal equality does not specify an operator domain. In a real inner product it is tautological for any operator on that vector, and must be replaced by the polarized symmetry identity.}
\GSfield{Fate}{Standard elliptic theory is developed in \citet[Secs.~1.3--1.6]{gilkey1984invariance}; the reconstruction supplies the assumptions needed for all subsequent heat and index statements.}
\GSfield{Present status}{As of November 2026, the stated, qualified result is established; the notebook wording is interpreted only with the hypotheses specified above.}
\GSfield{Assessment}{literal statement: correct with missing hypotheses; publication match: exact; verification: read the result; development: annotated; confidence: high}
\end{GSitem}
\begin{GSitem}{1979-465.10}{E}{The scalar specialization}
\GSfield{Kind; source}{claim; \GSpage{P000851}}
\GSfield{The notes}{\(\Delta\phi(0,0)=-\Box\phi\), with covariant derivatives and a spin connection.}
\GSfield{Statement}{On scalar functions with the trivial coefficient bundle, the positive Riemannian Laplacian is \(\Delta_0=d^*d=-\nabla^a\nabla_a\) under the stated sign convention. A nontrivial twisting connection or a scalar potential produces a different operator.}
\GSfield{Reconstruction}{The spin connection acts trivially on untwisted scalars. The factor in a two-spinor contraction depends on the \(1/\sqrt2\) normalization of the soldering matrices.}
\GSfield{Fate}{Preexisting Hodge theory \citep[Sec.~1.5]{gilkey1984invariance}.}
\GSfield{Present status}{As of November 2026, the stated, qualified result is established; the notebook wording is interpreted only with the hypotheses specified above.}
\GSfield{Assessment}{literal statement: correct with missing hypotheses; publication match: exact; verification: read the result; development: annotated; confidence: high}
\end{GSitem}
\begin{GSitem}{1979-465.11}{E}{The spinor Lichnerowicz formula}
\GSfield{Kind; source}{theorem; \GSpage{P000851}}
\GSfield{The notes}{Chiral Dirac spinors are written \(\Psi_\pm=(1\pm\gamma_5)\Psi\), and \(\Delta\Psi_\pm=(-\Box+R/4)\Psi_\pm\).}
\GSfield{Statement}{On a Riemannian spin manifold, the untwisted Dirac operator satisfies \(D^2=\nabla^*\nabla+R/4\). In even dimension the chiral projections are \(P_\pm=(1\pm\gamma_5)/2\), provided \(\gamma_5^2=1\).}
\GSfield{Reconstruction}{The displayed \(1\pm\gamma_5\) yields twice the projected spinor; it is not an idempotent projection. Additional gauge twisting adds the corresponding curvature Clifford term to \(D^2\).}
\GSfield{Fate}{The classical spin-complex identity is treated in \citet[Sec.~3.4]{gilkey1984invariance}. The factor \(1/2\) is checked directly from \(P_\pm^2=P_\pm\).}
\GSfield{Present status}{As of November 2026, the stated, qualified result is established; the notebook wording is interpreted only with the hypotheses specified above.}
\GSfield{Assessment}{literal statement: correct with missing hypotheses; publication match: exact; verification: read the result; development: annotated; confidence: high}
\end{GSitem}
\begin{GSitem}{1979-465.12}{E}{The Maxwell/Hodge operator in Lorenz gauge}
\GSfield{Kind; source}{theorem; \GSpage{P000851}}
\GSfield{The notes}{\(\Delta A_a=(-\Box g_{ab}+R_{ab})A^b,\ \mathcal D^aA_a=0\).}
\GSfield{Statement}{For a real one-form on a Riemannian manifold, \((d\delta+\delta d)A=-\nabla^b\nabla_bA+\operatorname{Ric}(A)\). Under the Lorenz gauge \(\delta A=0\), the source-free Maxwell potential equation \(\delta dA=0\) becomes \(\Delta_1A=0\).}
\GSfield{Reconstruction}{The sign of curvature depends on the Riemann-tensor convention, and Lorentz signature gives the corresponding wave operator rather than a positive elliptic Laplacian.}
\GSfield{Fate}{Classical Weitzenböck/Hodge theory \citep[Sec.~1.5]{gilkey1984invariance}; it is a low-spin instance of the framework cited by Kallosh.}
\GSfield{Present status}{As of November 2026, the stated, qualified result is established; the notebook wording is interpreted only with the hypotheses specified above.}
\GSfield{Assessment}{literal statement: correct with missing hypotheses; publication match: exact; verification: read the result; development: annotated; confidence: high}
\end{GSitem}
\begin{GSitem}{1979-465.13}{F}{The spin-three-halves example}
\GSfield{Kind; source}{claim; \GSpage{P000851}}
\GSfield{The notes}{For \(s=3/2\), the notes say somebody and Schwinger, without an equation.}
\GSfield{Statement}{The natural identification is the Rarita--Schwinger vector-spinor equation, but the notebook does not specify its mass, gauge symmetry, constraints, background or squared operator.}
\GSfield{Reconstruction}{This identification repairs a name, not a complete field equation. A unique higher-spin Laplacian cannot be reconstructed merely from the spin value.}
\GSfield{Fate}{Rarita--Schwinger fields and supergravity are covered in the checked primary discussion \citet{duff1982ultraviolet}; the precise notebook assertion remains unspecified.}
\GSfield{Present status}{As of November 2026, this is an identifiable preexisting example, with insufficient detail for an exact operator comparison.}
\GSfield{Assessment}{literal statement: not determinable; publication match: thematic only; verification: indirect; development: annotated; confidence: high}
\end{GSitem}
\begin{GSitem}{1979-465.14}{F}{The spin-two Einstein example}
\GSfield{Kind; source}{claim; \GSpage{P000851}}
\GSfield{The notes}{The \(s=2\) line contains only Einstein.}
\GSfield{Statement}{Linearized gravity acts on symmetric two-tensors modulo infinitesimal diffeomorphisms; a Laplace-type operator emerges after a background and gauge are fixed. The name alone does not choose that operator.}
\GSfield{Reconstruction}{Ricci-flat, Einstein-with-cosmological-constant and general off-shell backgrounds have different lower-order curvature terms.}
\GSfield{Fate}{The spin-two heat-kernel setting is discussed in \citet[Sec.~4]{duff1982ultraviolet}, but no exact numerical operator is claimed as recovered from this one-word line.}
\GSfield{Present status}{As of November 2026, the physical example is clear but its mathematical operator is underdetermined.}
\GSfield{Assessment}{literal statement: not determinable; publication match: thematic only; verification: indirect; development: annotated; confidence: high}
\end{GSitem}
\begin{GSitem}{1979-465.15}{E}{Differential forms and the de Rham Laplacian}
\GSfield{Kind; source}{construction; \GSpage{P000851}}
\GSfield{The notes}{\(w=(1/p!)\omega_{\mu_1\ldots\mu_p}dx^{\mu_1}\wedge\cdots\wedge dx^{\mu_p}\), \(\Delta=d\delta+\delta d\), with scalar, vector and two-form spin types listed.}
\GSfield{Statement}{On a closed oriented Riemannian manifold, \(d\) and its \(L^2\)-adjoint \(\delta\) define a nonnegative elliptic Laplacian on forms; its kernel in degree \(p\) represents \(H^p_{\mathrm{dR}}\). In dimension four the two-form representation splits into \((1,0)\) and \((0,1)\).}
\GSfield{Reconstruction}{The notebook's language that forms describe bosons concerns transformation type. Form degree, Grassmann parity of field coefficients and physical statistics are different gradings.}
\GSfield{Fate}{This is preexisting Hodge theory \citep[Sec.~1.5]{gilkey1984invariance}; the displayed coordinate expression is correct.}
\GSfield{Present status}{As of November 2026, the stated, qualified result is established; the notebook wording is interpreted only with the hypotheses specified above.}
\GSfield{Assessment}{literal statement: correct with missing hypotheses; publication match: exact; verification: read the result; development: annotated; confidence: high}
\end{GSitem}
\begin{GSitem}{1979-465.16}{F}{Mixed-parity differential forms and an index theorem}
\GSfield{Kind; source}{problem; \GSpage{P000851}; Renata Kallosh and I. M. Gelfand}
\GSfield{The notes}{Needed: anticommuting \(dx\), then forms with commuting and anticommuting differentials, an operator \(\Delta\), and an index theorem; Gelfand says it need not be \(d\delta+\delta d\).}
\GSfield{Statement}{On a supermanifold, the parity-shifted tangent construction has odd \(dx\) for even coordinates \(x\) and even \(d\theta\) for odd coordinates \(\theta\). This gives a differential calculus, but no particular Fredholm complex or index theorem is determined by the notebook.}
\GSfield{Reconstruction}{Ordinary exterior differentials already anticommute. The substantive requested extension is the mixed parity calculus and a suitable analytical index, with domains, compactness and symbol assumptions still absent.}
\GSfield{Fate}{The Berezin--Leites supermanifold framework predates the lecture \citep{berezinleites1975supermanifolds}; its existence does not resolve the unspecified analytic request. This repeats the programme of 1979-462.10.}
\GSfield{Present status}{As of November 2026, the construction language is established but no unique super-index problem can be inferred.}
\GSfield{Assessment}{literal statement: not determinable; publication match: thematic only; verification: indirect; development: annotated; confidence: high}
\end{GSitem}
\begin{GSitem}{1979-465.17}{E}{The spin-five-halves gravitational-coupling obstruction}
\GSfield{Kind; source}{claim; \GSpage{P000851}}
\GSfield{The notes}{For spin \(5/2\), none of the equations will do because an interaction cannot be constructed; for spin \(3/2\) it has just been done in supergravity.}
\GSfield{Statement}{Minimal coupling of an isolated massless spin-\(5/2\) field to a general gravitational background introduces Weyl-curvature terms in the gauge consistency conditions. This obstructs the naive coupling; it is not a theorem excluding every nonminimal, self-dual, lower-dimensional or infinite-tower theory.}
\GSfield{Reconstruction}{The scope restriction matters even in the 1979 source: the identified paper explicitly discusses self-dual exceptions.}
\GSfield{Fate}{\citet{aragonedeser1979hypergravity}, published 24 September 1979, is the directly relevant preexisting result. Its primary publisher abstract was checked, but the full three-page proof was not retrieved.}
\GSfield{Present status}{As of November 2026, the stated, qualified result is established; the notebook wording is interpreted only with the hypotheses specified above.}
\GSfield{Assessment}{literal statement: false as written; publication match: partial; verification: indirect; development: annotated; confidence: high}
\end{GSitem}
\begin{GSitem}{1979-465.18}{E}{Reducible tensor powers versus irreducible spin labels}
\GSfield{Kind; source}{theorem; \GSpage{P000852}}
\GSfield{The notes}{\([1,0]=(1,0)\oplus(0,0)\), with a two-index spinor split into symmetric and epsilon-trace parts.}
\GSfield{Statement}{For a two-dimensional spin module \(S\), \(S\otimes S=\operatorname{Sym}^2S\oplus\Lambda^2S\), with dimensions \(4=3+1\). Thus brackets can consistently denote an unsymmetrized tensor power, while parentheses denote an irreducible spin type.}
\GSfield{Reconstruction}{The square brackets are essential; reading both sides as irreducible labels would create a false equation. The trace sign in the indexed formula depends on the epsilon raising convention.}
\GSfield{Fate}{This is elementary preexisting representation theory \citep[Sec.~3.2]{gilkey1984invariance}. It explains why the powers \(4^S\) in the later shorthand need not be dimensions of irreducible \((A,B)\) modules.}
\GSfield{Present status}{As of November 2026, the stated, qualified result is established; the notebook wording is interpreted only with the hypotheses specified above.}
\GSfield{Assessment}{literal statement: correct with missing hypotheses; publication match: exact; verification: read the result; development: annotated; confidence: high}
\end{GSitem}
\begin{GSitem}{1979-465.19}{F}{Compatibility and spectral intertwiners}
\GSfield{Kind; source}{claim; \GSpage{P000852-P000853}}
\GSfield{The notes}{Symmetry and the eigenvalue equation imply a compatibility identity; \([\nabla,\Delta]\phi=0\), and differentiating changes \([A,B]\) to \([A-\frac12,B+\frac12]\), preserving nonzero modes.}
\GSfield{Statement}{An operator \(Q:E\to F\) intertwines two Laplacians if \(QP_E=P_FQ\). To obtain an isomorphism of positive-eigenvalue spaces one also needs an inverse relation such as \(Q^*Q=P_E\) and \(QQ^*=P_F\), or its appropriately normalized counterpart. Then \(Q/\sqrt\lambda\) is an isomorphism on each \(\lambda>0\) eigenspace.}
\GSfield{Reconstruction}{Covariant derivatives do not commute with Laplacians on arbitrary curved backgrounds. The special curvature/compatibility assumptions are indispensable; an intertwiner alone may have a nontrivial kernel.}
\GSfield{Fate}{The identified higher-spin/self-dual work concerns such relations \citep{christensenduff1979indices,hawkingpope1978symmetry}; exact agreement with equations~(27), (33)--(36) was not verified. The rigorous conditional operator lemma above is checked directly.}
\GSfield{Present status}{As of November 2026, the mechanism is established but the notebook does not supply enough hypotheses to certify its particular differential intertwiner.}
\GSfield{Assessment}{literal statement: correct with missing hypotheses; publication match: partial; verification: indirect; development: annotated; confidence: medium}
\end{GSitem}
\begin{GSitem}{1979-465.20}{E}{Heat-kernel expansion and its spectral realization}
\GSfield{Kind; source}{theorem; \GSpage{P000852-P000853}}
\GSfield{The notes}{The heat equation has delta initial data, a four-dimensional expansion \(\sum b_k t^{(k-4)/2}\), and a sum over orthonormal eigenfunctions.}
\GSfield{Statement}{For a nonnegative self-adjoint Laplace-type operator \(P\) on a Hermitian bundle over a closed four-manifold, \(K_t(x,x')=\sum_j e^{-t\lambda_j}\phi_j(x)\otimes\phi_j(x')^*\). Its trace has asymptotic expansion \(\sum_{k\ge0}B_k t^{(k-4)/2}\), with \(B_k=\int\operatorname{tr}b_k(x)\,d\operatorname{vol}\); odd-index coefficients vanish in the closed case.}
\GSfield{Reconstruction}{The adjoint on the second eigenfunction and fibre trace are required for complex vector bundles. The sum converges smoothly for \(t>0\), and the delta initial condition is distributional.}
\GSfield{Fate}{Standard spectral/heat theory, fully treated in \citet[Secs.~1.6--1.7]{gilkey1984invariance}. The malformed exponent in (43) is fixed by (29) and the dimension.}
\GSfield{Present status}{As of November 2026, the stated, qualified result is established; the notebook wording is interpreted only with the hypotheses specified above.}
\GSfield{Assessment}{literal statement: correct with missing hypotheses; publication match: exact; verification: read the result; development: annotated; confidence: high}
\end{GSitem}
\begin{GSitem}{1979-465.21}{E}{The regularized nonzero-mode contribution need not be an integer}
\GSfield{Kind; source}{claim; \GSpage{P000852-P000853}}
\GSfield{The notes}{\(B_4=n+m\); the notebook first conjectures a common multiplicity, explicitly marks it incorrect, then calls \(m\) a regularized sum and nevertheless an integer.}
\GSfield{Statement}{For a nonnegative self-adjoint Laplace-type operator \(P\) on a closed smooth four-dimensional Riemannian manifold, form \(\zeta_P\) from its positive eigenvalues and let \(B_4\) denote the constant term in the integrated heat expansion. Then
\[
m=\zeta_P(0)=B_4-\dim\ker P.
\]
This value need not be an integer. For the scalar Laplacian on the unit round \(S^4\), the standard coefficient gives \(B_4=29/90\) and \(m=-61/90\), since the kernel consists of constants.}
\GSfield{Reconstruction}{\oursComputed\ For that sphere \(R=12\), \(|\operatorname{Ric}|^2=36\), \(|\operatorname{Riem}|^2=24\), and \(\operatorname{vol}=8\pi^2/3\). Substituting in \((360(4\pi)^2)^{-1}\int(5R^2-2|\operatorname{Ric}|^2+2|\operatorname{Riem}|^2)\) yields \(29/90\). An index supertrace may be integral even when individual regularized counts are not.}
\GSfield{Fate}{The heat/zeta relation and coefficient formula are in \citet[Secs.~1.10 and~4.8]{gilkey1984invariance}. The notebook's own rejection of the multiplicity interpretation is retained, and the later integrality claim is independently disproved.}
\GSfield{Present status}{As of November 2026, the stated, qualified result is established; the notebook wording is interpreted only with the hypotheses specified above.}
\GSfield{Assessment}{literal statement: false as written; publication match: partial; verification: read the result; development: annotated; confidence: high}
\end{GSitem}
\begin{GSitem}{1979-465.22}{E}{Heat coefficients through order six and the boundary caveat}
\GSfield{Kind; source}{claim; \GSpage{P000852}}
\GSfield{The notes}{Gilkey has computed through \(B_6\), with 1975 and a correction mentioning 1979; boundary \(b_4\) has recently been written.}
\GSfield{Statement}{For Laplace-type operators on closed manifolds, local heat coefficients through order six are universal curvature/endomorphism expressions. With boundary, a coefficient is not determined until an elliptic boundary condition is specified; it includes boundary invariants as well as interior terms.}
\GSfield{Reconstruction}{The exact 1979 reference is recoverable. The anonymous boundary claim is not: Dirichlet, Robin, mixed and spectral boundary conditions yield different statements.}
\GSfield{Fate}{\citet{gilkey1979recursion}, Compositio 38, pp.~201--240, is the verified 1979 paper; the later author monograph \citep[Sec.~4.8]{gilkey1984invariance} gives explicit formulae. The unspecified boundary contribution remains only a thematic match.}
\GSfield{Present status}{As of November 2026, the stated, qualified result is established; the notebook wording is interpreted only with the hypotheses specified above.}
\GSfield{Assessment}{literal statement: correct with missing hypotheses; publication match: partial; verification: read the result; development: annotated; confidence: high}
\end{GSitem}
\begin{GSitem}{1979-465.23}{F}{Index differences after spectral cancellation}
\GSfield{Kind; source}{claim; \GSpage{P000853}}
\GSfield{The notes}{If \(m[S,T]\) is independent of \(T\), then \(n[S,T]-n[S,T']=B_4[S,T]-B_4[S,T']\), followed by formula~(39).}
\GSfield{Statement}{For two operators satisfying equality of their regularized nonzero-mode contributions, subtracting \(B_4=\dim\ker+\zeta(0)\) gives the displayed kernel-dimension difference. To turn it into a topological formula one needs a verified elliptic-complex or spectral-pairing construction and the correct heat coefficients.}
\GSfield{Reconstruction}{This conditional subtraction is valid, but the numerical factor in (39) is not independently verified and disagrees with P000841: \(s-\frac23\) versus \(s-\frac13\). The distinction between tensor powers and irreducible bundles is unresolved.}
\GSfield{Fate}{The general mechanism is established heat-index theory \citep[Sec.~1.7]{gilkey1984invariance}; the exact proposed higher-spin family matches the topic of \citet{christensenduff1979indices}, not a checked coefficient-by-coefficient theorem.}
\GSfield{Present status}{As of November 2026, the algebraic subtraction is verified, while the notebook's numerical higher-spin identity remains uncertain.}
\GSfield{Assessment}{literal statement: not determinable; publication match: partial; verification: indirect; development: annotated; confidence: medium}
\end{GSitem}
\begin{GSitem}{1979-465.24}{E}{Euler and Pontryagin integrals}
\GSfield{Kind; source}{claim; \GSpage{P000853}}
\GSfield{The notes}{\(\chi=\int {}^*R_{abcd}\,{}^*R^{abcd}\,d\operatorname{vol}\) and \(p=\int{}^*R_{abcd}R^{abcd}\,d\operatorname{vol}\), without constants.}
\GSfield{Statement}{On a closed oriented Riemannian four-manifold,
\[
\chi(M)=\frac1{32\pi^2}\int_M
\bigl(|\operatorname{Riem}|^2-4|\operatorname{Ric}|^2+R^2\bigr)\,d\operatorname{vol}.
\]
With curvature two-form matrix \(\Omega\),
\[
p_1(TM)=-\frac1{8\pi^2}\operatorname{tr}(\Omega\wedge\Omega)
\]
in the stated Chern--Weil convention.}
\GSfield{Reconstruction}{Applying the same Hodge star to both factors preserves their inner product and gives \(|R|^2\), which is not the general Euler density. Ricci-flatness simplifies the Euler formula but does not remove its normalization. The Pontryagin scalar contraction also needs normalization and orientation convention.}
\GSfield{Fate}{These are classical Chern--Weil and Gauss--Bonnet formulae \citep[Secs.~2.2, 2.4 and~3.1]{gilkey1984invariance}. The PNG confirms both stars in the first printed formula.}
\GSfield{Present status}{As of November 2026, the stated, qualified result is established; the notebook wording is interpreted only with the hypotheses specified above.}
\GSfield{Assessment}{literal statement: false as written; publication match: partial; verification: read the result; development: annotated; confidence: high}
\end{GSitem}

\GStalk{466}{1979-10-08}{I. M. Gelfand}{P000854}{Reading programme and exercises}
The actual contribution is the opening of P000854: a reading programme and exercises on Gaussian and quartic Fourier integrals and on self-adjoint operators for a symplectic form. The exercises are definite mathematical items, although the contribution is primarily a programme. Their solutions must distinguish the self-adjoint relation for an alternating form from the Hamiltonian, skew-adjoint relation.

\begin{GSitem}{1979-466.1}{E}{The two Gaussian Fourier exercises}
\GSfield{Kind; source}{problem; \GSpage{P000854}}
\GSfield{The notes}{Approximate \(\int e^{-x^2}\cos1000x\,dx\) and \(\int e^{-1000x^2}\cos x\,dx\), over the full real line.}
\GSfield{Statement}{For \(a>0\), \(\int_{\mathbb R}e^{-ax^2}\cos(bx)\,dx=\sqrt{\pi/a}\,e^{-b^2/(4a)}\). Thus the answers are \(\sqrt\pi\,e^{-250000}\) and \(\sqrt{\pi/1000}\,e^{-1/4000}\), respectively.}
\GSfield{Reconstruction}{\ours\ The exact Fourier formula follows by differentiating in \(b\), integrating by parts to obtain \(I'=-bI/(2a)\), and using the Gaussian integral at \(b=0\). The first integral is positive but extraordinarily small; coarse oscillatory quadrature is unsuitable.}
\GSfield{Fate}{Elementary preexisting Fourier analysis; the transform normalization is consistent with \citet[Sec.~1.1]{gilkey1984invariance}. The exercise is solved explicitly here, with no research-problem status.}
\GSfield{Present status}{As of November 2026, the stated, qualified result is established; the notebook wording is interpreted only with the hypotheses specified above.}
\GSfield{Assessment}{literal statement: correct; publication match: exact; verification: read the result; development: annotated; confidence: high}
\end{GSitem}
\begin{GSitem}{1979-466.2}{E}{The narrow quartic Fourier exercise}
\GSfield{Kind; source}{problem; \GSpage{P000854}}
\GSfield{The notes}{Approximate \(\int_{\mathbb R}e^{-1000x^4}\cos x\,dx\).}
\GSfield{Statement}{For \(a>0\), \(I(a,b)=\frac1{2a^{1/4}}\sum_{k\ge0}\frac{(-1)^k\Gamma((2k+1)/4)}{(2k)!}(ba^{-1/4})^{2k}\). At \(a=1000,b=1\), this gives approximately \(0.32064797070715458\).}
\GSfield{Reconstruction}{\ours\ Termwise integration follows from domination by \(e^{-ax^4+|b||x|}\). The series converges for every finite \(b\). The value was checked independently by direct high-precision quadrature; the small scaled frequency makes this stable.}
\GSfield{Fate}{An elementary scaling and Taylor-series exercise, fully resolved by the displayed derivation. No exact historical publication is required for the calculation.}
\GSfield{Present status}{As of November 2026, the stated, qualified result is established; the notebook wording is interpreted only with the hypotheses specified above.}
\GSfield{Assessment}{literal statement: correct; publication match: none found; verification: indirect; development: annotated; confidence: high}
\end{GSitem}
\begin{GSitem}{1979-466.3}{E}{The highly oscillatory quartic Fourier exercise}
\GSfield{Kind; source}{problem; \GSpage{P000854}}
\GSfield{The notes}{Approximate \(\int_{\mathbb R}e^{-x^4}\cos1000x\,dx\).}
\GSfield{Statement}{Put \(r=(b/4)^{1/3}\). As \(b\to+\infty\), \(I(1,b)=\sqrt{2\pi/3}\,r^{-1}e^{-3r^4/2}\{\cos(3\sqrt3\,r^4/2-\pi/6)+O(r^{-4})\}\). At \(b=1000\), shifted-contour quadrature gives approximately \(1.7020109384\times10^{-1027}\); the leading saddle expression gives \(1.7019868243\times10^{-1027}\).}
\GSfield{Reconstruction}{\oursComputed\ The contributing saddles of \(-z^4+ibz\) are \(r e^{i\pi/6}\) and \(r e^{5i\pi/6}\). Shifting the real contour to height \(r/2\) makes their exponentially small contribution numerically accessible. The error is absolute inside the braces, not a uniform relative error near zeros of the cosine.}
\GSfield{Fate}{The exercise is resolved by elementary steepest descent and a separately saved numerical check. The convergent series in the preceding item is exact but suffers severe cancellation at this frequency.}
\GSfield{Present status}{As of November 2026, the stated, qualified result is established; the notebook wording is interpreted only with the hypotheses specified above.}
\GSfield{Assessment}{literal statement: correct; publication match: none found; verification: indirect; development: annotated; confidence: high}
\end{GSitem}
\begin{GSitem}{1979-466.4}{E}{Canonical form of a symplectically self-adjoint operator}
\GSfield{Kind; source}{problem; \GSpage{P000854}}
\GSfield{The notes}{Find the canonical form of a self-adjoint operator in a symplectic space; the adjective is abbreviated in the manuscript.}
\GSfield{Statement}{Over \(\mathbb C\), let \(V\) be a finite-dimensional vector space with nondegenerate alternating form \(\omega\), and let \(A:V\to V\) satisfy \(\omega(Au,v)=\omega(u,Av)\). There is a symplectic basis in which \(A=\operatorname{diag}(J,J^{\mathsf T})\), with \(J\) in Jordan normal form. Each Jordan block therefore occurs with its paired transpose.}
\GSfield{Reconstruction}{This is self-adjointness for the alternating form, not the skew-adjoint Hamiltonian condition. Over a perfect field of characteristic not two, an analogous polarization holds; a Jordan matrix over the base field additionally requires split characteristic polynomial.}
\GSfield{Fate}{The exact result is proved in \citet[Theorem~1 and Corollary~11]{sanabria2017symplectic}, which explicitly identifies earlier classification results. This is a standard linear-algebra exercise rather than a claim of a new 2017 theorem.}
\GSfield{Present status}{As of November 2026, the stated, qualified result is established; the notebook wording is interpreted only with the hypotheses specified above.}
\GSfield{Assessment}{literal statement: correct with missing hypotheses; publication match: exact; verification: read the result; development: annotated; confidence: high}
\end{GSitem}

\GStalk{467}{1979-10-08}{S. G. Gindikin}{P000854-P000858}{Twistors and the conformal compactification of Minkowski space}
Starting on P000854, Gindikin describes the Hermitian space of signature (2,2), its two-plane Grassmannian, tube domains, the real Minkowski chart and null incidence. This continues the geometry of 1978-455.1. The manuscript's unrestricted affine chart stabilizer includes dilations, and its final point/ray correspondence interchanges the two sides of the twistor incidence relation. The higher-dimensional challenge is a programme, not a unique conjecture.

\begin{GSitem}{1979-467.1}{E}{Projective twistors and the Hermitian null cone}
\GSfield{Kind; source}{definition; \GSpage{P000854-P000856}}
\GSfield{The notes}{\(\mathbb C^4\setminus0\to\mathbb CP^3\), a Hermitian form of signature \((2,2)\), positive/negative/null twistors and a seven-dimensional cone \(N\).}
\GSfield{Statement}{Let \(T=\mathbb C^4\) have nondegenerate Hermitian form \(h\) of signature \((2,2)\). A projective twistor is a complex line \([z]\in\mathbb P(T)\), whose positive, negative or null type is well defined because \(h(cz,cz)=|c|^2h(z,z)\). The nonzero null cone is a smooth real seven-manifold; its projectivization has real dimension five.}
\GSfield{Reconstruction}{The origin is a singular point of the cone and is excluded in the dimension statement. The dual projective space describes the dual twistor construction, not an extra choice of four-dimensional spacetime coordinates.}
\GSfield{Fate}{This is standard preexisting twistor geometry, checked in the classical-geometry sections of \citet[Secs.~3.1 and~3.3.2]{woit2021euclidean}. No claim from that paper's unification proposal is used.}
\GSfield{Present status}{As of November 2026, the stated, qualified result is established; the notebook wording is interpreted only with the hypotheses specified above.}
\GSfield{Assessment}{literal statement: correct with missing hypotheses; publication match: exact; verification: read the result; development: annotated; confidence: high}
\end{GSitem}
\begin{GSitem}{1979-467.2}{E}{The Grassmannian and its frame bundle}
\GSfield{Kind; source}{construction; \GSpage{P000855}}
\GSfield{The notes}{Projective lines in \(\mathbb CP^3\) form \(G_{4,2}^{\mathbb C}\); ordered two-frames form a principal bundle \(E\to G\).}
\GSfield{Statement}{The Grassmannian \(G=\operatorname{Gr}(2,T)\) has complex dimension four. The space \(E\) of ordered linearly independent pairs in \(T\) maps to \(G\) by span, with free transitive right \(GL_2(\mathbb C)\) action on each fibre. It is the frame bundle of the tautological rank-two bundle.}
\GSfield{Reconstruction}{The two-frames are bases of variable two-planes; they need not be orthonormal for the indefinite Hermitian form.}
\GSfield{Fate}{Classical Grassmannian geometry; the twistor incidence interpretation is checked in \citet[Sec.~3.3]{woit2021euclidean}. This continues 1978-455.1.}
\GSfield{Present status}{As of November 2026, the stated, qualified result is established; the notebook wording is interpreted only with the hypotheses specified above.}
\GSfield{Assessment}{literal statement: correct; publication match: exact; verification: read the result; development: annotated; confidence: high}
\end{GSitem}
\begin{GSitem}{1979-467.3}{E}{The graph coordinate chart}
\GSfield{Kind; source}{construction; \GSpage{P000855-P000856}}
\GSfield{The notes}{After choosing transverse two-planes \(L,L'\), P000855 says any other is a graph; P000856 restores \(\widetilde L\cap L'=0\).}
\GSfield{Statement}{For \(T=L\oplus L'\), the open set \(G_0=\{W:W\cap L'=0\}\) is identified with \(\operatorname{Hom}(L,L')\) by \(U\mapsto\operatorname{graph}U\). Every such graph has basis \(e_A+\sum_{B'}u_{AB'}e^{B'}\), with row/column conventions fixed.}
\GSfield{Reconstruction}{The transversality condition excludes a Schubert divisor. The later page supplies it explicitly, so the earlier universal wording is corrected internally.}
\GSfield{Fate}{Elementary Grassmannian coordinates, checked directly by projection \(W\to L\); the same chart construction is used in \citet[Sec.~3.3]{woit2021euclidean}.}
\GSfield{Present status}{As of November 2026, the stated, qualified result is established; the notebook wording is interpreted only with the hypotheses specified above.}
\GSfield{Assessment}{literal statement: correct with missing hypotheses; publication match: exact; verification: read the result; development: annotated; confidence: high}
\end{GSitem}
\begin{GSitem}{1979-467.4}{E}{Complementary isotropic planes and the GL2 subgroup}
\GSfield{Kind; source}{problem; \GSpage{P000855-P000856}}
\GSfield{The notes}{Prove the existence of transverse isotropic \(L,L'\), then prove that \(U(2,2)\) contains a copy of \(GL_2(\mathbb C)\).}
\GSfield{Statement}{In a Hermitian space of signature \((2,2)\), choose positive orthonormal vectors \(p_i\) and negative orthonormal vectors \(q_i\). Their sums and suitable imaginary multiples of their differences give complementary isotropic two-planes with form matrix \(H=\left(\begin{smallmatrix}0&iI\\-iI&0\end{smallmatrix}\right)\). Then \(g\mapsto\operatorname{diag}(g,(g^*)^{-1})\) is an injective homomorphism into \(U(H)\).}
\GSfield{Reconstruction}{Direct block multiplication gives \(G^*HG=H\). The mixed pairings determine the conjugate inverse transformation of the second basis. Their reversed pairings are \(-i\), so the phrase all other scalar products zero means only the independent pairings.}
\GSfield{Fate}{The exercise is solved by elementary Witt decomposition and the displayed block check; it belongs to preexisting indefinite-unitary linear algebra \citep[Sec.~3.3.2]{woit2021euclidean}.}
\GSfield{Present status}{As of November 2026, the stated, qualified result is established; the notebook wording is interpreted only with the hypotheses specified above.}
\GSfield{Assessment}{literal statement: correct with missing hypotheses; publication match: exact; verification: read the result; development: annotated; confidence: high}
\end{GSitem}
\begin{GSitem}{1979-467.5}{E}{Tube domains, real boundary and boundary strata}
\GSfield{Kind; source}{construction; \GSpage{P000856-P000857}}
\GSfield{The notes}{\(S_\pm\) are positive/negative two-planes, with \((U-U^*)/i\) definite; \(M_0=\{U=U^*\}\) is their common skeleton.}
\GSfield{Statement}{In a compatible graph convention, positivity of the restricted Hermitian form is equivalent to positive definiteness of \(\operatorname{Im}U=(U-U^*)/(2i)\); negativity is analogous. The real chart of maximal isotropic planes is \(\operatorname{Herm}_2(\mathbb C)\), of real dimension four. The finite boundary of the positive tube has \(\operatorname{Im}U\ge0\) singular, including rank-one and rank-zero strata.}
\GSfield{Reconstruction}{The Hermitian locus is the distinguished real boundary, not the entire topological boundary. Positive and negative planes are automatically transverse to the fixed isotropic complement; maximal isotropic planes need not be.}
\GSfield{Fate}{Classical matrix-domain geometry underlying the Minkowski real form \citep[Sec.~3.3.2]{woit2021euclidean}; the definiteness criterion follows by restricting \(H\) to a graph.}
\GSfield{Present status}{As of November 2026, the stated, qualified result is established; the notebook wording is interpreted only with the hypotheses specified above.}
\GSfield{Assessment}{literal statement: correct with missing hypotheses; publication match: exact; verification: read the result; development: annotated; confidence: high}
\end{GSitem}
\begin{GSitem}{1979-467.6}{E}{Transitivity and conformal compactification}
\GSfield{Kind; source}{theorem; \GSpage{P000857}}
\GSfield{The notes}{\(U(2,2)\) preserves \(S_+,S_-,M\) and acts transitively on each; \(M\) is the conformal compactification of Minkowski space.}
\GSfield{Statement}{The group \(U(2,2)\) acts transitively on maximal positive, maximal negative and maximal isotropic complex two-planes. The last space is compact: in the positive/negative decomposition, each maximal isotropic plane is the graph of a unitary map, giving \(M\simeq U(2)\). Its dense Hermitian chart carries the Minkowski conformal class.}
\GSfield{Reconstruction}{Transitivity follows by extending adapted bases to Hermitian isometries. \(U(2)\) is \((S^1\times SU(2))/\{\pm1\}\), so the compactification has the usual finite-cover convention; no unqualified \(S^1\times S^3\) identification is imposed.}
\GSfield{Fate}{The compactification is preexisting twistor theory and is checked in \citet[Sec.~3.3.2]{woit2021euclidean}.}
\GSfield{Present status}{As of November 2026, the stated, qualified result is established; the notebook wording is interpreted only with the hypotheses specified above.}
\GSfield{Assessment}{literal statement: correct with missing hypotheses; publication match: exact; verification: read the result; development: annotated; confidence: high}
\end{GSitem}
\begin{GSitem}{1979-467.7}{E}{Fractional-linear action on matrix coordinates}
\GSfield{Kind; source}{construction; \GSpage{P000857}}
\GSfield{The notes}{For \(g=\left(\begin{smallmatrix}A&B\\C&D\end{smallmatrix}\right)\), the action is \(U\mapsto(AU+B)(CU+D)^{-1}\).}
\GSfield{Statement}{If a plane is represented by columns \(\binom{U}{I}\), block multiplication and renormalization give exactly \(g\cdot U=(AU+B)(CU+D)^{-1}\), on the open set where \(CU+D\) is invertible. Using columns \(\binom{I}{U}\) instead gives \((C+DU)(A+BU)^{-1}\).}
\GSfield{Reconstruction}{The earlier graph-basis notation appears to use the opposite ordering. Both formulas express the same Grassmannian action after a consistent chart convention; the rational map is not defined on the whole affine chart for every \(g\).}
\GSfield{Fate}{Elementary projective/Grassmannian linear algebra, directly checked by block multiplication; compare \citet[Sec.~3.3.2]{woit2021euclidean}.}
\GSfield{Present status}{As of November 2026, the stated, qualified result is established; the notebook wording is interpreted only with the hypotheses specified above.}
\GSfield{Assessment}{literal statement: correct with missing hypotheses; publication match: exact; verification: read the result; development: annotated; confidence: high}
\end{GSitem}
\begin{GSitem}{1979-467.8}{E}{The affine chart stabilizer includes dilations}
\GSfield{Kind; source}{claim; \GSpage{P000857}}
\GSfield{The notes}{The subgroup preserving \(G_0\), acting by \(U\mapsto aUa^*+b\), \(a\in GL_2(\mathbb C)\), \(b=b^*\), is said to be isomorphic to the Poincaré group.}
\GSfield{Statement}{The displayed affine action preserves the Minkowski quadratic form up to the constant factor \(|\det a|^2\), because \(\det(a\,dU\,a^*)=|\det a|^2\det(dU)\). It is a similarity group including positive dilations. Restricting \(|\det a|=1\), and quotienting the scalar \(U(1)\) kernel, gives the connected Lorentz part together with translations.}
\GSfield{Reconstruction}{For example \(a=2I\) sends \(U\) to \(4U\), which is not a Minkowski isometry. The covering relation between \(SL_2(\mathbb C)\) and the proper connected Lorentz group must also be distinguished from isomorphism.}
\GSfield{Fate}{The classical conformal action is reviewed in \citet[Sec.~3.3.2]{woit2021euclidean}; the counterexample follows immediately from the notebook formula.}
\GSfield{Present status}{As of November 2026, the stated, qualified result is established; the notebook wording is interpreted only with the hypotheses specified above.}
\GSfield{Assessment}{literal statement: false as written; publication match: partial; verification: read the result; development: annotated; confidence: high}
\end{GSitem}
\begin{GSitem}{1979-467.9}{E}{The determinant is the Minkowski quadratic form}
\GSfield{Kind; source}{claim; \GSpage{P000857-P000858}}
\GSfield{The notes}{\(\det(dU)\) is the metric; the Pauli matrix uses coordinates \(u_0,u_1,u_3,u_4\), but the following line uses \(u_0,u_1,u_2,u_3\) and ends with a cube.}
\GSfield{Statement}{For \(U=\left(\begin{smallmatrix}x_0-x_1&x_2+ix_3\\x_2-ix_3&x_0+x_1\end{smallmatrix}\right)\), \(\det(dU)=dx_0^2-dx_1^2-dx_2^2-dx_3^2\). This is a quadratic form of signature \((1,3)\), whose polarization defines the Minkowski metric.}
\GSfield{Reconstruction}{The relabelling of spatial coordinates must be made consistently. The final exponent \(3\), visible in the original PNG, is a faithful error and must be \(2\).}
\GSfield{Fate}{This is an elementary determinant calculation in the standard spinor realization \citep[Sec.~3.1]{woit2021euclidean}.}
\GSfield{Present status}{As of November 2026, the stated, qualified result is established; the notebook wording is interpreted only with the hypotheses specified above.}
\GSfield{Assessment}{literal statement: false as written; publication match: exact; verification: read the result; development: annotated; confidence: high}
\end{GSitem}
\begin{GSitem}{1979-467.10}{E}{Intersecting twistor lines and the light cone}
\GSfield{Kind; source}{theorem; \GSpage{P000858}}
\GSfield{The notes}{All projective lines meeting a fixed line \(\ell_0\) satisfy \(\det(U^0-U)=0\); on the real slice this is the light cone.}
\GSfield{Statement}{Two graph two-planes \(W_U,W_V\) meet nontrivially if and only if \(\ker(U-V)\ne0\), equivalently \(\det(U-V)=0\). For Hermitian \(U,V\), this means their Minkowski separation is null.}
\GSfield{Reconstruction}{The statement concerns the graph chart; the Grassmannian incidence condition itself extends globally. The cone includes its vertex \(U=V\), where the intersection dimension jumps from one to two.}
\GSfield{Fate}{Classical twistor incidence, also present in the earlier seminar occurrence 1978-455.1; checked in \citet[Sec.~3.3.2]{woit2021euclidean} and directly by solving the graph intersection equations.}
\GSfield{Present status}{As of November 2026, the stated, qualified result is established; the notebook wording is interpreted only with the hypotheses specified above.}
\GSfield{Assessment}{literal statement: correct; publication match: exact; verification: read the result; development: annotated; confidence: high}
\end{GSitem}
\begin{GSitem}{1979-467.11}{E}{Correcting the point/ray correspondence}
\GSfield{Kind; source}{claim; \GSpage{P000858}}
\GSfield{The notes}{One sentence claims a bijection between points of \(M\) and rays in \(N\); the following line correctly relates light rays on \(M\) to isotropic projective twistors.}
\GSfield{Statement}{A point of compactified Minkowski space corresponds to a maximal isotropic two-plane in \(T\), hence to a projective line \(\mathbb P^1\) of null twistors. A projective null twistor corresponds to the null geodesic formed by all maximal isotropic two-planes containing that line in \(T\).}
\GSfield{Reconstruction}{The first claimed bijection is dimensionally impossible: \(M\) has real dimension four, while the projective null-twistor space has dimension five. The incidence relation, with its different fibres, is the correct replacement.}
\GSfield{Fate}{The standard correspondence is explicit in \citet[Sec.~3.3.2]{woit2021euclidean}. The contradiction is already visible between the two adjacent notebook sentences.}
\GSfield{Present status}{As of November 2026, the stated, qualified result is established; the notebook wording is interpreted only with the hypotheses specified above.}
\GSfield{Assessment}{literal statement: false as written; publication match: exact; verification: read the result; development: annotated; confidence: high}
\end{GSitem}
\begin{GSitem}{1979-467.12}{E}{The status of curved twistor theory in 1979}
\GSfield{Kind; source}{claim; \GSpage{P000856}}
\GSfield{The notes}{Twistor theory in curved space has yet to be worked out, although papers already exist.}
\GSfield{Statement}{A general curved-spacetime extension is a broad programme, but a precise self-dual curved construction already existed before this lecture. It is therefore inaccurate to read the sentence as denying every curved twistor correspondence.}
\GSfield{Reconstruction}{The nonlinear-graviton construction concerns complex geometry with self-dual curvature under additional constraints; it is not a solution for all real Lorentzian Einstein metrics.}
\GSfield{Fate}{Penrose's \emph{Nonlinear gravitons and curved twistor theory} appeared in January 1976 \citep{penrose1976nonlinear}. Publisher metadata was opened; the full original theorem was not retrieved, so no detailed theorem formulation is assigned to that citation.}
\GSfield{Present status}{As of November 2026, the stated, qualified result is established; the notebook wording is interpreted only with the hypotheses specified above.}
\GSfield{Assessment}{literal statement: correct with missing hypotheses; publication match: partial; verification: publisher metadata checked; development: annotated; confidence: high}
\end{GSitem}
\begin{GSitem}{1979-467.13}{F}{A higher-dimensional twistor analogue}
\GSfield{Kind; source}{problem; \GSpage{P000858}}
\GSfield{The notes}{Contest: formulate a higher-dimensional generalization, with the hint line \(\rightsquigarrow\) flag.}
\GSfield{Statement}{The instruction specifies neither the ambient dimension, real form, flag type, field equation nor the incidence property to preserve. One direct matrix generalization uses maximal isotropic \(m\)-planes in \(\mathbb C^{m,m}\), a Hermitian \(m\times m\) chart and \(U(m,m)\) fractional transformations.}
\GSfield{Reconstruction}{That chart has real dimension \(m^2\), and \(\det(dU)\) has degree \(m\); it is a quadratic Lorentz form only when \(m=2\). Thus merely replacing two by \(m\) does not reproduce all four-dimensional features.}
\GSfield{Fate}{The elementary extension explains one possible direction, but does not constitute a unique solution of the contest. No exact publication can be matched to the unspecified desired analogue.}
\GSfield{Present status}{As of November 2026, this remains an open-ended programme in the notebook, not a uniquely determined open conjecture.}
\GSfield{Assessment}{literal statement: not determinable; publication match: none found; verification: unresolved lead; development: annotated; confidence: high}
\end{GSitem}

\GStalk{468}{1979-10-15}{F. A. Berezin}{P000859-P000862}{Grassmann algebra, supermatrices and classification problems}
The dated lecture begins on P000859 and ends with two explicit matrix problems at the top of P000862. Gelfand proposes graph analogues of the even Grassmann algebra; Gindikin then recalls the BGG description of projective bundles and asks for actual maps and a Grassmannian counterpart. The repeated-body-spectrum problem is separated from the odd-entry commutant problem. Finite coefficient algorithms solve the literal finite-dimensional computational tasks; this chapter examines what additional structural information they provide.

\begin{GSitem}{1979-468.1}{E}{Parity in a graded vector space}
\GSfield{Kind; source}{definition; \GSpage{P000859}}
\GSfield{The notes}{\(L=L_0\oplus L_1\); the notes call elements in either summand homogeneous and use parity \(\alpha(x)\in\mathbb Z_2\).}
\GSfield{Statement}{For \(\mathbb K=\mathbb R\) or \(\mathbb C\), a super vector space is a direct sum \(L_{\bar0}\oplus L_{\bar1}\). A nonzero homogeneous vector in \(L_{\bar i}\) has parity \(i\); signs in tensor products are determined modulo two.}
\GSfield{Reconstruction}{The zero vector belongs to both summands and has no uniquely determined parity. This harmless convention is made explicit.}
\GSfield{Fate}{This is the preexisting language of superlinear algebra; compare \citet[Part I, pp.~90--130]{berezin1987superanalysis}. Publisher contents, rather than the full chapter, were checked.}
\GSfield{Present status}{As of November 2026, the stated, qualified result is established; the notebook wording is interpreted only with the hypotheses specified above.}
\GSfield{Assessment}{literal statement: correct with missing hypotheses; publication match: exact; verification: publisher metadata checked; development: annotated; confidence: high}
\end{GSitem}
\begin{GSitem}{1979-468.2}{E}{The exterior algebra and its monomial basis}
\GSfield{Kind; source}{definition; \GSpage{P000859-P000860}}
\GSfield{The notes}{\(\Lambda\) has generators \(1,\xi_1,\ldots,\xi_N\), relations \(\xi_i\xi_j+\xi_j\xi_i=0\), ordered monomial basis, dimension \(2^N\), and \(fg=(-1)^{\alpha(f)\alpha(g)}gf\).}
\GSfield{Statement}{Over a field of characteristic different from two, \(\Lambda=\bigwedge_{\mathbb K}\mathbb K^N\), with basis \(\xi_I=\xi_{i_1}\cdots\xi_{i_k}\) for ordered subsets \(I\), including \(\xi_\varnothing=1\). Its even and odd summands have dimensions \(2^{N-1}\) when \(N\ge1\); homogeneous elements supercommute.}
\GSfield{Reconstruction}{The restriction on characteristic is supplied; it holds for the two fields named in the notes.}
\GSfield{Fate}{Classical exterior algebra, already used throughout the earlier Berezin lectures and the 1978 BGG construction \citep[Sec.~2]{bernsteingelfandgelfand1978bundles}.}
\GSfield{Present status}{As of November 2026, the stated, qualified result is established; the notebook wording is interpreted only with the hypotheses specified above.}
\GSfield{Assessment}{literal statement: correct; publication match: exact; verification: read the result; development: annotated; confidence: high}
\end{GSitem}
\begin{GSitem}{1979-468.3}{E}{The even Grassmann envelope}
\GSfield{Kind; source}{construction; \GSpage{P000860}}
\GSfield{The notes}{The initial \(L\otimes\Lambda\) definition is cancelled. The replacement is \(L(\Lambda)=\{\sum a_ie_i:\alpha(a_i)=\alpha(e_i)\}=(\Lambda\otimes L)_{\bar0}\), a left \(\Lambda_0\)-module.}
\GSfield{Statement}{For finite graded \(L\), \(L(\Lambda)=(\Lambda_{\bar0}\otimes L_{\bar0})\oplus(\Lambda_{\bar1}\otimes L_{\bar1})\). This is the even part of the graded tensor product, and is closed under multiplication by \(\Lambda_{\bar0}\).}
\GSfield{Reconstruction}{The cancellation is substantive: the envelope is not the entire tensor product. Tensor multiplication for an algebra \(L\) requires the Koszul sign, or the equivalent coefficient-and-basis convention adopted in the lecture.}
\GSfield{Fate}{The standard functor-of-Grassmann-points construction predates this talk; see \citet[Part I, linear algebra]{berezin1987superanalysis}, and the preceding seminar occurrences 1976-406.5--9.}
\GSfield{Present status}{As of November 2026, the stated, qualified result is established; the notebook wording is interpreted only with the hypotheses specified above.}
\GSfield{Assessment}{literal statement: correct; publication match: exact; verification: publisher metadata checked; development: annotated; confidence: high}
\end{GSitem}
\begin{GSitem}{1979-468.4}{F}{Graph analogues of the even Grassmann algebra}
\GSfield{Kind; source}{conjecture; \GSpage{P000860}; I. M. Gelfand}
\GSfield{The notes}{Gelfand proposes an analogue for every Dynkin diagram, with the ordinary example attached to \(A_n\); connected generators have product zero, disconnected ones a new product. He cites Functional Analysis 1979, no.~3.}
\GSfield{Statement}{The recorded prescription is not a complete presentation: it omits whether generators commute, square to zero, obey further path relations, and what universal property the analogue must satisfy. It therefore does not determine a unique theorem or conjecture.}
\GSfield{Reconstruction}{The exact cited paper can be identified as Gelfand--Ponomarev's model algebras, but identifying a citation does not identify the additional proposed Grassmann analogue.}
\GSfield{Fate}{\citet[pp.~157--166]{gelfandponomarev1979model} is the exact volume/issue match, received 23 March 1979. Its graph-algebra setting supports the reference; no exact theorem equating the notebook prescription to its model algebras was verified.}
\GSfield{Present status}{As of November 2026, the all-Dynkin assertion remains underdetermined by these notes. It is not entered as an open problem.}
\GSfield{Assessment}{literal statement: not determinable; publication match: thematic only; verification: indirect; development: annotated; confidence: medium}
\end{GSitem}
\begin{GSitem}{1979-468.5}{E}{Even block supermatrices}
\GSfield{Kind; source}{construction; \GSpage{P000860-P000861}}
\GSfield{The notes}{\(L_{p,q}(\Lambda)\) has even entries in \(A_{11},A_{22}\), odd entries in \(A_{12},A_{21}\), and is an algebra.}
\GSfield{Statement}{Fix a homogeneous basis of \(\mathbb K^{p|q}\). The even endomorphisms over \(\Lambda\) are the matrices \(A=\left(\begin{smallmatrix}P&Q\\R&S\end{smallmatrix}\right)\) with \(P,S\) even and \(Q,R\) odd. With the lecture's left-module convention their composition is ordinary ordered matrix multiplication; parity is preserved.}
\GSfield{Reconstruction}{An even supermatrix need not have all entries even. This distinction is essential: Problem 2 later asks about all-odd-entry matrices, not general odd endomorphisms of a \(p|q\) module.}
\GSfield{Fate}{This is standard superlinear algebra; \citet[Part I, pp.~90--130]{berezin1987superanalysis} is a later systematic exposition, not evidence of a new 1979 theorem.}
\GSfield{Present status}{As of November 2026, the stated, qualified result is established; the notebook wording is interpreted only with the hypotheses specified above.}
\GSfield{Assessment}{literal statement: correct; publication match: exact; verification: publisher metadata checked; development: annotated; confidence: high}
\end{GSitem}
\begin{GSitem}{1979-468.6}{E}{Cyclicity of the supertrace}
\GSfield{Kind; source}{theorem; \GSpage{P000861}}
\GSfield{The notes}{\(\operatorname{str}A=\operatorname{tr}A_{11}-\operatorname{tr}A_{22}\) and \(\operatorname{str}(AB)=\operatorname{str}(BA)\).}
\GSfield{Statement}{For even supermatrices \(A,B\) over a supercommutative algebra, \(\operatorname{str}(AB)=\operatorname{str}(BA)\). For homogeneous superendomorphisms of arbitrary parity the correct general relation is \(\operatorname{str}(AB)=(-1)^{|A||B|}\operatorname{str}(BA)\).}
\GSfield{Reconstruction}{The even-envelope domain is retained; deleting it would make the unsigned cyclicity assertion false. Off-diagonal odd products change sign on interchange, cancelling the minus in the lower-block trace.}
\GSfield{Fate}{Already part of the seminar's 1976-406.9 exposition and of Berezin's superlinear algebra \citep[Part I]{berezin1987superanalysis}.}
\GSfield{Present status}{As of November 2026, the stated, qualified result is established; the notebook wording is interpreted only with the hypotheses specified above.}
\GSfield{Assessment}{literal statement: correct; publication match: exact; verification: publisher metadata checked; development: annotated; confidence: high}
\end{GSitem}
\begin{GSitem}{1979-468.7}{E}{The Berezinian and its multiplicativity}
\GSfield{Kind; source}{definition; \GSpage{P000861}}
\GSfield{The notes}{\(\operatorname{sdet}A=\det(P-QS^{-1}R)/\det S\) and \(\operatorname{sdet}(AB)=\operatorname{sdet}A\,\operatorname{sdet}B\).}
\GSfield{Statement}{For an invertible even matrix over a finite Grassmann algebra, the bodies of \(P,S\) are invertible, hence so are \(P,S\). Define \(\operatorname{Ber}A=\det(P-QS^{-1}R)(\det S)^{-1}\in\Lambda_{\bar0}^{\times}\). It is multiplicative.}
\GSfield{Reconstruction}{Invertibility, omitted at the definition, is necessary. This rational invariant differs from the ordinary determinant even when the matrix is block diagonal: it is \(\det P/\det S\).}
\GSfield{Fate}{The Berezinian already occurs in \citet[formula~(3)]{berezinleites1975supermanifolds}; it recurs here after 1976-406.8--10. The primary metadata was rechecked and the displayed formula was checked algebraically against the notebook.}
\GSfield{Present status}{As of November 2026, the stated, qualified result is established; the notebook wording is interpreted only with the hypotheses specified above.}
\GSfield{Assessment}{literal statement: correct with missing hypotheses; publication match: exact; verification: indirect; development: annotated; confidence: high}
\end{GSitem}
\begin{GSitem}{1979-468.8}{E}{The 1|1 multiplicativity check}
\GSfield{Kind; source}{claim; \GSpage{P000861}}
\GSfield{The notes}{For \(A=\left(\begin{smallmatrix}1&\xi_1\\\xi_2&1\end{smallmatrix}\right)\), the notes calculate \(\operatorname{Ber}A=1-\xi_1\xi_2\), \(A^2\), and \(\operatorname{Ber}(A^2)=1-2\xi_1\xi_2=(\operatorname{Ber}A)^2\).}
\GSfield{Statement}{For odd generators \(\xi_1,\xi_2\) in characteristic zero, put \(t=\xi_1\xi_2\), so \(t^2=0\). Then \(A^2=\left(\begin{smallmatrix}1+t&2\xi_1\\2\xi_2&1-t\end{smallmatrix}\right)\), and the Schur-complement formula gives \((1-3t)(1+t)=1-2t\).}
\GSfield{Reconstruction}{This verifies the particular example. It does not by itself prove multiplicativity for every \(1|1\) matrix.}
\GSfield{Fate}{A direct elementary instance of the preexisting Berezinian identity \citep[formula~(3)]{berezinleites1975supermanifolds}; the arithmetic is independently checked here.}
\GSfield{Present status}{As of November 2026, the stated, qualified result is established; the notebook wording is interpreted only with the hypotheses specified above.}
\GSfield{Assessment}{literal statement: correct; publication match: exact; verification: indirect; development: annotated; confidence: high}
\end{GSitem}
\begin{GSitem}{1979-468.9}{E}{The body map and invertibility}
\GSfield{Kind; source}{construction; \GSpage{P000861}}
\GSfield{The notes}{The map \(m\) is evaluation at \(\xi=0\); the body of an even supermatrix is block diagonal.}
\GSfield{Statement}{The augmentation \(m:\Lambda\to\mathbb K\) has nilpotent kernel \(J=(\xi_1,\ldots,\xi_N)\). An even supermatrix is invertible if and only if its body is invertible. Indeed, after factoring the body matrix, the remaining inverse is the finite Neumann series for \(1+K\), with \(K\in M_{p+q}(J)\).}
\GSfield{Reconstruction}{The invertibility consequence is reconstructed to clarify the allowed change-of-basis group in the following theorem and problems.}
\GSfield{Fate}{The body construction is standard in Berezin's superanalysis \citep[Part I, linear algebra]{berezin1987superanalysis}; the finite-series argument is given explicitly here.}
\GSfield{Present status}{As of November 2026, the stated, qualified result is established; the notebook wording is interpreted only with the hypotheses specified above.}
\GSfield{Assessment}{literal statement: correct; publication match: exact; verification: publisher metadata checked; development: annotated; confidence: high}
\end{GSitem}
\begin{GSitem}{1979-468.10}{E}{Diagonalization when body eigenvalues are distinct}
\GSfield{Kind; source}{theorem; \GSpage{P000861-P000862}}
\GSfield{The notes}{Over \(\mathbb C\), if \(m(A)\) has distinct eigenvalues, then an invertible \(B\) makes \(BAB^{-1}\) diagonal.}
\GSfield{Statement}{For a finite Grassmann algebra over \(\mathbb C\), let \(A\) be even and suppose all \(p+q\) eigenvalues of \(m(A)\), across both parity blocks, are distinct. There is an invertible even \(B\) with \(BAB^{-1}=\operatorname{diag}(c_1,\ldots,c_{p+q})\), where \(c_i\in\Lambda_{\bar0}\) has the corresponding eigenvalue as body.}
\GSfield{Reconstruction}{Diagonalize the two body blocks first. At each positive Grassmann degree, remove the off-diagonal \((i,j)\) coefficient by division by the nonzero scalar \(\lambda_i-\lambda_j\); the nilpotent filtration terminates. This supplies a finite proof and specifies the parity of the conjugator.}
\GSfield{Fate}{The statement belongs to standard supermatrix perturbation theory. The later Berezin book identifies the subject, but its full theorem was not retrieved; the proof here verifies the result independently \citep[Part I]{berezin1987superanalysis}.}
\GSfield{Present status}{As of November 2026, the stated, qualified result is established; the notebook wording is interpreted only with the hypotheses specified above.}
\GSfield{Assessment}{literal statement: correct; publication match: exact; verification: publisher metadata checked; development: annotated; confidence: high}
\end{GSitem}
\begin{GSitem}{1979-468.11}{B}{Canonical forms beyond simple body spectrum}
\GSfield{Kind; source}{problem; \GSpage{P000862}}
\GSfield{The notes}{\ru{Аналоги жорд. клеток здесь неизвестны. Задача 1. Сделать это (аналог норм. форм)!}}
\GSfield{Statement}{For fixed finite \(N,p,q\), classify even matrices \(A\in L_{p,q}(\Lambda_N)\) up to conjugation by invertible even matrices, without requiring distinct eigenvalues of the body. An effective canonical representative is one precise interpretation of an analogue of Jordan normal form.}
\GSfield{Reconstruction}{The notebook does not require a short list of blocks. The chapter separates existence of a canonical algorithm from manageable block descriptions, and does not confuse wildness with nonexistence of canonical forms.}
\GSfield{Fate}{\citet[Theorems~1.1 and~1.3]{sergeichuk2000canonical} gives canonical algorithms for linear matrix problems: a verified Bielefeld preprint appeared in 1999, the journal paper in 2000, and the repository deposit in 2007. The chapter supplies the application through the faithful regular representation of the finite-dimensional even matrix algebra; the application itself is a reconstruction, not an explicit supermatrix theorem claimed in that paper.}
\GSfield{Present status}{As of November 2026, the finite-\(N\) effective canonical-form interpretation is resolved by that reduction. The chapter also gives a small-\(N\) normal form and an embedding of simultaneous similarity, explaining why a Jordan-block-style classification is a stronger request.}
\GSfield{Assessment}{literal statement: not determinable; publication match: partial; verification: read the result; development: developed in this chapter; confidence: high; problem \hyperref[prob:GP-1979-01]{GP-1979-01}}
\end{GSitem}
\begin{GSitem}{1979-468.12}{E}{Even commutants of all-odd-entry matrices}
\GSfield{Kind; source}{problem; \GSpage{P000862}}
\GSfield{The notes}{Given \(A=(a_{ik})\) with every \(a_{ik}\in\Lambda_1\), describe all \(B=(b_{ik})\) with every \(b_{ik}\in\Lambda_0\) and \(AB=BA\). Four generators are related to Yang--Mills; independent entries are said to be tractable.}
\GSfield{Statement}{For fixed \(n,N\), a specified \(A\in M_n(\Lambda_N^{\bar1})\), and a field of characteristic zero, determine the even subalgebra \(C_{\bar0}(A)=\{B\in M_n(\Lambda_N^{\bar0}):AB=BA\}\), including dimensions, bases and multiplication if desired.}
\GSfield{Reconstruction}{Expand in the ordered monomial basis. This is the kernel of a finite \(\mathbb K\)-linear map \(B\mapsto AB-BA\), so the literal computation always has a terminating linear-algebra algorithm. A uniform structural description as \(A\) varies is a stronger question.}
\GSfield{Fate}{The literal finite problem is solved by classical coefficient linear algebra: expand \(A,B\) in the ordered Grassmann monomials and take the kernel of \(B\mapsto AB-BA\). No specific later publication is needed for existence of this terminating procedure. Searches did not identify the precise historical structural description Berezin may have wanted. The chapter proves a complete bounded \(2\times2,\Lambda_4\) parametrization including cubic entries; no claim of historical novelty is made.}
\GSfield{Present status}{As of November 2026, the literal finite computational task is resolved. A uniform structural normal form as \(A,n,N\) vary is a stronger interpretation, not an additional requirement specified in the manuscript, and is not labelled an open problem here.}
\GSfield{Assessment}{literal statement: correct; publication match: none found; verification: read the result; development: developed in this chapter; confidence: medium; problem \hyperref[prob:GP-1979-02]{GP-1979-02}}
\end{GSitem}
\begin{GSitem}{1979-468.13}{E}{Stable exterior modules and projective vector bundles}
\GSfield{Kind; source}{theorem; \GSpage{P000862}; S. G. Gindikin}
\GSfield{The notes}{Gindikin associates to a graded \(\Lambda(4)\)-module the differential \(a_1\xi_1+\cdots+a_4\xi_4\), whose cohomology is concentrated at zero, and says such modules modulo trivial ones classify instantons.}
\GSfield{Statement}{For \(E=\mathbb C^4\), a finite graded \(\bigwedge E\)-module \(V\) gives a linear complex of bundles on \(\mathbb P(E)\) with differential \(\sum a_i\xi_i\). Modules whose fibre complexes have homology only in degree zero correspond to algebraic vector bundles, up to the stable equivalence generated by free exterior modules.}
\GSfield{Reconstruction}{The anticommutation relations imply \(d^2=0\). A general module need not have the required fibre homology. Calling every resulting bundle an instanton drops the real structure and triviality on appropriate twistor lines required for the gauge-theoretic correspondence.}
\GSfield{Fate}{\citet[Secs.~3--5, Theorems~1--3]{bernsteingelfandgelfand1978bundles} already states the bundle correspondence, received 10 January 1978 and published in the 1978 volume. The recurrence is exact for bundles; the unqualified instanton identification is not. Compare 1977-437.6--12 and 1977-438.2--5.}
\GSfield{Present status}{As of November 2026, the stated, qualified result is established; the notebook wording is interpreted only with the hypotheses specified above.}
\GSfield{Assessment}{literal statement: correct with missing hypotheses; publication match: exact; verification: read the result; development: annotated; confidence: high; problem \hyperref[prob:GP-1977-11]{GP-1977-11}}
\end{GSitem}
\begin{GSitem}{1979-468.14}{B}{Construct the maps, not only the cohomology spaces}
\GSfield{Kind; source}{problem; \GSpage{P000862}; S. G. Gindikin}
\GSfield{The notes}{\ru{дать явную конструкцию. Не сделано (строятся пр-ва, но не отображения)}.}
\GSfield{Statement}{Given a coherent sheaf on \(\mathbb P^3\) by finite algebraic data, construct the differentials in a finite monad representing it, or the corresponding exterior-algebra resolution, rather than only the dimensions of its terms.}
\GSfield{Reconstruction}{This makes precise the surrounding BGG/Beilinson context; the notes do not specify an input encoding or require canonical choices of bases.}
\GSfield{Fate}{\citet[Theorems~4.1 and~6.1; Sec.~5]{eisenbudfloystadschreyer2003exterior} constructs the Tate resolution and the maps of the Beilinson monad. The joint preprint appeared 20 April 2001, checked version 8 June 2001; journal publication was 2003. This answers the explicit-morphism request, although it is not the earliest possible algorithmic implementation.}
\GSfield{Present status}{As of November 2026, the stated, qualified result is established; the notebook wording is interpreted only with the hypotheses specified above.}
\GSfield{Assessment}{literal statement: not determinable; publication match: exact; verification: read the result; development: annotated; confidence: high; problem \hyperref[prob:GP-1979-03]{GP-1979-03}}
\end{GSitem}
\begin{GSitem}{1979-468.15}{B}{A linear-algebra description for Grassmannians}
\GSfield{Kind; source}{problem; \GSpage{P000862}; S. G. Gindikin}
\GSfield{The notes}{\ru{сделать классиф. для грассманова мн-зия} follows the projective-space construction.}
\GSfield{Statement}{For a complex Grassmannian \(\operatorname{Gr}(k,n)\), seek a finite algebraic model for its bounded derived category of coherent sheaves, with complexes of explicit homogeneous bundles and morphisms, analogous to Beilinson's projective-space model.}
\GSfield{Reconstruction}{The derived-category formulation is the precise reconstruction of the requested analogue. It is not an assertion that all vector bundles admit a finite numerical list.}
\GSfield{Fate}{\citet[Theorem~3.4 and Secs.~3.1--3.5]{kapranov1988homogeneous} supplies an equivalence using Schur functors of the tautological bundle. The Grassmannian result appeared first in Russian in 1984 (English translation 1985); the checked 1988 paper gives the fuller framework. Different differentials may represent the same homotopy class.}
\GSfield{Present status}{As of November 2026, the stated, qualified result is established; the notebook wording is interpreted only with the hypotheses specified above.}
\GSfield{Assessment}{literal statement: not determinable; publication match: partial; verification: read the result; development: annotated; confidence: high; problem \hyperref[prob:GP-1979-04]{GP-1979-04}}
\end{GSitem}

\GStalk{469}{1979-10-15}{G. B. Shabat}{P000862-P000863}{Double translation hypersurfaces and the Lie--Wirtinger theorem}
The contribution starts at the bottom of P000862 and ends midway down P000863. The catalogue names G. B. Shabat; the manuscript initial looks like \ru{Ю.}, a retained attribution discrepancy. The discussion explains the theta-divisor example and reports the classical converse, followed by Gelfand's broad integrability speculation. Genericity and the distinction between smooth Jacobians, generalized Jacobians and degenerations are indispensable.

\begin{GSitem}{1979-469.1}{E}{Two distinct translation parametrizations}
\GSfield{Kind; source}{definition; \GSpage{P000862}}
\GSfield{The notes}{A hypersurface germ \(\Gamma\subset\mathbb C^n\) is written both as \(\gamma_1+\cdots+\gamma_{n-1}\) and as \(\delta_1+\cdots+\delta_{n-1}\); permutations and translations are excluded.}
\GSfield{Statement}{A translation parametrization of a smooth hypersurface germ is a local biholomorphism \((t_1,\ldots,t_{n-1})\mapsto\sum_i\gamma_i(t_i)\) onto the germ. Distinctness up to permutations and constant translations is a minimal condition; the classical converse needs stronger local nondegeneracy.}
\GSfield{Reconstruction}{Without the immersion and local uniqueness conditions, arbitrary sums of curve germs can have smaller dimension or singularities.}
\GSfield{Fate}{The classical subject and its precise genericity conditions are reviewed in \citet[Definitions~2.1--2.2]{little1983translation}.}
\GSfield{Present status}{As of November 2026, the stated, qualified result is established; the notebook wording is interpreted only with the hypotheses specified above.}
\GSfield{Assessment}{literal statement: correct with missing hypotheses; publication match: exact; verification: read the result; development: annotated; confidence: high}
\end{GSitem}
\begin{GSitem}{1979-469.2}{E}{Theta divisors as sums of Abel curves}
\GSfield{Kind; source}{construction; \GSpage{P000862-P000863}}
\GSfield{The notes}{For a smooth curve \(X\) of genus \(n\), the notes write \(\Theta=jX+\cdots+jX\), with \(n-1\) summands, and \(\Theta=K-\Theta\).}
\GSfield{Statement}{Let \(X\) be a smooth projective complex curve of genus \(g\ge2\), \(j:X\to J(X)\) an Abel map based at \(p_0\), and \(W_{g-1}\) the image of effective divisors of degree \(g-1\). Then \(W_{g-1}\) is a translate of the theta divisor and is invariant under \(x\mapsto [K_X-(2g-2)p_0]-x\).}
\GSfield{Reconstruction}{The symmetry follows from Riemann--Roch: an effective divisor \(D\) of degree \(g-1\) has \(h^0(K_X-D)=h^0(D)\). At regular points, two complementary effective divisors give the two local sums.}
\GSfield{Fate}{This is preexisting Riemann theta-divisor theory; \citet[Example~4.1]{little1983translation} gives the exact translation-manifold example.}
\GSfield{Present status}{As of November 2026, the stated, qualified result is established; the notebook wording is interpreted only with the hypotheses specified above.}
\GSfield{Assessment}{literal statement: correct; publication match: exact; verification: read the result; development: annotated; confidence: high}
\end{GSitem}
\begin{GSitem}{1979-469.3}{E}{When the two Abel curves are distinct}
\GSfield{Kind; source}{claim; \GSpage{P000863}}
\GSfield{The notes}{\(X\ne-jX\) is said to be equivalent to nonhyperellipticity.}
\GSfield{Statement}{For a smooth complex curve of genus \(g\ge2\), \(j(X)\) is a translate of \(-j(X)\) if and only if \(X\) is hyperelliptic. In the hyperelliptic case the two theta-divisor translation decompositions arising from inversion coincide.}
\GSfield{Reconstruction}{The left-hand \(X\) is shorthand for its Abel image, and equality must be considered up to translation. For the converse, an inversion symmetry yields a fixed degree-two divisor class \(p+\iota(p)\), hence a degree-two map to \(\mathbb P^1\).}
\GSfield{Fate}{Classical Abel-curve geometry; the hyperelliptic exception is explicit in \citet[Sec.~4, following Example~4.1]{little1983translation}.}
\GSfield{Present status}{As of November 2026, the stated, qualified result is established; the notebook wording is interpreted only with the hypotheses specified above.}
\GSfield{Assessment}{literal statement: correct with missing hypotheses; publication match: exact; verification: read the result; development: annotated; confidence: high}
\end{GSitem}
\begin{GSitem}{1979-469.4}{E}{The Lie--Wirtinger converse with nondegeneracy}
\GSfield{Kind; source}{theorem; \GSpage{P000863}}
\GSfield{The notes}{All double translation surfaces are said to be theta divisors and their degenerations; the parenthesis claims that a small deformation is a theta-divisor germ.}
\GSfield{Statement}{Take a smooth non-linear hypersurface germ in \(\mathbb C^n\) with two locally invertible sum parametrizations. Assume their \(2n-2\) tangent lines at the base point are distinct, every generating curve has acceleration transverse to the tangent hyperplane, and varying at least one first-system parameter while fixing the others changes every second-system parameter. Then the germ is a translate of the lift of an Abel-divisor locus in the generalized Jacobian of a reduced, connected canonical curve of degree \(2n-2\) and arithmetic genus \(n\).}
\GSfield{Reconstruction}{The first two extra assumptions are Little's Definition~2.2; the last is condition (*) in Proposition~3.7. Singular or reducible curves give generalized Jacobians. Smoothability of every possible degeneration is not asserted here.}
\GSfield{Fate}{Lie--Wirtinger is classical, not a new 1979 assertion. \citet[Corollary~3.5, Proposition~3.7, Theorem~3.9]{little1983translation} supplies a precise later proof and distinguishes the missing conditions. Its Theorem~5.1 recovers smooth nonhyperelliptic Jacobians when the ambient variety is principally polarized.}
\GSfield{Present status}{As of November 2026, the stated, qualified result is established; the notebook wording is interpreted only with the hypotheses specified above.}
\GSfield{Assessment}{literal statement: correct with missing hypotheses; publication match: partial; verification: read the result; development: annotated; confidence: high}
\end{GSitem}
\begin{GSitem}{1979-469.5}{F}{A universal explanation of integrability by translation geometry}
\GSfield{Kind; source}{conjecture; \GSpage{P000863}; I. M. Gelfand}
\GSfield{The notes}{Gelfand conjectures that all integrable systems have this theorem as their basis, and that it is the first of a large family of such theorems.}
\GSfield{Statement}{No class of systems, notion of integrability, construction of an associated hypersurface, or meaning of being a basis is specified. There is no uniquely recoverable quantified conjecture.}
\GSfield{Reconstruction}{The appearance of theta functions in algebraically integrable systems motivates the comment, but cannot turn a universal informal analogy into a theorem.}
\GSfield{Fate}{The checked Lie--Wirtinger source concerns translation hypersurfaces, not all integrable systems \citep{little1983translation}. Its theorem therefore supplies no exact publication match for this universal programme.}
\GSfield{Present status}{As of November 2026, the recorded programme remains underdetermined; it is not labelled an open theorem.}
\GSfield{Assessment}{literal statement: not determinable; publication match: thematic only; verification: read the result; development: annotated; confidence: high}
\end{GSitem}

\GStalk{470}{1979-10-15}{Raphael Høegh-Krohn}{P000863-P000865}{Energy representations of current groups and Brownian paths}
The actual talk occupies the lower part of P000863 and P000864--P000865, translated by Gelfand. It constructs the energy representation, states irreducibility in dimension at least three, asks about dimension two and explains the one-dimensional Brownian realization. The name bears a question mark in the notebook but is identified by the catalogue. The construction needs a Gaussian realization outside the Hilbert space and a consistent Maurer--Cartan convention. A later claimed proof in every dimension cannot settle the two-dimensional question: the primary literature explicitly identifies a gap.

\begin{GSitem}{1979-470.1}{E}{The current group and its one-form Hilbert space}
\GSfield{Kind; source}{construction; \GSpage{P000863-P000864}}
\GSfield{The notes}{\(G^X=C^\infty_0(X,G)\), with \(G\) compact simple; the one-form scalar product uses the Riemannian metric and the Killing form.}
\GSfield{Statement}{For a smooth Riemannian manifold \(X\) and compact semisimple \(G\), let \(\mathcal G=C^\infty_c(X,G)\), meaning maps equal to the identity off a compact set. Complete \(\Omega^1_c(X,\mathfrak g)\) for \(\int_X \langle\omega_1,\omega_2\rangle_{g,B}\,d\mu\), with \(B\) an \(\operatorname{Ad}G\)-invariant positive inner product, for example minus the Killing form.}
\GSfield{Reconstruction}{The Killing form itself is negative definite on a compact semisimple real Lie algebra. The sign is material to the Hilbert-space construction. The trace can be taken on either tangent or Lie-algebra endomorphisms because \(\operatorname{tr}(AB)=\operatorname{tr}(BA)\).}
\GSfield{Fate}{This setting predates the lecture in \citet[Secs.~2--3]{albeveriohoeghkrohn1978energy}; the general geometric formulation is explicit in \citet[Sec.~6]{wallach1987gauge}.}
\GSfield{Present status}{As of November 2026, the stated, qualified result is established; the notebook wording is interpreted only with the hypotheses specified above.}
\GSfield{Assessment}{literal statement: correct with missing hypotheses; publication match: exact; verification: read the result; development: annotated; confidence: high}
\end{GSitem}
\begin{GSitem}{1979-470.2}{E}{The infinite-dimensional Gaussian realization}
\GSfield{Kind; source}{construction; \GSpage{P000864}}
\GSfield{The notes}{The notation is \(L^2(\mathcal H,\mu)\), with a Gaussian analogous to \(e^{-|x|^2/2}dx\), invariant under rotations and quasi-invariant under translations.}
\GSfield{Statement}{On an infinite-dimensional separable real Hilbert space \(\mathcal H\), covariance \(I\) does not define a countably additive Gaussian Borel probability on \(\mathcal H\). Use an isonormal Gaussian process on a probability space, a nuclear dual containing \(\mathcal H\), or the symmetric Fock space over its complexification. Cameron--Martin translations are the vectors of \(\mathcal H\); abstract second quantization implements every orthogonal operator.}
\GSfield{Reconstruction}{If independent coordinates had law \(N(0,1)\), their squared sum would diverge almost surely, so a sample would not lie in \(\ell^2\). An arbitrary orthogonal map need not act pointwise on a chosen Banach realization; its Gaussian unitary is nevertheless defined.}
\GSfield{Fate}{\citet[Secs.~1 and~5]{wallach1987gauge} uses distribution spaces and Gaussian Hilbert spaces; \citet[Sec.~5]{albeverioetal2015brownian} gives the Fock/white-noise construction. The literal Hilbert-space probability notation is false.}
\GSfield{Present status}{As of November 2026, the stated, qualified result is established; the notebook wording is interpreted only with the hypotheses specified above.}
\GSfield{Assessment}{literal statement: false as written; publication match: exact; verification: read the result; development: annotated; confidence: high}
\end{GSitem}
\begin{GSitem}{1979-470.3}{E}{The Maurer--Cartan cocycle and affine action}
\GSfield{Kind; source}{claim; \GSpage{P000864}}
\GSfield{The notes}{\(\alpha(\varphi)\omega=\operatorname{Ad}\varphi\,\omega+\varphi^{-1}d\varphi\) is declared a group homomorphism.}
\GSfield{Statement}{With \(\pi(\varphi)\omega=\operatorname{Ad}\varphi\,\omega\), the correct left cocycle is \(\beta(\varphi)=d\varphi\,\varphi^{-1}\), since \(\beta(\varphi\psi)=\beta(\varphi)+\pi(\varphi)\beta(\psi)\). Thus \(\alpha(\varphi)\omega=\pi(\varphi)\omega+\beta(\varphi)\) is a left affine action.}
\GSfield{Reconstruction}{The displayed left Maurer--Cartan form instead obeys \((\varphi\psi)^{-1}d(\varphi\psi)=\operatorname{Ad}\psi^{-1}(\varphi^{-1}d\varphi)+\psi^{-1}d\psi\). It cannot be paired with \(\operatorname{Ad}\varphi\) in the stated left-action convention.}
\GSfield{Fate}{The corrected cocycle appears in the energy construction \citep[Sec.~5]{wallach1987gauge}. Direct differentiation of a product independently verifies the correction; the original PNG confirms the manuscript order, so this is a faithful mathematical error.}
\GSfield{Present status}{As of November 2026, the stated, qualified result is established; the notebook wording is interpreted only with the hypotheses specified above.}
\GSfield{Assessment}{literal statement: false as written; publication match: exact; verification: read the result; development: annotated; confidence: high}
\end{GSitem}
\begin{GSitem}{1979-470.4}{E}{The energy representation}
\GSfield{Kind; source}{construction; \GSpage{P000864}}
\GSfield{The notes}{The notes form \(V_\varphi=U_{\alpha(\varphi)}\) from the Gaussian representation of the Euclidean group of \(\mathcal H\).}
\GSfield{Statement}{For a real orthogonal representation \(\pi\) and continuous one-cocycle \(\beta\), the bosonic Fock representation satisfies \(V_g e(h)=\exp(-\|\beta(g)\|^2/2-\langle\pi(g)h,\beta(g)\rangle)e(\pi(g)h+\beta(g))\), with exponential vectors normalized by \(\langle e(h),e(k)\rangle=\exp\langle h,k\rangle\). Apply this to the current-group action and corrected Maurer--Cartan cocycle.}
\GSfield{Reconstruction}{The cocycle identity proves the representation law and the exponential-vector inner product proves unitarity. Inner-product conventions may conjugate the displayed scalar in a complex realization; the real Gaussian model fixes it unambiguously.}
\GSfield{Fate}{The energy representation was already constructed in \citet[Sec.~3]{albeveriohoeghkrohn1978energy}; its later Fock formulation is checked in \citet[Sec.~5]{albeverioetal2015brownian}.}
\GSfield{Present status}{As of November 2026, the stated, qualified result is established; the notebook wording is interpreted only with the hypotheses specified above.}
\GSfield{Assessment}{literal statement: correct with missing hypotheses; publication match: exact; verification: read the result; development: annotated; confidence: high}
\end{GSitem}
\begin{GSitem}{1979-470.5}{A}{Irreducibility in dimension at least three}
\GSfield{Kind; source}{theorem; \GSpage{P000864}}
\GSfield{The notes}{Theorem~1: \(V\) is irreducible for \(\dim X\ge3\), attributed to Albeverio and Høegh-Krohn after earlier bounds five and four.}
\GSfield{Statement}{For connected, orientable, paracompact smooth \(X\) of dimension at least three, compact semisimple \(G\), and the current-group energy representation from a Riemannian metric, smooth positive density and invariant positive Lie-algebra inner product, \(V\) is irreducible.}
\GSfield{Reconstruction}{The lecture omits topological and measure hypotheses. Compactly supported currents and the full current group, rather than an arbitrary proper subgroup, are required.}
\GSfield{Fate}{The soon-after joint paper \citet{albeveriohoeghkrohntestard1981irreducibility} includes Testard. Its result and earlier restrictions are documented and extended to this generality by the fully read \citet[Theorem~6.2]{wallach1987gauge}. The notebook is an October 1979 occurrence; a November 1979 conference presentation is not treated as the first circulation.}
\GSfield{Present status}{As of November 2026, the stated, qualified result is established; the notebook wording is interpreted only with the hypotheses specified above.}
\GSfield{Assessment}{literal statement: correct with missing hypotheses; publication match: exact; verification: read the result; development: annotated; confidence: high}
\end{GSitem}
\begin{GSitem}{1979-470.6}{A}{Reducibility in one dimension}
\GSfield{Kind; source}{theorem; \GSpage{P000864-P000865}}
\GSfield{The notes}{Theorem~2 states reducibility for \(\dim X=1\), followed by the Brownian path model.}
\GSfield{Statement}{For the one-dimensional based path-group energy representation on a finite interval, the Brownian realization has non-scalar commuting right-action operators as well as the left-action representation. Hence the left representation is reducible for nontrivial \(G\).}
\GSfield{Reconstruction}{Boundary and base-point conventions matter. This is the path model actually developed next in the notes, not an unqualified assertion for every subgroup or completion of a one-dimensional mapping group.}
\GSfield{Fate}{The 1981 Albeverio--Høegh-Krohn--Testard paper records reducibility. A directly checked later proof uses commuting left/right Brownian representations and their energy equivalence \citep[Theorems~4.6 and~5.8]{albeverioetal2015brownian}.}
\GSfield{Present status}{As of November 2026, the stated, qualified result is established; the notebook wording is interpreted only with the hypotheses specified above.}
\GSfield{Assessment}{literal statement: correct with missing hypotheses; publication match: partial; verification: read the result; development: annotated; confidence: high}
\end{GSitem}
\begin{GSitem}{1979-470.7}{C}{Irreducibility on a two-dimensional base}
\GSfield{Kind; source}{problem; \GSpage{P000864}}
\GSfield{The notes}{\(\dim X=2\) is underlined and followed by a question mark.}
\GSfield{Statement}{For a connected smooth Riemannian surface \(X\), compact semisimple \(G\), positive smooth measure density and positive invariant inner product on \(\mathfrak g\), decide irreducibility of the ordinary current-group energy representation for arbitrary such data.}
\GSfield{Reconstruction}{The parameters are part of the question: a result after rescaling the covariance or coupling is not an all-parameter solution. Projective positive-energy representations of a semidirect product with time evolution are a different representation class.}
\GSfield{Fate}{\citet[Theorem~6.2]{wallach1987gauge} proves irreducibility when \(d\mu=\rho\,d\operatorname{vol}\) and every root obeys \(\|\alpha\|_B>\sqrt{8\pi|\rho(x)|}\) at every \(x\). \citet{shimada2005irreducibility} claimed all dimensions, but \citet[p.~1]{hasebe2010white} identifies the invalid use of \(d\Phi=\beta(\exp\Phi)\) for nonabelian \(G\); \citet[Introduction]{ando2010local} also warns about the claim. The explicit \(SU(2)\) diagnostic in Section~\ref{dev:energy-diagnostic} illustrates the false differential identity; it is not a reducibility counterexample.}
\GSfield{Present status}{As of November 2026, the verified sources establish a genuine parameter-restricted theorem, while the unrestricted two-dimensional claim remains unresolved within the checked literature. Targeted recent searches found no verified repair; this is not an exhaustive assertion of current global openness.}
\GSfield{Assessment}{literal statement: not determinable; publication match: partial; verification: read the result; development: annotated; confidence: medium; problem \hyperref[prob:GP-1979-06]{GP-1979-06}}
\GSfield{Worth developing}{The remaining distinction is mathematically sharp: remove the root/covariance bound without importing a flawed nonabelian differential identity or changing the representation.}
\end{GSitem}
\begin{GSitem}{1979-470.8}{E}{Stochastic development to Brownian paths on a group}
\GSfield{Kind; source}{construction; \GSpage{P000864-P000865}}
\GSfield{The notes}{\(\eta^{-1}d\eta=db\); the notes say first to solve an ordinary equation and then transfer the measure. The displayed source space is accidentally \(C(\mathbb R^+,G)\).}
\GSfield{Statement}{Let \(b\) be Brownian motion in \(\mathfrak g\) for a positive invariant metric, with \(b_0=0\). On a finite interval, the Stratonovich equation \(\eta^{-1}\circ d\eta=db\), \(\eta_0=e\), defines group Brownian motion and a measurable stochastic-development map from Wiener paths in \(\mathfrak g\) to continuous paths in \(G\).}
\GSfield{Reconstruction}{Brownian paths are not differentiable; the equation needs a stochastic convention. The input is Lie-algebra path space, and the measure \(\mu_0\) belongs there. The group output measure is its pushforward.}
\GSfield{Fate}{The construction is used in \citet[Sec.~4]{albeveriohoeghkrohn1978energy}; the modern formulation and stochastic logarithms are explicit in \citet[Sec.~2]{albeverioetal2015brownian}.}
\GSfield{Present status}{As of November 2026, the stated, qualified result is established; the notebook wording is interpreted only with the hypotheses specified above.}
\GSfield{Assessment}{literal statement: false as written; publication match: exact; verification: read the result; development: annotated; confidence: high}
\end{GSitem}
\begin{GSitem}{1979-470.9}{E}{Quasi-invariance and the Brownian representation}
\GSfield{Kind; source}{theorem; \GSpage{P000865}}
\GSfield{The notes}{Left multiplication \(\eta\mapsto\varphi\eta\) is assigned a Radon--Nikodym derivative for \(\varphi\in C^\infty(\mathbb R^+,G)\), producing \(W_\varphi\).}
\GSfield{Statement}{On based group Wiener space over \([0,T]\), multiplication by an absolutely continuous based path \(h\) with \(\int_0^T\|h^{-1}\dot h\|^2dt<\infty\) is quasi-invariant. Pullback multiplied by the square root of the appropriate Radon--Nikodym derivative is unitary.}
\GSfield{Reconstruction}{The finite-energy condition is essential on an infinite interval. Smoothness alone on \(\mathbb R^+\) does not ensure finite energy. The initial value must preserve the based path space.}
\GSfield{Fate}{The precise finite-energy path group and its density are given in \citet[Sec.~2]{albeverioetal2015brownian}; this is the qualified version of the construction already used in 1978.}
\GSfield{Present status}{As of November 2026, the stated, qualified result is established; the notebook wording is interpreted only with the hypotheses specified above.}
\GSfield{Assessment}{literal statement: correct with missing hypotheses; publication match: exact; verification: read the result; development: annotated; confidence: high}
\end{GSitem}
\begin{GSitem}{1979-470.10}{E}{Equivalence of Brownian and energy representations}
\GSfield{Kind; source}{theorem; \GSpage{P000865}}
\GSfield{The notes}{Theorem: \(W_\varphi\cong V_\varphi\); the corresponding cocycle is the Maurer--Cartan form.}
\GSfield{Statement}{For a compact connected Lie group and the based finite-energy path group on a finite interval, the Brownian representation is unitarily equivalent to the energy representation with the same invariant metric. Left and right Brownian versions are related by inversion.}
\GSfield{Reconstruction}{The comparison is between matching path groups, interval conventions and covariance normalizations. Stochastic logarithm followed by Fourier--Wiener transform provides an explicit intertwiner.}
\GSfield{Fate}{The equivalence is already central to \citet[Sec.~4]{albeveriohoeghkrohn1978energy}. The full later construction is \citet[Theorem~5.8]{albeverioetal2015brownian}; its 2015 preprint is corroboration, not the first announcement.}
\GSfield{Present status}{As of November 2026, the stated, qualified result is established; the notebook wording is interpreted only with the hypotheses specified above.}
\GSfield{Assessment}{literal statement: correct with missing hypotheses; publication match: exact; verification: read the result; development: annotated; confidence: high}
\end{GSitem}
\begin{GSitem}{1979-470.11}{F}{The asserted type-III factor}
\GSfield{Kind; source}{claim; \GSpage{P000865}}
\GSfield{The notes}{\([L_\infty\cup W_\varphi]''\) is called a factor of type~3.}
\GSfield{Statement}{On \(L^2(Y,\mu)\) for a standard probability space, let \(L^\infty(Y,\mu)\) act by multiplication and let \(W\) be the quasi-regular unitaries of a nonsingular action. The commutant of \(\mathcal M=(L^\infty\cup W)''\) consists of invariant multiplication operators. If \(\mathcal M\) is a factor, that commutant is scalar, and \(\mathcal M=B(L^2(Y,\mu))\), a type-I factor.}
\GSfield{Reconstruction}{\ours\ Indeed the multiplication algebra is maximal abelian, so \(\mathcal M'\subset L^\infty\subset\mathcal M\). Factoriality forces \(\mathcal M'=\mathbb C I\). A group-measure-space crossed product acts on a different Hilbert space and must not be substituted silently.}
\GSfield{Fate}{The literal assertion is disproved by this elementary commutant argument. Later type-III factorial path-group representations \citep{albeverioetal1983factorial} concern a different generated algebra; that work is an indirect thematic lead, not an exact repair.}
\GSfield{Present status}{As of November 2026, the printed operator-algebra expression is false under its stated realization. Which different factor construction was intended is not recoverable from the notebook.}
\GSfield{Assessment}{literal statement: false as written; publication match: thematic only; verification: indirect; development: annotated; confidence: high}
\end{GSitem}

\GStalk{471}{1979-10-22}{I. A. Batalin}{P000866}{First-class constraints and structure functions}
The dated block on P000866 introduces first-class constraints, asks whether their structure functions can be made constant, and writes a multiplier action. The catalogue attributes it to Batalin; the notebook's tentative \ru{Некто (Батурин)} is retained as an attribution uncertainty. The catalogue's overlapping P000864--P000867 interval is not a literal boundary: the new block starts on P000866 and ends before Kallosh's heading on that page.

\begin{GSitem}{1979-471.1}{E}{First-class constraints on a symplectic phase space}
\GSfield{Kind; source}{definition; \GSpage{P000866}}
\GSfield{The notes}{The notes set \(g=(q,p)\), use the canonical antisymmetric matrix, and write \(\{T_\alpha,T_\beta\}=u_{\alpha\beta}^{\gamma}T\).}
\GSfield{Statement}{On a symplectic manifold \((M,\omega)\), let \(C=\{T_1=\cdots=T_m=0\}\) with independent differentials on \(C\). The constraints are first class when \(\{T_\alpha,T_\beta\}|_C=0\); locally regularity gives \(\{T_\alpha,T_\beta\}=u_{\alpha\beta}^{\gamma}T_\gamma\). Equivalently \(C\) is coisotropic. The coefficients depend on the chosen generators and need not be constants.}
\GSfield{Reconstruction}{Regularity and the subscript on the final \(T_\gamma\) are supplied; the geometric formulation is an editorial translation.}
\GSfield{Fate}{This is Dirac's pre-existing distinction: \citet[\S6, pp.~136--138]{dirac1950generalized}; the article derives from lectures given in August 1949 and was published in 1950.}
\GSfield{Present status}{As of 10 November 2026, this is the standard local definition; singular constraint sets require separate treatment.}
\GSfield{Assessment}{literal statement: correct with missing hypotheses; publication match: exact; verification: read the result; development: annotated; confidence: high}
\end{GSitem}
\begin{GSitem}{1979-471.2}{E}{Can structure functions be made constant?}
\GSfield{Kind; source}{question; \GSpage{P000866}}
\GSfield{The notes}{The explicit question is \ru{можно ли сделать заменами пер., чтобы было \(u_{\alpha\beta}^{\gamma}=\mathrm{const}\)}. It does not specify locality or whether the defining constraint functions may be replaced.}
\GSfield{Statement}{Locally near a regular point of a coisotropic submanifold of codimension \(m\), one can choose canonical coordinates and equivalent defining functions \(P_1,\ldots,P_m\) so that \(C=\{P_1=\cdots=P_m=0\}\) and \(\{P_a,P_b\}=0\). This affirmative statement allows a nonsingular change of constraint generators. A coordinate change alone, keeping a specified nonconstant coefficient function as the same scalar function, cannot make that function constant.}
\GSfield{Reconstruction}{The stated local regular question is a bounded reconstruction. No global abelianization, preservation of covariance, or uniqueness of structure functions is asserted. The notebook's unqualified wording has no single yes/no answer.}
\GSfield{Fate}{Equivalent constraints already occur in \citet[\S8, eq.~(39)]{dirac1950generalized}. \citet[\S1, eqs.~(1.8)--(1.10), and \S2.A]{gogilidze1996abelianization} explicitly study constructive abelianization (preprint 28 April 1995; journal 1996). That later paper does not establish priority for the elementary local normal-form assertion.}
\GSfield{Present status}{As of 10 November 2026, the regular local version with equivalent generators is settled; global, singular, and covariance-preserving versions are different questions and are not attributed to the speaker.}
\GSfield{Assessment}{literal statement: correct with missing hypotheses; publication match: partial; verification: read the result; development: annotated; confidence: medium; problem \hyperref[prob:GP-1979-20]{GP-1979-20}}
\end{GSitem}
\begin{GSitem}{1979-471.3}{E}{Multiplier action and Hamiltonian gauge directions}
\GSfield{Kind; source}{construction; \GSpage{P000866}}
\GSfield{The notes}{The action is \(S=\int[\frac12 g^a\epsilon_{ab}\dot g^b-H(g)-\lambda^\alpha T_\alpha(g)]\,dt\), followed by the statement that gauge transformations are generated only by first-class constraints.}
\GSfield{Statement}{For a regular first-class system with \(\{H,T_\alpha\}=V_\alpha{}^\beta T_\beta\), the characteristic directions of \(C\) are generated by the Hamiltonian fields \(X_{T_\alpha}|_C\). The action is invariant up to endpoint terms under the corresponding time-dependent gauge variations when the multipliers transform with the requisite time derivative and structure-function terms and the gauge parameter satisfies the chosen endpoint conditions.}
\GSfield{Reconstruction}{Preservation by \(H\), transformations of \(\lambda\), and endpoint conditions are supplied. The canonical one-form in the displayed symmetrized action differs from \(p\,dq\) by an exact term, depending on the matrix convention.}
\GSfield{Fate}{\citet[\S\S6--8]{dirac1950generalized} gives the constraint and multiplier framework before 1979. The stronger assertion that each arbitrary first-class combination by itself gives a full off-shell Lagrangian gauge transformation needs qualifications.}
\GSfield{Present status}{As of 10 November 2026, the regular characteristic-distribution statement is established; the displayed action alone does not prove an unrestricted Dirac conjecture.}
\GSfield{Assessment}{literal statement: correct with missing hypotheses; publication match: partial; verification: read the result; development: annotated; confidence: high}
\end{GSitem}

\GStalk{472}{1979-10-22}{Renata Kallosh}{P000866-P000870}{Background-field quantization, ghosts and gauge dependence}
Kallosh relates formal generating functionals and scattering amplitudes, then presents background-field gauge fixing for gravity, ghost terms, and an on-shell gauge-independence argument. Maxwell theory supplies a transparent example of gauge degeneracy. Several formulas suppress normalization, loop order, or curvature terms. The supergravity discussion concerns gauge-algebra and ghost modifications; it does not claim ultraviolet finiteness.

\begin{GSitem}{1979-472.1}{E}{Functional integral and perturbative scattering functional}
\GSfield{Kind; source}{construction; \GSpage{P000866-P000867}}
\GSfield{The notes}{\(e^{iW[\varphi]}=\int d\phi\,e^{iS[\phi,\varphi]}\) is related to an \(S\)-matrix functional after replacing \(\varphi\) by a solution expressed through an incoming free field.}
\GSfield{Statement}{After specifying a free propagator, asymptotic conditions, normalization and a perturbative prescription, a generating functional determines time-ordered Green functions, and amputated on-shell limits determine scattering amplitudes when the LSZ hypotheses hold. A bare oscillatory functional integral without a measure is only a formal notation.}
\GSfield{Reconstruction}{The source\textquotesingle s \(W\) is a generating functional. Its relation to an effective action requires a Legendre transform; connected functions and the effective action must be distinguished in the later argument.}
\GSfield{Fate}{\citet[\S1, pp.~778--779]{faddeevpopov1974gravity} reviews the earlier functional-integral and diagrammatic construction. The Russian review appeared in November 1973, the English translation in 1974.}
\GSfield{Present status}{As of 10 November 2026, the perturbative construction is standard; these notes do not construct a nonperturbative gravitational measure.}
\GSfield{Assessment}{literal statement: correct with missing hypotheses; publication match: partial; verification: read the result; development: annotated; confidence: high}
\end{GSitem}
\begin{GSitem}{1979-472.2}{E}{Iterating the classical field equation as a tree expansion}
\GSfield{Kind; source}{method; \GSpage{P000867}}
\GSfield{The notes}{The field satisfies a schematic equation \(\widetilde\varphi=\varphi_{\rm in}+D^c*(\delta L_{\rm int}/\delta\widetilde\varphi)\), illustrated by tree diagrams.}
\GSfield{Statement}{For a polynomial interaction and a fixed Green operator \(D\), the equation \(\phi=\phi_0+D N(\phi)\) can be solved recursively as a formal power series in a coupling. Each recursion inserts a vertex and propagator, giving rooted trees. Signs depend on the definition of the free operator; convergence and an incoming scattering interpretation require additional analytic and boundary assumptions.}
\GSfield{Reconstruction}{The tree recursion is separated from loop corrections and from rigorous existence of a nonlinear wave solution. The notebook's \(D^c\) designation by itself does not fix retarded versus Feynman boundary conditions.}
\GSfield{Fate}{\citet[\S1]{faddeevpopov1974gravity} supplies the pre-1979 diagrammatic setting. The precise rooted-tree recursion here is an elementary reconstruction, not an identified new theorem of this talk.}
\GSfield{Present status}{As of 10 November 2026, the formal iterative identity is settled; no convergence assertion is present.}
\GSfield{Assessment}{literal statement: correct with missing hypotheses; publication match: partial; verification: read the result; development: annotated; confidence: high}
\end{GSitem}
\begin{GSitem}{1979-472.3}{E}{Background split and harmonic gauge for the metric fluctuation}
\GSfield{Kind; source}{construction; \GSpage{P000867-P000868}}
\GSfield{The notes}{The gravitational integral uses \(S[g+h]\), with \(F_\nu=D^\mu[g]h_{\mu\nu}-\frac12D_\nu[g]h^\mu{}_\mu\) and a quadratic gauge-fixing action.}
\GSfield{Statement}{Fix a smooth nondegenerate background metric \(g\). In a local perturbative gauge, add \(\frac1{2\alpha}\int \sqrt{|g|}\,F_\nu F^\nu\) to the Einstein--Hilbert quadratic fluctuation operator, and include the determinant of the linearized variation of \(F\) along diffeomorphisms. Invertibility means invertibility with specified boundary conditions and after handling residual zero modes.}
\GSfield{Reconstruction}{The gauge parameter, boundary conditions and qualification concerning zero modes are explicit additions.}
\GSfield{Fate}{Covariant gravitational gauge fixing and its ghost determinant predate the talk: \citet[\S\S2--3]{faddeevpopov1974gravity}.}
\GSfield{Present status}{As of 10 November 2026, this remains the conventional perturbative background gauge, with global gauge-copy and measure issues outside the claim.}
\GSfield{Assessment}{literal statement: correct with missing hypotheses; publication match: exact; verification: read the result; development: annotated; confidence: high}
\end{GSitem}
\begin{GSitem}{1979-472.4}{E}{The gravitational ghost operator and perturbative unitarity}
\GSfield{Kind; source}{claim; \GSpage{P000867}}
\GSfield{The notes}{The notes write a ghost term schematically as \(\bar c^\nu D_\mu[g]D^\mu[g+h]c_\nu\), and attribute loss of unitarity without an added term to Feynman.}
\GSfield{Statement}{The Faddeev--Popov operator is \(M_\nu{}^\rho=\delta F_\nu/\delta\xi_\rho\). At \(h=0\), with \(\delta h_{\mu\nu}=\nabla_\mu\xi_\nu+\nabla_\nu\xi_\mu\), it is \(\nabla^2\delta_\nu{}^\rho+R_\nu{}^\rho\), up to curvature and action-sign conventions. Its Grassmann determinant implements cancellations needed for physical perturbative unitarity; gauge fixing alone is not that determinant.}
\GSfield{Reconstruction}{The Ricci term and the identification of the indispensable extra term as the ghost determinant are supplied. A positive physical Hilbert-space statement cannot be inferred just from the unconstrained integral.}
\GSfield{Fate}{\citet[introduction and \S3, eqs.~(3.20)--(3.21)]{faddeevpopov1974gravity} distinguishes Feynman's 1963 one-loop observation from later all-order ghost rules. This supplies an exact match for the mechanism, not for the abbreviated operator.}
\GSfield{Present status}{As of 10 November 2026, the curvature-complete operator and perturbative physical-state cancellation are standard.}
\GSfield{Assessment}{literal statement: false as written; publication match: partial; verification: read the result; development: annotated; confidence: high}
\end{GSitem}
\begin{GSitem}{1979-472.5}{F}{Additional ghost interactions in supergravity}
\GSfield{Kind; source}{remark; \GSpage{P000867}}
\GSfield{The notes}{The text says that one further term suffices in supergravity whereas an arbitrary theory requires infinitely many corrections.}
\GSfield{Statement}{For a gauge algebra that closes only modulo field equations, the naive determinant-only Faddeev--Popov rules need additional terms. Particular supergravity formulations admit quartic ghost interactions. No precise model, auxiliary-field content or bound on ghost degree is recorded, so the universal one-term/infinitely-many contrast is not a determinate theorem.}
\GSfield{Fate}{\citet[pp.~141--152]{kallosh1978modified} is the contemporaneous 1978 publication lead; its publisher abstract and metadata were found but the full article was inaccessible. \citet[\S2.2 and introduction]{fuster2005brst} later distinguishes open algebras and the antifield formalism. It does not turn this 1979 remark into a statement of the later BV master-equation theory.}
\GSfield{Present status}{As of 10 November 2026, the general open-algebra framework is established, but the notebook's asserted truncation cannot be verified without a specified supergravity model.}
\GSfield{Assessment}{literal statement: not determinable; publication match: partial; verification: indirect; development: annotated; confidence: medium}
\end{GSitem}
\begin{GSitem}{1979-472.6}{E}{Diffeomorphism invariance and background covariance}
\GSfield{Kind; source}{claim; \GSpage{P000868}}
\GSfield{The notes}{Equations (9)--(12) compare \(S[g+\delta g]\), \(W[g+\delta g]\), and the induced change of the fluctuation \(h\).}
\GSfield{Statement}{Einstein--Hilbert action is invariant under diffeomorphisms compatible with boundary data. Simultaneous background and fluctuation transformations can preserve a background-covariant gauge-fixed formal integral if the measure, determinant and regularization preserve that symmetry. Infinitesimal equality is understood to first order in the parameter.}
\GSfield{Reconstruction}{Background covariance is kept distinct from independence of the choice of gauge-fixing function.}
\GSfield{Fate}{\citet[\S\S2--3]{faddeevpopov1974gravity} provides the pre-existing covariant construction. Modern BRST qualifications concerning the measure and anomalies are explicit in \citet[\S\S7--9]{fuster2005brst}.}
\GSfield{Present status}{As of 10 November 2026, the classical identity is exact under its boundary hypotheses; a quantum version requires anomaly-free symmetry preservation.}
\GSfield{Assessment}{literal statement: correct with missing hypotheses; publication match: partial; verification: read the result; development: annotated; confidence: high}
\end{GSitem}
\begin{GSitem}{1979-472.7}{E}{On-shell gauge independence and its loop-order restriction}
\GSfield{Kind; source}{claim; \GSpage{P000868-P000869}}
\GSfield{The notes}{After subtracting \(S[g]\) and its linear fluctuation term, the variation of the radiative functional is proportional to \(\delta S/\delta g\); the notes conclude independence on the classical shell.}
\GSfield{Statement}{BRST variation yields a gauge-parameter identity of the form \(\delta\Gamma=\Gamma_{,i}K^i\) for the properly defined effective action. Consequently its value at a quantum stationary point is gauge independent. At one loop \(\Gamma=S+\hbar\Gamma_1+\cdots\), the corresponding variation of \(\Gamma_1\) vanishes on the classical shell \(S_{,i}=0\).}
\GSfield{Reconstruction}{The effective-action interpretation, loop expansion and distinction between the classical and quantum shell are supplied. The displayed notebook formula is treated as a schematic first-order identity; the later Higgs-background treatment checks the principle, not every gravity-specific detail.}
\GSfield{Fate}{\citet[pp.~173--188]{nielsen1975gauge} is the 1975 original publication lead, checked indirectly. The explicitly read primary later treatment \citet[\S2.1, eq.~(1)]{espinosa2015higgstory} states the effective-action identity and its quantum-stationary consequence. The one-loop classical-shell consequence follows by expanding that identity consistently in \(\hbar\). Neither source permits classical and quantum equations of motion to be identified at all orders.}
\GSfield{Present status}{As of 10 November 2026, on-shell gauge independence is established within the BRST perturbative hypotheses; off-shell effective actions generally depend on gauge.}
\GSfield{Assessment}{literal statement: correct with missing hypotheses; publication match: partial; verification: read the result; development: annotated; confidence: medium}
\end{GSitem}
\begin{GSitem}{1979-472.8}{E}{Changing a gauge-fixing functional by a change of variables}
\GSfield{Kind; source}{method; \GSpage{P000869}}
\GSfield{The notes}{The proof varies a quadratic gauge form \(F\) and compensates by \(h\mapsto h+\delta h\). An expectation of \(\delta h\) is written without its insertion, and equation (20) has an unmatched factor \(1/2\).}
\GSfield{Statement}{In a finite-dimensional regular gauge model, an admissible local gauge slice has an invertible derivative along gauge orbits. A change of slice can then be compensated locally, with its Jacobian included. In a formal field integral, the same manipulation requires a regulated measure, appropriate boundaries and consistent normalization. Invertibility of a quadratic fluctuation form alone does not prove global gauge independence.}
\GSfield{Reconstruction}{The elementary change-of-variables conditions replace the notebook's unrestricted functional-integral claim. The flawed equality is not used as a proof.}
\GSfield{Fate}{\citet[\S8]{fuster2005brst} states the BRST gauge-independence argument in a modern form; \citet[\S\S2--3]{faddeevpopov1974gravity} gives its earlier determinant setting. Neither validates the two literal misprints.}
\GSfield{Present status}{As of 10 November 2026, the local/formal argument is standard; a nonperturbative global statement is not established by these formulas.}
\GSfield{Assessment}{literal statement: false as written; publication match: partial; verification: read the result; development: annotated; confidence: high}
\end{GSitem}
\begin{GSitem}{1979-472.9}{E}{Maxwell gauge invariance and its Hessian kernel}
\GSfield{Kind; source}{definition; \GSpage{P000869-P000870}}
\GSfield{The notes}{\(S[A]=-\frac14\int(\partial_\mu A_\nu-\partial_\nu A_\mu)^2\) is invariant under \(A_\mu\mapsto A_\mu+\partial_\mu\xi\); the Hessian is degenerate.}
\GSfield{Statement}{For a flat background and suitable compact support or boundary conditions, \(F=dA\) is unchanged by \(A\mapsto A+d\xi\). The quadratic operator \(K_{\mu\nu}=\eta_{\mu\nu}\Box-\partial_\mu\partial_\nu\) satisfies \(K_{\mu\nu}\partial^\nu\xi=0\), so no inverse exists before removing gauge directions.}
\GSfield{Fate}{This is the Abelian special case of the earlier gauge-action construction; compare \citet[eqs.~(1)--(6)]{yangmills1954conservation}.}
\GSfield{Present status}{As of 10 November 2026, the algebraic kernel identity is elementary and unchanged; boundary conditions determine residual modes.}
\GSfield{Assessment}{literal statement: correct; publication match: exact; verification: read the result; development: annotated; confidence: high}
\end{GSitem}
\begin{GSitem}{1979-472.10}{E}{Lorenz and generalized linear gauges for electromagnetism}
\GSfield{Kind; source}{construction; \GSpage{P000870}}
\GSfield{The notes}{The two examples add \((\partial\cdot A)^2/(2\alpha)\) and \((\partial\cdot A+\beta n\cdot A)^2/(2\alpha)\).}
\GSfield{Statement}{A linear gauge functional \(F(A)=L^\mu A_\mu\) produces a quadratic gauge-fixing term. For Abelian transformations its Faddeev--Popov operator is \(L^\mu\partial_\mu\), independent of \(A\), and the ghosts decouple when that operator and its boundary prescription are fixed. Gauge parameters must avoid singular gauges or require a separate limiting prescription.}
\GSfield{Reconstruction}{The fixed vector \(n\), residual gauge kernel and parameter restrictions are specified; the notes do not prove equivalence for every \(\alpha,\beta\).}
\GSfield{Fate}{\citet[\S\S1--3]{faddeevpopov1974gravity} reviews the determinant construction for which Maxwell theory is the linear case.}
\GSfield{Present status}{As of 10 November 2026, the two expressions are standard admissible local gauge choices when the corresponding operator can be inverted.}
\GSfield{Assessment}{literal statement: correct with missing hypotheses; publication match: partial; verification: read the result; development: annotated; confidence: high}
\end{GSitem}

\GStalk{473}{1979-10-29}{I. A. Batalin}{P000871-P000877}{Faddeev--Popov reduction and BRST invariance}
The scan on P000871 dates this lecture \ru{29/X-79}, so 29 October is used here; the supplied catalogue's 29 September is recorded as a provenance discrepancy. After stationary phase and Maxwell gauge degeneracy, a finite-dimensional group-action model derives the gauge-slice determinant and ghost representation. BRST transformations, physical unitarity, gauge independence, open supergravity algebras and the Yang--Mills example follow. The notes are an exposition of methods circulating before the talk, not an instance of the later odd-antibracket BV formalism.

\begin{GSitem}{1979-473.1}{E}{Sources and time-ordered correlation functions}
\GSfield{Kind; source}{construction; \GSpage{P000871}}
\GSfield{The notes}{The notes define \(Z(J)=c\int e^{iW[\varphi]+iJ\varphi}\,d\mu(\varphi)\) and identify source derivatives with products of fields.}
\GSfield{Statement}{With \(\hbar=1\) and normalization \(Z(0)=1\), the formal time-ordered \(n\)-point function is \(i^{-n}\delta^n Z/\delta J(x_1)\cdots\delta J(x_n)|_{J=0}\). For non-normalized \(Z\), divide by \(Z(0)\). Derivatives of \(-i\log Z\) generate connected functions with the corresponding powers of \(i\).}
\GSfield{Reconstruction}{The missing factors, evaluation at zero source and normalization are stated explicitly. No measure-existence assertion is added.}
\GSfield{Fate}{The construction is reviewed before the seminar in \citet[\S1]{faddeevpopov1974gravity}, Russian 1973/English 1974.}
\GSfield{Present status}{As of 10 November 2026, this is a formal or suitably constructed generating-functional identity, not a proof of existence of an interacting quantum field theory.}
\GSfield{Assessment}{literal statement: correct with missing hypotheses; publication match: exact; verification: read the result; development: annotated; confidence: high}
\end{GSitem}
\begin{GSitem}{1979-473.2}{E}{Coupling expansion versus expansion around a classical solution}
\GSfield{Kind; source}{method; \GSpage{P000871}}
\GSfield{The notes}{A scalar \(\varphi^4\) example precedes two schemes: a coupling expansion and stationary phase near a solution. The notebook says the first is inapplicable to gravity and writes a dot on \(\varphi\) in the mass term.}
\GSfield{Statement}{For a scalar field the conventional action has a kinetic term, a mass term \(m^2\varphi^2/2\), and a quartic interaction. Coupling expansion and an \(\hbar\)-expansion about a saddle are distinct formal organizations; both can be used in perturbative gravity after gauge fixing. Nonrenormalizability by power counting does not make every coupling expansion undefined.}
\GSfield{Reconstruction}{The mass-term emendation and the distinction between perturbative definability and renormalizability are editorial.}
\GSfield{Fate}{\citet[\S\S1--2]{faddeevpopov1974gravity} explicitly presents perturbation theory for gravity before 1979 and therefore rules out the blanket prohibition.}
\GSfield{Present status}{As of 10 November 2026, the two expansions are standard; the printed mass term and absolute prohibition need correction.}
\GSfield{Assessment}{literal statement: false as written; publication match: partial; verification: read the result; development: annotated; confidence: high}
\end{GSitem}
\begin{GSitem}{1979-473.3}{E}{Nondegenerate stationary phase and its missing normalization}
\GSfield{Kind; source}{theorem; \GSpage{P000871-P000872}}
\GSfield{The notes}{The notes approximate an oscillatory integral by the inverse square root of the Hessian determinant and a value at a stationary point.}
\GSfield{Statement}{Let \(S\in C^\infty(\mathbb R^n;\mathbb R)\), \(a\in C_c^\infty(\mathbb R^n)\), and let \(x_0\) be the unique critical point of \(S\) in the relevant support, with nonsingular Hessian \(H=S''(x_0)\). As \(\hbar\downarrow0\),
\[
\begin{aligned}
\int_{\mathbb R^n}e^{iS(x)/\hbar}a(x)\,dx
&=e^{iS(x_0)/\hbar}(2\pi\hbar)^{n/2}
  e^{i\pi\operatorname{sgn}(H)/4}\\
&\qquad\times|\det H|^{-1/2}\bigl[a(x_0)+O(\hbar)\bigr].
\end{aligned}
\]
Multiple critical points contribute a sum; a critical manifold requires a different theorem.}
\GSfield{Reconstruction}{Amplitude, support, nondegeneracy, absolute determinant, signature phase and power of \(2\pi\hbar\) are supplied. The generic remainder is \(O(\hbar)\), not the notebook's stronger small-o assertion.}
\GSfield{Fate}{The finite-dimensional Gaussian basis of the earlier path-integral method is given in \citet[\S1]{faddeevpopov1974gravity}. The precise stationary-phase statement here is a classical mathematical reconstruction, not a new result identified with this lecture.}
\GSfield{Present status}{As of 10 November 2026, the finite-dimensional theorem is settled; it does not justify an infinite-dimensional gravity saddle expansion without further work.}
\GSfield{Assessment}{literal statement: false as written; publication match: partial; verification: read the result; development: annotated; confidence: high}
\end{GSitem}
\begin{GSitem}{1979-473.4}{E}{Gauge degeneracy obstructs the naive saddle formula}
\GSfield{Kind; source}{claim; \GSpage{P000872}}
\GSfield{The notes}{For Maxwell theory the Hessian annihilates gradients, and integration along a gauge orbit is described as producing a group-volume factor.}
\GSfield{Statement}{If \(S\) is invariant under a smooth group action, its Hessian at a critical point annihilates infinitesimal orbit directions. Thus the isolated nondegenerate stationary-phase formula cannot be applied in those directions. For Maxwell theory the identity is \((\eta_{\mu\nu}\Box-\partial_\mu\partial_\nu)\partial^\nu\xi=0\).}
\GSfield{Fate}{This is the obstruction addressed by the earlier Faddeev--Popov construction; \citet[\S2]{faddeevpopov1974gravity}.}
\GSfield{Present status}{As of 10 November 2026, the degeneracy statement is exact; an infinite group volume is a formal description unless a quotient measure is defined.}
\GSfield{Assessment}{literal statement: correct with missing hypotheses; publication match: exact; verification: read the result; development: annotated; confidence: high}
\end{GSitem}
\begin{GSitem}{1979-473.5}{E}{Infinitesimal group generators, Lie equations and invariant measures}
\GSfield{Kind; source}{construction; \GSpage{P000872-P000873}}
\GSfield{The notes}{A transformation \(g\mapsto f(g,\theta)\), generators \(R_\alpha^a\), structure constants and a matrix \(\lambda(\theta)\) are introduced. Divergence-free generators are associated with a \ru{унимодулярная} group.}
\GSfield{Statement}{For a smooth Lie-group action, the derivatives of \(f\) along group coordinates factor through fundamental vector fields and a Maurer--Cartan frame. A density \(\rho\,dg\) is invariant precisely when \(\partial_a(\rho R_\alpha^a)=0\). This is an action-and-density condition. Unimodularity of the abstract group means its left and right Haar measures agree and is not equivalent to \(\partial_aR_\alpha^a=0\) in arbitrary coordinates on the acted-on manifold.}
\GSfield{Reconstruction}{The density and the distinction between group unimodularity and volume preservation of a particular action are supplied.}
\GSfield{Fate}{The invariant-measure requirement is part of the pre-existing gauge argument; compare \citet[\S2]{faddeevpopov1974gravity}. \citet[\S~II]{bering1997geometry} later makes its geometric measure assumptions explicit (1997 preprint, revised 1998).}
\GSfield{Present status}{As of 10 November 2026, the Lie-action identities are standard; the notebook's terminology should not be used as an equivalence theorem.}
\GSfield{Assessment}{literal statement: false as written; publication match: partial; verification: read the result; development: annotated; confidence: high}
\end{GSitem}
\begin{GSitem}{1979-473.6}{E}{An admissible local gauge slice}
\GSfield{Kind; source}{definition; \GSpage{P000873}}
\GSfield{The notes}{The condition \(\psi^\alpha(f(g,\theta))=0\) is required to have one solution \(\theta\) for every \(g\); the notes call this a section.}
\GSfield{Statement}{For a locally free action of an \(m\)-dimensional Lie group, \(m\) gauge functions define a local slice when \(M^\alpha{}_\beta(g)=d\psi^\alpha_g(R_\beta(g))\) is invertible on the slice. The inverse-function theorem then gives local uniqueness near a chosen intersection. Global existence and uniqueness require additional assumptions and generally fail in non-Abelian gauge theory.}
\GSfield{Reconstruction}{Local uniqueness is substituted for the universal global section that the notes assume rather than prove.}
\GSfield{Fate}{The local determinant method predates 1979. The global obstruction is the topic of \citet[pp.~1--19]{gribov1978quantization}, published in 1978, but that original text was not recovered in this audit. The independently read later primary analysis \citet[\S~II]{bering1997geometry} explicitly excludes global Gribov problems from its gauge-independence argument.}
\GSfield{Present status}{As of 10 November 2026, local slices remain useful; a universal global slice is not available in general.}
\GSfield{Assessment}{literal statement: correct with missing hypotheses; publication match: partial; verification: indirect; development: annotated; confidence: medium}
\end{GSitem}
\begin{GSitem}{1979-473.7}{E}{Faddeev--Popov determinant identity}
\GSfield{Kind; source}{theorem; \GSpage{P000873}}
\GSfield{The notes}{Equations (21)--(27) derive \(1=\int\delta(\psi(f(g,\theta)))\Delta_\psi(f(g,\theta))\,d\sigma(\theta)\), with \(\Delta_\psi=\det(\psi_{,a}R^a)\). Equation (23) omits \(f\) in one argument.}
\GSfield{Statement}{If the orbit-to-gauge map has exactly one regular zero in the integration chart, the ordinary real delta-function change-of-variables formula gives \(1=\int \delta^m(\psi(f(g,\theta)))|\det \partial_\theta\psi(f(g,\theta))|\,d^m\theta\). With a consistently oriented formal measure one uses the signed determinant. The group-coordinate Jacobian converts this to the Haar-measure form.}
\GSfield{Reconstruction}{The absolute value for positive measures, possible orientation convention and unique-zero hypothesis are supplied. Equation (23)'s argument must be \(f(g,\theta)\).}
\GSfield{Fate}{\citet[\S2]{faddeevpopov1974gravity} gives the earlier determinant insertion. The local real-variable identity explains exactly which regularity and multiplicity conditions are needed.}
\GSfield{Present status}{As of 10 November 2026, the local change-of-variables identity is settled; multiple gauge intersections change the naive formula.}
\GSfield{Assessment}{literal statement: correct with missing hypotheses; publication match: exact; verification: read the result; development: annotated; confidence: high}
\end{GSitem}
\begin{GSitem}{1979-473.8}{E}{Factoring gauge volume from an invariant integral}
\GSfield{Kind; source}{method; \GSpage{P000873-P000874}}
\GSfield{The notes}{The invariant integrand and measure are changed by \(g\mapsto f(g,\theta)\), producing a common gauge-volume factor and a reduced delta-function integral.}
\GSfield{Statement}{For a free proper action with compatible measures and convergent integrals, disintegration along orbits yields a quotient integral and the orbit-volume factor. In the formal field-theoretic version, common infinite gauge volumes are canceled only after a specified normalization; insertion of the determinant and compatible gauge slice is essential.}
\GSfield{Reconstruction}{The honest finite-dimensional hypotheses are supplied. The notes' repeated \(d\sigma\) in an integral over field coordinates is a variable-label error, not a new measure.}
\GSfield{Fate}{The preceding reduction is the pre-1979 Faddeev--Popov method; \citet[\S2]{faddeevpopov1974gravity}.}
\GSfield{Present status}{As of 10 November 2026, the geometric quotient construction is standard where its hypotheses hold; cancellation of divergent factors alone is not a proof.}
\GSfield{Assessment}{literal statement: correct with missing hypotheses; publication match: partial; verification: read the result; development: annotated; confidence: high}
\end{GSitem}
\begin{GSitem}{1979-473.9}{E}{Gauge averaging and the indispensable normalization}
\GSfield{Kind; source}{construction; \GSpage{P000874-P000875}}
\GSfield{The notes}{The slice is moved to \(\psi(g)=c\), averaged with \(e^{i\Omega(c)}dc\), and equation (40) divides by \(\int e^{i\Omega(c)}dc\). The next displayed equality drops this denominator.}
\GSfield{Statement}{If \(I(c)\) has the same value \(I\) for every admissible slice parameter \(c\), then
\[
\frac{\int w(c)I(c)\,dc}{\int w(c)\,dc}=I,
\]
provided the denominator is nonzero and the integrals or their regularized meanings exist. Taking \(w=e^{i\Omega}\) gives a gauge-fixing action; normalization is still required for the stated independence.}
\GSfield{Reconstruction}{The distinction between the normalized quotient and the raw averaged integral is explicit.}
\GSfield{Fate}{This is the normalization step in the earlier covariant gauge procedure; \citet[\S3]{faddeevpopov1974gravity}.}
\GSfield{Present status}{As of 10 November 2026, the normalized identity is elementary; equation (41) needs the factor already present in equation (40).}
\GSfield{Assessment}{literal statement: false as written; publication match: partial; verification: read the result; development: annotated; confidence: high}
\end{GSitem}
\begin{GSitem}{1979-473.10}{E}{Grassmann representation of the gauge determinant}
\GSfield{Kind; source}{construction; \GSpage{P000874-P000875}}
\GSfield{The notes}{The determinant \(\det(\psi_{,a}R^a)\) is represented using anticommuting \(\xi,\bar\xi\), producing \(W_{\rm eff}=W+\Omega(\psi)+\bar\xi\,\psi_{,a}R^a\xi\).}
\GSfield{Statement}{For an \(m\times m\) matrix \(M\), a fixed Berezin integration order gives \(\int d\bar c\,dc\,e^{-\bar cMc}=\det M\). In an oscillatory convention \(e^{i\bar cMc}\), the harmless factor \(i^m\) is absorbed in normalization. Ghost parity is opposite to the gauge-parameter parity; fermionic gauge directions require commuting ghosts.}
\GSfield{Reconstruction}{Finite-dimensional normalization and the graded qualification are supplied; the general fermionic extension is not claimed to have been read in the inaccessible 1978 article.}
\GSfield{Fate}{\citet[\S3, following eq.~(3.21)]{faddeevpopov1974gravity} uses the determinant representation. The graded first/second-class extension is the pre-1979 lead \citet[pp.~343--348]{fradkinfradkina1978quantization}.}
\GSfield{Present status}{As of 10 November 2026, the finite-dimensional Berezin identity is exact; its functional-integral use remains subject to the chosen regularization.}
\GSfield{Assessment}{literal statement: correct with missing hypotheses; publication match: exact; verification: read the result; development: annotated; confidence: high}
\end{GSitem}
\begin{GSitem}{1979-473.11}{E}{BRST transformations and the antighost qualification}
\GSfield{Kind; source}{construction; \GSpage{P000875-P000876}}
\GSfield{The notes}{Equations (43)--(46) give \(\delta g^a=R_\alpha^a\xi^\alpha\mu\), \(\delta\xi^\alpha=\frac12t^\alpha_{\beta\gamma}\xi^\beta\xi^\gamma\mu\), and \(\delta\bar\xi_\alpha=\Omega_{,\alpha}(\psi)\mu\), with odd constant \(\mu\).}
\GSfield{Statement}{For a closed Lie gauge algebra, \(s\phi=R_\alpha c^\alpha\), \(sc^\alpha=-\frac12f^\alpha_{\beta\gamma}c^\beta c^\gamma\) define a nilpotent odd derivation with consistent left/right conventions. Add an auxiliary field \(b_\alpha\) and set \(s\bar c_\alpha=b_\alpha,\ sb_\alpha=0\) for off-shell nilpotence. Eliminating \(b\) by its equation yields the notebook's gauge-function variation, whose antighost nilpotence is generally only on shell.}
\GSfield{Reconstruction}{Signs are translated from the notebook's right odd parameter. The auxiliary field is a modern clarification, not a field visibly introduced in this passage.}
\GSfield{Fate}{The original development was publicly announced in 1974 and appears in the April 1975 report \citet[\S\S2--3]{becchirouetstora1975renormalization}; \citet[\S6, eqs.~(6.1)--(6.5)]{fuster2005brst} makes the auxiliary-field formulation explicit.}
\GSfield{Present status}{As of 10 November 2026, the BRST construction is established; nilpotence of the displayed antighost rule requires the stated on-shell qualification.}
\GSfield{Assessment}{literal statement: correct with missing hypotheses; publication match: partial; verification: read the result; development: annotated; confidence: high}
\end{GSitem}
\begin{GSitem}{1979-473.12}{E}{Chevalley--Eilenberg cancellation behind BRST nilpotence}
\GSfield{Kind; source}{claim; \GSpage{P000876}}
\GSfield{The notes}{The notes introduce \(\Gamma_\alpha=R_\alpha^a\partial_a\), compute \((\Gamma_\alpha\xi^\alpha)^2\), and cancel it using the ghost transformation.}
\GSfield{Statement}{If \([\Gamma_\alpha,\Gamma_\beta]=f_{\alpha\beta}^\gamma\Gamma_\gamma\), then \(s=c^\alpha\Gamma_\alpha-\frac12f_{\alpha\beta}^{\gamma}c^\alpha c^\beta\partial/\partial c^\gamma\) satisfies \(s^2=0\). The representation identity and Jacobi identity give the cancellations; this statement applies to a closed Lie algebra action, not automatically to field-dependent or on-shell closing algebras.}
\GSfield{Reconstruction}{The operator is written with a single left-derivative convention. Structure functions require the corresponding more general resolution.}
\GSfield{Fate}{The pre-1979 BRST report \citet[\S2]{becchirouetstora1975renormalization} presents this algebraic mechanism; \citet[\S4]{fuster2005brst} relates it to cohomology.}
\GSfield{Present status}{As of 10 November 2026, the nilpotence identity is exact under the closed-algebra assumptions.}
\GSfield{Assessment}{literal statement: correct with missing hypotheses; publication match: exact; verification: read the result; development: annotated; confidence: high}
\end{GSitem}
\begin{GSitem}{1979-473.13}{E}{Physical unitarity and cancellation of ghost states}
\GSfield{Kind; source}{claim; \GSpage{P000876}}
\GSfield{The notes}{The notes discuss a unitary \(S\)-operator, internal ghosts, and cancellations leaving physical particles.}
\GSfield{Statement}{In an anomaly-free perturbative BRST theory, the conserved nilpotent charge defines physical states through its cohomology. If the induced inner product is positive and unphysical states form canceling sectors, the \(S\)-operator induced on physical cohomology is unitary. The enlarged ghost state space itself need not have a positive-definite inner product.}
\GSfield{Reconstruction}{The physical quotient and positivity assumptions prevent the bare symbolic equality from being interpreted as a universal Hilbert-space theorem.}
\GSfield{Fate}{\citet[\S3, pp.~29--39 and abstract]{becchirouetstora1975renormalization} proves physical unitarity and gauge-parameter independence for its specified massive Higgs-model setting, not for every conceivable gauge theory.}
\GSfield{Present status}{As of 10 November 2026, physical perturbative unitarity is established for standard anomaly-free models; the notes do not settle global or nonperturbative unitarity of supergravity.}
\GSfield{Assessment}{literal statement: correct with missing hypotheses; publication match: partial; verification: read the result; development: annotated; confidence: high}
\end{GSitem}
\begin{GSitem}{1979-473.14}{E}{Field-dependent BRST parameters and gauge independence}
\GSfield{Kind; source}{method; \GSpage{P000876-P000877}}
\GSfield{The notes}{A field-dependent odd parameter involving \(\bar\xi\,\delta\psi\) is used to cancel a change of gauge; Slavnov identities are mentioned.}
\GSfield{Statement}{An infinitesimal change of gauge fermion changes the BRST-exact gauge-fixing action. A BRST change of integration variables, with its super-Jacobian included, leaves expectation values of BRST-closed observables unchanged when the measure is BRST invariant, boundary terms vanish, and the integration domain is admissible. A field-dependent parameter generally has a nontrivial Jacobian even if a constant BRST transformation preserves the measure.}
\GSfield{Reconstruction}{The measure/domain qualifications and the distinction between constant and field-dependent parameters are made explicit.}
\GSfield{Fate}{The earlier BRST report \citet[\S3, pp.~29--34]{becchirouetstora1975renormalization} provides the gauge-independence identities; \citet[\S V, eqs.~(V.8)--(V.9)]{bering1997geometry} explicitly treats the later geometric gauge-fixing argument and its endpoint condition.}
\GSfield{Present status}{As of 10 November 2026, the formal gauge-independence theorem is standard with these hypotheses; it is not a proof that gauge copies or anomalies cannot occur.}
\GSfield{Assessment}{literal statement: correct with missing hypotheses; publication match: exact; verification: read the result; development: annotated; confidence: high}
\end{GSitem}
\begin{GSitem}{1979-473.15}{E}{Open supergravity algebras}
\GSfield{Kind; source}{remark; \GSpage{P000877}}
\GSfield{The notes}{The lecture says supergravity is \ru{не группа}, in contrast with the preceding Yang--Mills derivation.}
\GSfield{Statement}{The identifiable mathematical content is that commutators of gauge transformations in some supergravity formulations close only modulo equations of motion, or have field-dependent coefficients. This is an open gauge algebra. It is not the assertion that no groupoid, integrated symmetry or auxiliary-field formulation can exist.}
\GSfield{Reconstruction}{Open-algebra terminology is supplied. The notes do not specify which supergravity or auxiliary fields are intended.}
\GSfield{Fate}{\citet[pp.~141--152]{kallosh1978modified} is a pre-talk lead; \citet[\S2.2]{fuster2005brst} verifies the later precise distinction. The original 1978 full text was not accessed.}
\GSfield{Present status}{As of 10 November 2026, the distinction is established, but no model-specific closure identity can be recovered from this sentence.}
\GSfield{Assessment}{literal statement: not determinable; publication match: partial; verification: indirect; development: annotated; confidence: medium}
\end{GSitem}
\begin{GSitem}{1979-473.16}{E}{Yang--Mills curvature, infinitesimal gauge generator and ghosts}
\GSfield{Kind; source}{construction; \GSpage{P000877}}
\GSfield{The notes}{The final example takes a compact semisimple group, \(F=dA+[A,A]\), a Lorenz gauge, and a ghost term \(-(\partial_\mu\bar\xi)D^\mu\xi\). Equation (60) abbreviates the functional generator.}
\GSfield{Statement}{For a connection \(A\) on a fixed bundle and invariant inner product on its compact semisimple Lie algebra, \(F_A=dA+\frac12[A,A]\) and \(\delta_\epsilon A=D_A\epsilon\). In a local Lorenz gauge the Faddeev--Popov operator is \(\partial^\mu D_\mu\), with ghost action represented after integration by parts. As a distributional kernel, \(R_{\mu\,b}^{a}(x,y)=D_{\mu\,b}^{a}(x)\delta(x-y)\); the delta function and spacetime index cannot be dropped in a literal kernel equation.}
\GSfield{Reconstruction}{Bundle, contraction, boundary and kernel conventions are supplied; the gauge parameter's reciprocal can be absorbed by renaming it.}
\GSfield{Fate}{\citet[eqs.~(1)--(11)]{yangmills1954conservation} supplies the classical 1954 model; \citet[\S2]{becchirouetstora1975renormalization} gives its pre-1979 ghost and BRST extension.}
\GSfield{Present status}{As of 10 November 2026, the classical and perturbative construction is standard; a rigorous four-dimensional continuum quantum theory is a separate problem.}
\GSfield{Assessment}{literal statement: correct with missing hypotheses; publication match: exact; verification: read the result; development: annotated; confidence: high}
\end{GSitem}

\GStalk{474}{1979-11-12}{I. M. Gelfand}{P000878-P000879}{Gravity, Einstein metrics and the gravitational boundary action}
The heading says \ru{По Хокингу}. Gelfand compares approaches to quantum gravity, writes the Einstein--Hilbert variational problem, mentions Yau, and motivates the gravitational boundary term and reference subtraction. The notebook's purported universal uniqueness theorem for Einstein metrics is substantially false. The corrected Kähler statements are identified with Yau's 1977 announcement/1978 proof; the boundary action matches Gibbons--Hawking 1977.

\begin{GSitem}{1979-474.1}{F}{Difficulties in operator and canonical approaches to gravity}
\GSfield{Kind; source}{remark; \GSpage{P000878}}
\GSfield{The notes}{A heading promises three approaches; the written enumeration discusses operator-valued metrics and a Hamiltonian splitting \(M^3\times\mathbb R\). Metric contraction and covariance/topology are raised as difficulties.}
\GSfield{Statement}{No theorem with quantified hypotheses is stated. A Hamiltonian formulation chooses suitable spacelike hypersurfaces and imposes gravitational constraints; this restricts the geometric setting and does not by itself prove incompatibility with diffeomorphism invariance. Operator-valued metric contractions require an ordering and regularization prescription.}
\GSfield{Reconstruction}{The unwritten third enumerated approach is not invented. The statement distinguishes methodological questions from an impossibility theorem.}
\GSfield{Fate}{\citet[\S1]{gibbonshawking1977action} supplies the path-integral context explicitly named by the notebook. It is a thematic match, not a publication of Gelfand's informal assessment.}
\GSfield{Present status}{As of 10 November 2026, the notes do not specify a uniquely identifiable obstruction theorem or research problem.}
\GSfield{Assessment}{literal statement: not determinable; publication match: thematic only; verification: read the result; development: annotated; confidence: medium}
\end{GSitem}
\begin{GSitem}{1979-474.2}{E}{Einstein equation from the action and the stress tensor normalization}
\GSfield{Kind; source}{construction; \GSpage{P000878-P000879}}
\GSfield{The notes}{The notes vary \((16\pi G)^{-1}\int(R-2\Lambda)\sqrt{-g}\,d^4x+\int L_m\sqrt{-g}\,d^4x\), but the displayed definition of \(T_{ab}\) has a factor \(-\frac12\sqrt{-g}\) and varies with respect to \(g_{ab}\).}
\GSfield{Statement}{Let \(S_m[g,\phi]\) be the matter action and define \(T_{ab}=-2(-g)^{-1/2}\delta S_m/\delta g^{ab}\). With fixed boundary data and the requisite gravitational boundary term, metric stationarity gives \(R_{ab}-\frac12Rg_{ab}+\Lambda g_{ab}=8\pi G\,T_{ab}\). Equivalently \(T^{ab}=2(-g)^{-1/2}\delta S_m/\delta g_{ab}\).}
\GSfield{Reconstruction}{The distinction between \(L_m\), its volume density and the integrated matter action, and the inverse-metric variation are supplied.}
\GSfield{Fate}{The variational action and its boundary qualification are explicit in \citet[\S II, eq.~(2.1)]{gibbonshawking1977action}, published 15 May 1977; receipt on 4 October 1976 is not evidence of a public preprint.}
\GSfield{Present status}{As of 10 November 2026, the corrected variational equation is standard; the notebook's stress-tensor normalization cannot be retained as a definition of the displayed Einstein equation.}
\GSfield{Assessment}{literal statement: false as written; publication match: partial; verification: read the result; development: annotated; confidence: high}
\end{GSitem}
\begin{GSitem}{1979-474.3}{E}{What Yau's theorem actually says}
\GSfield{Kind; source}{claim; \GSpage{P000879}}
\GSfield{The notes}{The page attributes to Yau existence and uniqueness of an Einstein metric on a compact Kähler manifold with \(\Lambda\ne0\).}
\GSfield{Statement}{On a compact Kähler manifold, any real closed \((1,1)\)-form representing \(2\pi c_1(X)\) is the Ricci form of a unique Kähler metric in each fixed Kähler class. For \(c_1(X)=0\) this yields a unique Ricci-flat metric in each class. If \(c_1(X)<0\), there is a uniquely normalized negative Kähler--Einstein metric. Compact Kähler manifolds with positive \(c_1\) need not admit a Kähler--Einstein metric; arbitrary compact Kähler manifolds do not satisfy the notebook's claim.}
\GSfield{Reconstruction}{The sign of the first Chern class, prescribed Ricci form, fixed Kähler class and normalization are all essential emendations.}
\GSfield{Fate}{\citet[Theorem~1, p.~1798]{yau1977calabi} announced the corrected results in May 1977; the full proof appeared in \citet[pp.~339--411]{yau1978ricci} in 1978. The announcement credits Aubin in the negative case. Neither asserts universal \(\Lambda\ne0\) existence.}
\GSfield{Present status}{As of 10 November 2026, these corrected theorems remain established. The notes' unrestricted uniqueness is not an open problem.}
\GSfield{Assessment}{literal statement: false as written; publication match: partial; verification: read the result; development: annotated; confidence: high}
\end{GSitem}
\begin{GSitem}{1979-474.4}{E}{Gibbons--Hawking boundary term and gluing}
\GSfield{Kind; source}{construction; \GSpage{P000879}}
\GSfield{The notes}{The action acquires \((8\pi)^{-1}\int_{\partial M}K\sqrt{-h}\), motivated by boundary variation and an additivity relation denoted \(I[g_1+g_2]=I[g_1]+I[g_2]\).}
\GSfield{Statement}{For a smooth non-null boundary and fixed induced metric, the Einstein--Hilbert action needs the extrinsic-curvature boundary term, with sign determined by signature and outward-normal conventions, to cancel normal derivatives of the metric variation. When two regions are glued smoothly along a common hypersurface, their opposite normals cancel the interface terms. Corners or nonsmooth gluing require additional analysis.}
\GSfield{Reconstruction}{\(g_1+g_2\) is interpreted as gluing regions carrying metrics, not adding two metric tensors. Smoothness and the absence of corners are supplied.}
\GSfield{Fate}{\citet[\S II, eq.~(2.1) and p.~2753]{gibbonshawking1977action} gives the 1977 Dirichlet boundary prescription.}
\GSfield{Present status}{As of 10 November 2026, the boundary term is standard; the literal tensor-addition interpretation of the notebook notation is invalid.}
\GSfield{Assessment}{literal statement: correct with missing hypotheses; publication match: exact; verification: read the result; development: annotated; confidence: high}
\end{GSitem}
\begin{GSitem}{1979-474.5}{E}{Reference subtraction for an asymptotically flat action}
\GSfield{Kind; source}{construction; \GSpage{P000879}}
\GSfield{The notes}{The boundary integrand is changed from \(K\) to \(K-K_0\), with \(K_0\) associated to a comparison geometry.}
\GSfield{Statement}{For a chosen asymptotically flat reference and compatible induced boundary geometry, subtracting its extrinsic-curvature contribution fixes the zero of the action and can render the limit finite. A function of the fixed induced boundary metric may be added without changing the bulk Dirichlet variational equation. Existence or uniqueness of a reference embedding is not automatic for every boundary.}
\GSfield{Reconstruction}{The asymptotically flat comparison and fixed-boundary conditions prevent a universal reference-subtraction theorem from being attributed to the notes.}
\GSfield{Fate}{\citet[\S II, p.~2753, before eq.~(2.2)]{gibbonshawking1977action} explicitly chooses the subtraction to make the flat-space action vanish.}
\GSfield{Present status}{As of 10 November 2026, reference subtraction is established in its stated setting; its geometric choices remain part of the prescription.}
\GSfield{Assessment}{literal statement: correct with missing hypotheses; publication match: exact; verification: read the result; development: annotated; confidence: high}
\end{GSitem}

\GStalk{475}{1979-11-12}{I. A. Batalin}{P000880-P000883}{Graded constraints, Dirac brackets and phase-space measures}
The lecture develops graded canonical brackets, first- and second-class constraints, reduction by Dirac brackets, and the phase-space path integral. Gelfand recognizes the Dirac bracket as the inverse of the induced symplectic form. Source-generating equations lead to a formal constrained integral with the square-root Berezinian; adding gauge conditions makes a regular first-class system second class. These are pre-existing Hamiltonian methods, with a contemporaneous 1977/1978 graded extension.

\begin{GSitem}{1979-475.1}{E}{Even graded canonical Poisson bracket}
\GSfield{Kind; source}{definition; \GSpage{P000880}}
\GSfield{The notes}{Coordinates are assigned parities, and equations (1)--(4) give a canonical bracket, graded antisymmetry and the graded Jacobi identity.}
\GSfield{Statement}{On an even symplectic supermanifold with homogeneous canonical coordinate pairs of equal parity, the Poisson bracket has parity \(\epsilon(f)+\epsilon(g)\), obeys \(\{f,g\}=-(-1)^{\epsilon(f)\epsilon(g)}\{g,f\}\), the graded Leibniz rule and graded Jacobi identity. Left/right derivatives must be used consistently. In particular the bracket of two odd functions is symmetric.}
\GSfield{Reconstruction}{The bracket here is even; it is not the later odd BV antibracket. Coordinate signs are restated invariantly rather than silently copied between incompatible derivative conventions.}
\GSfield{Fate}{The boson/fermion constraint extension was published before the seminar in \citet[pp.~343--348]{fradkinfradkina1978quantization}. Its detailed sign conventions remain an indirectly checked publication lead.}
\GSfield{Present status}{As of 10 November 2026, even symplectic supergeometry supplies the standard mathematical interpretation.}
\GSfield{Assessment}{literal statement: correct with missing hypotheses; publication match: partial; verification: indirect; development: annotated; confidence: medium}
\end{GSitem}
\begin{GSitem}{1979-475.2}{E}{First and second class: why determinant zero is insufficient}
\GSfield{Kind; source}{definition; \GSpage{P000880}}
\GSfield{The notes}{\(Q_{\alpha\beta}=\{T_\alpha,T_\beta\}\) is introduced; nonzero determinant is labeled second class, whereas determinant zero together with a Hamiltonian condition is labeled first class.}
\GSfield{Statement}{For independent constraints, second class means that the full constraint bracket matrix is invertible on the constraint surface; in the graded setting use the appropriate even supermatrix and Berezinian. A function is first class if its bracket with every constraint vanishes on that surface. A degenerate \(Q\) can contain both first- and second-class directions, so \(\det Q=0\) does not imply that all listed constraints are first class.}
\GSfield{Reconstruction}{Regularity, weak equality and the distinction between degeneracy and complete first-class closure are supplied.}
\GSfield{Fate}{\citet[\S\S6--8]{dirac1950generalized} already distinguishes mixed systems and changes of constraint basis; \citet[pp.~343--348]{fradkinfradkina1978quantization} is the graded extension lead.}
\GSfield{Present status}{As of 10 November 2026, the corrected classification is standard; a rank test must be supplemented by the full bracket structure.}
\GSfield{Assessment}{literal statement: false as written; publication match: partial; verification: read the result; development: annotated; confidence: high}
\end{GSitem}
\begin{GSitem}{1979-475.3}{E}{Constrained stationary action and Lagrange multipliers}
\GSfield{Kind; source}{theorem; \GSpage{P000880}}
\GSfield{The notes}{A constrained Hamiltonian action is replaced by an unconstrained action with multipliers \(\lambda^\alpha T_\alpha\).}
\GSfield{Statement}{For smooth independent constraints and variations with fixed compatible endpoints, stationarity on the constraint submanifold is equivalent locally to stationarity of the multiplier action together with \(T_\alpha=0\). Constraint preservation imposes \(\dot T_\alpha=\{T_\alpha,H\}+\lambda^\beta\{T_\alpha,T_\beta\}=0\); it may determine multipliers or generate further constraints.}
\GSfield{Reconstruction}{Independence and the possibility of secondary constraints are essential; singular constraints need not admit this elementary equivalence.}
\GSfield{Fate}{\citet[\S\S2--6]{dirac1950generalized} develops the multiplier consistency algorithm before 1979.}
\GSfield{Present status}{As of 10 November 2026, the regular multiplier statement is settled.}
\GSfield{Assessment}{literal statement: correct with missing hypotheses; publication match: exact; verification: read the result; development: annotated; confidence: high}
\end{GSitem}
\begin{GSitem}{1979-475.4}{E}{Dirac bracket for second-class constraints}
\GSfield{Kind; source}{construction; \GSpage{P000880-P000881}}
\GSfield{The notes}{The inverse matrix \(Q^{-1}\) is inserted into a corrected bracket, and constraints then bracket to zero.}
\GSfield{Statement}{In the ordinary even case, \(\{f,g\}_D=\{f,g\}-\{f,T_\alpha\}(Q^{-1})^{\alpha\beta}\{T_\beta,g\}\). It is a Poisson bracket wherever \(Q\) is invertible and satisfies \(\{T_\alpha,f\}_D=0\). Its restriction to the constraint surface is independent of extensions of the restricted functions. The super version uses the corresponding graded contractions.}
\GSfield{Reconstruction}{An unambiguous ordinary-even formula is stated first; graded sign factors depend on the source's derivative convention.}
\GSfield{Fate}{\citet[\S8, eqs.~(34)--(38), and appendix]{dirac1950generalized} gives the construction and Jacobi verification in 1950.}
\GSfield{Present status}{As of 10 November 2026, the Dirac reduction bracket is established on the regular second-class locus.}
\GSfield{Assessment}{literal statement: correct with missing hypotheses; publication match: exact; verification: read the result; development: annotated; confidence: high}
\end{GSitem}
\begin{GSitem}{1979-475.5}{E}{Gelfand's geometric interpretation of the Dirac bracket}
\GSfield{Kind; source}{remark; \GSpage{P000881}; I. M. Gelfand}
\GSfield{The notes}{Gelfand identifies the construction with restricting a nondegenerate skew form to the constraint surface and inverting it there.}
\GSfield{Statement}{For a symplectic manifold \((M,\omega)\), the regular second-class surface \(C\) is symplectic: \(\iota^*\omega\) is nondegenerate. Its inverse Poisson bivector computes \(\{f|_C,g|_C\}_{C}=\{f,g\}_D|_C\).}
\GSfield{Reconstruction}{The restriction is nondegenerate because the constraints are second class, not merely because the ambient form is nondegenerate.}
\GSfield{Fate}{This is the geometric content of the earlier reduction in \citet[\S8, pp.~138--140]{dirac1950generalized}; the notebook records Gelfand's explanatory intervention, not a new priority claim.}
\GSfield{Present status}{As of 10 November 2026, this intrinsic interpretation is standard.}
\GSfield{Assessment}{literal statement: correct with missing hypotheses; publication match: exact; verification: read the result; development: annotated; confidence: high}
\end{GSitem}
\begin{GSitem}{1979-475.6}{E}{Reduced equations and local canonical coordinates}
\GSfield{Kind; source}{construction; \GSpage{P000881}}
\GSfield{The notes}{The multipliers are solved using \(Q^{-1}\); the reduced phase space has variables \(q^*,p^*\), a physical Hamiltonian and a reduced canonical one-form.}
\GSfield{Statement}{For \(2m\) independent second-class constraints in a \(2n\)-dimensional phase space, the reduced space has dimension \(2(n-m)\). Hamilton's equations are \(\dot f=\{f,H\}_D|_C\), and Darboux coordinates exist locally on \(C\). The restricted canonical one-form differs from \(p^*dq^*\) by a locally exact form; an action equality therefore needs an endpoint term or fixed endpoint convention.}
\GSfield{Reconstruction}{The number of constraints is made even in the ordinary bosonic case, and the exact differential omitted by the notes is restored.}
\GSfield{Fate}{\citet[\S8, eqs.~(36)--(38)]{dirac1950generalized} gives the reduced bracket and dynamics. Its local canonical interpretation is classical; it asserts no global Darboux chart.}
\GSfield{Present status}{As of 10 November 2026, regular local reduction is settled; global coordinates and singular strata are separate issues.}
\GSfield{Assessment}{literal statement: correct with missing hypotheses; publication match: partial; verification: read the result; development: annotated; confidence: high}
\end{GSitem}
\begin{GSitem}{1979-475.7}{E}{Canonical commutators and the limits of quantization}
\GSfield{Kind; source}{method; \GSpage{P000881-P000882}}
\GSfield{The notes}{The reduced variables are quantized by canonical (super)commutators, and \(\dot f\) is replaced by a commutator with the Hamiltonian.}
\GSfield{Statement}{Canonical quantization prescribes the Heisenberg or Clifford relations for elementary even or odd variables and requires a representation, domains, ordering and a Hamiltonian operator. For a time-independent observable \(\widehat f\), \(\dot{\widehat f}=(i\hbar)^{-1}[\widehat f,\widehat H]\) in Heisenberg conventions. This does not imply a bracket-preserving quantization of every classical function.}
\GSfield{Reconstruction}{The limitation prevents a universal quantization theorem from being read into a formal replacement rule.}
\GSfield{Fate}{\citet[\S9, pp.~143--145]{dirac1950generalized} explicitly warns that the classical-to-operator correspondence is not uniquely defined because of ordering.}
\GSfield{Present status}{As of 10 November 2026, the elementary canonical algebra is standard; ordering, domains and consistency remain parts of each quantization problem.}
\GSfield{Assessment}{literal statement: correct with missing hypotheses; publication match: partial; verification: read the result; development: annotated; confidence: high}
\end{GSitem}
\begin{GSitem}{1979-475.8}{E}{Time-ordered source functional and Schwinger--Dyson equations}
\GSfield{Kind; source}{construction; \GSpage{P000882}}
\GSfield{The notes}{Sources \(J,I\) couple to the reduced coordinates and momenta; time ordering and a source-dependent Hamiltonian generate functional differential equations.}
\GSfield{Statement}{For a specified Hamiltonian operator, vacuum or boundary states and time-ordering convention, differentiating the source evolution operator yields time-ordered coordinate and momentum insertions. The operator equations of motion give functional Schwinger--Dyson relations, with contact terms at coincident times fixed by the ordering prescription.}
\GSfield{Reconstruction}{States, contact terms and operator ordering are specified rather than inferred to be automatic.}
\GSfield{Fate}{The functional-integral framework is reviewed in \citet[\S1]{faddeevpopov1974gravity}. The notebook itself explicitly flags the equal-time ordering issue; it is not resolved by writing a formal differential equation.}
\GSfield{Present status}{As of 10 November 2026, the formal source identities are standard; no convergence or unique time-slicing prescription is established here.}
\GSfield{Assessment}{literal statement: correct with missing hypotheses; publication match: partial; verification: read the result; development: annotated; confidence: high}
\end{GSitem}
\begin{GSitem}{1979-475.9}{E}{Phase-space functional Fourier representation}
\GSfield{Kind; source}{construction; \GSpage{P000882-P000883}}
\GSfield{The notes}{The source equations are represented by an integral with phase \(\int(p^*\dot q^*-H_{\rm phys}+Jq^*+Ip^*)dt\); a margin discusses products at equal times.}
\GSfield{Statement}{A time-sliced phase-space path integral can represent the evolution generated by a chosen ordered Hamiltonian, with discretization matched to its operator symbol. Different prescriptions can produce contact or ordering corrections. The continuous formula is shorthand for that chosen limit or for a formal perturbation series.}
\GSfield{Reconstruction}{The symbol and discretization are additional data; the notebook's Fourier transform does not remove their necessity.}
\GSfield{Fate}{\citet[\S5]{faddeevpopov1974gravity} derives the canonical functional-integral expression in the earlier gravity setting.}
\GSfield{Present status}{As of 10 November 2026, the representation is standard with a fixed prescription; the continuum integral is not a standalone definition for every Hamiltonian.}
\GSfield{Assessment}{literal statement: correct with missing hypotheses; publication match: partial; verification: read the result; development: annotated; confidence: high}
\end{GSitem}
\begin{GSitem}{1979-475.10}{E}{The second-class constraint measure}
\GSfield{Kind; source}{construction; \GSpage{P000883}}
\GSfield{The notes}{The unreduced phase-space integral contains \(\delta(T)\sqrt{\operatorname{sdet}\{T,T\}}\).}
\GSfield{Statement}{For regular second-class constraints \(T_\alpha\), the induced Liouville density can be represented in ambient coordinates by \(\delta(T)\sqrt{\det Q}\,dp\,dq\) in the ordinary even case, with an orientation or positive square-root convention. The super version uses the square root of the Berezinian and compatible odd delta conventions. The phase is the original Hamiltonian action restricted to the surface.}
\GSfield{Reconstruction}{The local regular density identity and branch/orientation conditions are specified. This geometric measure identity is distinct from proving existence of the resulting oscillatory path integral.}
\GSfield{Fate}{The ordinary construction is the publication lead \citet[pp.~227--261]{senjanovic1976path}; the graded extension is \citet[pp.~343--348]{fradkinfradkina1978quantization}. Both predate the lecture. The original full texts were not recovered in this audit, so no internal equation locator or independently read proof is claimed.}
\GSfield{Present status}{As of 10 November 2026, the regular measure formula is standard; the notebook does not establish a nonperturbative measure or cover singular constraints.}
\GSfield{Assessment}{literal statement: correct with missing hypotheses; publication match: partial; verification: indirect; development: annotated; confidence: medium}
\end{GSitem}
\begin{GSitem}{1979-475.11}{E}{Gauge fixing converts a regular first-class system to second class}
\GSfield{Kind; source}{construction; \GSpage{P000883}}
\GSfield{The notes}{The first-class constraints are supplemented by gauge conditions \(\Phi_\alpha=0\); the combined bracket matrix is asserted to be nondegenerate.}
\GSfield{Statement}{For \(m\) independent first-class constraints \(T_\alpha\), choose \(m\) local conditions \(\Phi_\alpha\) such that \(M_{\alpha\beta}=\{\Phi_\alpha,T_\beta\}|_{T=\Phi=0}\) is invertible. The combined \(2m\) constraints are then second class; on their surface the square root of the block determinant gives \(|\det M|\) in the bosonic positive-density convention. This realizes a local reduced phase-space integral.}
\GSfield{Reconstruction}{The nonzero mixed determinant is a hypothesis, not a consequence of arbitrarily naming conditions a gauge. No global cross-section is promised.}
\GSfield{Fate}{\citet[\S8]{dirac1950generalized} treats supplementary constraints; the path-integral determinant is part of \citet[\S5]{faddeevpopov1974gravity}.}
\GSfield{Present status}{As of 10 November 2026, the regular local construction is established; Gribov copies and singular actions require further treatment.}
\GSfield{Assessment}{literal statement: correct with missing hypotheses; publication match: exact; verification: read the result; development: annotated; confidence: high}
\end{GSitem}

\GStalk{477}{1979-11-19}{I. M. Gelfand}{P000884}{Dilogarithms and six proposed connections}
The dated introductory page asks where dilogarithms occur and lists six areas, then writes an integral and a differential five-term relation. The catalogue attributes this introduction to Gelfand, although the page itself says only Gelfand seminar. In the actual notebook it precedes Beilinson's P000885--P000886 contribution (record 476); the catalogue numbers reverse that order. The list is a research map, not six fully formulated conjectures.

\begin{GSitem}{1979-477.1}{E}{The notebook's dilogarithm integral}
\GSfield{Kind; source}{definition; \GSpage{P000884}}
\GSfield{The notes}{The displayed function is \(F(\xi)=\int_0^\xi\log x\,dx/(1-x)\).}
\GSfield{Statement}{For \(0<\xi<1\), \(F(\xi)=-\log\xi\log(1-\xi)-\operatorname{Li}_2(\xi)=\operatorname{Li}_2(1-\xi)-\pi^2/6\). Thus this is a standard dilogarithmic variant, not the usual defining integral \(-\int_0^\xi\log(1-x)\,dx/x\) for \(\operatorname{Li}_2\). Complex continuation needs a domain and logarithm branches.}
\GSfield{Reconstruction}{A real interval is chosen to make the literal integral and signs unambiguous; the equivalence follows by differentiation and the value at zero.}
\GSfield{Fate}{The classical identities are reviewed in \citet[Chapter I, \S1]{zagier2007dilogarithm}; this later exposition does not imply a 2007 discovery.}
\GSfield{Present status}{As of 10 November 2026, the integral identity is elementary and settled.}
\GSfield{Assessment}{literal statement: correct with missing hypotheses; publication match: exact; verification: read the result; development: annotated; confidence: high}
\end{GSitem}
\begin{GSitem}{1979-477.2}{E}{The differential five-term relation needs related arguments}
\GSfield{Kind; source}{theorem; \GSpage{P000884}}
\GSfield{The notes}{The page writes \(\sum_{j=1}^5[\log u_j/(1-u_j)+\log(1-u_j)/u_j]du_j=0\), without relations among the five arguments.}
\GSfield{Statement}{Set \(L(u)=\operatorname{Li}_2(u)+\tfrac12\log u\log(1-u)\). For \(0<x,y<1\), the five arguments \(x,y,(1-x)/(1-xy),1-xy,(1-y)/(1-xy)\) obey \(\sum L(u_j)=3L(1)=\pi^2/2\). Differentiating gives exactly the displayed one-form identity, since its summand is \(-2dL(u_j)\). It is false for five independently varying arguments.}
\GSfield{Reconstruction}{The two-parameter substitution, range and branches are supplied; they cannot be recovered uniquely from the written equation alone.}
\GSfield{Fate}{\citet[Chapter II, \S1.A, Rogers-dilogarithm five-term identity]{zagier2007dilogarithm} gives this precise classical identity.}
\GSfield{Present status}{As of 10 November 2026, the corrected Rogers identity is established; the omitted relations are a source limitation.}
\GSfield{Assessment}{literal statement: correct with missing hypotheses; publication match: partial; verification: read the result; development: annotated; confidence: high}
\end{GSitem}
\begin{GSitem}{1979-477.3}{F}{Dilogarithms in algebraic K-theory}
\GSfield{Kind; source}{remark; \GSpage{P000884}}
\GSfield{The notes}{The first entry is K-theory, with Bloch and Beilinson named.}
\GSfield{Statement}{The following contribution supplies a concrete instance: a dilogarithmic primitive proves the Steinberg relation for the real regulator of \(K_2\) of a complex curve. This makes one precise connection, without interpreting the list as a theorem about every K-group.}
\GSfield{Reconstruction}{The mathematical example is supplied from the immediately following talk; the introductory phrase has no independent quantifiers.}
\GSfield{Fate}{\citet[\S1, Lemma 1 and the paragraph after it]{beilinson1980regulators} identifies the complex regulator and its real part, referring explicitly to Bloch's 1978 Irvine lectures. The lecture's precise construction is catalogued separately in record 476.}
\GSfield{Present status}{As of 10 November 2026, this connection is well established; the bare list does not identify a further theorem or open problem.}
\GSfield{Assessment}{literal statement: not determinable; publication match: thematic only; verification: read the result; development: annotated; confidence: medium}
\end{GSitem}
\begin{GSitem}{1979-477.4}{E}{Combinatorial characteristic classes: the Gelfand--Gabrielov--Losik lead}
\GSfield{Kind; source}{remark; \GSpage{P000884}}
\GSfield{The notes}{The second entry names combinatorial characteristic classes and the initials G.-G.-L.}
\GSfield{Statement}{For a smooth manifold with a linearizable decomposition into simple cells satisfying the configuration-space connectivity condition (A), the Gelfand--Gabrielov--Losik construction produces a rational codimension-four cycle Poincar\'e dual to the first Pontryagin class. A real dilogarithmic variant enters this combinatorial construction. The notebook does not specify its normalization or a formula for a higher characteristic class.}
\GSfield{Reconstruction}{This is a literature-supported identification of the named connection, with the original theorem's hypotheses retained. The source does not record enough to choose the dilogarithm's normalization or infer a new general characteristic-class claim.}
\GSfield{Fate}{The precise cycle theorem is \citet[\S1, Proposition 10 and Theorem 1, Russian p.~18]{gabrielovgelfandlosik1975computation}, published in 1975; receipt on 2 December 1974 is not a verified public circulation date. The later local formulation is \citet[Theorems 1--3]{gabrielovgelfandlosik1976local}. \citet[\S3, pp.~115--116]{goncharov1993classical} explicitly identifies the dilogarithmic ingredient in the 1975 construction.}
\GSfield{Present status}{As of 10 November 2026, the named connection is established and predates the talk; no additional conjecture is recoverable from the list.}
\GSfield{Assessment}{literal statement: not determinable; publication match: partial; verification: read the result; development: annotated; confidence: medium}
\end{GSitem}
\begin{GSitem}{1979-477.5}{E}{Hyperbolic tetrahedral volume}
\GSfield{Kind; source}{remark; \GSpage{P000884}}
\GSfield{The notes}{The third entry mentions the volume of a non-Euclidean tetrahedron and Gauss.}
\GSfield{Statement}{For an oriented ideal hyperbolic tetrahedron whose vertices on the sphere at infinity have cross-ratio \(z\), its signed volume is \(D(z)=\operatorname{Im}\operatorname{Li}_2(z)+\log|z|\arg(1-z)\). Orientation reversal changes the sign. A general finite tetrahedron needs a more elaborate formula.}
\GSfield{Reconstruction}{Ideal vertices, orientation and cross-ratio are additional specifications. The source's unqualified tetrahedron is not silently identified with this special case.}
\GSfield{Fate}{\citet[Chapter I, \S4, eq.~(4.1)]{zagier2007dilogarithm} states the ideal-tetrahedron formula and its classical historical setting. This is a precise special case of the notebook's broad reference; no priority claim for Gauss alone is inferred.}
\GSfield{Present status}{As of 10 November 2026, the ideal-volume identity is established; this page does not select a general tetrahedral-volume formula.}
\GSfield{Assessment}{literal statement: correct with missing hypotheses; publication match: partial; verification: read the result; development: annotated; confidence: high}
\end{GSitem}
\begin{GSitem}{1979-477.6}{F}{Apéry's zeta(3) result and an unspecified dilogarithm link}
\GSfield{Kind; source}{remark; \GSpage{P000884}}
\GSfield{The notes}{The fourth entry is simply \(\zeta(3)\), with Apéry's name.}
\GSfield{Statement}{Apéry's theorem asserts \(\zeta(3)\notin\mathbb Q\). The notebook gives no integral, recurrence, special-value identity or reduction of that theorem to a dilogarithm; those are not reconstructed as an additional claim.}
\GSfield{Reconstruction}{The established theorem and the unidentified link are deliberately separated.}
\GSfield{Fate}{\citet[pp.~11--13]{apery1979irrationalite} is the primary 1979 publication from the 1978 arithmetic meeting. It confirms the irrationality result but does not identify the particular dilogarithmic connection intended by this list.}
\GSfield{Present status}{As of 10 November 2026, Apéry's irrationality theorem is settled; the content of the listed connection is not sufficiently specified to classify as an open problem.}
\GSfield{Assessment}{literal statement: not determinable; publication match: thematic only; verification: read the result; development: annotated; confidence: medium}
\end{GSitem}
\begin{GSitem}{1979-477.7}{E}{Web geometry and Bol}
\GSfield{Kind; source}{remark; \GSpage{P000884}}
\GSfield{The notes}{The fifth entry names webs and Bol, with no formula or web specified.}
\GSfield{Statement}{The classical Bol five-web is the germ, at a generic point of \(\mathbb C^2\), of the foliations with first integrals \(x,y,x/y,(1-y)/(1-x),x(1-y)/(y(1-x))\). It has six independent Abelian relations, the maximal rank for a planar five-web, and is not locally equivalent to an algebraic web. The additional relation beyond those generated by three-subwebs is the dilogarithmic five-term identity.}
\GSfield{Reconstruction}{The particular web and its generic-point hypothesis come from the checked literature. This is a strong identification of the named example, not evidence that all of its hypotheses or rank claims were spoken in the seminar.}
\GSfield{Fate}{\citet[\S6.1.4, p.~190; introduction p.~6]{pereirapirio2009web} states the foliations and relation and gives the nonalgebraizability argument; its bibliography attributes the example to Bol's 1936 paper. The original 1936 proof was not separately recovered. This explicit older result identifies the named connection, without converting the notebook list into a classification problem.}
\GSfield{Present status}{As of 10 November 2026, Bol's exceptional web is established; the entry records no definite new open question.}
\GSfield{Assessment}{literal statement: not determinable; publication match: partial; verification: read the result; development: annotated; confidence: medium}
\end{GSitem}
\begin{GSitem}{1979-477.8}{F}{Scattering theory: an unqualified connection}
\GSfield{Kind; source}{remark; \GSpage{P000884}}
\GSfield{The notes}{A sixth entry in the right margin reads scattering theory.}
\GSfield{Statement}{No scattering equation, dimension, observable, loop order or dilogarithmic identity is stated. Consequently the remark has no determinate mathematical proposition beyond suggesting an area of application.}
\GSfield{Fate}{Primary-source searches combining dilogarithms with scattering and amplitudes yield many inequivalent applications. None can be designated the fate of this particular remark without a model or formula.}
\GSfield{Present status}{As of 10 November 2026, the intended application remains unidentified; no open scattering problem is manufactured from the label.}
\GSfield{Assessment}{literal statement: not determinable; publication match: none found; verification: unresolved lead; development: annotated; confidence: low}
\end{GSitem}

\GStalk{476}{1979-11-19}{A. A. Beilinson}{P000885-P000886}{Dilogarithms, tame symbols and regulators of curves}
Beilinson explains the real dilogarithmic regulator, its Steinberg identity and tame-symbol periods, then refines it to a complex multiplicative regulator. The latter matches his short 1980 publication closely. The opening remark about zeta functions has no sufficiently precise arithmetic conjecture. Source errors include the sign of the zeta(2) integral, a missing period normalization, and calling a multiplicative skew symbol symmetric. The distinction between K-theory of the curve and its function field is essential.

\begin{GSitem}{1979-476.1}{E}{The zeta(2) integral has a missing minus sign}
\GSfield{Kind; source}{theorem; \GSpage{P000885}}
\GSfield{The notes}{Equation (3) states \(\int_0^1\log x\,dx/(1-x)=\sum 1/n^2\), with no visible minus sign in the scan.}
\GSfield{Statement}{The correct identity is \(\int_0^1\log x\,dx/(1-x)=-\sum_{n\ge1}n^{-2}=-\pi^2/6\). Expanding \((1-x)^{-1}\) and applying monotone convergence to the nonnegative function \(-\log x/(1-x)\) justifies termwise integration.}
\GSfield{Reconstruction}{The lower limit on the sum and the minus sign are supplied; the elementary justification is editorial.}
\GSfield{Fate}{This is a classical dilogarithm evaluation; compare \citet[Chapter I, \S1]{zagier2007dilogarithm}. Its sign is also forced by the negative integrand on \((0,1)\).}
\GSfield{Present status}{As of 10 November 2026, the corrected identity is settled; both transcriptions faithfully retain the notebook's sign error.}
\GSfield{Assessment}{literal statement: false as written; publication match: partial; verification: read the result; development: annotated; confidence: high}
\end{GSitem}
\begin{GSitem}{1979-476.2}{F}{Bloch's proposed replacement of zeta(2) by a curve zeta function}
\GSfield{Kind; source}{remark; \GSpage{P000885}}
\GSfield{The notes}{The introduction says generalization to curves, names Bloch, and replaces the right-hand zeta value by the curve's zeta function.}
\GSfield{Statement}{No base number field, arithmetic model, value of the complex parameter, order of vanishing, regulator determinant or rational factor is recorded. There is therefore no uniquely recoverable special-value formula or conjecture.}
\GSfield{Fate}{\citet[\S1 and bibliography item 1]{beilinson1980regulators} explicitly relates the real regulator to Bloch's 1978 Irvine lectures. Its later \S2 concerns higher regulators and L-values, but that is not evidence that all those conjectural quantifiers were stated here.}
\GSfield{Present status}{As of 10 November 2026, the historical regulator context is identifiable; the vague zeta remark remains unsuitable for an independent problem record.}
\GSfield{Assessment}{literal statement: not determinable; publication match: thematic only; verification: read the result; development: annotated; confidence: medium}
\end{GSitem}
\begin{GSitem}{1979-476.3}{E}{A closed real regulator one-form}
\GSfield{Kind; source}{construction; \GSpage{P000885}}
\GSfield{The notes}{For nowhere-zero holomorphic functions on a Riemann surface the notes define \(\omega_{f,g}=\log|f|d\arg g-\log|g|d\arg f\) and infer that integration along paths is well defined.}
\GSfield{Statement}{On a Riemann surface where \(f,g\) are holomorphic units, \(d\arg f=\operatorname{Im}(df/f)\) is globally defined and \(d\omega_{f,g}=0\). Hence its integrals depend only on the relative homology class of a path with fixed endpoints; closedness does not assert independence of every path between those endpoints.}
\GSfield{Reconstruction}{The global one-form is distinguished from a globally chosen argument function; path homotopy or homology is the required qualification.}
\GSfield{Fate}{The real regulator is attributed to Bloch's 1978 lectures in \citet[\S1, paragraph following Lemma 1]{beilinson1980regulators}. The displayed differential identity also follows directly from the Cauchy--Riemann equations.}
\GSfield{Present status}{As of 10 November 2026, the closed-form construction is established; nonzero periods are precisely part of its regulator content.}
\GSfield{Assessment}{literal statement: correct with missing hypotheses; publication match: partial; verification: read the result; development: annotated; confidence: high}
\end{GSitem}
\begin{GSitem}{1979-476.4}{E}{Bilinearity and alternation of the real regulator form}
\GSfield{Kind; source}{theorem; \GSpage{P000885}}
\GSfield{The notes}{Equations (6)--(7) give additivity in a product of functions and antisymmetry under interchanging them.}
\GSfield{Statement}{For holomorphic units, \(\omega_{f_1f_2,g}=\omega_{f_1,g}+\omega_{f_2,g}\), \(\omega_{f,g_1g_2}=\omega_{f,g_1}+\omega_{f,g_2}\), and \(\omega_{g,f}=-\omega_{f,g}\). These are equalities of differential forms, not merely of cohomology classes.}
\GSfield{Fate}{These elementary identities underlie Bloch's real regulator, recalled by \citet[\S1]{beilinson1980regulators}; they are already implicit in the logarithmic formula.}
\GSfield{Present status}{As of 10 November 2026, the identities are exact on the common domain of invertibility.}
\GSfield{Assessment}{literal statement: correct; publication match: exact; verification: read the result; development: annotated; confidence: high}
\end{GSitem}
\begin{GSitem}{1979-476.5}{E}{The single-valued dilogarithm proves the Steinberg relation}
\GSfield{Kind; source}{theorem; \GSpage{P000885}}
\GSfield{The notes}{The form \(\omega_{f,1-f}\) is asserted to be exact. Equation (9) writes its primitive using a logarithm term and the imaginary part of an integral, and emphasizes single-valuedness.}
\GSfield{Statement}{The Bloch--Wigner dilogarithm \(D(z)=\operatorname{Im}\operatorname{Li}_2(z)+\log|z|\arg(1-z)\) is single valued on \(\mathbb C\setminus\{0,1\}\), extends continuously by zero at \(0,1,\infty\), and satisfies \(dD(z)=\log|z|d\arg(1-z)-\log|1-z|d\arg z\). Consequently \(\omega_{f,1-f}=d(D\circ f)\) whenever \(f,1-f\) are holomorphic units.}
\GSfield{Reconstruction}{Individual logarithms and the complex dilogarithm require local branches; their stated combination is the global function. The lower endpoint zero is understood by its continuous limit.}
\GSfield{Fate}{\citet[Chapter I, \S3]{zagier2007dilogarithm} records the precise single-valued function and differential. The pre-1979 Bloch regulator context is confirmed in \citet[\S1]{beilinson1980regulators}.}
\GSfield{Present status}{As of 10 November 2026, this exactness identity and hence the real Steinberg relation are established.}
\GSfield{Assessment}{literal statement: correct with missing hypotheses; publication match: exact; verification: read the result; development: annotated; confidence: high}
\end{GSitem}
\begin{GSitem}{1979-476.6}{E}{David Wigner and continuous cohomology of SL(2)}
\GSfield{Kind; source}{remark; \GSpage{P000885}}
\GSfield{The notes}{The notebook attributes the single-valued function to D. Wigner and says he connected it to continuous cohomology of \(SL(2)\).}
\GSfield{Statement}{Fix \(x\in\mathbb P^1(\mathbb C)\). For generic quadruples, put
\[
c_x(g_0,g_1,g_2,g_3)
=D\bigl(\operatorname{cr}(g_0x,g_1x,g_2x,g_3x)\bigr).
\]
The five-term identity makes this a homogeneous measurable real three-cocycle for \(GL_2(\mathbb C)\), restricting to \(SL_2(\mathbb C)\). It represents the nonzero degree-three volume class under the comparison between measurable and continuous cohomology. Changing \(x\) gives a cohomologous cocycle.}
\GSfield{Reconstruction}{The complex group, real coefficients, degree, generic quadruples and comparison map are supplied from the literature. This identifies a standard precise version of the note; the notebook itself suppresses these data.}
\GSfield{Fate}{\citet[\S3, p.~117, eqs.~(10)--(12)]{goncharov1993classical} explicitly attributes this construction to Wigner and places it in the 1970s. The 1993 chapter is the checked exposition, not the first-circulation date of Wigner's observation; no original Wigner preprint date was verified.}
\GSfield{Present status}{As of 10 November 2026, the volume-cocycle interpretation is established; the source's omission of field and degree is recorded rather than treated as an open problem.}
\GSfield{Assessment}{literal statement: correct with missing hypotheses; publication match: partial; verification: read the result; development: annotated; confidence: medium}
\end{GSitem}
\begin{GSitem}{1979-476.7}{E}{A puncture period and the tame symbol}
\GSfield{Kind; source}{theorem; \GSpage{P000885}}
\GSfield{The notes}{Equation (10) equates a small-loop integral of \(\omega_{f,g}\) with \(\log|f^{v_x(g)}/g^{v_x(f)}|(x)\), without a factor \(2\pi\).}
\GSfield{Statement}{Let \(f=z^m a(z)\), \(g=z^n b(z)\) near \(x\), with \(a(0)b(0)\ne0\). For a positively oriented small circle \(\gamma_x\), \(\int_{\gamma_x}\omega_{f,g}=2\pi\log|a(0)^n/b(0)^m|=2\pi\log|\partial_x\{f,g\}|\), where \(\partial_x\{f,g\}=(-1)^{mn}a(0)^n/b(0)^m\).}
\GSfield{Reconstruction}{Orientation, local expansions and the tame sign are specified. The tame sign disappears under absolute value but not in the later complex symbol.}
\GSfield{Fate}{The complex local tame-symbol identity is \citet[\S1, Lemma 1(b)]{beilinson1980regulators}; taking logarithmic absolute values gives the corrected real normalization.}
\GSfield{Present status}{As of 10 November 2026, the corrected period identity is exact. The missing factor is a faithful notebook error, not a translation issue.}
\GSfield{Assessment}{literal statement: false as written; publication match: partial; verification: read the result; development: annotated; confidence: high}
\end{GSitem}
\begin{GSitem}{1979-476.8}{E}{Symbols present K2 of the function field, not automatically of the curve}
\GSfield{Kind; source}{definition; \GSpage{P000885-P000886}}
\GSfield{The notes}{The notes quotient symbols by bimultiplicativity, antisymmetry and the Steinberg relation, then call the result both \(K_2(X)\) and K2 of its rational function field.}
\GSfield{Statement}{For a field \(F\), Matsumoto's theorem identifies Quillen \(K_2(F)\) with the abelian group on \(\{f,g\}\), \(f,g\in F^\times\), subject to bimultiplicativity and \(\{f,1-f\}=0\) for \(f\ne0,1\). Skew symmetry follows from these relations. Taking \(F=\mathbb C(X)\) defines the function-field group; it does not by itself define \(K_2(X)\).}
\GSfield{Reconstruction}{The ambient free group is made abelian, and the curve and function-field notation are separated.}
\GSfield{Fate}{The result dates to Matsumoto's 1969 theorem, stated precisely in \citet[III.6.1]{weibel2013kbook}. The notebook's Milnor/Steinberg attribution is compatible with this older framework.}
\GSfield{Present status}{As of 10 November 2026, the field presentation is established; the missing passage to the curve is supplied by localization in the following item.}
\GSfield{Assessment}{literal statement: false as written; publication match: partial; verification: read the result; development: annotated; confidence: high}
\end{GSitem}
\begin{GSitem}{1979-476.9}{E}{The unique smooth projective completion of an algebraic curve}
\GSfield{Kind; source}{theorem; \GSpage{P000885}}
\GSfield{The notes}{A parenthesis says an algebraic curve is compactified uniquely, for example by embedding it in projective space and taking the closure.}
\GSfield{Statement}{A smooth connected algebraic curve over \(\mathbb C\) has a smooth projective completion unique up to the unique isomorphism extending the given open curve. The normalization of a projective closure gives that smooth completion. A raw projective closure can be singular; arbitrary open Riemann surfaces need not be algebraic or of finite type.}
\GSfield{Reconstruction}{Smoothness, algebraicity, normalization and uniqueness relative to the given open immersion are the missing hypotheses. The preceding general Riemann surface is not silently assumed algebraic.}
\GSfield{Fate}{This is the classical smooth-curve model used in the regulator construction of \citet[\S1]{beilinson1980regulators}; it is not a new theorem announced by this talk.}
\GSfield{Present status}{As of 10 November 2026, the corrected completion statement is standard; the unrestricted topological or analytic reading is false.}
\GSfield{Assessment}{literal statement: correct with missing hypotheses; publication match: partial; verification: read the result; development: annotated; confidence: high}
\end{GSitem}
\begin{GSitem}{1979-476.10}{E}{Localization and extension of the real regulator over punctures}
\GSfield{Kind; source}{construction; \GSpage{P000886}}
\GSfield{The notes}{Equation (12) writes \(K_2(X)\to K_2(\mathbb C(X))\to\bigoplus_x\mathbb C^*\). A sum of symbols with trivial local terms then gives a class in \(H^1(X,\mathbb R)\).}
\GSfield{Statement}{For a smooth algebraic curve, localization identifies the image of \(K_2(X)\to K_2(\mathbb C(X))\) with the kernel of the sum of tame symbols. If \(\xi=\sum_i n_i\{f_i,g_i\}\) has trivial tame symbol at every point of \(X\), the closed form \(\sum_i n_i\omega_{f_i,g_i}\), initially off its divisors, has zero small-loop periods and determines a unique class in \(H^1(X,\mathbb R)\). This defines the real regulator on \(K_2(X)\) by composition.}
\GSfield{Reconstruction}{The statement concerns an image/kernel exactness assertion. The notes do not supply a proof of injectivity of the first arrow, and no such assertion is needed for the regulator construction. Uniqueness is uniqueness of the extended cohomology class, not of an arbitrary differential form.}
\GSfield{Fate}{\citet[\S1, immediately after Lemma 1]{beilinson1980regulators} explicitly uses Quillen's localization theorem and identifies the resulting logarithmic absolute value with Bloch's regulator, already circulated in his 1978 lectures.}
\GSfield{Present status}{As of 10 November 2026, the regular-curve regulator construction is established; it makes no special-value conjecture without further arithmetic data.}
\GSfield{Assessment}{literal statement: correct with missing hypotheses; publication match: exact; verification: read the result; development: annotated; confidence: high}
\end{GSitem}
\begin{GSitem}{1979-476.11}{A}{Beilinson's complex multiplicative regulator}
\GSfield{Kind; source}{construction; \GSpage{P000886}}
\GSfield{The notes}{Equation (16) defines an exponentiated logarithmic integral in cases according to the winding number of \(g\), with a sign \((-1)^{v_\gamma(f)v_\gamma(g)}\). The symbol is said to be bimultiplicative, symmetric and Steinberg, and to recover the real construction.}
\GSfield{Statement}{For a loop \(\gamma\) based at \(a\), avoid the divisors and continue chosen logarithms along the loop. A convenient equivalent formula is \[r_\gamma(f,g)=\exp\!\left[\frac{1}{2\pi i}\left(\int_\gamma\log f\,d\log g-\log g(a)\int_\gamma d\log f\right)\right].\] It is independent of the initial logarithm choices, homotopy invariant, bimultiplicative, multiplicatively skew, and satisfies \(r_\gamma(f,1-f)=1\). On a puncture circle it is the tame symbol. Moreover \(\log|r_\gamma(f,g)|=(2\pi)^{-1}\int_\gamma\omega_{f,g}\), so unramified symbols yield the complex refinement of the real regulator.}
\GSfield{Reconstruction}{The basepoint formula makes branch normalization explicit and avoids the notebook's case division. Multiplicative skew symmetry means \(r_\gamma(g,f)=r_\gamma(f,g)^{-1}\); literal symmetry would generally be false. Taking logarithmic absolute value, with the stated factor, is the recovery of the real regulator.}
\GSfield{Fate}{This is an exact construction match to \citet[\S1, eq.~(1) and Lemma 1]{beilinson1980regulators}, published in 1980. The journal records receipt on 21 December 1979; no earlier public preprint date was verified. The notebook documents the seminar on 19 November 1979. The paper also acknowledges an independent unpublished construction by Deligne.}
\GSfield{Present status}{As of 10 November 2026, this regulator refinement was published soon after the seminar and is established. It is not merely a retrospective thematic link to later Beilinson conjectures.}
\GSfield{Assessment}{literal statement: false as written; publication match: exact; verification: read the result; development: annotated; confidence: high}
\end{GSitem}
\begin{GSitem}{1979-476.12}{E}{Stable Steinberg extension and second group homology}
\GSfield{Kind; source}{theorem; \GSpage{P000886}}
\GSfield{The notes}{For an affine curve the notes write \(K_2(X)=H_2(SL_n(A),\mathbb Z)\) and \(1\to K_2(X)\to St^X\to SL^X\to1\), where \(A\) is its ring of algebraic functions.}
\GSfield{Statement}{For a unital ring \(A\), the stable Steinberg group gives the universal central extension \(1\to K_2(A)\to\operatorname{St}(A)\to E(A)\to1\), and \(K_2(A)\cong H_2(E(A),\mathbb Z)\). If \(X=\operatorname{Spec}A\), its algebraic K2 is \(K_2(A)\). Replacing the stable elementary group by a particular \(SL_n(A)\) requires both the relevant equality with the elementary group and homology/K2 stability hypotheses.}
\GSfield{Reconstruction}{The stable rank and elementary-group qualifications are essential. They are supplied rather than choosing an unjustified finite rank for the source.}
\GSfield{Fate}{The Kervaire--Steinberg theorem is stated in \citet[III.5.5; stability discussion III.5.5.2]{weibel2013kbook}, which traces the older Milnor/Steinberg theory. It does not assert the displayed equality for arbitrary \(n\).}
\GSfield{Present status}{As of 10 November 2026, the stable statement is established; the source's unspecified finite-rank formulation is not a complete theorem.}
\GSfield{Assessment}{literal statement: correct with missing hypotheses; publication match: partial; verification: read the result; development: annotated; confidence: high}
\end{GSitem}

\GStalk{478}{1979-11-19}{I. A. Batalin}{P000887-P000888}{Hamiltonian BRST quantization in a relativistic gauge}
The second Batalin contribution on 19 November extends phase space by multiplier momenta and ghosts, introduces an odd nilpotent BRST charge and a gauge fermion, and sketches gauge independence by a field-dependent BRST change of variables. This is the contemporaneous Hamiltonian BFV framework, whose even Poisson bracket must be distinguished from the later BV antibracket. The statement about termination of supergravity ghost expansions remains model dependent.

\begin{GSitem}{1979-478.1}{E}{Promoting Lagrange multipliers to canonical variables}
\GSfield{Kind; source}{construction; \GSpage{P000887}}
\GSfield{The notes}{The multiplier \(\lambda\) is paired with a momentum \(\pi\), and the extended constraints are \(G_a=(\pi_\mu,T_\mu)\). Gauge conditions include \(-\dot\lambda^\mu+\chi^\mu\). The count is \((n+m)-2m=n-m\).}
\GSfield{Statement}{For a regular irreducible bosonic system with \(n\) canonical pairs and \(m\) independent first-class constraints, add \(m\) multiplier pairs. The extended phase space has dimension \(2(n+m)\), and its \(2m\) first-class constraints, when regularly reduced, leave dimension \(2(n-m)\). Differential gauge conditions for the multipliers can be encoded by a suitable extended Hamiltonian. This is a local degree-of-freedom count, not a proof that arbitrary gauges are admissible.}
\GSfield{Reconstruction}{The count refers to independent regular bosonic constraints. The overwritten sign and momentum index in source equation (4) are not silently repaired; the clear equations (6)--(7) determine the construction.}
\GSfield{Fate}{The contemporaneous publication leads are \citet[pp.~309--312]{batalinvilkovisky1977relativistic} and \citet[pp.~343--348]{fradkinfradkina1978quantization}. The later Hamiltonian framework is explicitly described in \citet[introduction and \S~II]{bering1997geometry}.}
\GSfield{Present status}{As of 10 November 2026, the regular extended-phase-space construction is established; the literal ambiguous formula (4) remains a transcription uncertainty.}
\GSfield{Assessment}{literal statement: correct with missing hypotheses; publication match: partial; verification: indirect; development: annotated; confidence: medium}
\end{GSitem}
\begin{GSitem}{1979-478.2}{E}{Ghost canonical pairs and the BFV phase-space action}
\GSfield{Kind; source}{construction; \GSpage{P000887}}
\GSfield{The notes}{Equation (8) integrates \((q^A,p_A,\eta^a,\mathcal P_a)\) with phase \(\int(p_A\dot q^A+\mathcal P_a\dot\eta^a-H_{\rm complete})dt\).}
\GSfield{Statement}{For each constraint \(G_a\), introduce a ghost \(\eta^a\) of opposite parity and its conjugate momentum \(\mathcal P_a\), with ghost numbers \(+1\) and \(-1\). The full symplectic structure is even. With a chosen time-slicing, boundary states and integration prescription the displayed phase-space expression is the formal BFV partition functional.}
\GSfield{Reconstruction}{Ghost number, parity and boundary prescriptions are supplied. Odd ghost variables do not make the ambient Poisson bracket odd, and no BV antifields have been introduced.}
\GSfield{Fate}{The boson/fermion extension belongs to the 1977/1978 Hamiltonian publications just cited. \citet[\S~II and \S~V.2, eq.~(V.8)]{bering1997geometry} writes the later precise geometric version and states its assumptions.}
\GSfield{Present status}{As of 10 November 2026, this is the standard Hamiltonian ghost extension; the displayed continuum integral remains formal without further analytic construction.}
\GSfield{Assessment}{literal statement: correct with missing hypotheses; publication match: partial; verification: read the result; development: annotated; confidence: high}
\end{GSitem}
\begin{GSitem}{1979-478.3}{E}{Nilpotent Hamiltonian charge, invariant Hamiltonian and gauge fermion}
\GSfield{Kind; source}{construction; \GSpage{P000887-P000888}}
\GSfield{The notes}{Equations (9)--(11) require \(\{\Omega,\Omega\}=0\), \(\{H_1,\Omega\}=0\), and \(H_{\rm complete}=H_1-\{\Psi,\Omega\}\), with \(\Psi=\mathcal P_a\chi^a\).}
\GSfield{Statement}{Let the odd charge \(\Omega\) have ghost number one and satisfy the even-bracket master equation \(\{\Omega,\Omega\}=0\). Let \(H_1\) be even, ghost number zero, and BRST invariant. For any admissible odd gauge fermion \(\Psi\) of ghost number minus one, the displayed gauge-fixed Hamiltonian is again BRST invariant by the graded Jacobi identity.}
\GSfield{Reconstruction}{Parities and ghost numbers make the equations precise; the sign of the gauge-fixing term is conventional once \(\Psi\) is fixed.}
\GSfield{Fate}{\citet[\S~II, eq.~(II.6), and \S~V.2]{bering1997geometry} gives this Hamiltonian structure and attributes it to the earlier BFV work. The original 1977/1978 articles are indirect publication checks, not a claim that the 1979 notes anticipated the later Lagrangian BV formalism.}
\GSfield{Present status}{As of 10 November 2026, the algebraic BRST-invariance statement is exact; admissibility and quantum gauge independence require the additional hypotheses in item 1979-478.7.}
\GSfield{Assessment}{literal statement: correct with missing hypotheses; publication match: exact; verification: read the result; development: annotated; confidence: high}
\end{GSitem}
\begin{GSitem}{1979-478.4}{E}{Solving the Hamiltonian BRST equations as a formal series}
\GSfield{Kind; source}{claim; \GSpage{P000888}}
\GSfield{The notes}{The speaker says the equations for \(\Omega\) and \(H_1\) can be solved as series in ghosts and ghost momenta whenever the constraints are first class.}
\GSfield{Statement}{For a regular irreducible first-class constraint system, with the usual acyclic constraint resolution locally, one can recursively construct a formal charge beginning \(\Omega=\eta^aG_a+\cdots\) and a BRST extension of a Hamiltonian preserving the constraints. Reducible systems need further ghosts; singular constraints need additional hypotheses. Formal solvability asserts neither convergence nor finite termination.}
\GSfield{Reconstruction}{Regularity, irreducibility or a suitable resolution, locality and formality are necessary additions to the notebook's all-first-class wording.}
\GSfield{Fate}{\citet[pp.~343--348]{fradkinfradkina1978quantization} is the contemporaneous graded-constraint publication lead. \citet[introduction]{bering1997geometry} confirms the later BFV framework but assumes, rather than proves, the existence of the charge used in that article. The original recursive proof was not recovered here.}
\GSfield{Present status}{As of 10 November 2026, the regular formal existence framework is established; an unrestricted singular-system theorem is not verified by the notes or the checked publication text.}
\GSfield{Assessment}{literal statement: correct with missing hypotheses; publication match: partial; verification: indirect; development: annotated; confidence: medium}
\end{GSitem}
\begin{GSitem}{1979-478.5}{F}{Termination of the supergravity ghost expansion}
\GSfield{Kind; source}{remark; \GSpage{P000888}}
\GSfield{The notes}{The series are said usually to terminate in supergravity.}
\GSfield{Statement}{A finite ghost-degree bound depends on the particular supergravity theory, formulation, auxiliary fields and constraint generators. None is specified, so the word usually does not define a universal termination theorem.}
\GSfield{Fate}{The contemporary modified-rule lead is \citet[pp.~141--152]{kallosh1978modified}. The later open-algebra discussion in \citet[\S2.2]{fuster2005brst} confirms the relevant phenomenon but supplies no model-independent bound matching this phrase.}
\GSfield{Present status}{As of 10 November 2026, model-specific constructions exist; this notebook remark remains too imprecise for a separate open-problem identifier.}
\GSfield{Assessment}{literal statement: not determinable; publication match: thematic only; verification: indirect; development: annotated; confidence: medium}
\end{GSitem}
\begin{GSitem}{1979-478.6}{E}{Hamiltonian BRST transformations and nilpotence}
\GSfield{Kind; source}{theorem; \GSpage{P000888}}
\GSfield{The notes}{Equation (12) writes \(\delta X=\{X,\Omega\}\mu\), with \(\mu\) an odd constant.}
\GSfield{Statement}{The differential \(sX=\{X,\Omega\}\), with a consistent right/left derivation convention, has odd parity and satisfies \(s^2=0\) when \(\{\Omega,\Omega\}=0\). A constant odd parameter gives an even infinitesimal change of integration variables. The symplectic potential changes by an endpoint term, so invariance of an action also depends on boundary conditions.}
\GSfield{Reconstruction}{Nilpotence is a property of the odd derivation itself, not merely a vacuous consequence of the square of one Grassmann parameter being zero.}
\GSfield{Fate}{\citet[\S~II, eq.~(II.6), and \S~V.2, eq.~(V.9)]{bering1997geometry} displays the nilpotent charge and the endpoint condition in the Hamiltonian setting.}
\GSfield{Present status}{As of 10 November 2026, the algebraic identity is exact, with boundary conditions needed for the integrated action.}
\GSfield{Assessment}{literal statement: correct with missing hypotheses; publication match: exact; verification: read the result; development: annotated; confidence: high}
\end{GSitem}
\begin{GSitem}{1979-478.7}{E}{Field-dependent BRST change of variables and gauge independence}
\GSfield{Kind; source}{method; \GSpage{P000888}}
\GSfield{The notes}{The last formula chooses \(\mu=i\int\mathcal P_a\,\delta\chi^a\,dt\) and concludes BRST invariance.}
\GSfield{Statement}{For an infinitesimal change of an admissible gauge fermion, a BRST change of path variables with a field-dependent odd parameter can compensate the change of gauge-fixing action through its Jacobian. This proves formal gauge independence only when the measure is BRST compatible, boundary terms vanish, the gauges are connected without a singularity, and no global orbit obstruction has been ignored.}
\GSfield{Reconstruction}{A field-dependent parameter does not generate an ordinary canonical transformation with unit Jacobian. The Jacobian, rather than constant-parameter invariance alone, is the crucial missing step.}
\GSfield{Fate}{\citet[\S~II, \S~V.2 and \S~V.3]{bering1997geometry} analyzes precisely this field-dependent path-space change and states explicitly that global Gribov problems, phase-space boundaries and existence of the path integral are excluded or assumed. Its arXiv preprint first circulated in September 1997; the mechanism itself is attributed there to the original BFV work.}
\GSfield{Present status}{As of 10 November 2026, the qualified formal BFV theorem is established; the notes do not prove a global, nonperturbative gauge-independence theorem.}
\GSfield{Assessment}{literal statement: correct with missing hypotheses; publication match: partial; verification: read the result; development: annotated; confidence: high}
\end{GSitem}

\GStalk{479}{1979-12-03}{I. Yu. Kobzarev}{P000889-P000893}{Elementary particles, gauge representations and the Higgs potential}
The speaker is called physicist in this block and identified as Kobzarev by the catalogue and the following lecture's tentative name. The talk surveys quark and lepton charges, three generations, colour and chiral weak representations, gauge-boson exchange and a gauge-invariant scalar potential. It is exposition of pre-existing physics rather than a new theorem announcement. Its massless-neutrino language and unobserved top quark are retained as 1979 context, with current status distinguished.

\begin{GSitem}{1979-479.1}{E}{First-generation charges and valence nucleon diagrams}
\GSfield{Kind; source}{construction; \GSpage{P000889-P000890}}
\GSfield{The notes}{The diagrams pair \((\nu,e)\) and \((u,d)\), assign charges \((0,-1)\) and \((2/3,-1/3)\), and draw proton \(uud\) and neutron \(udd\).}
\GSfield{Statement}{In units of the positron charge, the first-generation charge eigenvalues are \(Q(\nu_e)=0\), \(Q(e)=-1\), \(Q(u)=2/3\), \(Q(d)=-1/3\). The valence assignments give \(Q(uud)=1\) and \(Q(udd)=0\). The valence picture does not assert that a physical nucleon has only three constituents and no gluon or sea-quark degrees of freedom.}
\GSfield{Reconstruction}{The scan's circles are interpreted as valence diagrams, and charge eigenvalues are distinguished from an identity matrix acting on a doublet.}
\GSfield{Fate}{These are pre-1979 quark-model and electroweak assignments. Their current representation conventions are stated in \citet[\S10.1]{deblasdittmaierkogler2026electroweak}; the QCD field content is given in \citet[\S9.1]{hustonrabbertzzanderighi2026qcd}.}
\GSfield{Present status}{As of 10 November 2026, these charge assignments remain standard; the schematic nucleon drawing is not a full wavefunction.}
\GSfield{Assessment}{literal statement: correct with missing hypotheses; publication match: exact; verification: read the result; development: annotated; confidence: high}
\end{GSitem}
\begin{GSitem}{1979-479.2}{E}{Three quark and lepton generations}
\GSfield{Kind; source}{construction; \GSpage{P000890}}
\GSfield{The notes}{The page lists \((\nu_e,e),(\nu_\mu,\mu),(\nu_\tau,\tau)\) and \((u,d),(c,s),(t,b)\), marking the top quark as not yet found.}
\GSfield{Statement}{The three-generation model repeats the same gauge-representation pattern for three families, while masses and mixing parameters differ. The table is a field-content proposal; it is not a derivation of the number three or of the masses. Heavy charged leptons and quarks decay, but the blanket wording that every later-generation particle quickly disappears does not describe neutrino propagation.}
\GSfield{Reconstruction}{Historical non-observation is separated from present field content. Rounded handwritten masses are not used as current numerical data.}
\GSfield{Fate}{\citet[pp.~654--657]{kobayashimaskawa1973cp} already studies six-quark weak mixing in 1973. The current three-generation CKM construction is \citet[\S12.1]{ceccucciligetisakai2026ckm}; neutrino propagation and mixing are described in \citet[\S14.1]{gonzalezgarciawendell2026neutrinos}.}
\GSfield{Present status}{As of 10 November 2026, all three generations form the established phenomenological framework, including the top quark; neutrino masses require extending the massless-neutrino model described in 1979.}
\GSfield{Assessment}{literal statement: correct with missing hypotheses; publication match: partial; verification: read the result; development: annotated; confidence: high}
\end{GSitem}
\begin{GSitem}{1979-479.3}{E}{Colour triplets, singlets and eight adjoint gluons}
\GSfield{Kind; source}{definition; \GSpage{P000890-P000892}}
\GSfield{The notes}{Each quark is tripled into colour components, \(SU(3)_c\) acts by a unitary matrix, leptons are not tripled, and eight gluons are identified with the adjoint representation.}
\GSfield{Statement}{Each quark flavour transforms in the fundamental \(\mathbf3\) of \(SU(3)_c\); leptons are colour singlets. The gauge connection takes values in \(\mathfrak{su}(3)\), whose dimension is eight, and its homogeneous colour rotation is adjoint. A local gauge transformation of the connection also has an inhomogeneous derivative term.}
\GSfield{Reconstruction}{The ambiguous adjoint annotation on P000890 is attached to the gauge field as clarified on P000892; it is not used to place quarks in the eight-dimensional adjoint representation.}
\GSfield{Fate}{\citet[\S9.1, eq.~(9.1) and its definitions]{hustonrabbertzzanderighi2026qcd} states the QCD representation and curvature conventions. This is confirmation of the older framework, not a later resolution of a seminar conjecture.}
\GSfield{Present status}{As of 10 November 2026, the representation distinction is established and unchanged.}
\GSfield{Assessment}{literal statement: correct with missing hypotheses; publication match: exact; verification: read the result; development: annotated; confidence: high}
\end{GSitem}
\begin{GSitem}{1979-479.4}{E}{Chiral weak doublets, singlets and hypercharge}
\GSfield{Kind; source}{definition; \GSpage{P000891-P000892}}
\GSfield{The notes}{The weak group mixes each doublet synchronously and is written \(SU(2)\times U(1)\) on P000891, but \(SU(2)\times SU(1)\) on P000892. The right electron is called a singlet under the displayed product.}
\GSfield{Statement}{With \(Q=T_3+Y\), the lepton representations are \(L_L=(\nu_L,e_L)\sim(\mathbf2,-1/2)\) and \(e_R\sim(\mathbf1,-1)\) under \(SU(2)_L\times U(1)_Y\). The right electron is a singlet only of the \(SU(2)\) factor. Gauge transformations are continuous linear transformations, not merely permutations of particle labels. The minimal model has no right-handed neutrino field.}
\GSfield{Reconstruction}{The Abelian group and nontrivial right-electron hypercharge are restored. Spontaneous symmetry breaking, rather than unequal masses by themselves, explains why the vacuum need not exhibit the whole weak symmetry.}
\GSfield{Fate}{\citet[eqs.~(1)--(4)]{weinberg1967leptons} gives the earlier chiral model; \citet[\S10.1]{deblasdittmaierkogler2026electroweak} states modern normalizations.}
\GSfield{Present status}{As of 10 November 2026, the corrected representation assignment is standard; a right-electron singlet of the entire product would have zero electric charge.}
\GSfield{Assessment}{literal statement: false as written; publication match: partial; verification: read the result; development: annotated; confidence: high}
\end{GSitem}
\begin{GSitem}{1979-479.5}{E}{Dirac matrices, chirality and the massless limit}
\GSfield{Kind; source}{definition; \GSpage{P000891-P000892}}
\GSfield{The notes}{The talk writes the Dirac equation, the Clifford anticommutator and chiral projectors. Its displayed \(\gamma^5\) has the opposite sign to one common convention, so the projector labeled left is \((1+\gamma^5)/2\).}
\GSfield{Statement}{For \(\{\gamma^\mu,\gamma^\nu\}=2\eta^{\mu\nu}\) and \(\eta=\operatorname{diag}(1,-1,-1,-1)\), choose \(\gamma^5\) with \((\gamma^5)^2=1\) and \(\{\gamma^5,\gamma^\mu\}=0\). Then \(P_\pm=(1\pm\gamma^5)/2\) are complementary projectors. The massless Dirac equation separates into two Weyl equations; a mass term couples them. For positive-energy massless particles chirality matches helicity after fixing conventions, but for massive particles they are distinct notions.}
\GSfield{Reconstruction}{The sign of \(\gamma^5\) is treated as a convention, not a transcription mistake. A four-component notation for a Weyl neutrino does not decide whether actual massive neutrinos are Dirac or Majorana.}
\GSfield{Fate}{The chiral lepton construction is already used in \citet[eqs.~(1)--(4)]{weinberg1967leptons}. \citet[\S14.1]{gonzalezgarciawendell2026neutrinos} distinguishes the massless minimal model from the observed neutrino-mixing framework.}
\GSfield{Present status}{As of 10 November 2026, the Clifford and projection identities are exact; the notebook's massless-neutrino assumption is a historical model assumption.}
\GSfield{Assessment}{literal statement: correct with missing hypotheses; publication match: partial; verification: read the result; development: annotated; confidence: high}
\end{GSitem}
\begin{GSitem}{1979-479.6}{E}{Photon, charged-current and neutral-current exchange}
\GSfield{Kind; source}{construction; \GSpage{P000892}}
\GSfield{The notes}{Feynman sketches show photon, \(W^-\) and \(Z^0\) exchange. The photon is massless, the weak bosons massive, and approximate mass predictions are labeled about \(85m_p\) and \(95m_p\).}
\GSfield{Statement}{In the electroweak model a massless photon couples to the electromagnetic charge current, charged \(W^\pm\) bosons couple the two left-handed components of weak doublets, and a neutral \(Z\) couples to a weak neutral current. Gauge-boson lines in a perturbative diagram represent propagators carrying four-momentum. The 1979 masses are model predictions and rough numerical estimates, not measurements supplied by the source.}
\GSfield{Reconstruction}{The diagram's particle labels are read as examples of the interaction vertices; no exact cross section or claim of experimental discovery in 1979 is inferred.}
\GSfield{Fate}{\citet[eqs.~(5)--(14)]{weinberg1967leptons} derives the earlier mass and current structure. Current measured masses and electroweak fits are treated in \citet[\S10.4.1]{deblasdittmaierkogler2026electroweak}.}
\GSfield{Present status}{As of 10 November 2026, the weak gauge bosons are experimentally established; detailed current couplings and the tree-level mass relation are reconstructed in record 480.}
\GSfield{Assessment}{literal statement: correct with missing hypotheses; publication match: partial; verification: read the result; development: annotated; confidence: high}
\end{GSitem}
\begin{GSitem}{1979-479.7}{E}{Non-Abelian curvature and the gauge kinetic action}
\GSfield{Kind; source}{definition; \GSpage{P000893}}
\GSfield{The notes}{Equations (8)--(11) define adjoint-indexed vector fields, curvature with a quadratic term, a trace normalization and the Yang--Mills kinetic Lagrangian. Both curvature factors are printed with upper spacetime indices.}
\GSfield{Statement}{For a connection \(A_\mu=A_\mu^aT_a\), define \(D_\mu=\partial_\mu+igA_\mu\) and \([D_\mu,D_\nu]=igF_{\mu\nu}\). Contract the curvature using an invariant Lie-algebra form and the spacetime metric: \(\mathcal L_{\rm YM}=-\tfrac14F^a_{\mu\nu}F^{a\mu\nu}\) in the standard normalization. Matrix traces depend on the representation: for conventionally normalized \(SU(N)\), \(\operatorname{tr}_{\rm fund}T_aT_b=\tfrac12\delta_{ab}\) but \(\operatorname{tr}_{\rm adj}T_aT_b=N\delta_{ab}\).}
\GSfield{Reconstruction}{The metric contraction is restored. The notebook's adjoint trace value needs a stated rescaling of the invariant form, or a correction of the representation label; it is not a universal trace convention.}
\GSfield{Fate}{\citet[eqs.~(1)--(11)]{yangmills1954conservation} is the original non-Abelian construction; \citet[\S9.1]{hustonrabbertzzanderighi2026qcd} gives the current QCD normalization.}
\GSfield{Present status}{As of 10 November 2026, the invariant curvature action is established; the compressed indices and trace convention require the stated corrections.}
\GSfield{Assessment}{literal statement: correct with missing hypotheses; publication match: partial; verification: read the result; development: annotated; confidence: high}
\end{GSitem}
\begin{GSitem}{1979-479.8}{E}{Minimal gauge coupling of fermions}
\GSfield{Kind; source}{construction; \GSpage{P000893}}
\GSfield{The notes}{The matter term is \(\bar\Psi i\widehat D\Psi\), with \(D_\mu=\partial_\mu+igI^aA^a_\mu\) and \(\widehat D=\gamma^\mu D_\mu\). It is described as massless before adding scalar fields.}
\GSfield{Statement}{The generators in \(D_\mu\Psi\) must act in the representation carried by \(\Psi\), which need not be the adjoint representation of the gauge field. With that choice, \(\bar\Psi i\gamma^\mu D_\mu\Psi\) is gauge invariant. A bare Dirac mass is gauge invariant only if the left- and right-handed representations can be paired; the chiral electroweak electron needs a Higgs Yukawa coupling instead.}
\GSfield{Reconstruction}{The representation on matter is separated from the adjoint field index, avoiding the misleading implication that all matter generators are adjoint matrices.}
\GSfield{Fate}{\citet[eqs.~(1)--(4)]{weinberg1967leptons} is the relevant pre-existing chiral matter construction; \citet[eqs.~(1)--(6)]{yangmills1954conservation} gives the original covariant-derivative mechanism.}
\GSfield{Present status}{As of 10 November 2026, minimal coupling is standard; the bare-mass prohibition depends on the representation, not solely on the presence of a gauge field.}
\GSfield{Assessment}{literal statement: correct with missing hypotheses; publication match: exact; verification: read the result; development: annotated; confidence: high}
\end{GSitem}
\begin{GSitem}{1979-479.9}{E}{Gauge-invariant Higgs potential and power counting}
\GSfield{Kind; source}{construction; \GSpage{P000893}}
\GSfield{The notes}{A scalar multiplet and potential \(P(g\varphi)=P(\varphi)\) are added; the \(U(1)\) example is \(\lambda^2(\varphi_1^2+\varphi_2^2-\eta^2)^2\). The second power is called the maximum compatible with renormalizability.}
\GSfield{Statement}{For a scalar of canonical mass dimension one in four spacetime dimensions, a polynomial interaction of degree at most four has a coupling of nonnegative mass dimension. Thus the displayed potential is quadratic in \(|\varphi|^2\) and quartic in the real scalar fields. Gauge invariance restricts it to invariant polynomials. Power counting is a necessary organizational criterion here, not by itself a proof of renormalizability for an arbitrary gauge model.}
\GSfield{Reconstruction}{The spacetime dimension and the variable whose second power is meant are made explicit.}
\GSfield{Fate}{\citet[eq.~(1)]{higgs1964broken} already uses this Abelian scalar-gauge model; \citet[eqs.~(1)--(4)]{weinberg1967leptons} gives the weak-doublet potential.}
\GSfield{Present status}{As of 10 November 2026, the quartic power-counting statement is established; no new renormalization theorem is present.}
\GSfield{Assessment}{literal statement: correct with missing hypotheses; publication match: partial; verification: read the result; development: annotated; confidence: high}
\end{GSitem}
\begin{GSitem}{1979-479.10}{E}{Classical vacuum as a minimum-energy configuration}
\GSfield{Kind; source}{definition; \GSpage{P000893}}
\GSfield{The notes}{The last sentence identifies classical vacuum solutions with minimum energy and hence with minima of \(P(\varphi)\).}
\GSfield{Statement}{In a trivial topological sector with compatible boundary conditions, canonical positive kinetic energy and potential bounded below, constant scalar configurations at global minima of \(P\), with vanishing field strength and covariant derivative, minimize the classical energy. Nontrivial boundary data, topological sectors or gauge constraints can obstruct that simple description. A local minimum is a metastable classical state, not automatically the global ground state.}
\GSfield{Reconstruction}{The vanishing-gradient and boundary hypotheses needed to replace energy minimization by potential minimization are supplied.}
\GSfield{Fate}{The symmetry-breaking vacuum choice is explicit in \citet[eqs.~(1)--(2)]{higgs1964broken} and \citet[eqs.~(3)--(5)]{weinberg1967leptons}. The following lecture develops the Hessian and mass spectrum around it.}
\GSfield{Present status}{As of 10 November 2026, the qualified classical statement is standard; it makes no nonperturbative quantum-vacuum claim.}
\GSfield{Assessment}{literal statement: correct with missing hypotheses; publication match: exact; verification: read the result; development: annotated; confidence: high}
\end{GSitem}

\GStalk{480}{1979-12-10}{I. Yu. Kobzarev}{P000894-P000899}{Confinement heuristics and the electroweak mass mechanism}
The continuation first contrasts Coulomb fields with a proposed chromoelectric tube, with a short intervention attributed to Manin. It then derives the scalar Hessian, symmetry zero modes, gauge-boson masses, the Higgs-doublet breaking pattern, weak mixing, currents and fermion masses, ending with the quark mixing matrix. The confinement picture is distinguished from its still-open four-dimensional continuum formulation. The Higgs and weak-interaction material is principally exposition of older results, with explicit corrections to several notebook formulas.

\begin{GSitem}{1979-480.1}{E}{Coulomb interaction energy and divergent point-charge self-energies}
\GSfield{Kind; source}{claim; \GSpage{P000894}}
\GSfield{The notes}{The two-source field-line drawing accompanies \(\int E^2dV\sim e^2/r\).}
\GSfield{Statement}{For two static point charges in three spatial dimensions, subtracting the two separation-independent self-energies leaves an interaction \(V(r)=q_1q_2/(4\pi r)\) in rationalized units. The total \(\int|E|^2dV\) is divergent at the point charges unless they are regularized. Opposite charges have negative interaction energy even though the unsubtracted field-energy density is nonnegative.}
\GSfield{Reconstruction}{The source's schematic proportionality is read as the separation-dependent energy after self-energy subtraction, not the literal divergent total integral.}
\GSfield{Fate}{The Abelian Coulomb behavior is contrasted with area-law confinement in \citet[\S~II]{wilson1974confinement}, published 15 October 1974. The electrostatic scaling itself is much older.}
\GSfield{Present status}{As of 10 November 2026, the Coulomb comparison is established; the subtraction and charge-sign qualifications are essential.}
\GSfield{Assessment}{literal statement: correct with missing hypotheses; publication match: partial; verification: read the result; development: annotated; confidence: high}
\end{GSitem}
\begin{GSitem}{1979-480.2}{B}{A flux tube of distance-independent transverse size}
\GSfield{Kind; source}{conjecture; \GSpage{P000894}}
\GSfield{The notes}{The page explicitly labels as a hypothesis that the tube cross-section or diameter does not depend on the separation of the sources.}
\GSfield{Statement}{A fixed intrinsic core profile in a classical flux-tube model and the quantum mean-square transverse width of a fluctuating string are different observables. In the long-string effective theory at zero temperature, the latter has leading behavior \(w^2(r/2)=(d-2)(2\pi\sigma)^{-1}\log(r/r_0)\). Thus a strictly fixed quantum width is incompatible with that effective description even when the leading energy is \(\sigma r\).}
\GSfield{Reconstruction}{The notebook specifies neither a width observable nor a classical/quantum averaging prescription. The effective-theory statement is a precise later comparison; it is not silently substituted for an exact theorem about the source's unspecified diameter.}
\GSfield{Fate}{The later study \citet[eqs.~(1)--(3), with lattice comparison in Figs.~2--3]{gliozzipepewiese2010width} derives this broadening and tests it in \((2+1)\)-dimensional \(SU(2)\) lattice theory (preprint 26 February 2010; journal 2010). It cites the 1981 Lüscher--Münster--Weisz precursor. This is a developed correction of the naive picture, not a proof of the full four-dimensional continuum theory.}
\GSfield{Present status}{As of 10 November 2026, distance-independent width is not a universal confinement criterion. Quantum broadening and finite intrinsic thickness can coexist, and neither alone settles the continuum static-potential problem.}
\GSfield{Assessment}{literal statement: not determinable; publication match: partial; verification: read the result; development: annotated; confidence: medium}
\end{GSitem}
\begin{GSitem}{1979-480.3}{C}{Linear static energy and continuum confinement}
\GSfield{Kind; source}{conjecture; \GSpage{P000894}}
\GSfield{The notes}{Constant cross-section and energy density are used to infer a string with \(\nu=d\varepsilon/d\ell\sim E^2\sigma\), where the source's \(\sigma\) denotes cross-sectional area.}
\GSfield{Statement}{The elementary conditional inference is that constant energy per unit length gives \(V(r)=\tau r+O(1)\). A distinct sharp gauge-theory question is whether four-dimensional continuum pure \(SU(3)\) Yang--Mills theory at zero temperature has a positive asymptotic fundamental static string tension: \(\lim_{r\to\infty}V(r)/r=\tau>0\), where \(V(r)=-\lim_{T\to\infty}T^{-1}\log\langle W(r,T)\rangle\) after the appropriate static self-energy subtraction. Dynamical fundamental quarks are excluded from that asymptotic formulation.}
\GSfield{Reconstruction}{The continuum, infinite-distance, zero-temperature and pure-gauge quantifiers are supplied to match the earlier problem 1975-377.10. The fixed-width hypothesis in item 1979-480.2 is stronger and is not part of that reused problem. The notation \(\tau\) avoids confusing tension with the source's area \(\sigma\).}
\GSfield{Fate}{\citet[\S~IV, contrasted with \S~VI]{wilson1974confinement} establishes the lattice strong-coupling mechanism, not the continuum limit. \citet[\S4, eq.~(4.4) and Figs.~2--3]{bulava2019string} demonstrates string breaking in lattice QCD with light and strange dynamical quarks (preprint 11 February 2019; journal 2019). \citet[\S5]{jaffewitten2000yangmills} separates confinement from the existence-and-mass-gap problem; the latter remains officially unsolved in \citet[problem status]{clay2026yangmills}.}
\GSfield{Present status}{As of 10 November 2026, the conditional tube-energy calculation is elementary; a complete rigorous four-dimensional continuum construction proving this asymptotic potential was not found in the checked literature. Lattice evidence and effective strings do not complete that theorem, and string breaking invalidates the same asymptotic assertion with dynamical quarks.}
\GSfield{Assessment}{literal statement: correct with missing hypotheses; publication match: partial; verification: read the result; development: annotated; confidence: high; problem \hyperref[prob:GP-1975-10]{GP-1975-10}}
\GSfield{Worth developing}{A proof of a continuum positive string tension would be fundamental, but a bounded chapter calculation can only clarify the hypotheses and known controlled limits, not resolve confinement.}
\end{GSitem}
\begin{GSitem}{1979-480.4}{F}{Manin's semiclassical qualification}
\GSfield{Kind; source}{remark; \GSpage{P000894}; Yu. I. Manin}
\GSfield{The notes}{A remark attributed to Manin refers to quantizing Yang--Mills in a semiclassical approximation and says that no exact model yet exists.}
\GSfield{Statement}{The recorded fragment gives no background solution, expansion parameter, observable or proposition about semiclassical asymptotics. It can qualify the preceding heuristic, but cannot be reconstructed as a definite semiclassical theorem or a separately quantified open problem.}
\GSfield{Reconstruction}{No instanton, soliton or vacuum saddle is supplied merely because the word semiclassical occurs.}
\GSfield{Fate}{The distinction between a formal or approximate gauge calculation and a construction of continuum quantum Yang--Mills theory is made explicit in \citet[\S\S2, 5--6]{jaffewitten2000yangmills}. This is context, not an exact publication match for Manin's unfinished sentence.}
\GSfield{Present status}{As of 10 November 2026, the precise intended approximation remains unidentified; the broader construction problem is recorded separately.}
\GSfield{Assessment}{literal statement: not determinable; publication match: thematic only; verification: read the result; development: annotated; confidence: low}
\end{GSitem}
\begin{GSitem}{1979-480.5}{E}{The product gauge connection and electroweak kinetic terms}
\GSfield{Kind; source}{construction; \GSpage{P000894-P000895}}
\GSfield{The notes}{The weak Lagrangian combines three \(SU(2)\) fields \(W^a\), an Abelian field \(B\), matter multiplets and a scalar multiplet; \(D_\mu=\partial_\mu+igt^aW^a_\mu+ig'yB_\mu\).}
\GSfield{Statement}{For \(SU(2)_L\times U(1)_Y\), the connection has Lie-algebra components in \(\mathfrak{su}(2)\oplus\mathfrak u(1)\). Its adjoint representation is a triplet plus a trivial one-dimensional factor. The gauge kinetic terms are \(-\tfrac14W^a_{\mu\nu}W^{a\mu\nu}-\tfrac14B_{\mu\nu}B^{\mu\nu}\); the Abelian curvature has no commutator term. Each matter covariant derivative uses that matter field's own representation and hypercharge.}
\GSfield{Reconstruction}{The notebook's alternation between \(A\), \(W\) and \(H\) for fields or curvatures is disambiguated. A coupling omitted from a curvature can be absorbed by rescaling its connection; the choice must be consistent with the kinetic term.}
\GSfield{Fate}{\citet[eqs.~(1)--(4)]{weinberg1967leptons} gives the older model, and \citet[\S10.1]{deblasdittmaierkogler2026electroweak} states the normalized modern Lagrangian.}
\GSfield{Present status}{As of 10 November 2026, the product-connection construction is established; its three-plus-one adjoint decomposition is not the representation of every fermion or scalar.}
\GSfield{Assessment}{literal statement: correct with missing hypotheses; publication match: exact; verification: read the result; development: annotated; confidence: high}
\end{GSitem}
\begin{GSitem}{1979-480.6}{E}{Scalar energy, stationary points and local vacua}
\GSfield{Kind; source}{definition; \GSpage{P000895}}
\GSfield{The notes}{The scalar energy is \(T_{00}=\tfrac12\dot\varphi^2+\tfrac12|\nabla\varphi|^2+P(\varphi)\), called positive, and any local minimum of \(P\) is called a vacuum.}
\GSfield{Statement}{With positive scalar kinetic form and a potential shifted so that its global minimum is zero, the classical energy density is nonnegative. In a trivial sector with compatible constant boundary data a constant global minimum minimizes the energy. A local minimum satisfies \(\nabla P=0\) and gives a candidate metastable state; stationary points alone need not be minima, and noncompact field space need not contain a global minimizer.}
\GSfield{Reconstruction}{The energy shift, boundedness/existence assumptions and distinction between local and global minima are supplied. The word vacuum in the source is explicitly a local classical convention.}
\GSfield{Fate}{The earlier scalar symmetry-breaking examples are \citet[eqs.~(1)--(2)]{higgs1964broken} and \citet[eqs.~(3)--(5)]{weinberg1967leptons}.}
\GSfield{Present status}{As of 10 November 2026, the qualified energy argument is elementary; it does not establish a unique stable quantum vacuum.}
\GSfield{Assessment}{literal statement: correct with missing hypotheses; publication match: partial; verification: read the result; development: annotated; confidence: high}
\end{GSitem}
\begin{GSitem}{1979-480.7}{E}{Small oscillations and the scalar mass-squared Hessian}
\GSfield{Kind; source}{theorem; \GSpage{P000895-P000896}}
\GSfield{The notes}{Linearization at \(\lambda\) gives \(\Box\varphi^\alpha+P_{,\alpha\beta}(\lambda)\varphi^\beta=0\), but the subsequent plane-wave equation inserts an additional factor \(1/2\).}
\GSfield{Statement}{For canonical kinetic term \(\tfrac12\partial_\mu\phi^\alpha\partial^\mu\phi^\alpha\) and \(\nabla P(\lambda)=0\), the mass-squared matrix is \(H_{\alpha\beta}=P_{,\alpha\beta}(\lambda)\). A plane wave obeys \(Hc=m^2c\), with no factor \(1/2\). The half in the Taylor expansion \(P(\lambda+\phi)=P(\lambda)+\tfrac12\phi^TH\phi+\cdots\) disappears when the equation is varied. A noncanonical kinetic metric would give a generalized eigenvalue problem instead.}
\GSfield{Reconstruction}{The factor is corrected against both the scan and the Euler--Lagrange equation, not silently changed in the transcription.}
\GSfield{Fate}{The scalar and vector fluctuation masses are already derived in \citet[eqs.~(2a)--(4)]{higgs1964broken}; the Hessian statement here also follows directly from the notebook's previous equation.}
\GSfield{Present status}{As of 10 November 2026, the corrected classical mass matrix is settled; the notebook's extra half is inconsistent with its canonical normalization.}
\GSfield{Assessment}{literal statement: false as written; publication match: partial; verification: read the result; development: annotated; confidence: high}
\end{GSitem}
\begin{GSitem}{1979-480.8}{E}{Symmetry zero modes and the correct Goldstone count}
\GSfield{Kind; source}{theorem; \GSpage{P000896}}
\GSfield{The notes}{Differentiating \(P(g\varphi)=P(\varphi)\) gives zero modes \(I^a\lambda\). The text initially counts the generators but then observes that some vectors can be dependent.}
\GSfield{Statement}{Let a compact Lie group \(G\) act orthogonally on a real scalar space, let \(P\) be invariant, and let \(\lambda\) be a stationary point. Differentiating the invariance identity gives \(H(T\lambda)=0\) for every \(T\in\mathfrak g\). Thus the Hessian kernel contains the tangent space of the orbit and has dimension at least \(r=\dim G-\dim G_\lambda\). Equality requires nondegeneracy transverse to the orbit; accidental flat directions can give additional zeros.}
\GSfield{Reconstruction}{The rank of the orbit, rather than the dimension of the representation or all generators, is the exact symmetry count. Stationarity is used when the term containing \(\nabla P\) is dropped.}
\GSfield{Fate}{The gauge-model relation to the older Goldstone result is discussed in \citet[opening paragraphs and eqs.~(2a)--(2c)]{higgs1964broken}. The finite-dimensional Hessian identity is explicitly derived in this notebook; it does not by itself prove the full quantum Goldstone theorem.}
\GSfield{Present status}{As of 10 November 2026, the classical kernel statement is exact; the source's later dependence qualification repairs its initial overcount.}
\GSfield{Assessment}{literal statement: correct with missing hypotheses; publication match: partial; verification: read the result; development: annotated; confidence: high}
\end{GSitem}
\begin{GSitem}{1979-480.9}{E}{Unitary gauge, vector masses and conserved physical degrees of freedom}
\GSfield{Kind; source}{construction; \GSpage{P000896-P000897}}
\GSfield{The notes}{A gauge is chosen to remove scalar components along Goldstone directions; the connection then acquires a quadratic mass term. The page counts \(2n\to3n\) vector polarizations and \(N\to N-n\) scalar components.}
\GSfield{Statement}{Locally near a nonzero vacuum on a fixed orbit-type stratum, choose a slice transverse to the \(r\)-dimensional gauge orbit. The scalar kinetic term produces the positive semidefinite Gram mass matrix \((M^2)_{ab}=g_ag_b\langle T_a\lambda,T_b\lambda\rangle\), of rank \(r\). In four spacetime dimensions, \(r\) vector fields have three physical polarizations and \(n-r\) remain massless with two; \(N-r\) scalar components remain. The total count is \(2n+N=3r+2(n-r)+(N-r)\).}
\GSfield{Reconstruction}{The source's \(3n\) count assumes every generator is broken. Unitary gauge is local and can fail at zeros of the scalar field or across different orbit strata; the claimed choice at every point is not a global theorem.}
\GSfield{Fate}{\citet[eqs.~(2)--(4)]{higgs1964broken} gives the 1964 Abelian mechanism explicitly, including the new massive vector and remaining scalar; \citet[eqs.~(5)--(12)]{weinberg1967leptons} applies the partial breaking relevant here.}
\GSfield{Present status}{As of 10 November 2026, the local classical Higgs mechanism is established; it redistributes physical polarizations without increasing their total number.}
\GSfield{Assessment}{literal statement: correct with missing hypotheses; publication match: partial; verification: read the result; development: annotated; confidence: high}
\end{GSitem}
\begin{GSitem}{1979-480.10}{E}{Exercise: eliminate the Goldstone field in the planar scalar model}
\GSfield{Kind; source}{question; \GSpage{P000896-P000897}}
\GSfield{The notes}{The explicit exercise asks to eliminate the Goldstone mode for \(P=\tfrac12(\varphi_1^2+\varphi_2^2-\eta^2)^2\), with \(G=O(2)\), and suggests shifting \(A\) by a derivative of an angle.}
\GSfield{Statement}{On a patch with \(|\varphi|>0\), write \(\varphi=(\eta+h)(\cos\theta,\sin\theta)\). For the connected \(SO(2)\) gauge action and convention \(D=\partial+gAJ\), set \(B_\mu=A_\mu+g^{-1}\partial_\mu\theta\). Then \(\tfrac12|D\varphi|^2=\tfrac12(\partial h)^2+\tfrac12g^2(\eta+h)^2B_\mu B^\mu\). The angular field disappears locally; the quadratic masses are \(m_B^2=g^2\eta^2\) and \(m_h^2=4\eta^2\) for the notebook's potential normalization.}
\GSfield{Reconstruction}{The coupling convention, connected subgroup and nonvanishing-field patch are supplied. The disconnected reflection in \(O(2)\) supplies no additional infinitesimal Goldstone generator.}
\GSfield{Fate}{This is the elementary local calculation underlying \citet[eqs.~(1)--(4)]{higgs1964broken}, already published in 1964. The exercise is answered here as an annotation; it is not an unresolved problem or a new research result.}
\GSfield{Present status}{As of 10 November 2026, the exercise has the displayed standard solution; a global angle is not asserted in sectors with zeros or winding.}
\GSfield{Assessment}{literal statement: correct with missing hypotheses; publication match: exact; verification: read the result; development: annotated; confidence: high}
\end{GSitem}
\begin{GSitem}{1979-480.11}{E}{The Higgs doublet and the unbroken electromagnetic subgroup}
\GSfield{Kind; source}{construction; \GSpage{P000897-P000899}}
\GSfield{The notes}{The scalar is a two-component complex spinor with four real components. A gauge rotation puts it in \((0,\varphi_2)\), and one combination of \(t_3\) and hypercharge annihilates that vector. Three Goldstones are said to be eaten.}
\GSfield{Statement}{Take a complex \(SU(2)\) doublet \(\Phi\) with \(Y=1/2\), and \(\langle\Phi\rangle=(0,v/\sqrt2)\). Then \(Q=T_3+Y\) annihilates the vacuum, leaving \(U(1)_{\rm em}\) unbroken. The orbit has dimension three: three of the four real scalar components supply longitudinal vector polarizations and one physical radial scalar remains. Spinor here describes the weak fundamental representation; the Higgs field is a Lorentz scalar.}
\GSfield{Reconstruction}{The notebook's hypercharge symbol is rescaled and sign-changed on P000897; choosing \(Q=T_3+Y\) avoids treating that convention as a physical discrepancy. Gauging away three components is local near the chosen nonzero vacuum.}
\GSfield{Fate}{\citet[eqs.~(3)--(12)]{weinberg1967leptons} is the direct pre-1979 model; \citet[\S10.1]{deblasdittmaierkogler2026electroweak} fixes the current doublet normalization.}
\GSfield{Present status}{As of 10 November 2026, this symmetry-breaking pattern and particle count are established; the remaining scalar's mass is a potential parameter, not fixed by the symmetry alone.}
\GSfield{Assessment}{literal statement: correct with missing hypotheses; publication match: exact; verification: read the result; development: annotated; confidence: high}
\end{GSitem}
\begin{GSitem}{1979-480.12}{E}{Diagonalizing the weak neutral mass matrix}
\GSfield{Kind; source}{theorem; \GSpage{P000898}}
\GSfield{The notes}{The displayed four-by-four matrix has charged entries \(a^2,a^2\) and neutral block \(\left(\begin{smallmatrix}a^2&ab\\ab&b^2\end{smallmatrix}\right)\), with \(a\) proportional to \(g\) and \(b\) to \(g'\). It yields \(\tan\theta_w=g'/g\) and \(m_Z=m_W/\cos\theta_w\).}
\GSfield{Statement}{The neutral block has determinant zero, eigenvalues \(0,a^2+b^2\), and orthogonal eigenvectors \((-b,a)\) and \((a,b)\). With the standard doublet normalization, \(m_W^2=g^2v^2/4\), \(m_Z^2=(g^2+g'^2)v^2/4\), and \(m_\gamma=0\), so \(m_W=m_Z\cos\theta_w\) at tree level. Changing the sign of the Abelian connection changes the off-diagonal sign but not the spectrum.}
\GSfield{Reconstruction}{The vertical bars in equation (33) are read as delimiting the mass matrix, not merely as its determinant. The coupling relation requires the conventional hypercharge normalization and common vacuum scale.}
\GSfield{Fate}{\citet[eqs.~(5)--(12)]{weinberg1967leptons} derives the direct earlier result. Radiative corrections and the distinction between mixing-angle schemes and measured masses are explicit in \citet[\S\S10.1--10.2]{deblasdittmaierkogler2026electroweak}.}
\GSfield{Present status}{As of 10 November 2026, the tree-level eigenvalue identity is exact; physical precision predictions include radiative corrections.}
\GSfield{Assessment}{literal statement: correct with missing hypotheses; publication match: exact; verification: read the result; development: annotated; confidence: high}
\end{GSitem}
\begin{GSitem}{1979-480.13}{E}{Charged, electromagnetic and neutral currents}
\GSfield{Kind; source}{construction; \GSpage{P000898}}
\GSfield{The notes}{Equation (37) lists \(J^\pm,J^{\rm el},J^0\). Its neutral term has a prefactor read as \(q/\cos\theta_w\) and a single \(\sin\theta_w\) where the standard expression requires a square.}
\GSfield{Statement}{For fixed generator and field conventions, \[\mathcal L_{\rm int}=-\frac{g}{\sqrt2}(J^-_\mu W^{+\mu}+J^+_\mu W^{-\mu})-eJ^{\rm em}_\mu A^\mu-\frac{g}{\cos\theta_w}(J^3_\mu-\sin^2\theta_w J^{\rm em}_\mu)Z^\mu,\quad e=g\sin\theta_w.\] Here \(J^3\) contains left-handed isospin and \(J^{\rm em}\) contains both chiralities with their electric charges. Vertex factors follow by differentiating the interaction action; probabilities further require amplitudes, sums over states and phase space.}
\GSfield{Reconstruction}{The meanings of \(J^+\) and \(J^-\) are fixed with the chosen conjugate W fields. The neutral coupling and squared sine are emendations of the source; a charge operator cannot replace the universal weak prefactor.}
\GSfield{Fate}{\citet[eqs.~(13)--(14)]{weinberg1967leptons} gives the earlier current structure; \citet[\S10.1]{deblasdittmaierkogler2026electroweak} provides the modern normalization and chiral couplings.}
\GSfield{Present status}{As of 10 November 2026, the corrected electroweak current formula is established; the original two coefficient errors would change neutral-current amplitudes.}
\GSfield{Assessment}{literal statement: false as written; publication match: partial; verification: read the result; development: annotated; confidence: high}
\end{GSitem}
\begin{GSitem}{1979-480.14}{E}{Electron mass from a gauge-invariant Yukawa coupling}
\GSfield{Kind; source}{construction; \GSpage{P000898-P000899}}
\GSfield{The notes}{The left lepton doublet and right electron are paired with the Higgs vacuum vector; after choosing \((0,\lambda)\) the term is \(f\lambda\bar e_Re_L+\mathrm{h.c.}\).}
\GSfield{Statement}{With \(L_L\sim(\mathbf2,-1/2)\), \(e_R\sim(\mathbf1,-1)\), and \(\Phi\sim(\mathbf2,1/2)\), the invariant interaction is \(-y_e\bar L_L\Phi e_R+\mathrm{h.c.}\). At \(\Phi=(0,(v+h)/\sqrt2)\), it gives \(m_e=|y_e|v/\sqrt2\) after rephasing and the associated electron--Higgs coupling. The minimal field content leaves neutrinos massless; that is a limitation of this model, not an established property of actual neutrinos.}
\GSfield{Reconstruction}{The scalar or conjugate scalar in the invariant contraction depends on whether the term is written with \(\bar L_L\) or \(\bar e_R\); the precise invariant expression fixes the source's suppressed conjugation.}
\GSfield{Fate}{\citet[eqs.~(1), (3)--(5)]{weinberg1967leptons} gives the earlier electron Yukawa term. The modern distinction between the minimal model and neutrino mass extensions is \citet[\S\S14.1--14.2]{gonzalezgarciawendell2026neutrinos}.}
\GSfield{Present status}{As of 10 November 2026, the Yukawa mechanism is established, while neutrino masses and mixing require an extension of the massless minimal field content.}
\GSfield{Assessment}{literal statement: correct with missing hypotheses; publication match: exact; verification: read the result; development: annotated; confidence: high}
\end{GSitem}
\begin{GSitem}{1979-480.15}{E}{Quark Lagrangian and a unitary mixing matrix}
\GSfield{Kind; source}{question; \GSpage{P000899}}
\GSfield{The notes}{The final explicit question asks how to write the quark Lagrangian. Left-handed pairs use \((d',s',b')^T=V(d,s,b)^T\), and \(V\) is labeled orthogonal; right-handed quarks are singlets.}
\GSfield{Statement}{For each generation, \(Q_L\sim(\mathbf3,\mathbf2,1/6)\), \(u_R\sim(\mathbf3,\mathbf1,2/3)\), \(d_R\sim(\mathbf3,\mathbf1,-1/3)\). Add gauge kinetic couplings and \(-\bar Q_LY_d\Phi d_R-\bar Q_LY_u\widetilde\Phi u_R+\mathrm{h.c.}\), where \(\widetilde\Phi=i\sigma_2\Phi^*\). Biunitary diagonalization of \(Y_u,Y_d\) gives \(V=U_{uL}^\dagger U_{dL}\) in charged currents. In the three-generation minimal model \(V\) is unitary and has three physical angles and one CP-violating phase; it is real orthogonal only after imposing a CP-conserving restriction.}
\GSfield{Reconstruction}{The right-handed singlets refer only to weak \(SU(2)\); they still carry colour and hypercharge. The symbol resembling a superscript 1 on the down-type fields is interpreted through the clearly written primed mixing relation, not as a new generation label.}
\GSfield{Fate}{\citet[pp.~654--657]{kobayashimaskawa1973cp} introduced the six-quark CP-violating structure before this talk. \citet[\S12.1, eqs.~(12.1)--(12.3)]{ceccucciligetisakai2026ckm} gives the explicit Yukawa and unitary-mixing construction.}
\GSfield{Present status}{As of 10 November 2026, the question has the standard displayed construction. The unrestricted orthogonality claim omits the physical CP phase and is not an open problem.}
\GSfield{Assessment}{literal statement: false as written; publication match: partial; verification: read the result; development: annotated; confidence: high}
\end{GSitem}

\GStalk{481}{}{-}{P000900-P000901}{The catalog heading YEAR 1980}
This undated catalog record is a year divider, not a separate mathematical session. The supplied page window overlaps Migdal: the first visible heading on P000900 is explicitly 21 January 1980. No mathematical item is duplicated under this heading.


\GStalk{482}{1980-01-21}{A. A. Migdal}{P000900-P000904}{Yang--Mills theory in loop space}
The actual notes occupy P000900--P000904; P000905 starts the next session. Migdal replaces a gauge potential by traces of parallel transport, defines area and path derivatives, and converts the classical Yang--Mills equation and formal Schwinger--Dyson identity to loop equations. Color contraction produces a multiloop hierarchy; normalized large-rank factorization formally closes it. The final assertion about confinement is a research aspiration, not a proof recorded in these pages. Gelfand intervenes on characters of the path group. The principal equation had already appeared in December 1979; the talk is not its first publication.

\begin{GSitem}{1979-482.1}{E}{Gauge connections, ordered parallel transport and Wilson loops}
\GSfield{Kind; source}{construction; \GSpage{P000900}}
\GSfield{The notes}{The notes write \(A_\mu=\sum_aT_aA_\mu^a\), \(\Psi(C)=\operatorname{Tr}P\exp\oint_C A\,dx\), and describe cancellation of an outward path followed by its reverse.}
\GSfield{Statement}{For a smooth matrix connection and a piecewise smooth parametrized path \(\gamma\), parallel transport is the unique solution of the associated first-order matrix ODE, with prescribed initial identity. Reversing \(\gamma\) inverts transport, and concatenation multiplies transports in the chosen order. For a loop based at \(x\), gauge transformation conjugates \(U(C)\) by the gauge matrix at \(x\), so \(\operatorname{Tr}U(C)\) is gauge invariant.}
\GSfield{Reconstruction}{The derivative \(dx\) in (3) is interpreted as \(dx^\mu\) contracted with \(A_\mu\). The side of matrix multiplication, sign of the connection, coupling and factors of \(i\) require a convention; the notes do not fix all of them.}
\GSfield{Fate}{Wilson-loop gauge variables predate this talk: \citet[Secs.~II--III]{wilson1974confinement}; the corresponding loop notation is explicit in \citet[Sec.~4.1]{makeenko2000large}.}
\GSfield{Present status}{As of 10 November 2026: A standard construction, with a fixed connection/transport sign convention. It is not by itself a construction of a continuum quantum measure.}
\GSfield{Assessment}{literal statement: correct with missing hypotheses; publication match: exact; verification: read the result; development: annotated; confidence: high}
\end{GSitem}
\begin{GSitem}{1979-482.2}{E}{Vanishing of gauge-variant expectations needs a precise observable}
\GSfield{Kind; source}{claim; \GSpage{P000900}}
\GSfield{The notes}{\ru{Неинвар. величины в квант. теории получаются равными 0}: noninvariant quantities have zero quantum expectation.}
\GSfield{Statement}{For a gauge-invariant finite-lattice measure without gauge fixing, the expectation of an observable equals that of its average over local gauge transformations. Hence an observable whose gauge average vanishes has zero expectation. A general noninvariant observable need not have zero expectation: it may contain a nonzero invariant component.}
\GSfield{Reconstruction}{The gauge-average qualification is necessary. The notebook sentence is false if read as quantifying over every noninvariant function.}
\GSfield{Fate}{\citet[abstract]{elitzur1975local} supplies the earlier local-gauge-symmetry result. Its publisher abstract was checked, not the full proof. The more elementary finite-volume projection statement here is an editorial consequence of invariance of the measure, so the publication match is recorded as partial.}
\GSfield{Present status}{As of 10 November 2026: The corrected projection statement is established; gauge fixing changes the assertion about gauge-dependent quantities.}
\GSfield{Assessment}{literal statement: false as written; publication match: partial; verification: publisher metadata checked; development: annotated; confidence: high}
\end{GSitem}
\begin{GSitem}{1979-482.3}{E}{Gelfand on path-group characters and connections}
\GSfield{Kind; source}{remark; \GSpage{P000900-P000902}; I. M. Gelfand}
\GSfield{The notes}{Gelfand says that a representation of the group of paths is determined by its character, and later calls such a representation a connection; curvature is compared with a Lie algebra.}
\GSfield{Statement}{Finite-dimensional unitary representations of a group with identical characters are equivalent: restrict to the compact closure of their joint image and use compact-group character theory. A smooth connection determines a smooth transport functor on the thin path groupoid. Conversely, reconstructing a connection requires regularity and compatibility of transport, not merely an arbitrary algebraic representation of paths.}
\GSfield{Reconstruction}{The finite-dimensional/unitary hypothesis and thin-path smoothness distinguish two claims compressed in the intervention. The curvature analogy is not a literal identification of curvature with a Lie algebra.}
\GSfield{Fate}{Holonomy reconstruction is developed precisely in \citet[abstract]{barrett1991holonomy}. This later theorem clarifies the missing hypotheses; it is a partial match to Gelfand\textquotesingle s informal character argument, not evidence of a new 1980 theorem.}
\GSfield{Present status}{As of 10 November 2026: Established after specifying the category and regularity. The notes contain no unrestricted reconstruction theorem.}
\GSfield{Assessment}{literal statement: correct with missing hypotheses; publication match: partial; verification: publisher metadata checked; development: annotated; confidence: medium}
\end{GSitem}
\begin{GSitem}{1979-482.4}{E}{Stokes-type loop functionals and the area derivative}
\GSfield{Kind; source}{definition; \GSpage{P000900-P000901}}
\GSfield{The notes}{A small attached loop defines \(\delta\Psi/\delta\sigma_{\mu\nu}\); the notes explicitly warn that length is not a functional of this kind. Formula (7)/(10) inserts \(F_{\mu\nu}(x)\) inside the ordered trace.}
\GSfield{Statement}{For smooth transport and an infinitesimal oriented rectangle at a marked point \(x\), the first term in its area is the curvature insertion: \(\delta W(C)/\delta\sigma^{\mu\nu}(x)=N^{-1}\operatorname{Tr}(F_{\mu\nu}(x)U(C_x))\), up to the agreed connection sign. A Stokes functional is insensitive to retracing and has an area derivative with this scaling; perimeter is not such a functional at a newly attached small loop.}
\GSfield{Reconstruction}{Normalized \(W=N^{-1}\Psi\) is introduced for later consistency; the notes use unnormalized \(\Psi\). Orientation and the placement of the marked insertion must be fixed.}
\GSfield{Fate}{The identity was already printed in \citet[Eq.~(12)]{makeenkomigdal1979exact}, received 17 September and published 3 December 1979, with attribution to Mandelstam. The modern notation is \citet[Eq.~(147)]{makeenko2000large}.}
\GSfield{Present status}{As of 10 November 2026: A classical differential identity, not a license to differentiate an unrenormalized quantum Wilson loop without a regulator.}
\GSfield{Assessment}{literal statement: correct with missing hypotheses; publication match: exact; verification: read the result; development: annotated; confidence: high}
\end{GSitem}
\begin{GSitem}{1979-482.5}{E}{Curvature from an infinitesimal commutator}
\GSfield{Kind; source}{method; \GSpage{P000901}}
\GSfield{The notes}{Formula (13) writes a product of four factors \(1\pm dx_\mu\nabla_\mu\) as \(1+dx_\mu\wedge dx_\nu[\nabla_\mu,\nabla_\nu]\); (11)--(14) implement shifts by differential operators.}
\GSfield{Statement}{For two smooth infinitesimal transport directions, \(e^{aD_\mu}e^{bD_\nu}e^{-aD_\mu}e^{-bD_\nu}=1+ab[D_\mu,D_\nu]+O(|a|^2|b|+|a||b|^2)\). The mixed \(ab\) coefficient gives curvature. Replacing inverse exponentials by \(1-aD_\mu\) introduces unwanted pure quadratic terms unless only the mixed coefficient is retained.}
\GSfield{Reconstruction}{Equation (13) must be read as a mixed infinitesimal identity, not as equality of finite products.}
\GSfield{Fate}{This elementary small-rectangle calculation underlies the Mandelstam formula already used in \citet[Eq.~(12)]{makeenkomigdal1979exact} before the talk. The subsequent textbook presentation is \citet[Eq.~(147)]{makeenko2000large}.}
\GSfield{Present status}{As of 10 November 2026: Valid as a mixed derivative or a properly truncated holonomy expansion. A literal finite-product identity is false.}
\GSfield{Assessment}{literal statement: false as written; publication match: exact; verification: read the result; development: annotated; confidence: high}
\end{GSitem}
\begin{GSitem}{1979-482.6}{E}{Marked path derivatives become covariant derivatives}
\GSfield{Kind; source}{construction; \GSpage{P000901-P000902}}
\GSfield{The notes}{Formulas (15)--(18) move the marked point, add a short connector and shift both ends of a closed transport; the notes emphasize that the left-hand derivative is noncommutative.}
\GSfield{Statement}{For a curvature-inserted loop \(\phi_{\mu\nu}(C_x)=\operatorname{Tr}(F_{\mu\nu}(x)U(C_x))\), the derivative defined by translating the insertion with short parallel-transport connectors is \(\partial^x_\alpha\phi_{\mu\nu}=\operatorname{Tr}((D_\alpha F_{\mu\nu})(x)U(C_x))\). Iteration gives covariant derivatives in the specified order. This operator is not an ordinary partial derivative on a scalar function of \(x\) alone.}
\GSfield{Reconstruction}{We make explicit the connector-based definition that the notes state after using it. This resolves their warning about noncommutativity.}
\GSfield{Fate}{The covariant derivative at the insertion is explicit in \citet[Eq.~(14) and the following explanation]{makeenkomigdal1979exact}, published December 1979; \citet[Eqs.~(148)--(150)]{makeenko2000large} gives the fuller path-derivative dictionary.}
\GSfield{Present status}{As of 10 November 2026: Established for classical smooth connections and suitable loop functionals.}
\GSfield{Assessment}{literal statement: correct with missing hypotheses; publication match: exact; verification: read the result; development: annotated; confidence: high}
\end{GSitem}
\begin{GSitem}{1979-482.7}{E}{Translation of a contour and the area-derivative integral}
\GSfield{Kind; source}{claim; \GSpage{P000902}}
\GSfield{The notes}{Formula (19) decomposes whole-contour translation into endpoint derivatives and an integral of area derivatives; (20) specializes to translation-invariant closed-contour functionals.}
\GSfield{Statement}{For an open-path Stokes functional with differentiable endpoint and shape variations, a uniform translation is the sum of its endpoint variations and the swept-strip area variation. With antisymmetric area convention this has the form \(\partial_{C,\mu}=\partial_{x,\mu}+\partial_{y,\mu}+\int_C dz^\alpha\,\delta/\delta\sigma^{\alpha\mu}(z)\). Vanishing of the left side gives the identity used for translation-invariant functionals, with endpoint conventions retained.}
\GSfield{Reconstruction}{This is a variational identity on a specified domain, not an operator identity for arbitrary functions on curves. The sign in (20) uses antisymmetry of the two-form index.}
\GSfield{Fate}{This is an elementary shape-variation identity in the area/path calculus already used by \citet[Eqs.~(12)--(14)]{makeenkomigdal1979exact}. The notation in \citet[Sec.~4.3]{makeenko2000large} is a partial match; neither source was verified to print the exact open-path form (19).}
\GSfield{Present status}{As of 10 November 2026: A conditional shape-variation formula; no independent open problem is stated.}
\GSfield{Assessment}{literal statement: correct with missing hypotheses; publication match: partial; verification: read the result; development: annotated; confidence: medium}
\end{GSitem}
\begin{GSitem}{1979-482.8}{E}{The classical Yang--Mills equation in loop notation}
\GSfield{Kind; source}{theorem; \GSpage{P000900-P000903}}
\GSfield{The notes}{From \(D_\mu F_{\mu\nu}=0\), formulas (15), (15\textquotesingle) and (24) obtain \(\partial^x_\mu\delta\Psi/\delta\sigma_{\mu\nu}(x)=0\). Formulas (21)--(23) invoke the curvature-squared action.}
\GSfield{Statement}{For a smooth solution of the source-free Yang--Mills equation on a Euclidean domain, the marked path derivative of the curvature-inserted Wilson loop vanishes: \(\partial^x_\mu\delta W(C)/\delta\sigma_{\mu\nu}(x)=N^{-1}\operatorname{Tr}((D_\mu F_{\mu\nu})U(C_x))=0\). With \(S=\beta\int\operatorname{Tr}F_{\mu\nu}F_{\mu\nu}\) summed over ordered indices and trace-pairing convention, variation gives \(\delta S/\delta A_\nu=-4\beta D_\mu F_{\mu\nu}\); alternate action normalization changes the prefactor.}
\GSfield{Reconstruction}{The notebook (23) has mismatched free indices and suppresses the action-dependent prefactor and sign. Only the zero set is unaffected by an overall nonzero factor.}
\GSfield{Fate}{\citet[Eq.~(14), followed by Eqs.~(16)--(18)]{makeenkomigdal1979exact} identifies the classical loop derivative before January 1980. The modern account is \citet[Sec.~4.3]{makeenko2000large}.}
\GSfield{Present status}{As of 10 November 2026: Established. It does not identify the quantum large-rank equation with the classical Yang--Mills equation.}
\GSfield{Assessment}{literal statement: correct with missing hypotheses; publication match: exact; verification: read the result; development: annotated; confidence: high}
\end{GSitem}
\begin{GSitem}{1979-482.9}{E}{Normalized expectation and the weak Schwinger--Dyson identity}
\GSfield{Kind; source}{method; \GSpage{P000903}}
\GSfield{The notes}{(25) displays adjacent integrals for the expectation; (26)--(28) integrate a functional derivative and replace the classical equation by \(\hbar\delta/\delta A\).}
\GSfield{Statement}{For a finite-dimensional regulator with density \(Z^{-1}e^{-S(a)/\hbar}\,da\), integrable \(Q\) and vanishing boundary terms, \(\langle(\partial_jS)Q\rangle=\hbar\langle\partial_jQ\rangle\). The continuum expression is formal until a regulator, measure, gauge treatment and renormalization are specified. The equality with a differential operator holds inside expectations, not as equality of ordinary functions.}
\GSfield{Reconstruction}{The partition function in (25) must divide the numerator. The required decay/domain assumptions are not in the notebook.}
\GSfield{Fate}{The weak field-equation replacement and normalized lattice expectation already occur in \citet[Eqs.~(16), (21)--(22)]{makeenkomigdal1979exact}; \citet[Eqs.~(135)--(136)]{makeenko2000large} makes the integration-by-parts interpretation explicit.}
\GSfield{Present status}{As of 10 November 2026: A rigorous integration-by-parts identity for appropriate regulators, and a formal continuum identity unless independently justified.}
\GSfield{Assessment}{literal statement: false as written; publication match: exact; verification: read the result; development: annotated; confidence: high}
\end{GSitem}
\begin{GSitem}{1979-482.10}{E}{The color contraction is specific to the defining unitary representation}
\GSfield{Kind; source}{theorem; \GSpage{P000904}}
\GSfield{The notes}{(31)--(33) split the transport at an insertion and claim a two-trace contraction for an arbitrary semisimple group, with the correction \(-N^{-1}\operatorname{Tr}(UV)\).}
\GSfield{Statement}{For traceless Hermitian generators of \(\mathfrak{su}(N)\) normalized by \(\operatorname{Tr}(T^aT^b)=\delta^{ab}\), completeness gives
\[
\sum_a(T^a)_{ij}(T^a)_{kl}
=\delta_{il}\delta_{jk}-N^{-1}\delta_{ij}\delta_{kl}.
\]
Consequently,
\[
\sum_a\operatorname{Tr}(T^aUT^aV)
=\operatorname{Tr}U\operatorname{Tr}V-N^{-1}\operatorname{Tr}(UV).
\]
Adding the normalized identity generator gives the \(U(N)\) version without the subtraction.}
\GSfield{Reconstruction}{The group, defining representation, and trace normalization are essential. The formula is not valid unchanged for every semisimple group.}
\GSfield{Fate}{The defining-representation color projection is explicit in \citet[Eq.~(16)]{makeenkomigdal1979exact}, published December 1979, and in \citet[Sec.~4.2 and Eq.~(157)]{makeenko2000large}.}
\GSfield{Present status}{As of 10 November 2026: An exact finite-dimensional identity in the stated convention. The notebook	extquotesingle s unrestricted group quantifier is false.}
\GSfield{Assessment}{literal statement: false as written; publication match: exact; verification: read the result; development: annotated; confidence: high}
\end{GSitem}
\begin{GSitem}{1979-482.11}{E}{Finite-rank loop equations form a hierarchy}
\GSfield{Kind; source}{theorem; \GSpage{P000903-P000904}}
\GSfield{The notes}{Equation (34) relates the derivative of one loop expectation to a two-loop expectation. The notes correctly say that one loop is insufficient and include self-intersections.}
\GSfield{Statement}{Writing \(W(C)=N^{-1}\langle\operatorname{Tr}U(C)\rangle\), the formal \(SU(N)\) equation in a conventional coupling normalization is \[\partial^x_\mu\frac{\delta W(C)}{\delta\sigma_{\mu\nu}(x)}=\lambda\oint_Cdy_\nu\,\delta^{(d)}(x-y)\left(W_2(C_{xy},C_{yx})-\frac{1}{N^2}W(C)\right).\] Here \(W_2=\langle N^{-1}\operatorname{Tr}U(C_{xy})\,N^{-1}\operatorname{Tr}U(C_{yx})\rangle\). Higher products generate a splitting/joining hierarchy.}
\GSfield{Reconstruction}{The \(SU(N)\) subtraction, trace normalization and coupling are restored. Contact terms and coincident points require a regulator.}
\GSfield{Fate}{\citet[Eqs.~(16)--(19)]{makeenkomigdal1979exact}, received 17 September and published 3 December 1979, already gives the finite-rank two-trace equation and explains its lack of closure. \citet[Eqs.~(157), (161)]{makeenko2000large} supplies the higher hierarchy explicitly.}
\GSfield{Present status}{As of 10 November 2026: The formal hierarchy is established; rigorous results are available in specified regulated and lower-dimensional settings, not an unrestricted four-dimensional continuum theorem.}
\GSfield{Assessment}{literal statement: correct with missing hypotheses; publication match: exact; verification: read the result; development: annotated; confidence: high}
\end{GSitem}
\begin{GSitem}{1979-482.12}{E}{Gauge-covariant regularization must reconnect separated endpoints}
\GSfield{Kind; source}{method; \GSpage{P000904}}
\GSfield{The notes}{The notes warn that simply smoothing \(\delta(x-y)\) opens the two contours; smoothing must be accompanied by closing them.}
\GSfield{Statement}{When \(x\ne y\) in a point-split color contraction, open transports transform at different endpoints. To preserve gauge invariance one must include connecting parallel transports (and specify their paths), or use a gauge-invariant lattice regulator. A scalar mollifier alone cannot turn the separated open factors into gauge-invariant Wilson loops.}
\GSfield{Reconstruction}{We spell out the endpoint transformation argument. The notes give no complete regularization prescription or convergence theorem.}
\GSfield{Fate}{A gauge-invariant lattice equation is already given in \citet[Eqs.~(20)--(23)]{makeenkomigdal1979exact}, after the singular contact term is discussed. \citet[Sec.~4.6]{makeenko2000large} gives later continuum regularizations. The endpoint-covariance constraint is elementary; these later constructions are not dated to 1979.}
\GSfield{Present status}{As of 10 November 2026: A valid methodological constraint; existence of a continuum limit is a separate problem.}
\GSfield{Assessment}{literal statement: correct with missing hypotheses; publication match: partial; verification: read the result; development: annotated; confidence: high}
\end{GSitem}
\begin{GSitem}{1979-482.13}{E}{Large-rank factorization requires normalized traces}
\GSfield{Kind; source}{claim; \GSpage{P000904}}
\GSfield{The notes}{(35) states \(\langle\Psi(C_1)\cdots\Psi(C_m)\rangle=\prod_i\langle\Psi(C_i)\rangle+O(N^{-2})\), then observes \(\Psi(0)=N\) and redefines \(\Psi=N\bar\Phi\).}
\GSfield{Statement}{For fixed loops, normalized single-trace observables \(\Phi=N^{-1}\operatorname{Tr}U\), fixed \(\lambda=g^2N\), and a regime admitting the usual topological expansion, \(\langle\Phi_1\cdots\Phi_m\rangle-\prod_i\langle\Phi_i\rangle=O(N^{-2})\). For the unnormalized \(\Psi_i\), the corresponding absolute error is \(O(N^{m-2})\). Factorization is an additional dynamical large-rank assertion, not a consequence of gauge invariance alone.}
\GSfield{Reconstruction}{The order estimate in (35) is moved to the normalized variables introduced immediately below it. The topological-expansion regime is explicit.}
\GSfield{Fate}{\citet[Eq.~(19)]{makeenkomigdal1979exact} explicitly states factorization with normalized traces before the talk. Its perturbative topological counting is rooted in \citet[Sec.~3, Eqs.~(3.2)--(3.7)]{thooft1974planar}, received December 1973 and published April 1974. \citet[Sec.~3.4 and Eqs.~(158)--(160)]{makeenko2000large} is the later detailed account.}
\GSfield{Present status}{As of 10 November 2026: Standard in perturbation theory and specified lattice regimes; the literal normalization in (35) is inconsistent.}
\GSfield{Assessment}{literal statement: false as written; publication match: exact; verification: read the result; development: annotated; confidence: high}
\end{GSitem}
\begin{GSitem}{1979-482.14}{C}{The general confinement goal and the proposed large-rank route}
\GSfield{Kind; source}{claim; \GSpage{P000903-P000904}}
\GSfield{The notes}{The notes say the \(SU(\infty)\) theory is near an exact solution and conclude that the equation for \(\bar\Phi\) gives \ru{confinement}. No solution, positivity criterion, uniqueness argument, continuum limit or area-law proof is supplied.}
\GSfield{Statement}{A definite modern version asks for the physically selected, renormalized four-dimensional pure \(SU(N)\) Yang--Mills Wilson-loop expectations in the large-\(N\) limit, satisfying the loop equation, positivity and normalization, and a positive string tension in the large rectangular-loop limit. Closure of an equation alone neither constructs that solution nor proves its area law.}
\GSfield{Reconstruction}{The notebook\textquotesingle s broad assertion is not a fully formulated theorem. The explicit pure-gauge, dimension, limit and positivity conditions define the identifiable confinement programme, rather than claiming those words were all spoken.}
\GSfield{Fate}{\citet[abstract]{wilson1974confinement} proves a strong-coupling lattice mechanism; \citet[main theorem]{chatterjee2015string} proves a rigorous large-rank lattice/string correspondence for \(SO(N)\) at strong coupling. These do not settle the continuum assertion. The continuing four-dimensional construction problem is described in \citet[problem statement]{jaffewitten2000yangmills}.}
\GSfield{Present status}{As of 10 November 2026: Partly developed. Controlled lattice and two-dimensional results do not constitute a general solution of four-dimensional continuum confinement from the loop equations.}
\GSfield{Assessment}{literal statement: not determinable; publication match: partial; verification: read the result; development: annotated; confidence: medium; problem \hyperref[prob:GP-1975-10]{GP-1975-10}}
\GSfield{Worth developing}{A sharply regulated special case can be worthwhile, but the unrestricted continuum confinement programme is not a plausible bounded contribution for this chapter.}
\end{GSitem}

\GStalk{483}{1980-01-28}{A. V. Zelevinsky}{P000905-P000907}{Kazhdan--Lusztig polynomials and representations}
The notes actually continue through the upper part of P000907, before MacPherson begins. Zelevinsky explicitly discusses the 1979 paper of Kazhdan and Lusztig: Coxeter and Hecke algebra preliminaries, the bar-invariant basis, multiplication and cell representations. The Verma-module multiplicity formula and Vogan\textquotesingle s cohomological reformulations are distinguished as conjectures. Their later proofs must not be confused with the basis theorem already proved in 1979.

\begin{GSitem}{1979-483.1}{E}{Coxeter presentations, reduced words and Bruhat order}
\GSfield{Kind; source}{definition; \GSpage{P000905}}
\GSfield{The notes}{Equations (1)--(3) define a Coxeter group and reduced words, and define the elements \(<w\) by deleting any subset of a reduced expression. Examples are symmetric and Weyl groups; the symmetric-group discussion displays only the neighboring braid relation.}
\GSfield{Statement}{Let \((W,S)\) be a Coxeter system with finite \(S\), exponents \(m(s,s)=1\), and \(m(s,t)\in\{2,3,\ldots,\infty\}\) for \(s\ne t\); no relation is imposed for exponent \(\infty\). The length is the least number of generators. The Bruhat relation \(y\le w\) holds exactly when a reduced expression for \(w\) has a subword whose product is \(y\); this does not depend on the chosen reduced expression. In \(S_n\), also require \(s_i s_j=s_j s_i\) for \(|i-j|>1\).}
\GSfield{Reconstruction}{The empty deletion gives \(y=w\), so the general relation is \(\le\), not strict \(<\). The omitted distant commutations are needed for the symmetric-group presentation.}
\GSfield{Fate}{The talk\textquotesingle s named source is \citet[Sec.~1]{kazhdanlusztig1979representations}, received March and published June 1979.}
\GSfield{Present status}{As of 10 November 2026: Standard Coxeter theory; the two qualifications prevent literal misreadings of the compressed examples.}
\GSfield{Assessment}{literal statement: false as written; publication match: exact; verification: publisher metadata checked; development: annotated; confidence: high}
\end{GSitem}
\begin{GSitem}{1979-483.2}{E}{The Hecke algebra, specialization and its standard basis}
\GSfield{Kind; source}{construction; \GSpage{P000905-P000906}}
\GSfield{The notes}{Equations (4)--(6) impose \((T_s+1)(T_s-q)=0\) and braid relations, note specialization \(q=1\), identify the finite-field double-coset algebra, and define \(T_w\) using a reduced word.}
\GSfield{Statement}{Over \(A=\mathbb Z[q^{1/2},q^{-1/2}]\), the equal-parameter Hecke algebra \(\mathcal H(W)\) has basis \(T_w\), with products along reduced expressions independent of the expression. Under \(q^{1/2}\mapsto1\) it becomes \(\mathbb Z[W]\). For a split finite reductive group over \(\mathbb F_q\) with Borel subgroup \(B\), the complex specialization is the algebra of \(B\)-bi-invariant functions, with convolution normalized so \(1_B\) is the identity and the double-coset generators satisfy the stated quadratic relation.}
\GSfield{Reconstruction}{The tensor product with \(\mathbb C\) needs the specialization homomorphism. Splitness, equal parameters and convolution normalization are not written out in the notes.}
\GSfield{Fate}{\citet[Sec.~1]{kazhdanlusztig1979representations} is the explicit pre-talk source.}
\GSfield{Present status}{As of 10 November 2026: Established. Unequal-parameter or twisted finite groups require a modified statement.}
\GSfield{Assessment}{literal statement: correct with missing hypotheses; publication match: exact; verification: publisher metadata checked; development: annotated; confidence: high}
\end{GSitem}
\begin{GSitem}{1979-483.3}{E}{Bar involution and inversion of Hecke generators}
\GSfield{Kind; source}{definition; \GSpage{P000906}}
\GSfield{The notes}{(7)--(9) prescribe \(\overline{q^{1/2}}=q^{-1/2}\), \(\overline{T_w}=T_{w^{-1}}^{-1}\), and \(T_s^{-1}=q^{-1}(T_s+1-q)\).}
\GSfield{Statement}{There is a semilinear algebra involution of \(\mathcal H\), semilinear over the involution \(q^{1/2}\mapsto q^{-1/2}\), determined by \(\overline{T_s}=T_s^{-1}\). The quadratic relation implies the displayed inverse, and the braid relations ensure well-definedness. Thus \(\overline{T_w}=T_{w^{-1}}^{-1}\).}
\GSfield{Reconstruction}{It is semilinear, not \(A\)-linear; the inverse in the subscript formula is distinguished from inversion of \(w\).}
\GSfield{Fate}{\citet[Sec.~1]{kazhdanlusztig1979representations}; the generator calculation is supplied in the notebook itself.}
\GSfield{Present status}{As of 10 November 2026: An exact algebraic construction.}
\GSfield{Assessment}{literal statement: correct; publication match: exact; verification: publisher metadata checked; development: annotated; confidence: high}
\end{GSitem}
\begin{GSitem}{1979-483.4}{E}{Existence and uniqueness of the Kazhdan--Lusztig basis}
\GSfield{Kind; source}{theorem; \GSpage{P000906}}
\GSfield{The notes}{\ru{Теорема 1.1}: a unique bar-invariant element has the expansion
\[
C^{\prime}_w=q^{-\ell(w)/2}\sum_{y\le w}P_{y,w}T_y,
\]
with \(P_{w,w}=1\) and the strict degree bound (12).}
\GSfield{Statement}{For every Coxeter system and every \(w\in W\), there is a unique \(C^{\prime}_w=q^{-\ell(w)/2}\sum_{y\le w}P_{y,w}(q)T_y\) such that \(\overline{C^{\prime}_w}=C^{\prime}_w\), \(P_{w,w}=1\), \(P_{y,w}\in\mathbb Z[q]\), and \(\deg P_{y,w}\le(\ell(w)-\ell(y)-1)/2\) for \(y<w\). These elements form an \(A\)-basis.}
\GSfield{Reconstruction}{For \(y\not\le w\), set \(P_{y,w}=0\). The bound is interpreted integrally and the zero polynomial is allowed.}
\GSfield{Fate}{The source and theorem number match \citet[Thm.~1.1]{kazhdanlusztig1979representations}, published in 1979.}
\GSfield{Present status}{As of 10 November 2026: Established for arbitrary Coxeter systems. This algebraic existence theorem is separate from geometric positivity or character formulas.}
\GSfield{Assessment}{literal statement: correct; publication match: exact; verification: publisher metadata checked; development: annotated; confidence: high}
\end{GSitem}
\begin{GSitem}{1979-483.5}{E}{The simple-reflection polynomial and the second canonical basis}
\GSfield{Kind; source}{construction; \GSpage{P000906}}
\GSfield{The notes}{The calculation gives \(C^{\prime}_s=q^{-1/2}(T_s+1)\), \(P_{e,s}=P_{s,s}=1\); (16) defines \(C_w=\varepsilon_wj(C^{\prime}_w)\).}
\GSfield{Statement}{The semilinear involution \(j\) is specified by \(j(q^{1/2})=q^{-1/2}\) and \(j(T_w)=(-1)^{\ell(w)}q^{-\ell(w)}T_w\). It gives the alternative basis \(C_w=(-1)^{\ell(w)}j(C^{\prime}_w)\). In rank one, \(C_s=q^{-1/2}T_s-q^{1/2}\), and direct substitution verifies both bar invariance and \(T_sC_s=-C_s\).}
\GSfield{Reconstruction}{The last expression and direct check are elementary consequences clarifying the sign convention; they are not a claimed new result.}
\GSfield{Fate}{\citet[Sec.~1]{kazhdanlusztig1979representations} uses these two conventions.}
\GSfield{Present status}{As of 10 November 2026: Exact; the two bases must not be interchanged in the multiplication formula.}
\GSfield{Assessment}{literal statement: correct; publication match: exact; verification: publisher metadata checked; development: annotated; confidence: high}
\end{GSitem}
\begin{GSitem}{1979-483.6}{B}{Kazhdan--Lusztig multiplicities of Verma modules}
\GSfield{Kind; source}{conjecture; \GSpage{P000906}}
\GSfield{The notes}{(13), twice labelled \ru{гипотеза}, states \([M_w:L_y]=P_{w_0w,w_0y}(1)\), with \(L_w\) the simple quotient and \(y\le w\). The notes say it has been checked for \(\mathrm{GL}(4)\).}
\GSfield{Statement}{Let \(\mathfrak g\) be a finite-dimensional complex semisimple Lie algebra, choose a Borel and Cartan, and let \(\lambda\) be a regular integral antidominant weight for the dot action. Put \(M_w=M(w\cdot\lambda)\), and let \(L_w\) be its simple quotient, all in the corresponding block of category \(\mathcal O\). Then \([M_w:L_y]=P_{w_0w,w_0y}(1)\) for \(y\le w\), and the multiplicity is zero otherwise.}
\GSfield{Reconstruction}{The regular integral block and antidominant indexing are absent from the notebook and are essential to this exact indexing. The isolated \(\mathrm{GL}(4)\) check is reported as historical notebook evidence, not independently reconstructed.}
\GSfield{Fate}{Proved independently by \citet{brylinskikashiwara1981holonomic} and \citet{beilinsonbernstein1981localisation}. \citet[pp.~43--44]{brylinski1981differential}, reporting the 1980 Toru\'n lectures, explicitly records both proofs before their 1981 publication.}
\GSfield{Present status}{As of 10 November 2026: A theorem in the stated characteristic-zero category. The 1980 note accurately distinguishes the then-conjectural multiplicities from the 1979 basis theorem.}
\GSfield{Assessment}{literal statement: correct with missing hypotheses; publication match: exact; verification: read the result; development: annotated; confidence: high; problem \hyperref[prob:GP-1979-41]{GP-1979-41}}
\end{GSitem}
\begin{GSitem}{1979-483.7}{B}{Vogan parity and the Ext-polynomial refinement}
\GSfield{Kind; source}{conjecture; \GSpage{P000906}}
\GSfield{The notes}{(14)--(15) give \(P_{y,w}=\sum_{i\ge0}q^i\dim\operatorname{Ext}^{\ell(w)-\ell(y)-2i}(M_y,L_w)\) and vanishing in the other parity, attributed to Vogan.}
\GSfield{Statement}{With the same regular integral antidominant indexing in category \(\mathcal O\), for \(y\le w\) the dimensions satisfy \[P_{y,w}(q)=\sum_{i\ge0}q^i\dim\operatorname{Ext}^{\ell(w)-\ell(y)-2i}_{\mathcal O}(M_y,L_w),\] taking negative Ext degrees to be zero, and \(\operatorname{Ext}^j_{\mathcal O}(M_y,L_w)=0\) unless \(j\equiv\ell(w)-\ell(y)\pmod2\). Equivalently these are the appropriate weight spaces of nilradical cohomology.}
\GSfield{Reconstruction}{The category, block, and indexing convention must be fixed; unrestricted Ext over all \(U(\mathfrak g)\)-modules is not the stated formula.}
\GSfield{Fate}{\citet[Conj.~3.4 and Thm.~3.5]{vogan1979irreducible} states the equivalent nilradical-cohomology refinement; received 16 July and published December 1979. Its geometric proof belongs to the same 1980--1981 resolution recorded by \citet[pp.~43--44]{brylinski1981differential}. The exact category-\(\mathcal O\) Ext formula in this normalization is also printed in \citet[p.~262]{gabberjoseph1981towards}.}
\GSfield{Present status}{As of 10 November 2026: Established under these hypotheses. The coefficient formula gives more information than evaluating at \(q=1\); its equivalence to the character conjecture uses additional structure.}
\GSfield{Assessment}{literal statement: correct with missing hypotheses; publication match: exact; verification: read the result; development: annotated; confidence: high; problem \hyperref[prob:GP-1979-41]{GP-1979-41}}
\end{GSitem}
\begin{GSitem}{1979-483.8}{E}{Multiplication in the canonical basis and cell subquotients}
\GSfield{Kind; source}{theorem; \GSpage{P000907}}
\GSfield{The notes}{Equation (17) defines the leading coefficient \(\mu\); \ru{Теорема 1.3}, equation (18), gives the action of \(T_s\). The final sentence says the regular representation becomes block triangular.}
\GSfield{Statement}{Let \(\mu(z,v)\) be the coefficient of \(q^{(\ell(v)-\ell(z)-1)/2}\) in \(P_{z,v}\), set to zero when that degree is not a nonnegative integer or \(z\not<v\). Then \[T_sC_v=\begin{cases}-C_v,&sv<v,\\qC_v+q^{1/2}C_{sv}+q^{1/2}\sum_{z<v,\,sz<z}\mu(z,v)C_z,&sv>v.\end{cases}\] The associated preorder yields invariant filtrations and left-cell subquotients of the regular representation. Cell modules are not asserted to be irreducible in every Coxeter type.}
\GSfield{Reconstruction}{The zero convention for \(\mu\) makes the notebook sum well-defined. The phrase block triangular is interpreted as a filtration, not a canonical direct-sum decomposition into irreducibles.}
\GSfield{Fate}{\citet[Thm.~1.3 and the cell construction]{kazhdanlusztig1979representations}; this was already published in June 1979.}
\GSfield{Present status}{As of 10 November 2026: Established. The formula is for \(C_v\), not \(C^{\prime}_v\).}
\GSfield{Assessment}{literal statement: correct with missing hypotheses; publication match: exact; verification: publisher metadata checked; development: annotated; confidence: high}
\end{GSitem}

\GStalk{484}{1980-01-28}{R. D. MacPherson}{P000907}{Intersection homology and Schubert singularities}
The actual contribution begins at the MacPherson heading in the lower part of P000907 and ends on that page. It introduces intersection homology, its smooth-space normalization and rational duality, and its interpretation of Kazhdan--Lusztig polynomials. The question mark in the grading shift is preserved as incomplete notation. P000908 begins the 4 February session; it is not a continuation of the January intersection-homology talk. The notebook calls the speaker Bob; his name is Robert D. MacPherson.

\begin{GSitem}{1979-484.1}{E}{Allowable chains for middle intersection homology}
\GSfield{Kind; source}{definition; \GSpage{P000907}}
\GSfield{The notes}{(3)--(4) use an even-dimensional filtration and cycles meeting \(X_{n-2c}\) in dimension \(i-c-1\); bounding chains have the bound \(i-c\). The notes mention local topological triviality along strata.}
\GSfield{Statement}{Let \(X\) be a pure real \(n\)-dimensional stratified PL pseudomanifold with locally conical structure and no codimension-one stratum. For middle perversity \(\bar m(k)=\lfloor(k-2)/2\rfloor\), the allowable-chain complex consists of PL chains \(\xi\) satisfying \(\dim(|\xi|\cap X_{n-k})\le i-k+\bar m(k)\) and the analogous boundary bound \(\dim(|\partial\xi|\cap X_{n-k})\le i-1-k+\bar m(k)\). Its homology is \(IH_i^{\bar m}(X;\mathbb Q)\). For codimension \(2c\), the first bound is \(i-c-1\).}
\GSfield{Reconstruction}{The inequalities, boundary condition, coefficient field and conical pseudomanifold hypotheses are supplied. Equality of dimensions is neither required nor correct.}
\GSfield{Fate}{Already announced in \citet[definition and propositions]{goreskymacpherson1977dualite}; the full chain treatment appears in \citet[Sec.~1]{goreskymacpherson1980intersection}.}
\GSfield{Present status}{As of 10 November 2026: Established. The chain complex, not an unconstrained quotient of all bounding chains, is the precise object.}
\GSfield{Assessment}{literal statement: correct with missing hypotheses; publication match: exact; verification: read the result; development: annotated; confidence: high}
\end{GSitem}
\begin{GSitem}{1979-484.2}{E}{Smooth normalization and rational Poincare duality}
\GSfield{Kind; source}{theorem; \GSpage{P000907}}
\GSfield{The notes}{(1) gives \(IH_i(X)=H_i(X)\) when \(X\) is nonsingular; (2) pairs \(IH_i\) with \(IH_{n-i}\) into \(\mathbb Q\), with an unexplained superscript-like mark.}
\GSfield{Statement}{For a manifold the allowable conditions are vacuous, so intersection homology agrees with ordinary homology. For a compact oriented stratified pseudomanifold, complementary perversities \(\bar p(k)+\bar q(k)=k-2\) admit a nondegenerate rational pairing \(IH_i^{\bar p}(X;\mathbb Q)\otimes IH_{n-i}^{\bar q}(X;\mathbb Q)\to\mathbb Q\). If all strata have even codimension, middle perversity is self-complementary and gives the notebook\textquotesingle s same-perversity pairing.}
\GSfield{Reconstruction}{Compactness, orientation, rational coefficients and even codimensions are necessary qualifications. No interpretation is forced on the stray superscript. For noncompact spaces use the appropriate compact-support/Borel--Moore pairing.}
\GSfield{Fate}{\citet[duality theorem]{goreskymacpherson1977dualite}; detailed proof in \citet[Sec.~3.3]{goreskymacpherson1980intersection}.}
\GSfield{Present status}{As of 10 November 2026: Established with these hypotheses. Integral perfect duality does not follow in general.}
\GSfield{Assessment}{literal statement: correct with missing hypotheses; publication match: exact; verification: read the result; development: annotated; confidence: high}
\end{GSitem}
\begin{GSitem}{1979-484.3}{E}{Schubert closures supply the relevant even stratification}
\GSfield{Kind; source}{construction; \GSpage{P000907}}
\GSfield{The notes}{The note first says strange cohomology of Schubert cells, then explicitly writes \ru{шуберт. кл. (замкн.)} and uses smaller cells as the strata.}
\GSfield{Statement}{For a connected complex reductive group \(G\), Borel \(B\), Weyl group \(W\), and \(t\in W\), set \(X_t=\overline{BtB/B}\subset G/B\). Its Bruhat cells \(BsB/B\), \(s\le t\), are smooth affine complex spaces of dimensions \(\ell(s)\); they give the even real-dimensional stratification used for the local intersection-cohomology construction.}
\GSfield{Reconstruction}{The singular space is a Schubert variety, the closure of a cell. An open Bruhat cell itself is smooth and cannot carry the claimed singularity information.}
\GSfield{Fate}{The geometric theorem is \citet{kazhdanlusztig1980schubert}, published in 1980. A primary contemporary witness, \citet[p.~261, footnote; p.~301, reference~11]{gabberjoseph1981towards}, explicitly cites the Harvard 1979 preprint and distinguishes its solved geometric problem from the later character-formula proofs. Thus the 1980 publication date does not make this a development after the January talk. The original KL proof itself was not retrieved.}
\GSfield{Present status}{As of 10 November 2026: The closure/open-cell distinction is essential.}
\GSfield{Assessment}{literal statement: correct with missing hypotheses; publication match: exact; verification: publisher metadata checked; development: annotated; confidence: high}
\end{GSitem}
\begin{GSitem}{1979-484.4}{E}{Local intersection cohomology computes Kazhdan--Lusztig polynomials}
\GSfield{Kind; source}{theorem; \GSpage{P000907}}
\GSfield{The notes}{(5) writes \(P_{s,t}(q)=q^{?}\sum_i\dim IH_i(X_t,X_t-y_s)q^i\). The exponent is literally a question mark; the use of special characteristic-\(p\) cohomology is mentioned.}
\GSfield{Statement}{For \(s\le t\) in the Weyl group of a complex reductive group, let \(IC_{X_t}\) be the rational intersection-cohomology complex normalized to restrict to \(\mathbb Q[\ell(t)]\) on the open cell. At \(y_s\) in the \(s\)-cell, \[P_{s,t}(q)=\sum_{j\ge0}\dim\mathcal H^{2j-\ell(t)}(IC_{X_t})_{y_s}\,q^j.\] Stalk degrees of the opposite parity vanish. This is a local stalk formula, not the ordinary or global Poincare polynomial of \(X_t\).}
\GSfield{Reconstruction}{The displayed formula is an editorial modern normalization, replacing the visibly unfinished homological notation. No value is silently inserted for the original question mark; homological local groups require duality and indexing translations.}
\GSfield{Fate}{The geometric theorem is \citet{kazhdanlusztig1980schubert}, published in 1980. A primary contemporary witness, \citet[p.~261, footnote; p.~301, reference~11]{gabberjoseph1981towards}, explicitly cites the Harvard 1979 preprint and distinguishes its solved geometric problem from the later character-formula proofs. Thus the 1980 publication date does not make this a development after the January talk. The original KL proof itself was not retrieved.}
\GSfield{Present status}{As of 10 November 2026: An established characteristic-zero geometric interpretation. The notebook records the theorem with an unresolved grading convention.}
\GSfield{Assessment}{literal statement: not determinable; publication match: exact; verification: publisher metadata checked; development: annotated; confidence: high}
\end{GSitem}

\GStalk{485}{1980-02-04}{G. I. Olshanskii}{P000908}{Unitary representations of infinite classical groups}
The upper part of P000908 describes a multiparameter family for the infinite real general linear group, with finite-multiplicity discrete restriction to its orthogonal subgroup, and type-I/classification results for completed groups of fixed finite real rank. The completion topology and the distinction between the 1978 type-I theorem and the more complete 1980 construction are indispensable. MacPherson begins a separate contribution below the named heading on the same page.

\begin{GSitem}{1979-485.1}{E}{The Gaussian family with orthogonal-group labels}
\GSfield{Kind; source}{construction; \GSpage{P000908}}
\GSfield{The notes}{(1)--(3) display \(G=\bigcup GL(n,\mathbb R)\), \(K=\bigcup SO(n)\), and \(T_{\pi_1,s_1;\ldots;\pi_r,s_r}\), where \(\pi_i\in\widehat{O(n_i)}\) and \(s_i\in\mathbb R\).}
\GSfield{Statement}{For the Gaussian quasi-regular representation \(T_s\) of the finite-matrix group \(SL(\infty,\mathbb R)\) on \(L^2(\mathbb R^{\mathbb N},\mu)\), the commutant of \(T_s^{\otimes n}\) is generated by the natural \(O(n)\) action mixing the rows. Its \(O(n)\)-isotypic multiplicity spaces carry irreducible representations \(T_{\pi,s}\). For pairwise distinct real parameters \(s_i\), \(\bigotimes_iT_{\pi_i,s_i}\) is irreducible and determines the unordered parameter family. These representations extend to the polar completion \(GL(\infty,\mathbb R)=UR\), where \(U\) has the strong operator topology and \(R\) consists of positive invertible operators \(1+A\) with \(A\) Hilbert--Schmidt, with its Hilbert--Schmidt topology.}
\GSfield{Reconstruction}{The Gaussian model, extension topology and distinctness condition are reconstructed from the identified publication. The notebook gives only the parameter list, not a classification of all abstract representations of the algebraic union.}
\GSfield{Fate}{Exact family in \citet[Secs.~3 and 5, Thms.~5--7]{olshanskii1980construction}; presented 21 June and received 25 June 1979, published in volume 250(2), 1980. The exact issue-release day relative to 4 February was not verified. The result was already formulated in the note presented in June 1979; this is treated as contemporary exposition, not as a post-seminar first discovery. The precise public release day of the 1980 issue remains unverified.}
\GSfield{Present status}{As of 10 November 2026: A proved contemporary construction. The A label identifies the speaker	extquotesingle s 1980 publication; it does not imply that the manuscript first circulated after the seminar.}
\GSfield{Assessment}{literal statement: correct with missing hypotheses; publication match: exact; verification: read the result; development: annotated; confidence: high}
\end{GSitem}
\begin{GSitem}{1979-485.2}{E}{Discrete finite-multiplicity restriction and the maximal-compact analogy}
\GSfield{Kind; source}{claim; \GSpage{P000908}}
\GSfield{The notes}{The family restricts to \(K\) with \ru{конечнокр. дискр. спектр}; the notes call \(K\) analogous to a maximal compact subgroup and mention a unitary variant.}
\GSfield{Statement}{For these Gaussian finite-row representations, orthogonal transformations of the countably many columns preserve the Gaussian measure, so the restriction is independent of the parameters \(s_i\). Its Hilbert decomposition into continuous tensor-type \(SO(\infty)\) representations is discrete, with finite multiplicities. The complex construction replaces real Gaussian space and \(O(n_i)\) labels by their complex/unitary counterparts. The union \(SO(\infty)\) is not itself a compact group.}
\GSfield{Reconstruction}{The assertion is restricted to the family displayed in the notes and to its continuous tensor category; it is not a Peter--Weyl theorem for all representations of an arbitrary inductive-limit topology.}
\GSfield{Fate}{\citet[Sec.~3 and the proof of Thm.~6]{olshanskii1980construction} supplies the Gaussian model and its compact dual; compare the decomposition under restriction in \citet[Sec.~1]{olshanskii1978unitary}. The result was formulated in the note presented in June 1979, so the talk is treated as contemporary exposition. The precise public release day of the 1980 issue remains unverified.}
\GSfield{Present status}{As of 10 November 2026: Established in this representation class. The maximal-compact phrase is an analogy about restriction theory.}
\GSfield{Assessment}{literal statement: correct with missing hypotheses; publication match: partial; verification: read the result; development: annotated; confidence: medium}
\end{GSitem}
\begin{GSitem}{1979-485.3}{E}{Type I for the completed finite-real-rank groups}
\GSfield{Kind; source}{theorem; \GSpage{P000908}}
\GSfield{The notes}{(4)--(6) put closure bars over the unions \(SO(p,n)\), \(U(p,n)\), \(Sp(p,n)\), call them completions, and state that they are type I.}
\GSfield{Statement}{Fix finite \(p\). Use Olshanskii\textquotesingle s operator completions of \(SO_0(p,\infty)\), \(U(p,\infty)\), and \(Sp(p,\infty)\), not the discrete group or an unspecified final topology. Every continuous unitary representation generates a type-I von Neumann algebra; there are no factor representations of type II or III. In the norm/strong formulations specified by the paper, the continuous representation categories agree.}
\GSfield{Reconstruction}{The identity component \(SO_0\), Hilbert-space continuity and completion topology are supplied. Type I refers to the generated operator algebra of every continuous unitary representation.}
\GSfield{Fate}{\citet[Thm.~3.3]{olshanskii1978unitary}; received 20 October 1977, published 1978. The 1980 paper explicitly credits this earlier result.}
\GSfield{Present status}{As of 10 November 2026: Established for the stated completions. Changing the topology changes the assertion.}
\GSfield{Assessment}{literal statement: correct with missing hypotheses; publication match: exact; verification: read the result; development: annotated; confidence: high}
\end{GSitem}
\begin{GSitem}{1979-485.4}{E}{Exhaustion of finite-rank irreducibles by holomorphic constructions}
\GSfield{Kind; source}{theorem; \GSpage{P000908}}
\GSfield{The notes}{The same sentence says that all irreducible representations of the completed fixed-\(p\) groups have been described.}
\GSfield{Statement}{For the completed groups in the preceding item, the 1980 classification uses irreducible holomorphic representations of a specified \(\mathbb Z\)-cover of \(U(p,\infty)\). Their restrictions produce the irreducibles of \(SO_0(p,\infty)\) and, from \(U(2p,\infty)\), of \(Sp(p,\infty)\); holomorphic--antiholomorphic tensor products give the unitary case. The admissible highest-weight pairs in Theorems 1--2 determine the parameters. For the un-covered \(U(p,\infty)\), retain precisely the representations on which the covering kernel acts trivially.}
\GSfield{Reconstruction}{The cover, admissibility and descent condition are absent from the one-line notebook summary. We state the classification mechanism, not an unsupported claim for arbitrary representations of the algebraic union.}
\GSfield{Fate}{\citet[Thms.~1--4]{olshanskii1980construction} proves the exhaustion; manuscript June 1979, publication 1980. The 1978 paper proves type I and constructs many representations but does not claim this full exhaustion. The result was already formulated in the note presented in June 1979; this is treated as contemporary exposition, not as a post-seminar first discovery. The precise public release day of the 1980 issue remains unverified.}
\GSfield{Present status}{As of 10 November 2026: Established in the stated continuous category; the two chronology layers are kept separate.}
\GSfield{Assessment}{literal statement: correct with missing hypotheses; publication match: exact; verification: read the result; development: annotated; confidence: high}
\end{GSitem}

\GStalk{486}{1980-02-04}{R. D. MacPherson}{P000908-P000913}{Combinatorial Pontryagin classes and the dilogarithm}
The contribution begins below Olshanskii on P000908 and occupies P000908--P000913. Most of the construction is identifiable in MacPherson\textquotesingle s February 1977 Bourbaki exposition, published in 1978: locally linear stars, hypersimplices, configuration complexes, the compatible product complex and the signed \(1/48\) cocycle for \(p_1\). The end moves to Grassmannians and a dilogarithmic primitive, developed with Gelfand in their 1982 paper. The 11 February heading on P000914 begins the next talk. This is a return to the characteristic-class line in occurrences 1976-413.5--11, not a continuation of the January intersection-homology discussion.

\begin{GSitem}{1979-486.1}{E}{Locally linear flattenings are an additional triangulation hypothesis}
\GSfield{Kind; source}{definition; \GSpage{P000908}}
\GSfield{The notes}{\(X\) is a simplicial four-manifold, and for every simplex \(\sigma\) the notes assume \(\psi:\operatorname{St}\sigma\to\mathbb R^4\), linear on simplices and homeomorphic near the interior of \(\sigma\). A two-dimensional sketch illustrates the star.}
\GSfield{Statement}{Work with a finite triangulated PL four-manifold, with a local flattening of each simplex: a simplicewise affine map from its star to \(\mathbb R^4\) that is a homeomorphism on a neighborhood of the relative interior. Equivalently, the normal cone on the link of a codimension-\(k\) simplex has an appropriate linear flattening in \(\mathbb R^k\). This condition is imposed on the triangulation; it is not inferred merely from the word PL.}
\GSfield{Reconstruction}{The normal-space formulation explains the later dependence on codimension. A chosen triangulation of a PL manifold need not automatically satisfy this particular linear-star condition.}
\GSfield{Fate}{\citet[pp.~105--106]{macpherson1978combinatorial}; exposition February 1977, publication 1978.}
\GSfield{Present status}{As of 10 November 2026: A precise hypothesis of the construction, not a new 1980 assertion of universal linear realizability.}
\GSfield{Assessment}{literal statement: correct with missing hypotheses; publication match: exact; verification: read the result; development: annotated; confidence: high}
\end{GSitem}
\begin{GSitem}{1979-486.2}{E}{The hypersimplex over a simplex and its elementary models}
\GSfield{Kind; source}{construction; \GSpage{P000909}}
\GSfield{The notes}{(5)--(6) prescribe \(h(v)\in[0,1]\), \(h=1\) on vertices of \(\sigma\), and \(\sum_vh(v)=5\). Examples are a point, a segment, a triangle and an octahedron. Vertices are labelled by the places where \(h=1\).}
\GSfield{Statement}{For a simplex \(\sigma\) in a closed four-manifold, let \(V_\sigma\) be the vertex set of its star and put \[H_\sigma=\{h\in[0,1]^{V_\sigma}:h|_{V(\sigma)}=1,\ \sum h=5\}.\] If \(\sigma\) has codimension \(k\), the remaining coordinates sum to \(k\). All vertices are zero--one functions. Thus a top-dimensional simplex gives a point; a codimension-one simplex gives a segment. For codimension two, three remaining vertices give \(\Delta(2,3)\), a triangle, and four give \(\Delta(2,4)\), an octahedron.}
\GSfield{Reconstruction}{The notebook uses two-dimensional pictures for four-dimensional normal combinatorics. The exact cube slice is \(\sum h_i=2\), not an arbitrary equidistant cutting plane.}
\GSfield{Fate}{\citet[Sec.~1]{macpherson1978combinatorial}; compare the earlier seminar occurrence 1976-413.10.}
\GSfield{Present status}{As of 10 November 2026: An elementary convex-polytope construction. Boundary simplices require a relative variant of the segment example.}
\GSfield{Assessment}{literal statement: correct with missing hypotheses; publication match: exact; verification: read the result; development: annotated; confidence: high}
\end{GSitem}
\begin{GSitem}{1979-486.3}{E}{Gluing the hypersimplicial space over the original complex}
\GSfield{Kind; source}{construction; \GSpage{P000908-P000909}}
\GSfield{The notes}{For \(\sigma\subset\sigma^{\prime}\), the notes put \(H_{\sigma^{\prime}}\subset H_\sigma\), identify the common part \(H_{\sigma^{\prime}}\times\sigma\), and glue the products. They also draw the original and dual cell structures.}
\GSfield{Statement}{Extend a function on \(V_{\sigma^{\prime}}\) by zero to \(V_\sigma\); this embeds \(H_{\sigma^{\prime}}\) as a face of \(H_\sigma\). Glue the products \(H_\sigma\times\sigma\) along these face inclusions to obtain \(H\to X\). Use the barycentric/common subdivision \(X^{\prime}\) when comparing the primal simplex cells with the dual cells \(\sigma^D\); cellular means that each skeleton maps into the corresponding skeleton.}
\GSfield{Reconstruction}{The extension-by-zero rule and the common subdivision make the arrows in the notebook unambiguous.}
\GSfield{Fate}{\citet[Sec.~1]{macpherson1978combinatorial}.}
\GSfield{Present status}{As of 10 November 2026: Established; incidence reverses under passage from simplices to dual cells.}
\GSfield{Assessment}{literal statement: correct with missing hypotheses; publication match: exact; verification: read the result; development: annotated; confidence: high}
\end{GSitem}
\begin{GSitem}{1979-486.4}{E}{Configuration spaces and their degeneracy loci}
\GSfield{Kind; source}{definition; \GSpage{P000910}}
\GSfield{The notes}{\(C_\sigma\) first means all flattenings \(\psi\). General position requires every set consisting of all vertices of \(\sigma\) and \(k\) other star vertices, five points in total, to be independent. Speciality \(c\) counts the failures.}
\GSfield{Statement}{For a codimension-\(k\) simplex, pass to normal flattenings of the cone on its link in \(\mathbb R^k\), modulo \(GL(k,\mathbb R)\). A configuration is generic if each relevant \(k\)-subset of link vertices has nonzero determinant. The \(c\)-fold degeneracy locus has exactly \(c\) zero determinants. It is semialgebraic and has finitely many connected components. In the original four-space, these are the affine-independence tests on the displayed five points.}
\GSfield{Reconstruction}{The normal quotient and the word exactly are supplied. The notes subsequently reuse \(C_\sigma\) for a finite cell model; the raw realization space and that model are distinct.}
\GSfield{Fate}{\citet[Sec.~2]{macpherson1978combinatorial}.}
\GSfield{Present status}{As of 10 November 2026: A precise finite semialgebraic construction. No assertion that all its connected pieces are balls follows from the definition.}
\GSfield{Assessment}{literal statement: correct with missing hypotheses; publication match: exact; verification: read the result; development: annotated; confidence: high}
\end{GSitem}
\begin{GSitem}{1979-486.5}{F}{Does the degeneracy decomposition itself form a cell structure?}
\GSfield{Kind; source}{question; \GSpage{P000910}}
\GSfield{The notes}{After defining speciality, the recorder writes:\par\ru{Неизвестно: стратификация или клет. разбиение}.}
\GSfield{Statement}{The identifiable issue is whether the degeneracy pieces of the raw flattening quotient, or their connected components, have the topology and incidence properties of cells. The notes do not specify which decomposition, whether a refinement is allowed, or what notion of stratification is intended. Therefore they do not determine a unique yes/no theorem.}
\GSfield{Reconstruction}{This uncertainty is retained rather than converted into a blanket conjecture about all realization spaces. Semialgebraic refinement and the separately constructed finite complex answer different questions.}
\GSfield{Fate}{\citet[Sec.~2]{macpherson1978combinatorial} explicitly constructs a cell complex from components; it does not identify every raw component with an open cell.}
\GSfield{Present status}{As of 10 November 2026: No exact standalone problem can be recovered from this sentence. It is not evidence that the whole Pontryagin construction remained undeveloped.}
\GSfield{Assessment}{literal statement: not determinable; publication match: partial; verification: read the result; development: annotated; confidence: medium}
\end{GSitem}
\begin{GSitem}{1979-486.6}{E}{The low-codimension configuration cell models}
\GSfield{Kind; source}{construction; \GSpage{P000910-P000911}}
\GSfield{The notes}{For \(k=4\), one point is assigned to every generic component; for \(k=3\), walls of speciality one produce edges; for \(k=2\), speciality-two components give squares. The later point case mistakenly repeats \(k=2\).}
\GSfield{Statement}{Write \(K_\sigma\) for the finite model. For codimension four it is a discrete set. For codimension three it is the graph with vertices the generic components and an edge for each connected one-degeneracy wall, joining the two adjacent components. For codimension two fill the four-cycle around a two-degeneracy component by a square. For codimension one or zero put \(K_\sigma=\{*\}\). Crossing one determinant wall changes a sign, hence its two endpoints are distinct and no loop edge arises.}
\GSfield{Reconstruction}{The notation \(K_\sigma\) separates the model from the raw configuration space. The second point-case index must be zero, not two.}
\GSfield{Fate}{\citet[Sec.~2, pp.~109--111]{macpherson1978combinatorial}.}
\GSfield{Present status}{As of 10 November 2026: The stated finite construction is established. It is not a theorem that arbitrary higher-codimension realization spaces have such a simple cell decomposition.}
\GSfield{Assessment}{literal statement: false as written; publication match: exact; verification: read the result; development: annotated; confidence: high}
\end{GSitem}
\begin{GSitem}{1979-486.7}{E}{Connectivity of the codimension-three configuration graph}
\GSfield{Kind; source}{theorem; \GSpage{P000911}}
\GSfield{The notes}{The first line says \ru{Граф связен (трудная теорема!)}.}
\GSfield{Statement}{For a codimension-three simplex in the locally flattenable triangulated manifold under consideration, the configuration graph described in the preceding item is connected: generic local flattenings can be joined by passages through one-degeneracy walls.}
\GSfield{Reconstruction}{The assertion is about this local-star graph. It is not a general connectivity claim for all oriented-matroid realization spaces.}
\GSfield{Fate}{\citet[Note~2, p.~122]{macpherson1978combinatorial} explicitly supplies the connectivity justification via Cairns\textquotesingle s 1944 isotopy theorem.}
\GSfield{Present status}{As of 10 November 2026: Established in the indicated setting; the difficult step was already available before the talk.}
\GSfield{Assessment}{literal statement: correct with missing hypotheses; publication match: exact; verification: read the result; development: annotated; confidence: high}
\end{GSitem}
\begin{GSitem}{1979-486.8}{E}{Restriction maps glue the configuration models}
\GSfield{Kind; source}{construction; \GSpage{P000911}}
\GSfield{The notes}{The margin asks whether the \(C_\sigma\times\sigma^D\) glue naturally; the next lines restrict \(\psi\) from a larger star to a smaller one and draw \(C_{\sigma^{\prime}}\to C_\sigma\) when \(\sigma^{\prime}\subset\sigma\).}
\GSfield{Statement}{If \(\sigma^{\prime}\subset\sigma\), restriction followed by projection to the corresponding normal quotient induces a map of configuration spaces and of the finite cell models \(K_{\sigma^{\prime}}\to K_\sigma\). These compatible maps glue the products with dual cells to a cellular space \(C\to X^D\). Using \(X^{\prime}\), it can be combined with \(H\to X\).}
\GSfield{Reconstruction}{The normal projection, needed because codimension changes, is implicit in the notebook word restriction. The question is answered immediately by the construction rather than left as a newly open problem.}
\GSfield{Fate}{\citet[Sec.~2]{macpherson1978combinatorial}.}
\GSfield{Present status}{As of 10 November 2026: Established. This is related to the compatibility theme of 1976-413.11 but does not restate a new unsolved boundary-chain problem.}
\GSfield{Assessment}{literal statement: correct with missing hypotheses; publication match: exact; verification: read the result; development: annotated; confidence: high}
\end{GSitem}
\begin{GSitem}{1979-486.9}{E}{The compatible product complex Gamma}
\GSfield{Kind; source}{construction; \GSpage{P000908-P000911}}
\GSfield{The notes}{The notes draw \(\Gamma\subset H\times_XC\). The bottom of P000911 says a product is included when the five-point configurations specified by vertices of the hypersimplex face are generic.}
\GSfield{Statement}{Over the common subdivision, retain a product cell of the hypersimplicial and configuration factors only when every basis-sized subset indicated by a vertex of its hypersimplex face is nondegenerate for the corresponding flattening component. The resulting compatible cell subcomplex \(\Gamma\) has a projection \(\pi:\Gamma\to X\). Its allowed incidences make the characteristic cocycle well-defined.}
\GSfield{Reconstruction}{The handwritten membership description is compressed; the criterion is supplied in the vocabulary of the published construction. It does not imply that every product cell is admissible.}
\GSfield{Fate}{\citet[Sec.~3]{macpherson1978combinatorial}.}
\GSfield{Present status}{As of 10 November 2026: An established compatibility construction; the projection is not claimed to be an ordinary fiber bundle.}
\GSfield{Assessment}{literal statement: correct with missing hypotheses; publication match: exact; verification: read the result; development: annotated; confidence: high}
\end{GSitem}
\begin{GSitem}{1979-486.10}{E}{The signed one-forty-eighth cocycle on four-cells}
\GSfield{Kind; source}{construction; \GSpage{P000912}}
\GSfield{The notes}{The four face types have five variable coordinates and respectively three, two, one or zero additional coordinates fixed at one, above a vertex of \(X\). The two extreme types have value zero; the middle types have \(\widetilde P=\pm1/48\), with sign \((-1)^i\) from ten orientations.}
\GSfield{Statement}{For an oriented four-cell lying over a vertex, a four-dimensional hypersimplex face has five varying coordinates and \(b=0,1,2,3\) other coordinates fixed at one besides the base vertex. The varying coordinates sum to \(4-b\), giving \(\Delta(4,5),\Delta(3,5),\Delta(2,5),\Delta(1,5)\). The characteristic cochain vanishes on the extreme types and on the other prescribed four-cells. On the middle types project the five varying vertices along the fixed span to \(\mathbb R^2\) or \(\mathbb R^3\). With the cell orientation determined by their ordering, its value is \((-1)^i/48\), where \(i\) counts positively oriented ordered pairs or triples among the ten relevant minors.}
\GSfield{Reconstruction}{The base must be a vertex for this displayed fixed-coordinate count. Reversing the auxiliary orientation replaces \(i\) by \(10-i\), so the sign is unchanged; reversing the cell orientation reverses the value.}
\GSfield{Fate}{\citet[Sec.~4]{macpherson1978combinatorial}.}
\GSfield{Present status}{As of 10 November 2026: An explicit rational cochain in the specified orientation convention. The normalization is not a universal assertion about all Pontryagin cocycle models.}
\GSfield{Assessment}{literal statement: correct with missing hypotheses; publication match: exact; verification: read the result; development: annotated; confidence: high}
\end{GSitem}
\begin{GSitem}{1979-486.11}{E}{The combinatorial cocycle represents the first Pontryagin class}
\GSfield{Kind; source}{theorem; \GSpage{P000909-P000912}}
\GSfield{The notes}{(4) says \(\pi^*p^1(X)=[\widetilde P]\), and the intervening construction is meant to produce \(\widetilde P\in C^4(\Gamma)\).}
\GSfield{Statement}{Under the local flattening hypotheses, the prescribed cochain is a cocycle and \([\widetilde P]=\pi^*p_1(TX)\in H^4(\Gamma;\mathbb Q)\). A homological four-section \(s\) of \(\pi\) yields a representative \(s^*\widetilde P\) of \(p_1(X)\). Such a section is a chain map right-inverse through degree four; it is not necessarily a continuous section of \(\pi\).}
\GSfield{Reconstruction}{The rational coefficients and homological-section interpretation explain how the pulled-back identity computes a class on \(X\). These details come from the publication, not an additional theorem claimed in the notebook.}
\GSfield{Fate}{\citet[Sec.~3, Propositions~1--2; Sec.~4, theorem and Corollary~1]{macpherson1978combinatorial}.}
\GSfield{Present status}{As of 10 November 2026: Proved before 1980. The later all-Pontryagin formula of \citet{gelfandmacpherson1992combinatorial} is a broader development, not the publication date of this first-class formula.}
\GSfield{Assessment}{literal statement: correct with missing hypotheses; publication match: exact; verification: read the result; development: annotated; confidence: high}
\end{GSitem}
\begin{GSitem}{1979-486.12}{E}{Six coordinate degeneracy divisors on the real Grassmannian}
\GSfield{Kind; source}{claim; \GSpage{P000912}}
\GSfield{The notes}{The notes mention six Schubert-cell closures, six coordinate two-planes in \(\mathbb R^4\), and planes nontransverse to each of them.}
\GSfield{Statement}{For each coordinate two-plane \(E_I\subset\mathbb R^4\), the locus \(D_I=\{L\in\mathrm{Gr}(2,4):L\cap E_I\ne0\}\) is a Schubert degeneracy hypersurface for a flag containing \(E_I\). There are six such coordinate choices; their complement is the locus where all six Plucker coordinates are nonzero. Separately, the Schubert decomposition for one fixed full flag has six cells of dimensions \(0,1,2,2,3,4\). The six coordinate divisors are not those six cell closures.}
\GSfield{Reconstruction}{The source compresses two uses of six. The distinction is supplied without silently equating the coordinate hypersurfaces with all cells for one flag.}
\GSfield{Fate}{The Grassmannian/configuration setting is in \citet[Sec.~6]{macpherson1978combinatorial}; the precise generic-minor locus is also described in \citet[Secs.~3.1--3.3]{rojas2026explicit}.}
\GSfield{Present status}{As of 10 November 2026: Standard incidence geometry. The notebook wording does not determine which full flags were intended.}
\GSfield{Assessment}{literal statement: correct with missing hypotheses; publication match: partial; verification: read the result; development: annotated; confidence: medium}
\end{GSitem}
\begin{GSitem}{1979-486.13}{E}{Enhanced configurations of four rays and cross-ratio}
\GSfield{Kind; source}{construction; \GSpage{P000912-P000913}}
\GSfield{The notes}{The English phrase says four points in \(\mathbb RP^2\), but the next lines and figure say four rays in \(\mathbb R^2\), ordinary points in \(\mathbb RP^1\), and \(\mathbb R^4/L\) for a two-plane \(L\). Configurations are equivalent under \(GL(2,\mathbb R)\).}
\GSfield{Statement}{For \(L\in\mathrm{Gr}(2,4)\) with all relevant minors nonzero, the images of the four positive coordinate rays in \(\mathbb R^4/L\) are nonzero rays in a two-dimensional space. Up to linear isomorphism, they form an enhanced projective configuration: a projective-line configuration together with choices of lifts from lines to rays. For four distinct labelled projective points, an ordered cross-ratio is a complete continuous projective invariant; the ray choices add discrete data.}
\GSfield{Reconstruction}{Projectivization has dimension one. The genericity condition excludes zero quotient vectors and coincident projective points. Cross-ratio alone does not encode all enhancement signs.}
\GSfield{Fate}{\citet[Secs.~6--7]{macpherson1978combinatorial}.}
\GSfield{Present status}{As of 10 November 2026: Established. The configuration quotient may have non-Hausdorff boundary behavior; the generic interval components are sufficient for the cross-ratio discussion.}
\GSfield{Assessment}{literal statement: false as written; publication match: exact; verification: read the result; development: annotated; confidence: high}
\end{GSitem}
\begin{GSitem}{1979-486.14}{E}{A Pontryagin form on the Grassmannian}
\GSfield{Kind; source}{construction; \GSpage{P000913}}
\GSfield{The notes}{The notes label \(P\) a four-form representing the Pontryagin class of the standard rank-two bundle \(\xi\) on \(G_{2,2}(\mathbb R)\), and say that a connection is constructed.}
\GSfield{Statement}{Let \(\xi\) be the tautological real two-plane bundle on \(\mathrm{Gr}(2,4)\), with the metric and projected connection induced from Euclidean space. Its Chern--Weil representative is \(P=-\frac{1}{8\pi^2}\operatorname{tr}(\Omega\wedge\Omega)\), where \(\Omega\) is the curvature. The de Rham class is \(p_1(\xi)\), independent of the connection. On the oriented double cover it equals the square of the Euler class.}
\GSfield{Reconstruction}{The connection, curvature normalization and orientation-cover observation supply a precise version of the note. If instead one uses the quotient bundle under the complementary-plane identification, the corresponding sign/convention must be tracked.}
\GSfield{Fate}{\citet[Secs.~5 and 8]{macpherson1978combinatorial} uses the invariant Grassmannian characteristic form.}
\GSfield{Present status}{As of 10 November 2026: Classical Chern--Weil theory; its later fiber integral is a different operation.}
\GSfield{Assessment}{literal statement: correct with missing hypotheses; publication match: exact; verification: read the result; development: annotated; confidence: high}
\end{GSitem}
\begin{GSitem}{1979-486.15}{A}{Fiber integration of the Pontryagin form and the Rogers dilogarithm}
\GSfield{Kind; source}{claim; \GSpage{P000913}}
\GSfield{The notes}{The scan, more clearly than the linearized transcription, gives \(\bar\pi_*(P)=d\Phi\) and labels \(\Phi\) a function, \ru{дилогарифм}. The projection sends a two-plane to its enhanced four-ray configuration.}
\GSfield{Statement}{On a generic component of the configuration quotient of \(\mathrm{Gr}(2,4)\), integration of the invariant first-Pontryagin four-form along the three-dimensional diagonal-scaling fibers, with the hypersimplicial compactification and an orientation convention, gives a one-form proportional to \(dL_R(x)\), where \(x\) is cross-ratio and \[L_R(x)=\operatorname{Re}\operatorname{Li}_2(x)+\tfrac12\log|x|\log|1-x|,\qquad x\in\mathbb R\setminus\{0,1\}.\] Thus locally the fiber integral has a dilogarithmic primitive. Its scalar and additive constants depend on fiber/sign and normalization choices.}
\GSfield{Reconstruction}{The source supplies neither the compactification nor a scalar normalization. The statement uses the real Rogers dilogarithm, not the Bloch--Wigner function; it does not silently assert an equality for an unspecified arbitrary connection.}
\GSfield{Fate}{Developed by \citet[Secs.~0 and 3]{gelfandmacpherson1982geometry}, published June 1982. The full modern computation was read in \citet[Sec.~3.7 and Thm.~3.7, case \(k=1\)]{rojas2026explicit}, preprint 26 October 2026. The 1982 paper\textquotesingle s publisher metadata, not its full proof, was checked.}
\GSfield{Present status}{As of 10 November 2026: The first-class relation is established. The 2026 preprint supplies a new explicit treatment of powers of \(p_1\); its broader claims are not substituted for the notebook\textquotesingle s elementary case.}
\GSfield{Assessment}{literal statement: correct with missing hypotheses; publication match: exact; verification: read the result; development: annotated; confidence: high}
\end{GSitem}

\GStalk{487}{1980-02-11}{R. D. MacPherson}{P000914}{Grassmannians, hypersimplices and enhanced configurations}
The actual source is the short, separately dated fragment P000914; P000915 starts Zakharov on 18 February. It states the idea that a Grassmannian point is determined by a hypersimplicial datum and an enhanced projective configuration, with \(n+k\) labelled vectors in a \(k\)-dimensional space. The projected-space dimension is written incorrectly, and the bottom of the fragment is clipped. The finite statement is already present in the 1977/1978 Bourbaki exposition; no missing continuation is invented.

\begin{GSitem}{1979-487.1}{E}{From a Grassmannian point to a vector and ray configuration}
\GSfield{Kind; source}{construction; \GSpage{P000914}}
\GSfield{The notes}{\(G_{n,k}\) consists of unoriented \(n\)-planes in \(\mathbb R^{n+k}\); the text defines enhanced configurations of \(n+k\) points in \(k\)-space, parenthetically calling its projective dimension \(k+1\).}
\GSfield{Statement}{For \(L\in\mathrm{Gr}(n,n+k)\), take the images of the coordinate vectors in \(V=\mathbb R^{n+k}/L\), of dimension \(k\), modulo linear isomorphisms of \(V\). The resulting spanning labelled vector configuration determines \(L\) as the kernel of the corresponding surjection \(\mathbb R^{n+k}\to V\). Passing to the images of positive coordinate rays forgets their positive lengths and gives an enhanced projective configuration; zero quotient vectors need a separate boundary stratum. Nonzero projective points lie in \(\mathbb RP^{k-1}\).}
\GSfield{Reconstruction}{The kernel construction, generic/nonzero qualification, and dimension \(k-1\) are editorial clarification. The printed \(k+1\) is not retained as a mathematical claim.}
\GSfield{Fate}{\citet[Sec.~6]{macpherson1978combinatorial}, exposition February 1977 and publication 1978.}
\GSfield{Present status}{As of 10 November 2026: An exact quotient/configuration construction with the stated boundary qualification.}
\GSfield{Assessment}{literal statement: false as written; publication match: exact; verification: read the result; development: annotated; confidence: high}
\end{GSitem}
\begin{GSitem}{1979-487.2}{E}{A configuration plus a hypersimplex point recovers the plane}
\GSfield{Kind; source}{claim; \GSpage{P000914}}
\GSfield{The notes}{The fragment\textquotesingle s central line reads:\par\ru{Элементу} \(G_{n,k}\) \ru{соотв. пара гиперсимплекс + конфигурация}.}
\GSfield{Statement}{Fix a spanning rank-\(k\) vector configuration \(w_1,\ldots,w_{n+k}\) and denote by \(\Delta_I\) its maximal minors. A convenient hypersimplicial coordinate is \[h_i=\frac{\sum_{I\ni i}|\Delta_I|}{\sum_{|I|=k}|\Delta_I|},\qquad0\le h_i\le1,\quad\sum_i h_i=k.\] On each effective orbit of positive independent coordinate rescalings, quotienting by its full stabilizer, this map is a homeomorphism onto the relative interior of the polytope spanned by the indicator vectors of subsets with \(\Delta_I\ne0\); its closure maps to that polytope. Thus the enhanced projective configuration and this point recover the Grassmannian datum. For a generic configuration the polytope is the full \(\Delta(k,n+k)\).}
\GSfield{Reconstruction}{The explicit map, absolute-minor convention and non-generic face polytope come from the identified source, not the clipped continuation. Alternative moment-map normalizations using squared minors give related coordinates and are not conflated with this formula. In a nongeneric disconnected configuration, the stabilizer may be larger than common scalar multiplication; the assertion concerns the effective orbit in the Grassmannian, not the unreduced parameter torus.}
\GSfield{Fate}{\citet[Sec.~6, proposition]{macpherson1978combinatorial}; the further Grassmannian development is \citet{gelfandmacpherson1982geometry}.}
\GSfield{Present status}{As of 10 November 2026: Established in the specified configuration/orbit formulation. The notebook does not formulate a new open question here.}
\GSfield{Assessment}{literal statement: correct with missing hypotheses; publication match: exact; verification: read the result; development: annotated; confidence: high}
\end{GSitem}

\GStalk{488}{1980-02-18}{V. E. Zakharov}{P000915-P000916}{Benney equations, conservation laws and Hamiltonian variables}
The actual session occupies P000915--P000916. Starting from an incompressible fluid with a free surface, Zakharov takes the hydrostatic long-wave approximation and announces additional conservation laws, simpler Hamiltonian variables, and an extension to two horizontal dimensions. Gelfand interprets the model as collective dynamics through a common field. The contemporary papers supply precise versions of the short announcements; no theorem of complete integrability is attributed to these notes. P000917 explicitly starts 22 September 1980, with Fuchs on intersection homology; P000918 continues that discussion and turns to Graev on integral geometry. Those next-year items are inventoried separately in \texttt{code/boundary\_inventory.md} and are not occurrences of this record.

\begin{GSitem}{1979-488.1}{E}{Hydrostatic reduction of the incompressible free-surface equations}
\GSfield{Kind; source}{method; \GSpage{P000915-P000916}}
\GSfield{The notes}{The fluid occupies \(0<z<h(r,t)\), with horizontal velocity \(u\), vertical velocity \(w\), and \(w_z+\operatorname{div}u=0\). The vertical equation on P000915 has \(-p_z+1\), while P000916 writes \(p_z+1=0\) and obtains \(p=h-z\).}
\GSfield{Statement}{For an inviscid incompressible fluid of constant density, a flat impermeable bed, gravitational acceleration scaled to one, and zero pressure at the surface, neglecting the entire vertical material acceleration gives \(p_z=-1\), hence \(p(r,z,t)=h(r,t)-z\). The horizontal equation is \[u_t+(u\cdot\nabla_r)u+w u_z+\nabla_r h=0,\qquad w_z+\operatorname{div}_r u=0.\] This is the hydrostatic long-wave model, not the full free-surface Euler system.}
\GSfield{Reconstruction}{The pressure sign, surface-pressure condition and scaling assumptions are made explicit. The bracket marking terms to neglect in the notebook does not include every vertical-acceleration term; \(w_t\) must also be negligible for the displayed reduction.}
\GSfield{Fate}{The long-wave system is the starting point, credited to Benney, in \citet[Sec.~1, Eqs.~(1)--(4)]{zakharov1980benney} and \citet[Sec.~1]{zakharov1981benney}. Their introduction identifies the previously studied model, while their new claims concern its structure.}
\GSfield{Present status}{As of 10 November 2026: The formal reduction is established under the stated approximation. Its validity over a specified error scale and time interval requires a separate asymptotic justification.}
\GSfield{Assessment}{literal statement: false as written; publication match: exact; verification: read the result; development: annotated; confidence: high}
\end{GSitem}
\begin{GSitem}{1979-488.2}{E}{The height equation is integrated mass conservation}
\GSfield{Kind; source}{theorem; \GSpage{P000916}}
\GSfield{The notes}{The third displayed equation is \(h_t+\operatorname{div}\int_0^h u\,dz=0\); the notes say that the one-dimensional system had been studied by Manin and others.}
\GSfield{Statement}{If \(h>0\) and \(u,w,h\) are smooth, impose \(w(r,0,t)=0\) and the kinematic free-surface condition \(h_t+u(r,h,t)\cdot\nabla_r h=w(r,h,t)\). Integrating incompressibility and differentiating the moving upper limit yields \[h_t+\operatorname{div}_r\int_0^{h(r,t)}u(r,z,t)\,dz=0.\] Together with the horizontal equation and incompressibility, this is the displayed Benney system.}
\GSfield{Reconstruction}{The bed and free-surface boundary conditions are required to derive the equation; they are not all written in the fragment.}
\GSfield{Fate}{\citet[Sec.~1]{zakharov1980benney} and \citet[Eqs.~(1)--(4)]{zakharov1981benney} explicitly start from this pre-existing long-wave system.}
\GSfield{Present status}{As of 10 November 2026: An exact mass balance within the smooth hydrostatic model. It does not require any assertion of complete integrability.}
\GSfield{Assessment}{literal statement: correct with missing hypotheses; publication match: exact; verification: read the result; development: annotated; confidence: high}
\end{GSitem}
\begin{GSitem}{1979-488.3}{A}{Additional branches generate conservation laws beyond the usual moment hierarchy}
\GSfield{Kind; source}{claim; \GSpage{P000916}}
\GSfield{The notes}{Zakharov announces \ru{счетное число законов сохранения (гораздо больше, чем было)}: a countable collection, substantially more than previously known. No densities are written.}
\GSfield{Statement}{A precise finite-layer reconstruction is \(\rho_{a,t}+(\rho_a u_a)_x=0\), \(u_{a,t}+u_a u_{a,x}+h_x=0\), \(h=\sum_{a=1}^{N}\rho_a\), with smooth fields and distinct velocities on the region considered. Set \[R(p)=p+\sum_{a=1}^{N}\frac{\rho_a}{p-u_a}.\] Each formal inverse branch \(R(p)=\lambda\) satisfies \(p_t+\partial_x(p^2/2+h)=0\). Expansion of the branch at infinity yields the standard moment densities; the additional branches approaching the \(u_a\) give further conservation-law sequences. Their integrals are conserved for periodic fields or suitable decay.}
\GSfield{Reconstruction}{The finite-layer system and formal expansions are reconstructed from the identified papers. The continuum-label limit needs additional analytic or formal-series hypotheses and, for the displayed non-polynomial formulas in the paper, a specified principal-value prescription. The notebook does not state a completeness theorem for these integrals.}
\GSfield{Fate}{\citet[Sec.~4, Eqs.~(48)--(52)]{zakharov1980benney}, received 17 December 1979 and published in 1980, develops the announcement; \citet[Sec.~4, Eqs.~(35)--(39)]{zakharov1981benney} displays the extra branches and involution argument. The latter explicitly says that completeness, needed for full integrability, is not proved.}
\GSfield{Present status}{As of 10 November 2026: The formal conservation hierarchies are established for the stated reductions. They do not by themselves prove global solvability or complete integrability of every smooth Benney solution.}
\GSfield{Assessment}{literal statement: correct with missing hypotheses; publication match: partial; verification: read the result; development: annotated; confidence: high}
\end{GSitem}
\begin{GSitem}{1979-488.4}{A}{Lagrangian layers give a simpler one-dimensional Hamiltonian bracket}
\GSfield{Kind; source}{construction; \GSpage{P000916}}
\GSfield{The notes}{Zakharov announces \ru{более простая гамильтонова структура}, and says that the idea is a reinterpretation of the equation.}
\GSfield{Statement}{Use a material vertical label \(a\) and a smooth monotone height coordinate \(z=z(x,a,t)\) with \(\rho=z_a>0\). In a fixed label interval, let \(h=\int\rho\,da\). The equations become \(\rho_t+\partial_x(\rho u)=0\), \(u_t+u u_x+h_x=0\). With periodic or decaying fields, \[E=\tfrac12\int\!\left(\int\rho u^2\,da+h^2\right)dx,\qquad \binom{\rho_t}{u_t}=-\begin{pmatrix}0&\partial_x\\\partial_x&0\end{pmatrix}\binom{\delta E/\delta\rho}{\delta E/\delta u},\] where the bracket includes \(\delta(a-b)\) in the label variable. Thus the bracket has constant coefficients in these variables.}
\GSfield{Reconstruction}{The precise fields, energy, boundary conditions and bracket are absent from the notes. The coordinate change is valid only while the layers remain smooth and the vertical parametrization stays invertible.}
\GSfield{Fate}{\citet[Sec.~2]{zakharov1980benney}, manuscript received December 1979, and \citet[p.~195, Eqs.~(11)--(12)]{zakharov1981benney} give this Hamiltonian formulation.}
\GSfield{Present status}{As of 10 November 2026: Established locally in these smooth variables. This simple bracket must not be confused with a proof of existence of global smooth solutions.}
\GSfield{Assessment}{literal statement: correct with missing hypotheses; publication match: exact; verification: read the result; development: annotated; confidence: high}
\end{GSitem}
\begin{GSitem}{1979-488.5}{A}{Two horizontal dimensions retain a Hamiltonian formulation under appropriate variables}
\GSfield{Kind; source}{claim; \GSpage{P000916}}
\GSfield{The notes}{The third announcement is only \ru{что-то для двумерной зависимости (с гамильтоновостью)}: something for two-dimensional dependence, with Hamiltonian structure.}
\GSfield{Statement}{One definite potential-flow version has horizontal \(r\in\mathbb R^2\), layer density \(\rho(r,a,t)>0\), potential \(\phi(r,a,t)\), \(u=\nabla_r\phi\), \(h=\int\rho\,da\), and \[E=\tfrac12\int\!\left(\int\rho|\nabla_r\phi|^2\,da+h^2\right)dr.\] The canonical equations \(\rho_t=\delta E/\delta\phi\), \(\phi_t=-\delta E/\delta\rho\) give \(\rho_t+\operatorname{div}_r(\rho u)=0\) and \(u_t+(u\cdot\nabla_r)u+\nabla_r h=0\). General vortical flows require additional Clebsch or noncanonical variables and their own local/topological hypotheses.}
\GSfield{Reconstruction}{This is a precise identified special case of an imprecise announcement, not a claim that the notebook proposed only potential flows. It does not assert two-dimensional complete integrability.}
\GSfield{Fate}{\citet[Secs.~2 and 5]{zakharov1980benney} develops multidimensional variables; \citet[Sec.~2]{zakharov1981benney} supplies the many-fluid interpretation. The match to the compressed third announcement is partial.}
\GSfield{Present status}{As of 10 November 2026: A Hamiltonian structure exists in the stated variables and domain; the notes contain no definite stronger two-dimensional theorem to mark solved or open.}
\GSfield{Assessment}{literal statement: correct with missing hypotheses; publication match: partial; verification: read the result; development: annotated; confidence: medium}
\end{GSitem}
\begin{GSitem}{1979-488.6}{A}{The many-fluid variables also give a kinetic reinterpretation}
\GSfield{Kind; source}{method; \GSpage{P000916}}
\GSfield{The notes}{The general idea is to reinterpret the equation, followed immediately by Gelfand\textquotesingle s description of collective evolution through a common field.}
\GSfield{Statement}{For smooth layer fields satisfying \(\rho_t+\operatorname{div}_r(\rho u)=0\), \(u_t+(u\cdot\nabla_r)u=-\nabla_r h\), define the distribution in velocity \[f(r,v,t)=\int\rho(r,a,t)\,\delta(v-u(r,a,t))\,da,\qquad h=\int f\,dv.\] Then \(f_t+v\cdot\nabla_r f-\nabla_r h\cdot\nabla_v f=0\) distributionally. This gives a Vlasov-type collective-field equation. Its equivalence with a smooth hydrostatic velocity profile requires an invertible layer coordinate and, when solving for the velocity label, an appropriate monotonicity condition.}
\GSfield{Reconstruction}{The distributional formula makes the identified reinterpretation explicit. In two velocity dimensions a one-parameter family of layers gives a special class of velocity distributions; it is not every Vlasov solution.}
\GSfield{Fate}{\citet[Sec.~2]{zakharov1980benney} and \citet[p.~195, Eqs.~(13)--(15)]{zakharov1981benney} state the kinetic reformulation of this Benney model.}
\GSfield{Present status}{As of 10 November 2026: Established for the stated smooth layer solutions, with the distribution interpreted weakly. The converse needs the additional invertibility hypotheses.}
\GSfield{Assessment}{literal statement: correct with missing hypotheses; publication match: partial; verification: read the result; development: annotated; confidence: high}
\end{GSitem}
\begin{GSitem}{1979-488.7}{E}{Gelfand interprets the system as collective differential equations}
\GSfield{Kind; source}{remark; \GSpage{P000916}; I. M. Gelfand}
\GSfield{The notes}{Gelfand describes particles that do not interact directly, but jointly produce a field (an electric field is mentioned) which acts on every particle; he proposes studying such collective equations after ordinary differential equations.}
\GSfield{Statement}{A precise schematic model is a kinetic transport equation \(\partial_t f+v\cdot\nabla_xf+F[f]\cdot\nabla_vf=0\), where the common field \(F[f]\) is determined by the distribution through a separately specified field equation. The Benney reinterpretation uses \(F[f]=-\nabla_x\int f\,dv\); an electrostatic plasma uses a Poisson field instead. These are different closures of the same collective-dynamics idea.}
\GSfield{Reconstruction}{The schematic equation and distinction between closures explain the analogy; no unique field equation or new well-posedness problem can be extracted from the intervention alone.}
\GSfield{Fate}{The general collective-field concept predates the talk: \citet[Secs.~1--2]{vlasov1968vibrational} is the English version of the 1938 paper, as the publisher\textquotesingle s Russian reprint record explicitly confirms. The specific Benney closure is given in \citet[Eq.~(15)]{zakharov1981benney}.}
\GSfield{Present status}{As of 10 November 2026: An established interpretation, not a newly stated precise conjecture. The original Vlasov publisher record and table of contents were checked; its full proof was not needed for the displayed Benney identity.}
\GSfield{Assessment}{literal statement: correct with missing hypotheses; publication match: thematic only; verification: publisher metadata checked; development: annotated; confidence: high}
\end{GSitem}

\section{Developments}\label{dev:1979}
\subsection{Berezin's normal-form question: an algorithmic answer and the limit of the Jordan analogy}
\label{dev:berezin-normal}

\paragraph{What is being asked.}
On P000861--P000862 Berezin passes from diagonalization of an even supermatrix with distinct body eigenvalues to \ru{Аналоги жорд. клеток здесь неизвестны} and asks for an analogue of normal forms. The strongest precise reading is conjugacy of the even algebra
\[
 \mathcal A=M_{p|q}(\Lambda_N)_{\bar0},\qquad
 \Lambda_N=\bigwedge_{\mathbb C}(\xi_1,\ldots,\xi_N),
\]
under its invertible even elements, for fixed finite \(p,q,N\). The diagonal blocks have even coefficients and the off-diagonal blocks odd coefficients. We do not extend the notebook to an infinite Grassmann algebra, nor require eigenvalue and block-size data alone to classify every matrix. This subsection gives a retrospective resolution and explanatory examples for GP-1979-01, not a newly announced open problem.

\paragraph{Verified literature and chronology.}
\citet[Theorems~1.1 and~1.3; \S1.4]{sergeichuk2000canonical} supplies canonical representatives under the units of a finite-dimensional complex matrix algebra. Its bibliography identifies the 1999 Bielefeld preprint 99-070; the journal publication is 2000 and the arXiv deposit 2007. The reduction below makes its application here explicit. The meaning of ``wild'' used below is containment of simultaneous similarity of arbitrary matrix pairs; see \citet[\S1 and Theorem~1.1]{belitskiisergeichuk2003complexity}, published in 2003, with an author copy deposited in 2007. Neither term implies the nonexistence of canonical-form algorithms. No claim is made that these publications identify a direct historical response to Berezin's seminar question.

\begin{proposition}[Reduction to an ordinary matrix-algebra action \ours]
Let \(\mathcal A\) be a finite-dimensional unital complex algebra. Internal conjugacy in \(\mathcal A\) is exactly similarity of its left regular matrices under the units of the matrix algebra \(\Gamma=L(\mathcal A)\). Thus the cited canonical-form theorem applies to the notebook's finite supermatrix problem.
\end{proposition}
\begin{proof}
Choose a complex basis of \(\mathcal A\) and let \(L_a(x)=ax\). The map \(a\mapsto L_a\) is an injective unital algebra homomorphism, since \(L_a(1)=a\). If \(a\) is a unit then \(L_a\) is invertible. Conversely, if \(L_a\) is invertible, there is \(b\) with \(ab=1\); injectivity of \(L_a\) and \(a(ba-1)=0\) give \(ba=1\). Consequently \(\Gamma^\times=L(\mathcal A^\times)\), and
\[
 L_bL_aL_b^{-1}=L_{bab^{-1}}.
\]
Both implications of the conjugacy equivalence follow by injectivity. The algebra of even supermatrices is finite dimensional and closed under its ordinary matrix multiplication with Grassmann coefficients, so it satisfies the hypothesis. Fixing a basis and applying the cited matrix-algebra reduction and canonical algorithm gives representatives; this is an existence and reduction argument, not an implementation of that algorithm in the accompanying code.
\end{proof}

\paragraph{A tractable repeated-eigenvalue case.}
For the pure even subcase \(q=0\), \((\Lambda_2)_{\bar0}=\mathbb C[\varepsilon]/(\varepsilon^2)\), with \(\varepsilon=\xi_1\xi_2\). This already contains nontrivial repeated-body phenomena.

\begin{proposition}[Semisimple body over the dual numbers \ours]
Let \(A=A_0+\varepsilon A_1\in M_n(\mathbb C[\varepsilon]/(\varepsilon^2))\), and suppose \(A_0\) is diagonalizable. Write its distinct eigenvalues as \(\lambda_1,\ldots,\lambda_s\), with multiplicities \(m_i\). Up to ordering the eigenvalue groups and Jordan blocks, a complete normal form is
\[
 \bigoplus_{i=1}^s\bigl(\lambda_i I_{m_i}+\varepsilon J_i\bigr),
\]
where \(J_i\) is the ordinary Jordan form of the restriction of \(A_1\) to the \(\lambda_i\)-eigenspace, after projection along the other eigenspaces. These Jordan forms, together with the eigenvalues, are complete invariants in this subcase.
\end{proposition}
\begin{proof}
A constant change of basis puts \(A_0=D=\bigoplus\lambda_iI_{m_i}\). For any matrix \(S\), exact multiplication in the dual numbers gives
\[
 (I+\varepsilon S)(D+\varepsilon A_1)(I-\varepsilon S)
 =D+\varepsilon(A_1+[S,D]).
\]
In the \((i,j)\)-block, \([S,D]_{ij}=(\lambda_j-\lambda_i)S_{ij}\). Choose \(S_{ij}=-(A_1)_{ij}/(\lambda_j-\lambda_i)\) when \(i\ne j\), and \(S_{ii}=0\). The off-diagonal blocks disappear. Constant block-diagonal changes then put the surviving blocks into ordinary Jordan form. For completeness, any similarity between two such matrices has body commuting with \(D\), after the same ordering of distinct eigenvalues. Its body is therefore block diagonal. Its nilpotent part changes the first-order coefficient only by a commutator with \(D\), whose diagonal blocks are zero. Hence each surviving block changes only by the corresponding constant similarity. Conversely those block similarities provide an equivalence.
\end{proof}

\begin{example}[A diagonal body need not suffice for diagonalization \ours]
The matrix \(\varepsilon J_2(0)\) has zero, hence diagonalizable, body but is not diagonalizable over the dual numbers. If an invertible \(P=P_0+\varepsilon P_1\) diagonalized it, comparison of the \(\varepsilon\)-coefficient would diagonalize the nonzero nilpotent matrix \(J_2(0)\) over \(\mathbb C\). A diagonal nilpotent matrix is zero, a contradiction. This does not contradict the notebook's theorem, which requires \emph{distinct} body eigenvalues.
\end{example}

\begin{proposition}[Arbitrary matrix pairs occur already with three generators \ours]
For every \(n\), define
\[
 A(X,Y)=\lambda I_n+\xi_1\xi_2X+\xi_1\xi_3Y
 \quad\hbox{in }M_n((\Lambda_3)_{\bar0}).
\]
Then \(A(X,Y)\) and \(A(X',Y')\) are similar by an invertible matrix over \((\Lambda_3)_{\bar0}\) if and only if the ordinary pairs \((X,Y)\) and \((X',Y')\) are simultaneously similar.
\end{proposition}
\begin{proof}
The positive-degree ideal \(J\) of \((\Lambda_3)_{\bar0}\) is spanned by the three degree-two monomials and has \(J^2=0\). Any invertible matrix is \(P=P_0+R\), with \(P_0\in GL_n(\mathbb C)\) and \(R\in M_n(J)\). In \(PAP^{-1}\), the terms involving both \(R\) and \(A-\lambda I\) vanish, and conjugation of \(\lambda I\) changes nothing. Comparison of the independent coefficients \(\xi_1\xi_2\) and \(\xi_1\xi_3\) gives \(X'=P_0XP_0^{-1}\) and \(Y'=P_0YP_0^{-1}\). Conversely that constant matrix conjugates the two Grassmann matrices.
\end{proof}

\paragraph{What is settled and what remains.}
As of 10 November 2026, the finite algorithmic interpretation has the preceding answer; it should not be advertised as open. A short uniform Jordan-block list is a different demand and faces the explicit matrix-pair reduction. No impossibility of every conceivable canonical parametrization is claimed. One further bounded run could implement the reduced-algebra algorithm for a chosen \((p,q,N)\), or classify a specified small rank and body type. It could not justify a general ``no normal forms'' statement from wildness. The dual-number proof was checked a second way by expanding the three factors coefficient by coefficient, and the exact rational off-block computation in \texttt{code/verify\_berezin.py} agrees. The pair reduction was independently checked through the equivalent intertwining equation \(PA(X,Y)=A(X',Y')P\): its two degree-two equations are exactly \(P_0X=X'P_0\) and \(P_0Y=Y'P_0\).

\subsection{Berezin's even commutant problem: all coefficients and a complete four-generator case}
\label{dev:berezin-commutant}

\paragraph{Notebook formulation, scope, and literature limit.}
P000862 asks to describe every \(B=(b_{ik})\), \(b_{ik}\in\Lambda_{\bar0}\), satisfying \(AB=BA\) for a given \(n\times n\) matrix with \(a_{ik}\in\Lambda_{\bar1}\). It explicitly mentions four generators and says that independent entries can be handled. No additional differential equation, instanton constraint, gauge quotient, or prescribed format of the answer is given. We therefore fix \(\mathbb K=\mathbb R\) or \(\mathbb C\), finite \(N\), and arbitrary odd coefficients; the object is the actual even centralizer, with ordinary matrix multiplication. An all-odd \(n\times n\) coefficient matrix here must not be confused with an even \((p|q)\) supermatrix from the preceding problem.

The search record distinguishes this exact question from general polynomial-identity results and from centralizers of ordinary matrices. No primary source giving the following complete \(n=2,N=4\) stratification was located. That negative retrieval result is \emph{not} a priority claim or evidence that the finite equation was globally unsolved. Coefficient expansion already solves its effective finite-dimensional content by elementary linear algebra; the occurrence is therefore classified E. The bounded explanatory contribution is to give that expansion, prove a uniform description in the stated small case, and supply independently checked code. This is the precise scope of development for GP-1979-02. A stronger compact structural account for arbitrary \(n,N\) is not specified by the notebook, and we do not invent its requirements or declare it open. The historical attribution of Berezin's intended form of an answer was not recovered.

\begin{proposition}[Exact finite coefficient description \ours]
For increasing subsets \(I\subseteq\{1,\ldots,N\}\), put \(\xi_I=\prod_{i\in I}\xi_i\), and write
\[
 A=\sum_{|I|\ {\rm odd}}\xi_I A_I,\qquad
 B=\sum_{|J|\ {\rm even}}\xi_J B_J,
 \qquad A_I,B_J\in M_n(\mathbb K).
\]
For disjoint \(I,J\), let \(\epsilon(I,J)=(-1)^{\#\{(i,j)\in I\times J:i>j\}}\). Then the entire even commutant is the kernel of the linear map with equations
\begin{equation}
 \sum_{\substack{I\sqcup J=K\\|I|\ {\rm odd},\ |J|\ {\rm even}}}
 \epsilon(I,J)[A_I,B_J]=0
 \quad\text{for each odd subset }K.
 \label{eq:commutant-coefficients}
\end{equation}
For \(N\geq1\), there are \(2^{N-1}n^2\) scalar unknowns and the same number of scalar equations before redundancies. Solving this system gives a basis, and multiplication of basis elements gives the full algebra multiplication table.
\end{proposition}
\begin{proof}
Repeated generators give zero. For disjoint sets, sorting the product gives the displayed sign. Since \(|J|\) is even, \(\xi_J\xi_I=\xi_I\xi_J\). The coefficient of \(\xi_K\) in \(AB-BA\) is therefore exactly the left side of \eqref{eq:commutant-coefficients}. Linear independence of the exterior monomials proves necessity and sufficiency, with no genericity assumption. The kernel is closed under multiplication because \([A,BC]=[A,B]C+B[A,C]\), and contains the identity. Coordinates of the products in a kernel basis determine its algebra structure.
\end{proof}

\begin{remark}[Two useful boundary cases \ours]
For \(N=1\), \(A=\xi_1X\), and the commutant is the ordinary centralizer of \(X\). For \(N=2\), write \(A=\xi_1X+\xi_2Y\). Then
\[
 C_{\bar0}(A)=\{B_0+\xi_1\xi_2D:[X,B_0]=[Y,B_0]=0,
 \ D\in M_n(\mathbb K)\}.
\]
Its dimension is \(\dim Z(X,Y)+n^2\), and its positive-degree ideal is square zero. These are direct specializations, not asymptotic estimates.
\end{remark}

\paragraph{Four generators and matrix size two.}
Here \(M_2((\Lambda_4)_{\bar0})\) has dimension 32. Subtracting any odd scalar matrix \(sI_2\) from \(A\) does not change its even commutant. Thus replace every coefficient \(A_I\) by its traceless part. The division by two is legitimate over \(\mathbb R\) and \(\mathbb C\). Write
\[
 A=\sum_{i=1}^4\xi_i X_i+\sum_{i<j<k}\xi_i\xi_j\xi_k Z_{ijk},
 \qquad X_i,Z_{ijk}\in\mathfrak{sl}_2(\mathbb K),
\]
and let \(r=\dim_{\mathbb K}\operatorname{span}(X_1,X_2,X_3,X_4)\). An invertible linear change of the four Grassmann generators preserves the problem. In the descriptions below the generators have been changed to an adapted basis, and the cubic coefficients have been transformed by the same exterior-algebra automorphism. Put
\[
 H=\begin{pmatrix}1&0\\0&-1\end{pmatrix},\quad
 E=\begin{pmatrix}0&1\\0&0\end{pmatrix},\quad
 F=\begin{pmatrix}0&0\\1&0\end{pmatrix},\qquad
 B=B_0+\sum_{i<j}\xi_i\xi_jB_{ij}+\xi_1\xi_2\xi_3\xi_4D.
\]

\begin{theorem}[All even commutants for \(n=2,N=4\) \ours]
For every odd \(A\) as above, \(D\) is arbitrary and the remaining coefficients are exactly as follows.
\begin{enumerate}[label=(\roman*)]
\item If \(r=3\), choose the linear part \(\xi_1H+\xi_2E+\xi_3F\). Then \(B_0=cI_2\), the three coefficients \(B_{14},B_{24},B_{34}\) are arbitrary scalar matrices, and \(B_{12},B_{13},B_{23}\) are arbitrary subject only to
\[
 [H,B_{23}]-[E,B_{13}]+[F,B_{12}]=0.
\]
The dimension is 17, independently of the cubic part.
\item If \(r=2\), choose the linear part \(\xi_1X+\xi_2Y\), with \(X,Y\) independent in \(\mathfrak{sl}_2\). Then \(B_0\) and \(B_{34}\) are scalar, \(B_{12}\) is arbitrary, and the only other conditions are
\[
 [X,B_{23}]=[Y,B_{13}],\qquad [X,B_{24}]=[Y,B_{14}].
\]
The dimension is 20, independently of the cubic part.
\item If \(r=1\), choose the linear part \(\xi_1X\), \(X\ne0\). Then \(B_0=cI_2+dX\), the coefficients \(B_{12},B_{13},B_{14}\) are arbitrary, and
\[
 B_{jk}=dZ_{1jk}+u_{jk}I_2+v_{jk}X\quad(2\leq j<k\leq4),
 \qquad d[Z_{234},X]=0.
\]
The dimension is 24 if \([Z_{234},X]=0\), and 23 otherwise. This includes nilpotent \(X\).
\item If \(r=0\), every \(B_{ij}\) is arbitrary and \(B_0\) is in the joint ordinary centralizer of the cubic coefficients. The dimension is
\[
 28+\dim Z(Z_{123},Z_{124},Z_{134},Z_{234}),
\]
namely 32, 30, or 29 according as the span of the traceless cubic coefficients has dimension zero, one, or at least two.
\end{enumerate}
\end{theorem}

\begin{proof}
The degree-one equations are \([X_i,B_0]=0\). The degree-three equations are
\begin{equation}
 [X_i,B_{jk}]-[X_j,B_{ik}]+[X_k,B_{ij}]+[Z_{ijk},B_0]=0
 \quad(i<j<k).
 \label{eq:four-generator-equations}
\end{equation}
No term can involve \(D\), since its product with an odd positive-degree monomial has degree greater than four. These are all the equations.

We first record elementary matrix facts used in the case analysis. Every nonzero traceless two-by-two matrix \(X\) is nonscalar and cyclic; choose a vector \(v\) with \(v,Xv\) independent. A commuting matrix is determined by its value \(av+bXv\) at \(v\), and is \(aI_2+bX\). Thus the centralizer of \(X\) is \(\mathbb K I_2+\mathbb K X\), and two independent traceless matrices have joint centralizer \(\mathbb K I_2\). Moreover the trace pairing on \(M_2(\mathbb K)\) is nondegenerate and
\[
 \operatorname{tr}(T[X,U])=\operatorname{tr}([T,X]U).
\]
Consequently the orthogonal complement of \([X,M_2]+[Y,M_2]\) is the joint centralizer of \(X,Y\). For independent traceless \(X,Y\), the image is therefore the three-dimensional traceless space.

When \(r\geq2\), the first equations force \(B_0=cI_2\), so the cubic term in \eqref{eq:four-generator-equations} vanishes. For \(r=3\), a change of generator basis puts the three independent coefficients into the prescribed \(H,E,F\) basis. The equations involving the fourth generator are
\[
 [H,B_{24}]=[E,B_{14}],\quad
 [H,B_{34}]=[F,B_{14}],\quad
 [E,B_{34}]=[F,B_{24}].
\]
Subtract scalar parts and write
\[
 B_{14}=aH+bE+cF,\quad B_{24}=dH+eE+fF,
 \quad B_{34}=gH+hE+iF.
\]
Using \([H,E]=2E\), \([H,F]=-2F\), \([E,F]=H\), the first two equations give \(c=f=b=h=0\) and \(e=i=-a\); the third gives \(d=g=a=0\). Hence all three traceless parts vanish. The remaining equation on \(B_{12},B_{13},B_{23}\) maps a 12-dimensional space onto \(\mathfrak{sl}_2\): commutators with \(H\) give \(E,F\), and with \(E\) give \(H\). Its kernel has dimension nine. The total is \(1+3+9+4=17\).

For \(r=2\), the equations involving \(B_{34}\) say that it commutes with \(X\) and \(Y\), hence is scalar. Each displayed equation for a pair of other coefficients is a surjective map \(M_2\oplus M_2\to\mathfrak{sl}_2\), by the trace-pairing argument. Each solution space has dimension \(8-3=5\). Together with arbitrary \(B_{12}\) and \(D\), this gives \(1+1+4+5+5+4=20\).

For \(r=1\), the degree-one equation says \(B_0=cI_2+dX\). If \(1<j<k\), equation \eqref{eq:four-generator-equations} becomes \([X,B_{jk}-dZ_{1jk}]=0\), giving the asserted expression. The triple \(234\) gives \(d[Z_{234},X]=0\). The three \(B_{1j}\) contribute 12 parameters, the other pairs contribute six, and \(D\) contributes four. The constant term contributes two parameters when the last commutator vanishes and one otherwise. This proves the two dimensions and necessity and sufficiency of the parametrization.

For \(r=0\), equations \eqref{eq:four-generator-equations} reduce to commutation of \(B_0\) with every cubic coefficient. The other seven even coefficients are free, contributing 28 parameters. The elementary centralizer facts give dimensions four, two, and one for the possible joint centralizers. In every case we have used all and only the coefficient equations, so the list is exhaustive.
\end{proof}

\paragraph{Independent verification and exact output.}
\oursComputed\ The accompanying program uses two separately implemented constructions: exterior products by adjacent swaps of ordered tuples, followed by exact SymPy nullspaces; and bitmask coefficient signs, followed by independent Gaussian elimination with rational \texttt{Fraction} arithmetic. It checks 36 cases, including both semisimple and nilpotent rank-one coefficients, every dimension stratum, nonzero cubic terms, and invertible changes of the four generators. Every nullspace basis vector is substituted into the original product \(AB-BA\). All 1,706 ordered basis products in three representative commutants are checked for closure. The dimensions from both constructions agree in every case. The tests for \(N=1,2,3\) and the scalar \(n=1,N=4\) case also agree with direct formulas. The complete inputs and results, with fixed seed 19791015, are in \texttt{code/verify\_berezin\_output.json}. These finite computations check signs and degeneracies; the theorem's universal validity comes from its proof.

The proof was re-derived from \eqref{eq:four-generator-equations}, independently of the rank computations: in particular the rank-one transverse cubic term removes exactly one constant parameter, whereas cubic terms cannot affect either \(r\geq2\) stratum. This check rules out extending the generic dimension 17 to all four-generator matrices. No conjecture is used in the theorem.

\paragraph{Exact obstruction and a next bounded question.}
There is no obstruction to computing a specified finite commutant by \eqref{eq:commutant-coefficients}. The remaining structural task is to describe its parameter strata and multiplication uniformly as matrix size grows; the simplification \(Z(X,Y)=\mathbb K I_2\) for every independent traceless pair is special to size two and fails in higher rank. A realistic next run would fix \(n=3,N=4\), first assume the degree-one coefficients have scalar joint centralizer, and study the ranks of the explicit map
\[
 (B_{ij})_{i<j}\longmapsto
 \bigl([X_i,B_{jk}]-[X_j,B_{ik}]+[X_k,B_{ij}]\bigr)_{i<j<k}.
\]
Its domain and codomain have dimensions 54 and 36. Rank loci are determinantal algebraic sets, but a finite sample of ranks would not classify them. The same run could check Berezin's posthumous writings and specialist bibliographies for an exact historical attribution. The present result establishes neither new instanton data nor a general Yang--Mills statement; the notebook does not provide hypotheses permitting that inference.

\subsection{The surface energy representation: checking a documented proof failure}
\label{dev:energy-diagnostic}

\paragraph{Question and the bounded diagnostic.}
Item 1979-470.7 and GP-1979-06 concern irreducibility of the ordinary energy representation of a current group on a two-dimensional base. The precise all-parameter question is the one in the item: a connected smooth Riemannian surface, a compact semisimple group, a positive density and an invariant positive inner product. \citet[Theorem~6.2]{wallach1987gauge} gives a parameter-restricted theorem, not that unrestricted assertion. The later all-dimensional claim in \citet{shimada2005irreducibility} has a documented flaw: \citet[p.~1, revised version of 15 January 2011]{hasebe2010white} identifies the false nonabelian differential identity used in Lemma~4.6. The following explicit diagnostic confirms that flaw. It does not prove reducibility, repair the argument, or claim a new discovery.

\begin{example}[The logarithmic derivative of an exponential \ours]
There is a smooth compactly supported map \(\Phi:\mathbb R^2\to\mathfrak{su}(2)\) for which
\[
 d\Phi\ne\beta(\exp\Phi),\qquad
 \beta(\psi)=d\psi\,\psi^{-1}.
\]
The failure persists locally on any smooth surface containing a coordinate disk.
\end{example}
\begin{proof}
Put
\[
 X=\begin{pmatrix}0&1\\-1&0\end{pmatrix},\qquad
 Y=\begin{pmatrix}i&0\\0&-i\end{pmatrix}.
\]
Both matrices are traceless and anti-Hermitian, \(X^2=Y^2=-I\), and \(XY=-YX\). First take \(\Phi(x,y)=xX+yY\) near \(p=(\pi/2,0)\). With \(r=(x^2+y^2)^{1/2}\), the exponential is
\[
 e^{\Phi}=\cos r\,I+\frac{\sin r}{r}(xX+yY),
\]
where the apparent singularity at zero is removable. Along \(y=0\), for positive \(x\), direct differentiation and multiplication give
\begin{align*}
 \bigl(\partial_y e^{\Phi}\bigr)e^{-\Phi}
 &=\frac{\sin x}{x}Y(\cos x\,I-\sin x\,X)\\
 &=\frac{\sin(2x)}{2x}Y+\frac{\sin^2x}{x}XY.
\end{align*}
At \(p\), this is \((2/\pi)XY\), whereas \(\partial_y\Phi=Y\). The two matrices are different: \(Y\) is nonzero diagonal and \(XY\) is off diagonal. Now multiply \(xX+yY\) by a smooth compactly supported bump function that is identically one near \(p\). This leaves the value and all derivatives used above unchanged and makes \(\Phi\) compactly supported. Its exponential lies in \(SU(2)\) everywhere and is the identity off that compact support. The same local construction in a surface coordinate disk, extended by zero, proves the last assertion.
\end{proof}

\paragraph{Independent check and the remaining obstruction.}
\oursComputed\ A second derivation uses
\[
 (d e^{\Phi})e^{-\Phi}
 =\int_0^1 e^{t\Phi}(d\Phi)e^{-t\Phi}\,dt,
\]
obtained by differentiating the exponential power series and integrating the resulting ordered products. For \(\Phi=xX\) and the tangent \(Y\), conjugation gives \(\cos(2tx)Y+\sin(2tx)XY\); integration reproduces both coefficients above. The symbolic program \texttt{code/verify\_energy\_identity.py} differentiates the two-by-two exponential directly and, independently, integrates the conjugated tangent matrix. It checks the points \(x=\pi/2\) and \(x=\pi/4\), the Lie-algebra hypotheses, and the expansion \(Y+xXY+O(x^2)\) at zero. All checks agree; the exact output is included.

The valid derivative has commutator corrections. Any attempted repair of the cited proof must control those terms in its actual limiting argument; replacing them by zero is invalid even for the compact semisimple group \(SU(2)\). A further bounded run could re-examine that specific estimate with the integral derivative, or seek an explicit intertwiner in a specified surface/coupling case. Neither task is completed here, and the checked literature does not justify an all-parameter resolution as of November 2026. This subsection is an annotation supporting the unresolved-status assessment, not a theorem about the representation's commutant.

\section{Problem index}\label{problems:1979}
Earlier-year occurrence identifiers are retained as text; current-year identifiers link to their entries.

\subsection{GP-1979-01: Canonical forms of even Grassmann matrices}\label{prob:GP-1979-01}
Classify \(L_{p,q}(\Lambda_N)\) under its group of even units, for fixed finite \(N,p,q\), without the simple-body-spectrum hypothesis.
\par\noindent\textit{Occurrences:} \hyperref[item:1979-468.11]{1979-468.11}.
\par\noindent\textit{Status (2026-11):} As of November 2026, canonical algorithms follow from finite-dimensional algebra similarity and Sergeichuk's linear-matrix-problem theorem. A short Jordan-block classification is a stronger demand; the chapter supplies special cases and a simultaneous-similarity embedding.

\subsection{GP-1979-02: Even commutants of odd Grassmann matrices}\label{prob:GP-1979-02}
Given arbitrary \(A\in M_n(\Lambda_N^{\bar1})\), describe \(C_{\bar0}(A)\). The notebook specifically highlights \(N=4\).
\par\noindent\textit{Occurrences:} \hyperref[item:1979-468.12]{1979-468.12}.
\par\noindent\textit{Status (2026-11):} As of November 2026, the literal finite problem is solved by expanding the commutator in Grassmann monomials and computing a finite linear kernel. The chapter gives a complete structural \(n=2,N=4\) result. Any stronger all-size normal-form request is an interpretation not specified in the manuscript; no global openness or novelty claim is made.

\subsection{GP-1977-11: Describe projective bundles by finite exterior modules}\label{prob:GP-1977-11}
Describe algebraic vector bundles on complex projective space by proper finite graded exterior-algebra modules, up to addition of free modules.
\par\noindent\textit{Occurrences:} 1977-437.2, 1977-437.11, 1977-437.12, 1977-438.2, 1977-438.3, 1977-438.4, 1977-438.5, \hyperref[item:1979-468.13]{1979-468.13}.
\par\noindent\textit{Status (2026-11):} Established by BGG in 1978; the 1979 occurrence is exposition. Instanton bundles require extra reality and twistor-line conditions.

\subsection{GP-1979-03: Explicit maps in the BGG and Beilinson constructions}\label{prob:GP-1979-03}
Construct actual differentials from finite sheaf data, not just the terms of the associated monad or exterior resolution.
\par\noindent\textit{Occurrences:} \hyperref[item:1979-468.14]{1979-468.14}.
\par\noindent\textit{Status (2026-11):} Developed explicitly by Eisenbud--Fløystad--Schreyer, preprint 2001, publication 2003; see Theorem 6.1 and the construction preceding it.

\subsection{GP-1979-04: A Grassmannian counterpart of the projective-space linear model}\label{prob:GP-1979-04}
Construct a finite algebraic or homogeneous-bundle model for \(D^b(\operatorname{Coh}\operatorname{Gr}(k,n))\) over \(\mathbb C\).
\par\noindent\textit{Occurrences:} \hyperref[item:1979-468.15]{1979-468.15}.
\par\noindent\textit{Status (2026-11):} Resolved in the reconstructed derived-category sense by Kapranov, 1984/1985 and the fuller 1988 Theorem 3.4. This does not classify every bundle by a finite list of invariants.

\subsection{GP-1979-06: The two-dimensional energy representation}\label{prob:GP-1979-06}
Is the ordinary energy representation of \(C^\infty_c(X,G)\) irreducible for every connected smooth surface \(X\), compact semisimple \(G\), and invariant positive covariance data?
\par\noindent\textit{Occurrences:} \hyperref[item:1979-470.7]{1979-470.7}.
\par\noindent\textit{Status (2026-11):} As of November 2026, a root/covariance inequality gives irreducibility (Wallach 1987, Theorem 6.2). The claimed all-dimensional proof by Shimada has an explicitly documented gap. No unrestricted solution was verified in targeted searches; this is a bounded literature-status conclusion, not an assertion of exhaustive current openness.

\subsection{GP-1979-05: Universality of the difference-curve resonance construction}\label{prob:GP-1979-05}
Does every regular local two-dimensional three-wave dispersion law with infinitely many independent collision invariants admit the difference-curve parametrization on compatible resonance branches? This reconstructs the immediately preceding infinitely-many-first-integrals programme; it is not the earlier, broader request for one additional invariant.
\par\noindent\textit{Occurrences:} \hyperref[item:1979-463.3]{1979-463.3}.
\par\noindent\textit{Status (2026-11):} As of November 2026, Balk--Ferapontov's 1994 publication explicitly gives a classification of three-wave interactions with infinitely many invariants. Its full theorem was not retrieved, so exact local and branch hypotheses are unverified. The one-extra-invariant reading is distinct: finite-invariant examples were announced in 1991 and their uniqueness is discussed in the 1998 web-geometry development and Balk2014, Sec.2.1. This is a developed historical question with partial publication matching, not a claim of current openness.

\subsection{GP-1979-20: Local abelianization of regular first-class constraints}\label{prob:GP-1979-20}
Given independent first-class constraints defining a regular coisotropic submanifold, can local canonical coordinates and an invertible change of constraint generators make all constraint brackets zero? This specifies the local reading of the notebook's question about constant structure functions.
\par\noindent\textit{Occurrences:} \hyperref[item:1979-471.2]{1979-471.2}.
\par\noindent\textit{Status (2026-11):} As of November 2026 the stated local regular version is affirmative. The notebook does not specify the allowed transformations or any global quantifier; it must not be converted into a universal global abelianization conjecture. Dirac's generator equivalence is older; the constructive 1995 preprint/1996 article is a later explicit treatment.

\subsection{GP-1975-10: Continuum confinement and the static potential}\label{prob:GP-1975-10}
For four-dimensional continuum pure \(SU(N)\) Yang--Mills theory at zero temperature, with fixed finite \(N\ge2\) and external fundamental static sources, establish a well-defined subtracted static potential with positive asymptotic string tension:
\[
\lim_{r\to\infty}\frac{V(r)}r=\tau>0.
\]
There are no dynamical fundamental quarks in this formulation. Occurrence 1979-480.3 specializes to \(SU(3)\). Occurrence 1979-482.14 returns to the general confinement goal and proposes normalized large-rank loop equations as a route; that limiting method is not identified with the fixed-rank statement. The stronger constant-width flux-tube hypothesis is excluded from this problem identity.
\par\noindent\textit{Occurrences:} 1975-377.10, \hyperref[item:1979-480.3]{1979-480.3}, \hyperref[item:1979-482.14]{1979-482.14}.
\par\noindent\textit{Status (2026-11):} As of 10 November 2026, the verified results are partial. Wilson's strong-coupling lattice area law \citep[Secs.~IV and VI]{wilson1974confinement} and the controlled large-rank lattice expansion \citep[Theorem~3.1 and Corollaries~3.2--3.3]{chatterjee2015string} do not provide the asserted continuum proof. With light dynamical quarks, string breaking changes the asymptotic static-energy question \citep[Sec.~4]{bulava2019string}. The continuum existence-and-mass-gap problem and confinement are separately stated in \citet[Sec.~5]{jaffewitten2000yangmills}; the former is still listed as unsolved by \citet{clay2026yangmills}. The reused older identifier follows the static-potential goal, while the distinct proposed methods and physical regimes remain explicit.

\subsection{GP-1979-41: Kazhdan--Lusztig multiplicity and cohomological conjectures}\label{prob:GP-1979-41}
For a regular integral block of category \(\mathcal O\), recover Verma composition multiplicities and the graded Ext dimensions from the Kazhdan--Lusztig polynomials, with the precise antidominant indexing in occurrences 1979-483.6--7.
\par\noindent\textit{Occurrences:} \hyperref[item:1979-483.6]{1979-483.6}, \hyperref[item:1979-483.7]{1979-483.7}.
\par\noindent\textit{Status (2026-11):} As of 10 November 2026: solved. The independent geometric proofs circulated by the 1980 Toru\'n meeting and were published in 1981; see \citet[pp.~43--44]{brylinski1981differential}, \citet{brylinskikashiwara1981holonomic}, and \citet{beilinsonbernstein1981localisation}. This status concerns complex characteristic-zero regular integral category \(\mathcal O\), not unrestricted modular analogues.

\section{Corrections to the transcription and translation}\label{corr:1979}
The evidence distinguishes a transcription correction from a mathematical emendation of wording faithfully preserved in the notebook.
\subsection*{P000840 (both; provenance): carry-in before dated 24 September heading}
\noindent\textit{Transcription reads.} Catalogue record 462 page interval begins at P000840\par
\noindent\textit{Critical reading.} Retain the vertex-model carry-in as prior occurrence 1978-460.12--20; do not attribute it to Kallosh.\par
\noindent\textit{Evidence.} Both languages and original pages; preceding index already contains these items. New Kallosh heading occurs on P000841.\par
\subsection*{P000841 (both; ambiguity): formula (2)}
\noindent\textit{Transcription reads.} \(s-\frac13\)\par
\noindent\textit{Critical reading.} Unresolved: P000853 (39) instead has \(s-\frac23\); do not silently choose.\par
\noindent\textit{Evidence.} Original P000841 and P000853 PNGs inspected; operator and Pontryagin normalization are not recovered from inaccessible original paper.\par
\subsection*{P000842 (both; faithful error): exceptional symmetry}
\noindent\textit{Transcription reads.} \(E_7\to SU(8)\); \(E_7\) possibly called compact.\par
\noindent\textit{Critical reading.} For the identified maximal-supergravity scalar geometry use \(E_{7(7)}/SU(8)\), with appropriate finite quotient; the arrow is not a stated group map.\par
\noindent\textit{Evidence.} PNG and primary Duff (1982) discussion; real-form distinction is material.\par
\subsection*{P000842 (both; ambiguity): formula (3) shifted spin}
\noindent\textit{Transcription reads.} \(A\overset?{-}\frac12\)\par
\noindent\textit{Critical reading.} Retain uncertainty of the overwritten sign.\par
\noindent\textit{Evidence.} Original PNG shows overinking; self-duality hypothesis is explicit below. Neither transcription resolves it.\par
\subsection*{P000842 (both; faithful error): nonzero modes and B4}
\noindent\textit{Transcription reads.} m is a nonzero-mode count\par
\noindent\textit{Critical reading.} Interpret m as the regularized contribution \(\zeta_P(0)\) once the elliptic setup is fixed.\par
\noindent\textit{Evidence.} P000852 rejects the common-multiplicity conjecture; P000853 explicitly calls m a regularized sum. Gilkey, Sec. 1.10.\par
\subsection*{P000843 (both; faithful error): nontrivial collision invariant}
\noindent\textit{Transcription reads.} f does not differ from omega by a linear function\par
\noindent\textit{Critical reading.} Exclude all \(f=a\omega+b\cdot k\).\par
\noindent\textit{Evidence.} Otherwise \(f=2\omega\) is an immediate trivial example. Zakharov--Schulman (1988), eq. (3.9).\par
\subsection*{P000843 (both; faithful error): resonance surface}
\noindent\textit{Transcription reads.} omega nonlinear implies Gamma is a surface\par
\noindent\textit{Critical reading.} Assume a nonempty regular level-set component.\par
\noindent\textit{Evidence.} A nonlinear function can have empty or singular resonance sets; regularity follows from the implicit-function theorem only at noncritical zeros.\par
\subsection*{P000844 (both; faithful error): GLM equation and next line}
\noindent\textit{Transcription reads.} K=-F\par
\noindent\textit{Critical reading.} K=-F only at linear order in F, unless the composition term vanishes.\par
\noindent\textit{Evidence.} Original PNG shows F+K+K circled-plus F=0; substitution leaves the quadratic composition term.\par
\subsection*{P000844 (both; faithful error): evolution equations versus exponential}
\noindent\textit{Transcription reads.} t paired with L0, y with A0; exponential uses ell y and a t\par
\noindent\textit{Critical reading.} Restore a consistent assignment and sign convention; no unique emendation certified.\par
\noindent\textit{Evidence.} Original PNG confirms the mismatch; the text later calls ell and a symbols of L0 and A0.\par
\subsection*{P000844 (both; faithful error): last line}
\noindent\textit{Transcription reads.} Quadratic and cubic scalar symbols give KdV\par
\noindent\textit{Critical reading.} With both y and t present they give the KP construction; the y-independent reduction gives KdV.\par
\noindent\textit{Evidence.} P000845 explicitly names KP and computes its two-variable dispersion.\par
\subsection*{P000845 (both; ambiguity): KP evolution}
\noindent\textit{Transcription reads.} Expression ends without relation\par
\noindent\textit{Critical reading.} Read as an evolution equation equal to zero after fixing the uncertain sign.\par
\noindent\textit{Evidence.} Original PNG and both transcriptions; the final equality is absent.\par
\subsection*{P000845 (both; faithful error): complex Hamilton equation}
\noindent\textit{Transcription reads.} \(i a_t=\delta H/\delta a\)\par
\noindent\textit{Critical reading.} \(i a_t=\delta H/\delta\bar a\)\par
\noindent\textit{Evidence.} No bar visible in original PNG. Variation of \(H_0=\int\omega|a|^2\) and Zakharov--Schulman (1988) eq. (1.1) fix the standard canonical convention.\par
\subsection*{P000845-P000846 (both; faithful error): four-wave coefficient and vanishing claim}
\noindent\textit{Transcription reads.} Printed two fractions give w|Gamma=0\par
\noindent\textit{Critical reading.} Printed w does not vanish: the stated resonant quartet gives 6/5. Intended coefficient not recovered.\par
\noindent\textit{Evidence.} Original P000845 confirms second denominator pluses; exact rational check in annotation\_checks.json. Merely changing those pluses gives 2, not 0.\par
\subsection*{P000846 (both; faithful error): Poincare resonance criterion}
\noindent\textit{Transcription reads.} \(h_k\) restricted to resonance=0; infinitely many failures exclude an integral\par
\noindent\textit{Critical reading.} The first-order condition is \((k\cdot\nabla F_0)h_k=0\) on \(k\cdot\nabla H_0=0\); nonintegrability requires additional resonance richness.\par
\noindent\textit{Evidence.} Original PNG; direct expansion of {H,F}=0.\par
\subsection*{P000848 (both; faithful error): Hamilton square-root paragraph}
\noindent\textit{Transcription reads.} Anticommuting a,b,c with squares−1 plus unit give quaternions\par
\noindent\textit{Critical reading.} These relations give Cl(0,3)=H direct-sum H; add ab=c or a volume-element quotient to obtain H.\par
\noindent\textit{Evidence.} Original PNG inspected; universal monomial basis has dimension 8, not 4.\par
\subsection*{P000849 (both; provenance): bibliography Kallosh}
\noindent\textit{Transcription reads.} JETP Letters 29 (1979), 192; 493\par
\noindent\textit{Critical reading.} These are Russian page references to two papers, not an erratum; English pages 172 and 449 respectively.\par
\noindent\textit{Evidence.} Duff (1982) bibliography distinguishes 1979a Super-self-duality and 1979b Structure of divergences in supergravitation.\par
\subsection*{P000849 (both; faithful error): equation (3)}
\noindent\textit{Transcription reads.} \(\sigma_i\sigma_j=\delta_{ij}+\epsilon_{ijk}\sigma_k\)\par
\noindent\textit{Critical reading.} For standard Hermitian Pauli matrices add i; the scaled matrices on P000850 also require factors \(1/2\) and \(1/\sqrt2\).\par
\noindent\textit{Evidence.} PNG confirms missing i; direct multiplication.\par
\subsection*{P000849 (both; faithful error): equation (4)}
\noindent\textit{Transcription reads.} \(\sigma_a\bar\sigma_b-\bar\sigma_b\sigma_a\)\par
\noindent\textit{Critical reading.} Use a conventionally normalized antisymmetrization \(\sigma_a\bar\sigma_b-\sigma_b\bar\sigma_a\).\par
\noindent\textit{Evidence.} PNG confirms the displayed order; spacetime-label antisymmetry is required.\par
\subsection*{P000850 (both; faithful error): Hodge eigenvalues}
\noindent\textit{Transcription reads.} Eigenvalues of star are±1\par
\noindent\textit{Critical reading.} Require Euclidean signature; in Lorentz signature they are±i on complexified two-forms.\par
\noindent\textit{Evidence.} The notes alternate Lorentz and \(SO(4)\) conventions; star squared has the signature-dependent sign.\par
\subsection*{P000850 (both; faithful error): equation (10) index ordering}
\noindent\textit{Transcription reads.} \(F_{\alpha\beta\dot\gamma\dot\delta}\) on the left; \(\alpha,\gamma\) and \(\dot\beta,\dot\delta\) are paired on the right.\par
\noindent\textit{Critical reading.} Use \(F_{\alpha\dot\beta\gamma\dot\delta}\) for the indicated vector-index pairing.\par
\noindent\textit{Evidence.} Original PNG confirms mismatch; marginal invariant tensor decomposition fixes the intended identity.\par
\subsection*{P000850 (both; faithful error): self-dual curvature}
\noindent\textit{Transcription reads.} Self-duality implies Einstein equations\par
\noindent\textit{Critical reading.} Require self-duality of the full Riemann curvature; Weyl self-duality alone is insufficient.\par
\noindent\textit{Evidence.} PNG formula(13) acts on full Riemann tensor; curvature-block decomposition verifies Ricci-flat conclusion.\par
\subsection*{P000851 (both; faithful error): equation (18)}
\noindent\textit{Transcription reads.} \(\Psi_\pm=(1\pm\gamma_5)\Psi\)\par
\noindent\textit{Critical reading.} The idempotent projections are \((1\pm\gamma_5)/2\); the printed vectors are twice projected spinors.\par
\noindent\textit{Evidence.} PNG confirms no factor \(1/2\); direct idempotence check.\par
\subsection*{P000851 (both; faithful error): mixed differential forms programme}
\noindent\textit{Transcription reads.} Needed: anticommuting dx\par
\noindent\textit{Critical reading.} Ordinary exterior dx already anticommute; the intended extra calculus must include even differentials of odd coordinates.\par
\noindent\textit{Evidence.} PNG inspected; Berezin--Leites supermanifold convention. No unique analytic index theorem specified.\par
\subsection*{P000851 (both; faithful error): spin 5/2 interaction}
\noindent\textit{Transcription reads.} No interaction can be constructed\par
\noindent\textit{Critical reading.} Restrict to the naive minimal coupling on general gravitational backgrounds; do not assert universal impossibility.\par
\noindent\textit{Evidence.} Aragone--Deser (1979) primary abstract itself records self-dual exceptions.\par
\subsection*{P000852 (both; provenance): Gilkey date}
\noindent\textit{Transcription reads.} 1975;1979 exists\par
\noindent\textit{Critical reading.} Identify the checked 1979 Recursion relations paper without deleting the older-date note.\par
\noindent\textit{Evidence.} The original scan contains both dates. Numdam metadata confirms volume 38, number 2 (1979), pages 201--240.\par
\subsection*{P000852 (both; provenance): multiplicity conjecture beside formula (31)}
\noindent\textit{Transcription reads.} m is the common multiplicity of every nonzero eigenvalue\par
\noindent\textit{Critical reading.} Retain the manuscript's explicit rejection: incorrect. Use regularized spectral contribution.\par
\noindent\textit{Evidence.} The original scan has an arrow from the underlined Russian word for incorrect to the conjecture attributed to S.~Gelfand.\par
\subsection*{P000853 (both; faithful error): Euler integral (44)}
\noindent\textit{Transcription reads.} Integral of starR contracted with starR\par
\noindent\textit{Critical reading.} Use normalized Euler density \(|\operatorname{Riem}|^2-4|\operatorname{Ric}|^2+R^2\); same-star contraction alone is \(|\operatorname{Riem}|^2\).\par
\noindent\textit{Evidence.} Original PNG has both stars. General Chern--Weil formula and normalization in Gilkey, Ch. 2.\par
\subsection*{P000853 (both; faithful error): Pontryagin integral (45)}
\noindent\textit{Transcription reads.} p=int starR R without normalization\par
\noindent\textit{Critical reading.} State the p1 convention and include its normalization; equivalently \(p_1=-\operatorname{tr}(\Omega\wedge\Omega)/(8\pi^2)\).\par
\noindent\textit{Evidence.} Original PNG; absent factor affects numerical index formula comparison.\par
\subsection*{P000853 (both; ambiguity): equation (43)}
\noindent\textit{Transcription reads.} \(t^{(k-4/2}\), with the malformed exponent as transcribed.\par
\noindent\textit{Critical reading.} \(t^{(k-4)/2}\)\par
\noindent\textit{Evidence.} Earlier formula(29) gives the four-dimensional exponent unambiguously; scan's closing parenthesis is unclear.\par
\subsection*{P000853 (both; faithful error): heat-kernel spectral sum (40)}
\noindent\textit{Transcription reads.} \(\phi_i(x)\phi_i(x')\)\par
\noindent\textit{Critical reading.} For complex bundle sections use \(\phi_i(x)\otimes\phi_i(x')^*\) and a fibre trace on the diagonal.\par
\noindent\textit{Evidence.} PNG lacks the adjoint; general spectral-kernel formula.\par
\subsection*{P000853 (both; faithful error): regularized-count claim}
\noindent\textit{Transcription reads.} m is an integer on deep topological grounds\par
\noindent\textit{Critical reading.} Individual \(\zeta_P(0)\) need not be integral; for scalar S4 it is−61/90.\par
\noindent\textit{Evidence.} Gilkey heat coefficient; exact rational substitution saved in annotation\_checks.json. An index combination can be integral.\par
\subsection*{P000855 (both; faithful error): graph chart}
\noindent\textit{Transcription reads.} Any other two-plane is a graph\par
\noindent\textit{Critical reading.} Only two-planes transverse to Lprime are graphs.\par
\noindent\textit{Evidence.} P000856 explicitly restores this hypothesis; graph projection proof.\par
\subsection*{P000855-P000856 (both; faithful error): mixed Hermitian pairings}
\noindent\textit{Transcription reads.} All other scalar products zero after \(h(e_A,e^{A'})=i\)\par
\noindent\textit{Critical reading.} The reversed pairings are−i by Hermitian symmetry; only remaining independent pairings vanish.\par
\noindent\textit{Evidence.} Displayed Hermitian matrix(6) already has lower block−iI.\par
\subsection*{P000857 (both; faithful error): equation (21) group identification}
\noindent\textit{Transcription reads.} Affine action with arbitrary \(GL_2(\mathbb C)\) is the Poincare group\par
\noindent\textit{Critical reading.} It includes dilations; impose \(|\det a|=1\) and quotient scalar \(U(1)\) for the isometric connected Lorentz part.\par
\noindent\textit{Evidence.} Original PNG; \(\det(a\,dU\,a^*)=|\det a|^2\det(dU)\).\par
\subsection*{P000857-P000858 (both; faithful error): Pauli-coordinate labels}
\noindent\textit{Transcription reads.} u0,u1,u3,u4 then u0,u1,u2,u3\par
\noindent\textit{Critical reading.} Relabel the three spatial coordinates consistently.\par
\noindent\textit{Evidence.} Original PNGs; determinant uses all three spatial squares.\par
\subsection*{P000858 (both; faithful error): metric (25) last exponent}
\noindent\textit{Transcription reads.} \((du_3)^3\)\par
\noindent\textit{Critical reading.} \((du_3)^2\)\par
\noindent\textit{Evidence.} Original PNG plainly has 3; determinant calculation requires a quadratic form.\par
\subsection*{P000858 (both; faithful error): point/ray sentence}
\noindent\textit{Transcription reads.} Points of M biject with rays of N\par
\noindent\textit{Critical reading.} Points of M correspond to isotropic 2-planes; null projective twistors correspond to null geodesics of M.\par
\noindent\textit{Evidence.} Original PNG; following line correctly gives light rays as isotropic projective points. Dimension 4 versus 5 precludes the preceding bijection.\par
\subsection*{P000862 (both; faithful error): BGG conclusion}
\noindent\textit{Transcription reads.} Such exterior modules classify instantons\par
\noindent\textit{Critical reading.} They classify projective bundles under the proper-module condition and stable equivalence; instantons require extra reality and twistor-line conditions.\par
\noindent\textit{Evidence.} BGG (1978), Theorems 1--3; manuscript's preceding heading says vector bundles.\par
\subsection*{P000862 (both; ambiguity): speaker heading}
\noindent\textit{Transcription reads.} Yu.Shabbat or Yu. Shabat\par
\noindent\textit{Critical reading.} Preserve manuscript initial Yu. and catalogue G. B. Shabat as an attribution discrepancy.\par
\noindent\textit{Evidence.} Original PNG gives Russian Yu.; exact catalogue record 469 supplies G. B.; no silent identification.\par
\subsection*{P000863 (both; faithful error): Abel image comparison}
\noindent\textit{Transcription reads.} X differs from−jX\par
\noindent\textit{Critical reading.} Compare j(X) with translates of−j(X).\par
\noindent\textit{Evidence.} Original PNG lacks j on left; the assertion is about embedded Abel curves, not abstractX.\par
\subsection*{P000863 (both; faithful error): Lie--Wirtinger assertion}
\noindent\textit{Transcription reads.} All translation surfaces and small deformations are theta divisors\par
\noindent\textit{Critical reading.} Require double translation, explicit nondegeneracy and mixing condition(*); allow generalized Jacobians. Universal smoothability is unverified.\par
\noindent\textit{Evidence.} Little (1983), Definition 2.2, Proposition 3.7 and Theorem 3.9; source omits hypotheses.\par
\subsection*{P000863-P000864 (both; faithful error): Killing scalar product}
\noindent\textit{Transcription reads.} Killing form used as positive Hilbert metric\par
\noindent\textit{Critical reading.} Use minus Killing form for compact semisimple real Lie algebra, or any positive Ad-invariant inner product.\par
\noindent\textit{Evidence.} Original PNG and Albeverio et al. (2015), Sec. 2.\par
\subsection*{P000864 (both; faithful error): Gaussian Hilbert-space measure}
\noindent\textit{Transcription reads.} Covariance-I Gaussian probability on infinite-dimensional H\par
\noindent\textit{Critical reading.} Use a Gaussian distribution-space realization or bosonic Fock space; H is the Cameron--Martin space.\par
\noindent\textit{Evidence.} Wallach (1987), Secs. 1 and 5; independent coordinate-square sum diverges almost surely.\par
\subsection*{P000864 (both; faithful error): affine action (2)}
\noindent\textit{Transcription reads.} \(\operatorname{Ad}(\phi)\omega+\phi^{-1}d\phi\)\par
\noindent\textit{Critical reading.} For a left action use \(\operatorname{Ad}(\phi)\omega+d\phi\,\phi^{-1}\); alternatively change the action convention consistently.\par
\noindent\textit{Evidence.} Original PNG confirms order; product differentiation contradicts the stated homomorphism with the printed mixture.\par
\subsection*{P000865 (both; faithful error): last factor claim}
\noindent\textit{Transcription reads.} \((L^\infty\cup W_\phi)^{\prime\prime}\) on the stated \(L^2\) is a type III factor\par
\noindent\textit{Critical reading.} If this algebra is factorial, it is \(B(L^2)\), hence typeI. A different representation or crossed product is needed for a type III claim.\par
\noindent\textit{Evidence.} Multiplication by the bounded functions is a MASA; its invariant multiplication commutant lies inside generated algebra. Original PNG confirms expression.\par
\subsection*{P000865 (both; faithful error): path transformations}
\noindent\textit{Transcription reads.} Every smooth path on \(\mathbb R_+\) gives quasi-invariance\par
\noindent\textit{Critical reading.} Require based finite-energy paths, or work on a finite interval with compatible initial point.\par
\noindent\textit{Evidence.} Albeverio et al. (2015), Sec. 2.4; smoothness on an infinite interval does not imply finite energy.\par
\subsection*{P000865 (both; faithful error): stochastic-development domain}
\noindent\textit{Transcription reads.} \((C(\mathbb R_+,G),\mu_0)\) maps to \(C(\mathbb R_+,G)\).\par
\noindent\textit{Critical reading.} The source is Lie-algebra Wiener paths \((C_0(\mathbb R_+,\mathfrak g),\mu_0)\), suitably based.\par
\noindent\textit{Evidence.} Original PNG conflicts with its own preceding line defining \(\mu_0\); stochastic logarithm/development maps.\par
\subsection*{P000866 (both; faithful error): constraint closure formula}
\noindent\textit{Transcription reads.} \(\{T_\alpha,T_\beta\}=u_{\alpha\beta}^{\gamma}T\)\par
\noindent\textit{Critical reading.} \(\{T_\alpha,T_\beta\}=u_{\alpha\beta}^{\gamma}T_\gamma\)\par
\noindent\textit{Evidence.} Scan P000866 has the missing final subscript; the correction is mathematical, not a transcription repair.\par
\subsection*{P000866 (both; provenance): speaker heading}
\noindent\textit{Transcription reads.} Somebody (Baturin) / Некто (Батурин)\par
\noindent\textit{Critical reading.} Catalogue attribution: I. A. Batalin; retain the notebook's tentative name as provenance.\par
\noindent\textit{Evidence.} The original scan retains the tentative parenthesis; the supplied talk catalogue assigns record 471 to Batalin.\par
\subsection*{P000867 (both; ambiguity): Feynman unitarity remark}
\noindent\textit{Transcription reads.} without \(S_{\rm cal}\) there is no unitarity\par
\noindent\textit{Critical reading.} The historical cancellation concerns inclusion of the ghost contribution \(S_{\rm ph.ch}\); keep the scan's subscript in diplomatic text.\par
\noindent\textit{Evidence.} P000867 visibly reads the gauge-fixing subscript. Faddeev--Popov review introduction distinguishes gauge fixing from the required ghost loops.\par
\subsection*{P000869 (both; faithful error): equation (20)}
\noindent\textit{Transcription reads.} \(\tfrac12 h Fh=(h+\delta h)(F+\delta F)(h+\delta h)\)\par
\noindent\textit{Critical reading.} Use the same factor 1/2 on both quadratic forms, and specify equality to the intended variation order.\par
\noindent\textit{Evidence.} Scan confirms the unmatched factor. Setting both variations to zero already exposes the inconsistency.\par
\subsection*{P000869 (both; faithful error): expectation after equation (18)}
\noindent\textit{Transcription reads.} \(\langle\delta h_{\mu\nu}\rangle=\int dh\,e^{i\widetilde S}\)\par
\noindent\textit{Critical reading.} \(\langle\delta h_{\mu\nu}\rangle=Z^{-1}\int dh\,\delta h_{\mu\nu}\,e^{i\widetilde S}\)\par
\noindent\textit{Evidence.} Scan confirms missing insertion and normalization; these cannot be inferred as a transcription loss.\par
\subsection*{P000871 (both; provenance): date heading / catalogue record 473}
\noindent\textit{Transcription reads.} Catalogue: 29-Sep-79 (1979-09-29)\par
\noindent\textit{Critical reading.} Notebook date: 29/X-79, i.e. 1979-10-29; preserve catalogue discrepancy.\par
\noindent\textit{Evidence.} Direct inspection of P000871 shows X, not IX. The block lies after 22 October and before 12 November.\par
\subsection*{P000871 (both; faithful error): scalar example mass term}
\noindent\textit{Transcription reads.} \(m^2\dot\varphi^2\)\par
\noindent\textit{Critical reading.} \(m^2\varphi^2\)\par
\noindent\textit{Evidence.} The dot is present on the scan. This is a mathematical emendation to the intended massive scalar action.\par
\subsection*{P000873 (both; faithful error): equation (23)}
\noindent\textit{Transcription reads.} \(\det(\partial_\theta\psi)=\Delta_\psi(g)\det\lambda(\theta)\)\par
\noindent\textit{Critical reading.} \(\det(\partial_\theta\psi)=\Delta_\psi(f(g,\theta))\det\lambda(\theta)\)\par
\noindent\textit{Evidence.} Scan confirms the shortened argument; equation (22) immediately above contains the transformed point.\par
\subsection*{P000875 (both; faithful error): equation (41)}
\noindent\textit{Transcription reads.} Averaged ghost integral multiplied only by gauge-group volume\par
\noindent\textit{Critical reading.} Also divide by the gauge-average normalization integral already displayed in equation (40).\par
\noindent\textit{Evidence.} P000875 scan shows the denominator in (40) and its absence in (41).\par
\subsection*{P000879 (both; faithful error): Yau theorem}
\noindent\textit{Transcription reads.} Compact Kähler manifold: unique Einstein metric, Lambda nonzero\par
\noindent\textit{Critical reading.} Separate the Calabi prescribed-Ricci theorem in a fixed Kähler class, the Ricci-flat case c1=0 and the negative-c1 Kähler--Einstein theorem.\par
\noindent\textit{Evidence.} Scan has the unrestricted assertion. Yau (1977), Theorem 1 explicitly distinguishes these cases; positive \(c_1\) is not included.\par
\subsection*{P000879 (both; faithful error): stress tensor}
\noindent\textit{Transcription reads.} \(T_{ab}=-\tfrac12\sqrt{-g}\,\delta L_m/\delta g_{ab}\)\par
\noindent\textit{Critical reading.} \(T_{ab}=-2(-g)^{-1/2}\delta S_m/\delta g^{ab}\), or \(T^{ab}=2(-g)^{-1/2}\delta S_m/\delta g_{ab}\).\par
\noindent\textit{Evidence.} Direct scan inspection confirms the factor and covariant metric variation; the matter action/density distinction must also be supplied.\par
\subsection*{P000880 (both; faithful error): first-class definition}
\noindent\textit{Transcription reads.} det Q=0, together with a Hamiltonian condition, defines a first-class system\par
\noindent\textit{Critical reading.} Require every listed constraint bracket to vanish weakly; determinant zero alone permits a mixture of first- and second-class constraints.\par
\noindent\textit{Evidence.} P000880 scan verifies the determinant criterion. A degenerate skew matrix need not vanish.\par
\subsection*{P000881 (both; ambiguity): reduced canonical one-form}
\noindent\textit{Transcription reads.} p dq = p* dq* on the constrained surface\par
\noindent\textit{Critical reading.} Locally p dq|C = p* dq* + dF, with an endpoint term in the action, unless a special compatible primitive has been chosen.\par
\noindent\textit{Evidence.} The written action reduction needs a primitive convention; the symplectic forms agree, but their primitives can differ by an exact differential.\par
\subsection*{P000884 (both; ambiguity): five-term relation}
\noindent\textit{Transcription reads.} Five arguments uj with no relations specified\par
\noindent\textit{Critical reading.} Add a Rogers five-term substitution or an equivalent cyclic pentagon relation before claiming the differential sum is zero.\par
\noindent\textit{Evidence.} The original scan confirms that the argument relations are absent. With independent uj the displayed one-form cannot vanish identically.\par
\subsection*{P000884-P000886 (both; provenance): catalogue order}
\noindent\textit{Transcription reads.} Record 476 Beilinson precedes record 477 Gelfand in the catalogue\par
\noindent\textit{Critical reading.} Preserve catalogue IDs, but place record 477 (P000884) before record 476 (P000885-P000886) in source order.\par
\noindent\textit{Evidence.} P000884 has the dated introductory list; P000885 starts Beilinson and continues its equation numbering. Both source languages follow this physical order.\par
\subsection*{P000885 (both; faithful error): equation (10)}
\noindent\textit{Transcription reads.} Integral around puncture = \(\log|f^{v(g)}/g^{v(f)}|\)\par
\noindent\textit{Critical reading.} For ordinary positive winding, integral = 2 pi \(\log|f^{v(g)}/g^{v(f)}|\)\par
\noindent\textit{Evidence.} Original PNG has no coefficient. Substitution \(f=z^ma\), \(g=z^nb\) gives the factor 2pi; neither the preceding darg nor the loop integral is normalized in the notes.\par
\subsection*{P000885 (both; faithful error): equation (3)}
\noindent\textit{Transcription reads.} \(\int_0^1\log x/(1-x)\,dx=\sum_{n\ge1}1/n^2\)\par
\noindent\textit{Critical reading.} \(\int_0^1\log x/(1-x)\,dx=-\sum_{n\ge1}1/n^2=-\pi^2/6\)\par
\noindent\textit{Evidence.} The original PNG has no minus sign, and the integrand is negative throughout (0,1).\par
\subsection*{P000885-P000886 (both; ambiguity): K2 notation}
\noindent\textit{Transcription reads.} The symbol group is K2(X) of the rational function field\par
\noindent\textit{Critical reading.} Write K2(C(X)) for the field group and distinguish K2(X), whose image is characterized by tame-symbol localization.\par
\noindent\textit{Evidence.} Both source languages switch between curve and function-field notation. Equation (12) itself distinguishes them.\par
\subsection*{P000886 (both; faithful error): symbol symmetry}
\noindent\textit{Transcription reads.} symm. / симм.\par
\noindent\textit{Critical reading.} Multiplicatively skew-symmetric: \(r(g,f)=r(f,g)^{-1}\).\par
\noindent\textit{Evidence.} Scan P000886 confirms симм.; Beilinson (1980), section 1, Lemma 1 and the exponentiated formula give the inverse upon interchange.\par
\subsection*{P000887 (both; ambiguity): equation (4)}
\noindent\textit{Transcription reads.} An overwritten sign and momentum subscript precede the gauge condition\par
\noindent\textit{Critical reading.} Leave the sign/index unresolved; interpret the extended canonical construction using the unambiguous equations (6)-(9).\par
\noindent\textit{Evidence.} Original scan inspection confirms overwriting and does not uniquely determine the literal sign. Both transcriptions flag it.\par
\subsection*{P000892 (both; faithful error): right electron}
\noindent\textit{Transcription reads.} singlet under SU(2) times SU(1)\par
\noindent\textit{Critical reading.} Singlet under SU(2), with nonzero U(1) hypercharge\par
\noindent\textit{Evidence.} The handwritten caption includes the whole product, but the right electron's nonzero electric charge forbids a trivial action of both factors.\par
\subsection*{P000892 (both; faithful error): weak gauge group}
\noindent\textit{Transcription reads.} SU(2) times SU(1)\par
\noindent\textit{Critical reading.} SU(2) times U(1)\par
\noindent\textit{Evidence.} P000892 scan really writes SU(1); P000891 and P000894 give U(1). The electroweak Abelian factor is U(1), while SU(1) is trivial.\par
\subsection*{P000893 (both; faithful error): equations (11) and (13)}
\noindent\textit{Transcription reads.} \(F_a^{ik}F_a^{ik}\)\par
\noindent\textit{Critical reading.} \(F^a_{ik}F^{aik}\), with the chosen spacetime metric.\par
\noindent\textit{Evidence.} The scan confirms both upper pairs; a Lorentz-scalar kinetic term requires contraction with the metric.\par
\subsection*{P000893 (both; ambiguity): trace normalization}
\noindent\textit{Transcription reads.} Adjoint generators have \(\operatorname{Sp}(I_aI_b)=\tfrac12\delta_{ab}\)\par
\noindent\textit{Critical reading.} State the chosen invariant form; with conventional SU(N) generators the ordinary fundamental trace is \(\tfrac12\delta_{ab}\) and the adjoint trace is \(N\delta_{ab}\).\par
\noindent\textit{Evidence.} The original scan explicitly says adjoint. Rescaling a chosen Lie-algebra form is possible, so this is a normalization ambiguity rather than an irreparable construction.\par
\subsection*{P000894 (both; provenance): speaker attribution}
\noindent\textit{Transcription reads.} The same physicist; Igor Yuryevich Kobzar-prime (?)\par
\noindent\textit{Critical reading.} Catalogue: I. Yu. Kobzarev; retain the notebook's tentative spelling in provenance.\par
\noindent\textit{Evidence.} The scan has the first name and patronymic and an uncertain surname ending. It explicitly links this speaker with the preceding physicist of record 479.\par
\subsection*{P000896 (both; ambiguity): Goldstone count}
\noindent\textit{Transcription reads.} N zero modes, then marginal n=dim G\par
\noindent\textit{Critical reading.} The symmetry-enforced subspace has dimension dim G - dim stabilizer; total Hessian nullity can be larger.\par
\noindent\textit{Evidence.} The notebook itself subsequently notes dependent vectors Ia lambda. Thus neither N nor all generators gives the count without a stabilizer qualification.\par
\subsection*{P000896 (both; faithful error): equation (18)}
\noindent\textit{Transcription reads.} \(m^2c=\tfrac12\operatorname{Hess}(P)c\)\par
\noindent\textit{Critical reading.} \(m^2c=\operatorname{Hess}(P)c\) for the displayed canonical kinetic term.\par
\noindent\textit{Evidence.} The PNG verifies the factor 1/2; the previous equation (16) has no such factor and its plane-wave substitution yields the corrected equation.\par
\subsection*{P000897 (both; ambiguity): polarization count}
\noindent\textit{Transcription reads.} 2n becomes 3n; N becomes N-n\par
\noindent\textit{Critical reading.} For r broken generators: 2n+N = 3r+2(n-r)+(N-r). The source's count assumes r=n.\par
\noindent\textit{Evidence.} The same page immediately exhibits one unbroken weak generator, so the fully broken count cannot be applied to all four electroweak vectors.\par
\subsection*{P000898 (both; faithful error): neutral-current prefactor}
\noindent\textit{Transcription reads.} \(q/\cos\theta_W\)\par
\noindent\textit{Critical reading.} \(g/\cos\theta_W\)\par
\noindent\textit{Evidence.} The original image's letter is read as q in both transcriptions. The charge operator is already inside the current; Weinberg (1967), equations 13--14 and PDG (2026), section 10.1 give the weak coupling g.\par
\subsection*{P000898 (both; faithful error): neutral-current sine}
\noindent\textit{Transcription reads.} \(\sin\theta_W\) times the electromagnetic current.\par
\noindent\textit{Critical reading.} \(\sin^2\theta_W\) times the electromagnetic current.\par
\noindent\textit{Evidence.} No exponent is visible in the original PNG. Diagonalizing the neutral gauge fields supplies the square, as in the primary electroweak formulas.\par
\subsection*{P000899 (both; faithful error): mixing matrix}
\noindent\textit{Transcription reads.} V is an orthog. matrix\par
\noindent\textit{Critical reading.} V is unitary; an orthogonal matrix is a CP-conserving special case.\par
\noindent\textit{Evidence.} The scan explicitly reads orthog. A general three-generation quark mixing matrix contains a physical complex CP phase, already present in Kobayashi--Maskawa (1973).\par
\subsection*{P000900 (both; faithful error): Noninvariant expectations}
\noindent\textit{Transcription reads.} All noninvariant quantities have zero expectation.\par
\noindent\textit{Critical reading.} Only observables with zero gauge average must have zero expectation.\par
\noindent\textit{Evidence.} The scan and both transcriptions agree; an observable can have both invariant and noninvariant components. This is a mathematical qualification, not a paleographic correction.\par
\subsection*{P000901 (both; faithful error): Equation (13)}
\noindent\textit{Transcription reads.} \((1+aD_\mu)(1+bD_\nu)(1-aD_\mu)(1-bD_\nu)=1+ab[D_\mu,D_\nu]\).\par
\noindent\textit{Critical reading.} Read only the mixed infinitesimal coefficient, or use inverse exponential transports with remainder terms.\par
\noindent\textit{Evidence.} The scan confirms the four linear factors. Direct expansion leaves pure quadratic terms, so finite equality is false.\par
\subsection*{P000903 (both; faithful error): Equation (23)}
\noindent\textit{Transcription reads.} \(\delta S/\delta A_\mu=\beta D_\mu F_{\mu\nu}\).\par
\noindent\textit{Critical reading.} The variational denominator must have free index \(\nu\); the overall coefficient depends on the action and trace convention.\par
\noindent\textit{Evidence.} The scan explicitly shows the inconsistent free index. Varying the displayed action gives a nonzero scalar multiple of the Yang--Mills equation; it does not give the displayed exact equality.\par
\subsection*{P000903 (both; faithful error): Equation (25)}
\noindent\textit{Transcription reads.} The numerator path integral is multiplied by the partition function.\par
\noindent\textit{Critical reading.} Use the numerator divided by \(Z=\int\mathcal DA\,e^{-S}\).\par
\noindent\textit{Evidence.} The scan has two adjacent parenthesized factors and no visible inverse; normalization requires the second factor to be inverted. Faithful notebook omission.\par
\subsection*{P000903 (both; faithful error): Equation (28)}
\noindent\textit{Transcription reads.} \(\delta S/\delta A=\hbar\delta/\delta A\).\par
\noindent\textit{Critical reading.} An equality inside expectations after integration by parts, not an equality of functions and operators.\par
\noindent\textit{Evidence.} The preceding notebook equation (27) is the valid weak formulation; the displayed operator shorthand suppresses this domain.\par
\subsection*{P000904 (both; faithful error): Equation (33), group qualification}
\noindent\textit{Transcription reads.} For an arbitrary semisimple group.\par
\noindent\textit{Critical reading.} The displayed completeness formula is for \(\mathfrak{su}(N)\) in its defining representation, with specified generator normalization; \(U(N)\) has a related formula.\par
\noindent\textit{Evidence.} The handwritten Russian confirms the broad group qualification. Other representations and groups have different invariant contractions.\par
\subsection*{P000904 (both; faithful error): Equation (35)}
\noindent\textit{Transcription reads.} Factorization error \(O(N^{-2})\) for \(\Psi=\operatorname{Tr}U\).\par
\noindent\textit{Critical reading.} The error \(O(N^{-2})\) belongs to \(\bar\Phi=N^{-1}\Psi\); for a product of \(m\) unnormalized traces it scales as \(O(N^{m-2})\).\par
\noindent\textit{Evidence.} The scan immediately below (35) writes \(\Psi(0)=N\) and proposes the normalization. The mismatch is in the notebook itself.\par
\subsection*{P000905 (both; faithful error): Bruhat subword definition}
\noindent\textit{Transcription reads.} Deleting any subset gives \(w\prime<w\).\par
\noindent\textit{Critical reading.} Use \(w\prime\le w\); require a proper deletion for strict inequality.\par
\noindent\textit{Evidence.} The source permits any subset, including the empty subset; this yields w itself.\par
\subsection*{P000905 (both; faithful error): Symmetric-group presentation}
\noindent\textit{Transcription reads.} The neighboring braid relations generate all the remaining relations.\par
\noindent\textit{Critical reading.} Also include \(s_i s_j=s_j s_i\) for \(|i-j|>1\).\par
\noindent\textit{Evidence.} Both source texts omit the distant commuting relations. They do not follow from involutions and the neighboring braid relations for n at least 4.\par
\subsection*{P000907 (both; provenance): Boundary between records 483 and 484}
\noindent\textit{Transcription reads.} The catalog page window for Zelevinsky ends at P000906.\par
\noindent\textit{Critical reading.} Theorems (17)--(18) and the cell-representation sentence at the top of P000907 still belong to Zelevinsky; MacPherson starts at his named heading.\par
\noindent\textit{Evidence.} The original P000907 scan shows the continuation above the underlined McPherson heading.\par
\subsection*{P000907 (both; faithful error): Intersection-chain dimension conditions}
\noindent\textit{Transcription reads.} The intersection has dimension exactly \(i-c-1\), or \(i-c\) for a bounding chain.\par
\noindent\textit{Critical reading.} Read \(\le i-c-1\) and \(\le i-c\), and require the boundary to be allowable.\par
\noindent\textit{Evidence.} The scan contains no inequality sign. The allowable-chain definition in Goresky--MacPherson 1977/1980 uses upper bounds; disjoint intersections are allowed.\par
\subsection*{P000907 (both; ambiguity): Local polynomial formula}
\noindent\textit{Transcription reads.} \(q^{?}\sum_i\dim IH_i(X_t,X_t-y_s)q^i\).\par
\noindent\textit{Critical reading.} Leave the notebook exponent unresolved; use an explicitly labelled normalized intersection-cohomology stalk formula in the companion.\par
\noindent\textit{Evidence.} The question mark is visible in the original scan. Its repair requires a grading convention, not paleographic substitution.\par
\subsection*{P000907-P000908 (both; provenance): MacPherson date boundary}
\noindent\textit{Transcription reads.} Record 484 nominal window P000906-P000909.\par
\noindent\textit{Critical reading.} The 28 January contribution occupies the lower P000907; P000908 explicitly starts 4 February with Olshanskii, followed by a separate MacPherson talk.\par
\noindent\textit{Evidence.} Date and speaker headings in the scans determine the occurrence attribution.\par
\subsection*{P000908-P000914 (both; provenance): MacPherson February contribution boundaries}
\noindent\textit{Transcription reads.} Attached catalog record 486 has the overlapping range P000907-P000915.\par
\noindent\textit{Critical reading.} The 4 February contribution starts at the MacPherson heading on P000908, ends on P000913, and is followed by a separately dated 11 February contribution on P000914.\par
\noindent\textit{Evidence.} Original dated headings and speaker change on P000908/P000914.\par
\subsection*{P000911 (both; faithful error): Point cases of the configuration complex}
\noindent\textit{Transcription reads.} \(k=1,k=2\Rightarrow C_\sigma\) is a point.\par
\noindent\textit{Critical reading.} \(k=1,k=0\Rightarrow K_\sigma\) is a point.\par
\noindent\textit{Evidence.} The original scan clearly repeats k=2, immediately after defining a two-dimensional complex for that case. MacPherson 1978, Sec.2, says codimension 1 or 0. The transcription is faithful; the mathematics is emended.\par
\subsection*{P000912 (both; faithful error): Enhanced projective configuration dimension}
\noindent\textit{Transcription reads.} Four points in \(\mathbb RP^2\).\par
\noindent\textit{Critical reading.} Four rays in \(\mathbb R^2\), or enhanced points of \(\mathbb RP^1\).\par
\noindent\textit{Evidence.} The exponent 2 is visible in the scan; its own next line and P000913 give the quotient R4/L with dim L=2 and explicitly draw four rays in R2.\par
\subsection*{P000914 (both; faithful error): Projectivization dimension}
\noindent\textit{Transcription reads.} \(k\)-dimensional vector space or \((k+1)\)-dimensional projective space.\par
\noindent\textit{Critical reading.} A \(k\)-dimensional vector space projectivizes to dimension \(k-1\).\par
\noindent\textit{Evidence.} The scan clearly reads k+1, and the accompanying count uses a k-dimensional quotient. This is a faithful notebook dimensional error.\par
\subsection*{P000914-P000915 (both; provenance): Record 487 page window}
\noindent\textit{Transcription reads.} Attached catalog record 487 has the overlapping range P000913-P000916.\par
\noindent\textit{Critical reading.} The 11 February MacPherson fragment is P000914; P000915 explicitly names Zakharov and 18 February 1980.\par
\noindent\textit{Evidence.} The original page headings establish both dates and speakers.\par
\subsection*{P000915 (both; faithful error): Vertical Euler equation and the hydrostatic approximation}
\noindent\textit{Transcription reads.} \(w_t+(u\cdot\nabla)w+w w_z-p_z+1=0\), with only part of the vertical acceleration bracketed for neglect.\par
\noindent\textit{Critical reading.} The pressure term must be \(+p_z\) in the convention used for the horizontal equation; neglect the whole vertical material acceleration to obtain \(p_z+1=0\) and \(p=h-z\).\par
\noindent\textit{Evidence.} The minus sign is present in the original scan, so this is a faithful mathematical error. P000916 explicitly gives \(p_z+1=0\) and \(p=h-z\), fixing the intended convention.\par
\subsection*{P000915-P000918 (both; provenance): Record 488 page envelope and next-year spillover}
\noindent\textit{Transcription reads.} Attached catalog record 488 has the overlapping range P000914-P000918.\par
\noindent\textit{Critical reading.} The Zakharov session dated 18 February 1980 occupies P000915--P000916. P000917 begins 22 September 1980 with Fuchs; P000918 continues that discussion and introduces Graev. The latter pages are carried out to the boundary inventory.\par
\noindent\textit{Evidence.} Original dates and speaker headings; no continuous argument links the Benney fragment to the later intersection-homology and integral-geometry notes.\par
\subsection*{P000917 (both; faithful error): Carry-out: allowable-chain equation}
\noindent\textit{Transcription reads.} \(\dim(\sigma_j\cap X_{2n-i})=j-2i\).\par
\noindent\textit{Critical reading.} \(\dim(|c|\cap X_{2n-i})\le j-i+\bar p(i)\), with the separate boundary condition \(\dim(|\partial c|\cap X_{2n-i})\le j-1-i+\bar p(i)\).\par
\noindent\textit{Evidence.} The scan clearly has equality and \(j-2i\). The definitions in \citet{goreskymacpherson1977dualite,goreskymacpherson1980intersection} give the corrected perversity bounds. This is a faithful mathematical error in the next-year spillover.\par
\subsection*{P000917 (both; faithful error): Carry-out: stratification subscripts}
\noindent\textit{Transcription reads.} \(X_{2n}\supset X_{2n-2}\supset X_{n-3}\supset X_{n-4}\).\par
\noindent\textit{Critical reading.} Keep consistent real-dimensional subscripts, such as \(X_{2n-i}\); for a complex stratification odd levels may repeat.\par
\noindent\textit{Evidence.} Both transcriptions and original scan reproduce the abrupt 2n to n change. This belongs to Fuchs on 22 September 1980, separately inventoried, not to Zakharov record 488.\par
\subsection*{P000917 (both; ambiguity): Carry-out: unspecified L2 comparison}
\noindent\textit{Transcription reads.} Intersection homology coincides with \(L_2\)-cohomology.\par
\noindent\textit{Critical reading.} Specify the metric, regular locus, coefficients and differential domain; the arithmetic Zucker theorem is one precise partial match, not a metric-free universal statement.\par
\noindent\textit{Evidence.} The scan gives none of those choices. Looijenga (1988), p. 3 supplies the arithmetic metric/local-system hypotheses. Boundary inventory B05 preserves the unresolved identification.\par
\subsection*{P000917-P000918 (both; ambiguity): Carry-out: isolated singularity middle degree}
\noindent\textit{Transcription reads.} Only degrees strictly above or below half the dimension are described.\par
\noindent\textit{Critical reading.} For real dimension \(2n\) and field coefficients, the missing homological degree is \(\operatorname{im}(H_n(X_{\rm reg})\to H_n(X))\), under the stated compact conical hypotheses.\par
\noindent\textit{Evidence.} The omission is visible in both transcriptions and scan; the displayed formula is editorial reconstruction and is recorded with its assumptions in boundary inventory B06.\par
\subsection*{P000918 (both; provenance): Carry-out: authors of the second cited paper}
\noindent\textit{Transcription reads.} The FA13:2 and FA14:4 references are placed under Gelfand--Gindikin--Shapiro.\par
\noindent\textit{Critical reading.} FA13:2(1979) is Gelfand--Gindikin--Shapiro; FA14:4(1980) is Gelfand--Graev, Admissible n-dimensional complexes of curves in Rn.\par
\noindent\textit{Evidence.} Primary Math-Net records faa1896 and faa1851 and both full original papers verify titles, authors, volumes and pages. Received dates 25 December 1978 and 16 June 1980.\par
\subsection*{P000918 (both; ambiguity): Carry-out: completeness of the closedness condition}
\noindent\textit{Transcription reads.} For \(k>1\), closedness is the complete admissibility definition.\par
\noindent\textit{Critical reading.} Closedness implies the locality/period condition only under the required size and general-position hypotheses; these are not supplied by the inequality \(k>1\) alone.\par
\noindent\textit{Evidence.} The 1979 primary paper Sec. 1.3 includes both conditions, with a qualifying remark on p. 14. The notebook compresses that qualification.\par
\subsection*{P000918 (both; faithful error): Carry-out: coordinate and parameter counts}
\noindent\textit{Transcription reads.} \(f(x_1,\varphi_1,\ldots,\varphi_n)\), initially \(\xi=(\xi_1,\ldots,\xi_n)\).\par
\noindent\textit{Critical reading.} \(f(x_1,\varphi_2,\ldots,\varphi_n)\), \(\xi\in\Gamma^{n+k-1}\); curve graph functions have indices \(2,\ldots,n\).\par
\noindent\textit{Evidence.} The scan writes an n+1-entry argument list for a function on Rn and separately corrects the parameter count below. Exact conventions are in Gelfand--Gindikin--Shapiro (1979), Sec. 1.1, and Gelfand--Graev (1980), Eq. 4.\par
\subsection*{P000918 (both; faithful error): Carry-out: real admissibility versus local inversion}
\noindent\textit{Transcription reads.} Admissible real curve families are described as families with local inversion formulas.\par
\noindent\textit{Critical reading.} The technical real closed-form/admissibility identities need not give nonzero reconstruction periods; the cited papers obtain genuine local inversion in the complex version.\par
\noindent\textit{Evidence.} \citet[introduction and Sec.~1.3]{gelfandgindikinshapiro1979local} and \citet[Sec.~2, following Definition~1]{gelfandgraev1980admissible} explicitly state that the real coefficient is zero.\par
